\xchapter{Advertisement}




T\textsc{he} following work was written in the early
part of last year, for Messrs. Rivington's ``Theological Library;'' but as it seemed, on its completion, little fitted for the
objects with which that publication has been undertaken, it makes its
appearance in an independent form. Some apology is due to the reader
for the length of the introductory chapter, but it was intended as the
opening of a more extensive undertaking. It may be added, to prevent
mistake, that the theological works cited at the foot of the page, are
referred to for the facts, rather than the opinions they contain;
though some of them, as the ``Defensio Fidei Nicenæ,'' evince
gifts, moral and intellectual, of so high a cast, as to render it a
privilege to be allowed to sit at the feet of their authors, and to
receive the words, which they have been, as it were, commissioned to
deliver.


[\emph{October}, 1833.]




\xchapter{Advertisement to 3rd Edition}




A VERY few words will suffice for the purpose
of explaining in what respects the Third Edition of this Volume
differs from those which preceded it.


Its text has been relieved of some portion of the literary
imperfections necessarily incident to a historical sketch, its author's
first work, and written against time.


Also, some additions have been made to the footnotes. These are
enclosed in brackets, many of them being merely references (under the
abbreviation ``Ath. Tr.'') to his annotations on those
theological Treatises of Athanasius, which he translated for the
Oxford Library of the Fathers.


A few longer Notes, for the most part extracted from other
publications of his, form an Appendix.


The Table of Contents, and the Chronological Table have both been
enlarged.


No change has been made any where affecting the opinions,
sentiments, or speculations contained in the original edition,—though
they are sometimes expressed with a boldness or decision which now
displeases him;—except that two sentences, which needlessly
reflected on the modern Catholic Church, have, without hurting the
context, been relegated to a place by themselves at the end of the Appendix.


\emph{April}, 1871.

\xchapter{Chapter 1. Schools and Parties in and about the Ante-Nicene Church, Considered in their relation to the Arian Heresy}

\xsection{Section 1. The Church of Antioch}

IT is proposed in the following pages to
trace the outlines of the history of Arianism, between the first and
the second General Councils. These are its natural chronological
limits, whether by Arianism we mean a heresy or a party in the Church.
In the Council held at Nicæa, in Bithynia, A.D.
325, it was formally detected and condemned. In the subsequent years
it ran its course, through various modifications of opinion, and with
various success, till the date of the second General Council, held A.D.
381, at Constantinople, when the resources of heretical subtilty being
at length exhausted, the Arian party was ejected from the Catholic
body, and formed into a distinct sect, exterior to it. It is during
this period, while it still maintained its hold upon the creeds and
the government of the Church, that it especially invites the
attention of the student in ecclesiastical history. Afterwards,
Arianism presents nothing new in its doctrine, and is only remarkable
as becoming the animating principle of a second series of
persecutions, when the barbarians of the North, who were infected with
it, possessed themselves of the provinces of the Roman Empire.

The line of history which is thus limited by the two first
Ecumenical Councils, will be found to pass through a variety of
others, provincial and patriarchal, which form easy and intelligible
breaks in it, and present the heretical doctrine in the various stages
of its impiety. These, accordingly, shall be taken as cardinal points
for our narrative to rest upon;—and it will matter little in the
result, whether it be called a history of the Councils, or of Arianism,
between the eras already marked out.

However, it is necessary to direct the reader's attention in the
first place, to the state of parties and schools, in and about the
Church, at the time of its rise, and to the sacred doctrine which it
assailed, in order to obtain a due insight into the history of the
controversy; and the discussions which these subjects involve, will
occupy a considerable portion of the volume. I shall address myself
without delay to this work; and, in this chapter, propose first to
observe upon the connexion of Arianism with the Church of Antioch, and
upon the state and genius of that Church in primitive times. This
shall be the subject of the present section: in those which follow, I
shall consider its relation towards the heathen philosophies and
heresies then prevalent; and towards the Church of Alexandria, to
which, though with very little show of reasoning, it is often
referred. The consideration of the doctrine of the Trinity shall form
the second chapter.

1.

During the third century, the Church of Antioch was more or less
acknowledged as the metropolis of Syria, Cilicia, Phœnicia, Comagene,
Osrhoene, and Mesopotamia, in which provinces it afterwards held
patriarchal sway \footnote{Bingham, Antiq. ix. i.}. It had
been the original centre of Apostolical missions among the heathen \footnote{Acts xi., xiii., xiv.}; and claimed St. Peter himself for its first bishop, who had
been succeeded by Ignatius, Theophilus, Babylas, and others of sacred
memory in the universal Church, as champions and martyrs of the faith
\footnote{Vide Tillemont, Mem. vol. i. \&c.}. The secular importance
of the city added to the influence which accrued to it from the
religious associations thus connected with its name, especially when
the Emperors made Syria the seat of their government. This ancient and
celebrated Church, however, is painfully conspicuous in the middle of
the century, as affording so open a manifestation of the spirit of
Antichrist, as to fulfil almost literally the prophecy of the Apostle
in his second Epistle to the Thessalonians \footnote{Vide Euseb. vii. 30.}. Paulus, of Samosata, who was raised to the see of Antioch not
many years after the martyrdom of Babylas, after holding the
episcopate for ten years, was deposed by a Council of eastern bishops,
held in that city A.D. 272, on the ground of his
heretical notions concerning the nature of Christ. His original
calling seems to have been that of a sophist \footnote{Mosheim, de Reb. ante Constant. sæc. iii.
§ 35.}; how he obtained admittance into the clerical order is
unknown; his elevation, or at least his continuance in the see, he
owed to the celebrated Zenobia \footnote{He was raised to the episcopate at the
commencement of Odenatus's successes against Sapor (Tillemont, Mem.
vol. iv. Chronol.). In the years which followed, he held a civil
magistracy with his ecclesiastical dignity; in the temporalities of
which, moreover, he was upheld by Zenobia, some years after his formal
deposition by the neighbouring bishops. (Basnag. Annal. A.D.
269, § 6.)},
to whom his literary attainments, and his political talents, may be
supposed to have recommended him. Whatever were the personal virtues
of the Queen of the East, who is said to have been a Jewess by birth
or creed, it is not surprising that she was little solicitous for the
credit or influence of the Christian Church within her dominions. The
character of Paulus is consigned to history in the Synodal Letter of
the bishops, written at the time of his condemnation \footnote{Euseb. Hist. vii. 30.}; which, being circulated through the Church, might fairly be
trusted, even though the high names of Gregory of Neocæsarea and
Firmilian were not found in the number of his judges. He is therein
charged with a rapacity, an arrogance, a vulgar ostentation and desire
of popularity, an extraordinary profaneness, and a profligacy, which
cannot but reflect seriously upon the Church and clergy which elected,
and so long endured him. As to his heresy, it is difficult to
determine what were his precise sentiments concerning the Person of
Christ, though they were certainly derogatory of the doctrine of His
absolute divinity and eternal existence. Indeed, it is probable that
he had not any clear view on the solemn subject on which he allowed
himself to speculate; nor had any wish to make proselytes, and form a
party in the Church \footnote{Mosheim, de Reb. ante Const. § 35, n. 1.
[For the opinions of Paulus, vide Athan., Tr. p. 175.]}.
Ancient writers inform us that his heresy was a kind of Judaism in
doctrine, adopted to please his Jewish patroness \footnote{Athan. Epist. ad Monachos, § 71. Theod.
Hær. ii. 8. Chrysost. in Joann. Hom. 7, but Philastr. Hær. § 64,
says that Paulus docuit Zenobiam judaizare.}; and, if originating in this motive, it was not likely to be
very systematic or profound. His habits, too, as a sophist, would
dispose him to employ himself in attacks upon the Catholic doctrine,
and in irregular discussion, rather than in the sincere effort to
obtain some definite conclusions, to satisfy his own mind or convince
others. And the supercilious spirit, which the Synodal letter
describes as leading him to express contempt for the divines who
preceded him at Antioch, would naturally occasion incaution in his
theories, and a carelessness about guarding them from inconsistencies,
even where he perceived them. Indeed, the Primate of Syria had already
obtained the highest post to which ambition could aspire, and had
nothing to labour for; and having, as we find, additional engagements
as a civil magistrate, he would still less be likely to covet the
barren honours of an heresiarch. A sect, it is true, was formed upon
his tenets, and called after his name, and has a place in
ecclesiastical history till the middle of the fifth century; but it
never was a considerable body, and even as early as the date of the
Nicene Council had split into parties, differing by various shades of
heresy from the orthodox faith \footnote{Tillemont Mem. vol. iv. p. 126. Athan. in
Arianos, iv. 30.}. We shall have a more correct notion, then, of the heresy of
Paulus, if we consider him as the founder of a school rather than
of a sect, as encouraging in the Church the use of those disputations
and sceptical inquiries, which belonged to the Academy and other
heathen philosophies, and as scattering up and down the seeds of
errors, which sprang up and bore fruit in the generation after him. In
confirmation of this view, which is suggested by his original
vocation, by the temporal motives which are said to have influenced
him, and by his inconsistencies, it may be observed, that his intimate
friend and fellow-countryman, Lucian, who schismatized or was
excommunicated on his deposition, held heretical tenets of a
diametrically opposite nature, that is, such as were afterwards called
Semi-Arian, Paulus himself advocating a doctrine which nearly
resembled what is commonly called the Sabellian.

More shall be said concerning Paulus of Samosata presently; but now
let us advance to the history of this Lucian, a man of learning \footnote{He was distinguished in biblical
literature, as being the author of a third edition of the Septuagint.
Vide Tillemont, Mem. vol. v. p. 202, 203. Du Pin, cent. iii.}, and at length a martyr, but who may almost be considered the
author of Arianism. It is very common, though evidently illogical, to
attribute the actual rise of one school of opinion to another, from
some real or supposed similarity in their respective tenets. It is
thus, for instance, Platonism, or again, Origenism, has been assigned
as the actual source from which Arianism was derived. Now, Lucian's
doctrine is known to have been precisely the same as that species of
Arianism afterwards called Semi-Arianism \footnote{Bull, Baronius, and others, maintain his
orthodoxy. The Semi-Arians adopted his creed, which is extant. Though
a friend, as it appears, of Paulus, he opposed the Sabellians (by one
of whom he was at length betrayed to the heathen persecutors of the
Church), and this opposition would lead him to incautious statements
of an Arian tendency. Vide below, Section v. Epiphanius (Ancor. 33)
tells us, that he considered the Word in the Person of Christ as the
substitute for a human soul.}; but it is not on that account that I here trace the rise of
Arianism to Lucian. There is an historical, and not merely a doctrinal
connexion between him and the Arian party. In his school are found, in
matter of fact, the names of most of the original advocates of
Arianism, and all those who were the most influential in their
respective Churches throughout the East:—Arius himself, Eusebius of
Nicomedia, Leontius, Eudoxius, Asterius, and others, who will be
familiar to us in the sequel; and these men actually appealed to him
as their authority, and adopted from him the party designation of
Collucianists \footnote{Theod. Hist. i. 5. Epiph. Hær. lxix. 6.
Cave, Hist. Literar. vol. i.}. In
spite of this undoubted connexion between Lucian and the Arians, we
might be tempted to believe, that the assertions of the latter
concerning his heterodoxy, originated in their wish to implicate a man
of high character in the censures which the Church directed against
themselves, were it not undeniable, that during the times of the three
bishops who successively followed Paulus, Lucian was under
excommunication. The Catholics too, are silent in his vindication, and
some of them actually admit his unsoundness in faith \footnote{Theod. Hist. i. 4.}. However, ten or fifteen years before his martyrdom, he was
reconciled to the Church; and we may suppose that he then recanted
whatever was heretical in his creed: and his glorious end was allowed
to wipe out from the recollection of Catholics of succeeding times
those passages of his history, which nevertheless were so miserable in
their results in the age succeeding his own. Chrysostom's panegyric on
the festival of his martyrdom is still extant, Ruffinus mentions him
in honourable terms, and Jerome praises his industry, erudition, and
eloquence in writing \footnote{Vide Tillemont, Mem. vol. v.}.

Such is the historical connexion at the very first sight between
the Arian party and the school of Antioch \footnote{[Vide Appendix, \emph{Syrian School}.]}: corroborative evidence will hereafter appear, in the
similarity of character which exists between the two bodies. At
present, let it be taken as a confirmation of a fact, which Lucian's
history directly proves, that Eusebius the historian, who is suspected
of Arianism, and his friend Paulinus of Tyre, one of its first and
principal supporters, though not pupils of Lucian, were more or less
educated, and the latter ordained at Antioch \footnote{Vales. de Vit. Euseb. et ad Hist. x. i.}; while in addition to the Arian bishops at Nicæa already
mentioned, Theodotus of Laodicea, Gregory of Berytus, Narcissus of
Neronias, and two others, who were all supporters of Arianism at the
Council, were all situated within the ecclesiastical influence, and
some of them in the vicinity of Antioch \footnote{Tillemont, Mem. vol. vi. p. 276.}; so that (besides Arius himself), of thirteen, who according
to Theodoret, arianized at the Council, nine are referable to the
Syrian patriarchate. If we continue the history of the controversy, we
have fresh evidence of the connexion between Antioch and Arianism.
During the interval between the Nicene Council and the death of
Constantius (A.D. 325-361), Antioch is the
metropolis of the heretical, as Alexandria of the orthodox party. At
Antioch, the heresy recommenced its attack upon the Church after the
decision at Nicæa. In a Council held at Antioch, it first showed
itself in the shape of Semi-Arianism, when Lucian's creed was
produced. There, too, in this and subsequent Councils, negotiations on
the doctrine in dispute were conducted with the Western Church. At
Antioch, lastly, and at Tyre, a suffragan see, the sentence of
condemnation was pronounced upon Athanasius.

2.

Hitherto I have spoken of individuals as the authors of the
apostasy which is to engage our attention in the following chapters;
but there is reason to fear that men like Paulus were but symptoms of
a corrupted state of the Church. The history of the times gives us
sufficient evidence of the luxuriousness of Antioch; and it need
scarcely be said, that coldness in faith is the sure consequence of
relaxation of morals \footnote{[Vide a remarkable passage in Origen, on
the pomp of the Bishops of his day, quoted by Neander, Hist. vol. ii.
p. 330, Bohn.]}.
Here, however, passing by this consideration, which is too obvious to
require dwelling upon, I would rather direct the reader's attention to
the particular form which the Antiochene corruptions seem to have
assumed, viz., that of Judaism \footnote{[Lengerke, de Ephræm. Syr. p. 64. traces
the literal interpretation, which was the characteristic of the school
of Antioch, to the example of the Jews.]}; which at that time, it must be recollected, was the
creed of an existing nation, acting upon the Church, and not merely,
as at this day, a system of opinions more or less discoverable among
professing Christians.

The fortunes of the Jewish people had experienced a favourable
change since the reign of Hadrian. The violence of Roman persecution
had been directed against the Christian Church; while the Jews,
gradually recovering their strength, and obtaining permission to
settle and make proselytes to their creed, at length became an
influential political body in the neighbourhood of their ancient home,
especially in the Syrian provinces which were at that time the chief
residence of the court. Severus (A.D. 194) is
said to have been the first to extend to them the imperial favour,
though he afterwards withdrew it. Heliogabalus, and Alexander, natives
of Syria, gave them new privileges; and the latter went so far as to
place the image of Abraham in his private chapel, among the objects of
his ordinary worship. Philip the Arabian continued towards them a
countenance, which was converted into an open patronage in the reign
of Zenobia. During the Decian persecution, they had been sufficiently
secure at Carthage, to venture to take part in the popular ridicule
which the Christians excited; and they are even said to have
stimulated Valerian to his cruelties towards the Church \footnote{Basnage, Hist. des Juifs, vi. 12.
Tillemont, Hist. des Emper. iii. iv.}.

But this direct hostility was not the only, nor the most formidable
means of harassing their religious enemies, which their improving
fortunes opened upon them. With their advancement in wealth and
importance, their national character displayed itself under a new
exterior. The moroseness for which they were previously notorious, in
great measure disappears with their dislodgment from the soil of their
ancestors; and on their re-appearance as settlers in a strange land,
those festive, self-indulgent habits, which, in earlier times, had but
drawn on them the animadversion of their Prophets, became their
distinguishing mark in the eyes of external observers \footnote{Vide Gibbon, Hist. ch. xvi. note 6.
Chrysost. in Judæos, i. p. 386-388, \&c.}. Manifesting a rancorous malevolence towards the zealous
champions of the Church, they courted the Christian populace by arts
adapted to captivate and corrupt the unstable and worldly-minded.
Their pretensions to magical power gained them credit with the
superstitious, to whom they sold amulets for the cure of diseases;
their noisy spectacles attracted the curiosity of the idle, who
weakened their faith, while they disgraced their profession, by
attending the worship of the Synagogue. Accordingly there was formed
around the Church a mixed multitude, who, without relinquishing their
dependence on Christianity for the next world, sought in Judaism the
promise of temporal blessings, and a more accommodating rule of life
than the gospel revealed. Chrysostom found this evil so urgent at
Antioch in his day, as to interrupt his course of homilies on the
heresy of the Anomœans, in order to direct his preaching against the
seductions to which his hearers were then exposed, by the return of
the Jewish festivals \footnote{Chrysost. in Judæos, i. p. 389, \&c.
[Jerome speaks of a law of Valens:—``ne quis vitulorum carnibus
vesceretur, utilitati agriculturæ providens, et pessimam \emph{judaizantis}
\emph{vulgi} emendans consuetudinem.'' Adv. Jovinian. ii. 7.]}.
In another part of the empire, the Council of Illiberis found it
necessary to forbid a superstitious custom, which had been introduced
among the country people, of having recourse to the Jews for a
blessing on their fields. Afterwards, Constantine made a law against
the inter-marriage of Jews and Christians; and Constantius confiscated
the goods of Christians who lapsed to Judaism \footnote{Bingham, Antiq. xvi. 6. Basnage, Hist. des
Juifs, vi. 14.}. These successive enactments may be taken as evidence of the
view entertained by the Church of her own danger, from the artifices
of the Jews. Lastly, the attempt to rebuild the temple in Julian's
reign, was but the renewal of a project on their part, which
Constantine had already frustrated, for reinstating their religion in
its ancient ritual and country \footnote{Chrysost, in Judæos, iii. p. 435. [Vide
Chrysost. in Matth. Hom. 43, where he says that in Julian's time,
``they ranged themselves with the heathen and courted their
party.'' He proceeds to say that ``in all their other evil
works they surpass their predecessors, in sorceries, magic arts,
impurities.'' Oxford Transl.]}.

Such was the position of the Jews towards the primitive Church,
especially in the patriarchate of Antioch; which, I have said, was
their principal place of settlement, and was at one time under the
civil government of a Judaizing princess, the most illustrious
personage of her times, who possessed influence enough over the
Christian body to seduce the Metropolitan himself from the orthodox
faith.

3.

But the evidence of the existence of Judaism, as a system, in the
portion of Christendom in question, is contained in a
circumstance which deserves our particular attention; the adoption, in
those parts, of the quarto-deciman rule of observing Easter, when it
was on the point of being discontinued in the Churches of Proconsular
Asia, where it had first prevailed.

It is well known that at the close of the second century, a
controversy arose between Victor, Bishop of Rome, and Polycrates,
Bishop of Ephesus, concerning the proper time for celebrating the
Easter feast, or rather for terminating the ante-paschal fast. At that
time the whole of Christendom, with the exception of Proconsular Asia
(a district of about two hundred miles by fifty), and its immediate
neighbourhood \footnote{Euseb. Hist. v. 23-25, and Vales. ad loc.},
continued the fast on to the Sunday after the Jewish Passover, which
they kept as Easter Day as we do now, in order that the weekly and
yearly commemorations of the Resurrection might coincide. But the
Christians of the Proconsulate, guided by Jewish custom, ended the
fast on the very day of the paschal sacrifice, without regarding the
actual place held in the week by the feast, which immediately
followed; and were accordingly called Quarto-decimans \footnote{Exod. xii. 6. Vide Tillemont, Mem. vol.
iii. p. 629, \&c.}. Victor felt the inconvenience of this want of uniformity in
the celebration of the chief Christian festival; and was urgent, even
far beyond the bounds of charity, and the rights of his see, in his
endeavour to obtain the compliance of the Asiatics. Polycrates, who
was primate of the Quarto-deciman Churches, defended their peculiar
custom by a statement which is plain and unexceptionable. They had
received their rule, he said, from St. John and St. Philip the
Apostles, Polycarp of Smyrna, Melito of Sardis, and others; and deemed
it incumbent on them to transmit as they had received. There was
nothing Judaistic in this conduct; for, though the Apostles intended
the Jewish discipline to cease with those converts who were born under
it, yet it was by no means clear, that its calendar came under the
proscription of its rites. On the other hand, it was natural that the
Asian Churches should be affectionately attached to a custom which
their first founders, and they inspired teachers, had sanctioned.

But the case was very different, when Churches, which had for
centuries observed the Gentile rule, adopted a custom which at the
time had only existence among the Jews. The Quarto-decimans of the
Proconsulate had come to an end by A.D. 276;
and, up to that date, the Antiochene provinces kept their Easter feast
in conformity with the Catholic usage \footnote{Tillemont, Mem. vol. iii. p. 48, who
conjectures that Anatolius of Laodicea was the author of the change.
But changes require predisposing causes.}; yet, at the time of the Nicene Council (fifty years
afterwards), we find the Antiochenes the especial and solitary
champions of the Jewish rule \footnote{Athan. ad Afros, § 2.}. We can scarcely doubt that they adopted it in imitation of
the Jews who were settled among them, who are known to have influenced
them, and who about that very date, be it observed, had a patroness in
Zenobia, and, what was stranger, had almost a convert in the person of
the Christian Primate. There is evidence, moreover, of the actual
growth of the custom in the Patriarchate at the end of the third
century; which well agrees with the hypothesis of its being an
innovation, and not founded on ancient usage. And again (as was
natural, supposing the change to begin at Antioch), at the date of the
Nicene Council, it was established only in the Syrian Churches, and
was but making its way with incomplete success in the extremities of
the Patriarchate. In Mesopotamia, Audius began his schism with the
characteristic of the Quarto-deciman rule, just at the date of the
Council \footnote{Epiph. Hær. lxx. § 1.}; and about the
same time, Cilicia was contested between the two parties, as I gather
from the conflicting statements of Constantine and Athanasius, that it
did, and that it did not, conform to the Gentile custom \footnote{Athan. ad Afros, supra. Socr. Hist. i. 9,
where, by the bye, the Proconsulate is spoken of as conforming to the
general usage; so as clearly to distinguish between the two Quarto-deciman
schools.}. By the same time, the controversy had reached Egypt also.
Epiphanius refers to a celebrated contest, now totally unknown,
between one Crescentius and Alexander, the first defender of the
Catholic faith against Arianism \footnote{Epiph. ibid. § 9.}.

It is true that there was a third Quarto-deciman school, lying
geographically between the Proconsulate and Antioch, which at first
sight might seem to have been the medium by which the Jewish custom
was conveyed on from the former to the latter; but there is no
evidence of its existence till the end of the fourth century. In order
to complete my account of the Quarto-decimans, and show more fully
their relation to the Judaizers, I will here make mention of it;
though, in doing so, I must somewhat digress from the main subject
under consideration.

The portion of Asia Minor, lying between the Proconsulate and the
river Halys, may be regarded, in the Ante-Nicene times, as one
country, comprising the provinces of Phrygia, Galatia, Cappadocia, and
Paphlagonia, afterwards included within the Exarchate of Cæsarea; and
was then marked by a religious character of a peculiar cast. Socrates,
speaking of this district, informs us, that its inhabitants were
distinguished above other nations by a strictness and seriousness of
manners, having neither the ferocity of the Scythians and Thracians,
nor the frivolity and sensuality of the Orientals \footnote{Socrat. Hist. iv. 28, cf, Epiph. Hær.
xlviii. 14 [and xlvii. 1].}. The excellent qualities, however, implied in this
description, were tarnished by the love of singularity, the spirit of
insubordination and separatism, and the gloomy spiritual pride which
their history evidences. St. Paul's Epistle furnishes us with the
first specimen of this unchristian temper, as evinced in the conduct
of the Galatians \footnote{[Jerome calls the Galatians ``ad
intelligentiam tardiores, vecordes,'' and speaks of their ``stoliditas
barbara,'' in Galat. lib. ii. præf.]}, who,
dissatisfied with the exact evangelical doctrine, aspired to some
higher and more availing system than the Apostle preached to them.
What the Galatians were in the first century, Montanus and Novatian
became in the second and third; both authors of a harsh and arrogant
discipline, both natives of the country in question \footnote{Vales, ad loc. Socr. [Philostorg. viii.
15.]}, and both meeting with special success in that country,
although the schism of the latter was organized at Rome, of which
Church he was a presbyter. It was, moreover, the peculiarity, more or
less, of both Montanists and Novatians in those parts, to differ from
the general Church as to the time of observing Easter \footnote{Socrat. Hist. v. 22. Sozom. Hist. vii. 18.}; whereas, neither in Africa nor in Rome did the two sects
dissent from the received rule \footnote{Tertull. de Jejun. 14. Vales, ad Sozom.
vii. 18. Socrat. Hist. v. 21.}. What was the principle or origin of this irregularity, does
not clearly appear; unless we may consider as characteristic, what
seems to be the fact, that when their neighbours of the Proconsulate
were Quarto-decimans, they (in the words of Socrates) ``shrank
from feasting on the Jewish festival \footnote{Valesius ad loc. applies this differently.},'' and after the others had conformed to the Gentile rule,
they, on the contrary, openly judaized \footnote{Socrat. Hist. v. 21.}. This change in their practice, which took place at the end of
the fourth century, was mainly effected by a Jew, of the name of
Sabbatius, who becoming a convert to Christianity, rose to the
episcopate in the Novatian Church. Sozomen, in giving an account of
the transaction, observes that it was a national custom with the
Galatians and Phrygians to judaize in their observance of Easter.
Coupling this remark with Eusebius's mention of Churches in the
neighbourhood of the Proconsulate, as included among the Quarto-decimans
whom Victor condemned \footnote{Eusb. Hist. ut supra.},
we may suspect that the perverse spirit which St. Paul reproves in his
Epistle, and which we have been tracing in its Montanistic and
Novatian varieties, still lurked in those parts in its original
judaizing form, till after a course of years it was accidentally
brought out by circumstances upon the public scene of ecclesiastical
history. If further evidence of the connexion of the Quarto-deciman
usage with Judaism be required, I may refer to Constantine's
Nicene Edict, which forbids it, among other reasons, on the ground of
its being Jewish \footnote{Theod. Hist. i. 10.}.

4.

The evidence, which has been adduced for the existence of Judaism
in the Church of Antioch, is not without its bearing upon the history
of the rise of Arianism. I will not say that the Arian doctrine is the
direct result of a judaizing practice; but it deserves consideration
whether a tendency to derogate from the honour due to Christ, was not
created by an observance of the Jewish rites, and much more, by that
carnal, self-indulgent religion, which seems at that time to have
prevailed in the rejected nation. When the spirit and morals of a
people are materially debased, varieties of doctrinal error spring up,
as if self-sown, and are rapidly propagated. While Judaism inculcated
a superstitious, or even idolatrous dependence on the mere casualties
of daily life, and gave license to the grosser tastes of human nature,
it necessarily indisposed the mind for the severe and unexciting
mysteries, the large indefinite promises, and the remote sanctions, of
the Catholic faith; which fell as cold and uninviting on the depraved
imagination, as the doctrines of the Divine Unity and of implicit
trust in the unseen God, on the minds of the early Israelites. Those
who were not constrained by the message of mercy, had time attentively
to consider the intellectual difficulties which were the medium of its
communication, and heard but ``a hard saying'' in what was
sent from heaven as ``tidings of great joy.'' ``The
mind,'' says Hooker, ``feeling present joy, is always
marvellously unwilling to admit any other cogitation, and in that
case, casteth off those disputes whereunto the intellectual part at
other times easily draweth ... The people that are said in the sixth
of John to have gone after our Lord to Capernaum ... leaving Him on
the one side of the sea of Tiberias, and finding Him again as soon as
they themselves by ship were arrived on the contrary side ... as they
wondered, so they asked also, 'Rabbi, when camest Thou hither?' The
Disciples, when Christ appeared to them in a far more strange and
miraculous manner, moved no question, but rejoiced greatly in what
they saw ... The one, because they enjoyed not, disputed; the other
disputed not, because they enjoyed \footnote{Eccles. Pol. v. 67.}.''

It is also a question, whether the mere performance of the rites of
the Law, of which Christ came as anti-type and repealer, has not a
tendency to withdraw the mind from the contemplation of the more
glorious and real images of the Gospel; so that the Christians of
Antioch would diminish their reverence towards the true Saviour of
man, in proportion as they trusted to the media of worship provided
for a time by the Mosaic ritual. It is this consideration which
accounts for the energy with which the great Apostle combats the
adoption of the Jewish ordinances by the Christians of Galatia, and
which might seem excessive, till vindicated by events subsequent to
his own day \footnote{[Eusebius says, that St. Paul detected
humanitarianism in the Galatian Judaism. Contr. Marcell. i. 1, p. 7.]}. In the
Epistle addressed to them, the Judaizers are described as men
labouring under an irrational fascination, fallen from grace, and
self-excluded from the Christian privileges \footnote{Socrat. Hist. v. 22.}; when in appearance they were but using, what on the one hand
might be called mere external forms, and on the other, had actually
been delivered to the Jews on Divine authority. Some light is thrown
upon the subject by the Epistle to the Hebrews, in which it is implied
throughout, that the Jewish rites, after their Antitype was come, did
but conceal from the eye of faith His divinity, sovereignty, and
all-sufficiency. If we turn to the history of the Church, we seem to
see the evils in actual existence, which the Apostle anticipated in
prophecy; that is, we see, that in the obsolete furniture of the
Jewish ceremonial, there was in fact retained the pestilence of Jewish
unbelief, tending (whether directly or not, at least eventually) to
introduce fundamental error respecting the Person of Christ.

Before the end of the first century, this result is disclosed in
the system of the Cerinthians and the Ebionites. These sects, though
more or less infected with Gnosticism, were of Jewish origin, and
observed the Mosaic Law; and whatever might be the minute
peculiarities of their doctrinal views, they also agreed in
entertaining Jewish rather than Gnostic conceptions of the Person of
Christ \footnote{Burton, Bamp. Lect., notes 74, 82.}. Ebion,
especially, is characterised by his Humanitarian creed; while on the
other hand, his Judaism was so notorious, that Tertullian does not
scruple to describe him as virtually the object of the Apostle's
censure in his Epistle to the Galatians \footnote{Tertull. de Præscript. Hæret. c. 33, p.
243.}.

The Nazarenes are next to be noticed;—not for the influence they
exercised on the belief of Christians, but as evidencing, with the
sects just mentioned, the latent connection between a judaizing
discipline and heresy in doctrine. Who they were, and what their
tenets, has been a subject of much controversy. It is sufficient for
my purpose—and so far is undoubted—that they were at the same time
``zealous of the Law'' and unsound in their theology \footnote{Burton, Bampt. Lect., note 84.}; and this without being related to the Gnostic families: a
circumstance which establishes them as a more cogent evidence of the
real connexion of ritual with doctrinal Judaism than is furnished by
the mixed theologies of Ebion and Cerinthus \footnote{For the curious in ecclesiastical
antiquity, Mosheim has elicited the following account of their name
and sect (Mosheim de Reb. Christ. ante Constant. Sæcul. ii. § 38,
39). The title of Nazarene he considers to have originally belonged to
the body of Jewish converts, taken by them with a reference to Matt.
ii. 23, while the Gentiles at Antioch assumed the Greek appellation of
Christians. As the Mosaic ordinances gradually fell into disuse among
the former, in process of time it became the peculiar designation of
the Church of Jerusalem; and that Church in turn throwing off its
Jewish exterior in the reign of Hadrian, on being unfairly subjected
to the disabilities then laid upon the rebel nation, it finally
settled upon the scanty remnant, who considered their ancient
ceremonial to be an essential part of their present profession. These
Judaizers, from an over-attachment to the forms, proceeded in course
of time, to imbibe the spirit of the degenerate system; and ended in
doctrinal views not far short of modern Socinianism.}. It is worth observing that their declension from orthodoxy
appears to have been gradual; Epiphanius is the first writer who
includes them by name in the number of heretical sects \footnote{Burton, Bampt. Lect., note 84. Considering
the Judaism of the Quarto-decimans after Victor's age, is it
impossible that he may have suspected that the old leaven was
infecting the Churches of Asia? This will explain and partly excuse
his earnestness in the controversy with them. It must be recollected
that he witnessed, in his own branch of the Church, the rise of the
first simply humanitarian school which Christianity had seen, that of
Theodotus, Artemas, \&c. (Euseb. Hist. v. 28), the latter of whom
is charged by Alexander with reviving the heresy of the judaizing
Ebion (Theod. Hist., i. 4); [while at the same time at Rome Blastus
was introducing the Quarto-deciman rule]. Again, Theodotus, Montanus,
and Praxeas, whose respective heresies he was engaged in combating,
all belonged to the neighbourhood of the Proconsulate, where there
seems to have been a school, from which Praxeas derived his heresy (Theod.
Hær. iii. 3); while Montanism, as its after history shows, contained
in it the seeds, both of the Quarto-deciman and Sabellian errors (Tillemont,
Mem. vol. ii. p. 199, 205. Athan. in Arian. ii. 43). It may be added
that the younger Theodotus is suspected of Montanism (Tillemont. Mem.
vol. iii. p. 277).}.

5.

Such are the instances of the connexion between Judaism and
theological error, previously to the age of Paulus, who still more
strikingly exemplifies it. First, we are in possession of his
doctrinal opinions, which are grossly humanitarian; next we find, that
in early times they were acknowledged to be of Jewish origin; further,
that his ceremonial Judaism also was so notorious that one author even
affirms that he observed the rite of circumcision \footnote{Philastr. Hær. § 64. [Epiphanius denies
that the Paulianists circumcised. Hær. lxv. 2. It is remarkable that
the Arian Whiston looked favourably on the rite. Biograph. Brit. p.
4213.]}: and lastly, just after his day we discover the rise of a
Jewish usage, the Quarto-deciman, in the provinces of Christendom,
immediately subjected to his influence.

It may be added that this view of the bearing of Judaism upon the
sceptical school afterwards called Arian is countenanced by frequent
passages in the writings of the contemporary Fathers, on which no
stress, perhaps, could fairly be laid, were not their meaning
interpreted by the above historical facts \footnote{Athan. de Decret. 2. 27; Sentent. Dionys.
3, 4; ad Episc. Æg. 13; de fug, 2; in Arian. iii. 27, and \emph{passim}.
Chrysost. Hom. in Anomœos and in Judæos. Theod. Hist. i. 4. Epiphan.
Hær. lxix. 79.}. Moreover, in the popular risings which took place in Antioch
and Alexandria in favour of Arianism, the Jews sided with the
heretical party \footnote{Basnage, Hist. des Juifs. vi. 41.};
evincing thereby, not indeed any definite interest in the subject of
dispute, but a sort of spontaneous feeling, that the side of heresy
was their natural position; and further, that its spirit, and the
character which it created, were congenial to their own. Or, again, if
we consider the subject from a different point of view, and omitting
dates and schools, take a general survey of Christendom during the
first centuries, we shall find it divided into the same two parties,
both on the Arian and the Quarto-deciman questions; Rome and
Alexandria with their dependencies being the champions of the Catholic
tradition in either controversy, and Palestine, Syria, and Asia Minor,
being the strongholds of the opposition. And these are the two
questions which occasioned the deliberations of the Nicene Fathers.

However, it is of far less consequence, as it is less certain,
whether Arianism be of Jewish origin, than whether it arose at
Antioch: which is the point principally insisted on in the foregoing
pages. For in proportion as it is traced to Antioch, so is the charge
of originating it removed from the great Alexandrian School, upon
which various enemies of our Apostolical Church have been eager to
fasten it. In corroboration of what has been said above on this
subject, I here add the words of Alexander, in his letter to the
Church of Constantinople, at the beginning of the controversy;
which are of themselves decisive in evidence of the part, which
Antioch had, in giving rise to the detestable blasphemy which he was
combating.

``Ye are not ignorant,'' he writes to the
Constantinopolitan Church concerning Arianism, ``that this
rebellious doctrine belongs to Ebion and Artemas, and is in imitation
of Paulus of Samosata, Bishop of Antioch, who was excommunicated by
the sentence of the Bishops assembled in Council from all quarters.
Paulus was succeeded by Lucian, who remained in separation for many
years during the time of three bishops ... Our present heretics have
drunk up the dregs of the impiety of these men, and are their secret
offspring; Arius and Achillas, and their party of evil-doers, incited
as they are to greater excesses by three Syrian prelates, who agree
with them … Accordingly, they have been expelled from the Church, as
enemies of the pious Catholic teaching; according to St. Paul's
sentence, 'If any man preach any other Gospel unto you than that ye
have received, let him be anathema.''' \footnote{Theod. Hist. i. 4. [Simeon, Bishop of
Beth-Arsam, in Persia, A.D. 510-525, traces the
genealogy of Paulianism and Nestorianism from Judaism thus:—Caiaphas
to Simon Magus; Simon to Ebion; Ebion to Artemon; Artemon to Paul of
Somosata; Paul to Diodorus; Diodorus to Theodore; Theodore to
Nestorius. Asseman. Bibl. Orient. t. i. p. 347.]}

\xsection{Section 2. The Schools of the Sophists}

A\textsc{s} Antioch was the birthplace, so were the Schools of the
Sophists the place of education of the heretical spirit which we are
considering. In this section, I propose to show its disputatious
character, and to refer it to these Schools as the source of it.

The vigour of the first movement of the heresy, and the rapid
extension of the controversy which it introduced, are among the more
remarkable circumstances connected with its history. In the course of
six years it called for the interposition of a General Council; though
of three hundred and eighteen bishops there assembled, only
twenty-two, on the largest calculation and, as it really appears, only
thirteen, were after all found to be its supporters. Though thus
condemned by the whole Christian world, in a few years it broke out
again; secured the patronage of the imperial court, which had recently
been converted to the Christian faith; made its way into the highest
dignities of the Church; presided at her Councils, and tyrannized over
the majority of her members who were orthodox believers.

Now, doubtless, one chief cause of these successes is found in the
circumstance, that Lucian's pupils were brought together from so
many different places, and were promoted to posts of influence in so
many parts of the Church. Thus Eusebius, Maris, and Theognis, were
bishops of the principal sees of Bithynia; Menophantes was exarch of
Ephesus; and Eudoxius was one of the Bishops of Comagene. Other causes
will hereafter appear in the secular history of the day; but here I am
to speak of their talent for disputation, to which after all they were
principally indebted for their success.

It is obvious, that in every contest, the assailant, as such, has
the advantage of the party assailed; and that, not merely from the
recommendation which novelty gives to his cause in the eyes of
bystanders, but also from the greater facility in the nature of
things, of finding, than of solving objections, whatever be the
question in dispute. Accordingly, the skill of a disputant mainly
consists in securing an offensive position, fastening on the weaker
points of his adversary's case, and then not relaxing his hold till
the latter sinks under his impetuosity, without having the opportunity
to display the strength of his own cause, and to bring it to bear upon
his opponent; or, to make use of a familiar illustration, in causing a
sudden run upon his resources, which the circumstances of time and
place do not allow him to meet. This was the artifice to which
Arianism owed its first successes \footnote{[\emph{anap\={e}d\={o}si gar h\={o}s luss\={e}t\={e}res
kunes eis echthr\={o}n amunan}]. Epiph. Hær. lxix. 15. Vide the
whole passage.}. It owed them to the
circumstance of its being (in its original form) a sceptical rather
than a dogmatic teaching; to its proposing to inquire into and
reform the received creed, rather than to hazard one of its own. The
heresies which preceded it, originating in less subtle and dexterous
talent, took up a false position, professed a theory, and sunk under
the obligations which it involved. The monstrous dogmas of the various
Gnostic sects pass away from the scene of history as fast as they
enter it. Sabellianism, which succeeded, also ventured on a creed; and
vacillating between a similar wildness of doctrine, and a less
imposing ambiguity, soon vanished in its turn \footnote{Vide § 5, infra. [Gregory Naz. speaks of a
[\emph{gal\={e}n\={e}}] after these heresies, and before
Arianism. Orat. xxv. 8.]}. But the
Antiochene School, as represented by Paulus of Samosata and Arius,
took the ground of an assailant, attacked the Catholic doctrine, and
drew the attention of men to its difficulties, without attempting to
furnish a theory of less perplexity or clearer evidence.

The arguments of Paulus (which it is not to our purpose here to
detail) seem fairly to have overpowered the first of the Councils
summoned against him (A.D. 264), which dissolved
without coming to a decision \footnote{Euseb. Hist. vii. 28. Cave, Hist. Literar. vol. 1. p. 158.}. A second, and (according to
some writers) a third, were successfully convoked, when at length his
subtleties were exposed and condemned; not, however, by the reasonings
of the Fathers of the Council themselves, but by the instrumentality
of one Malchion, a presbyter of Antioch, who, having been by
profession a Sophist, encountered his adversary with his own arms \footnote{[[\emph{sphodra katapolemountai hoi polemioi,
hotan tois aut\={o}n hoplois chr\={o}metha kat' aut\={o}n}]. Socr. iii. 16.]}. Even in yielding, the arts of Paulus secured from his
judges an ill-advised concession, the abandonment of the celebrated
word \emph{homoüsion} (\emph{consubstantial}), afterwards adopted as
the test at Nicæa; which the orthodox had employed in the
controversy, and to which Paulus objected as open to a
misinterpretation \footnote{Bull. Defens. Fid. Nic. ii. i. § 9-14.}. Arius followed in the track thus marked
out by his predecessor. Turbulent by character, he is known in history
as an offender against ecclesiastical order, before his agitation
assumed the shape which has made his name familiar to posterity \footnote{Epiph. Hær. lxix. 3.}. When he betook himself to the doctrinal controversy, he chose for
the first open avowal of his heterodoxy the opportunity of an attack
upon his diocesan, who was discoursing on the mystery of the Trinity
to the clergy of Alexandria. Socrates, who is far from being a
partisan of the Catholics, informs us that Arius being well skilled in
dialectics sharply replied to the bishop, accused him of Sabellianism,
and went on to \emph{argue} that ``\emph{if} the Father begat the
Son, certain conclusions would follow,'' and so proceeded. His
heresy, thus founded in a syllogism, spread itself by instruments of a
kindred character. First, we read of the excitement which his
reasonings produced in Egypt and Lybia; then of his letters addressed
to Eusebius and to Alexander, which display a like pugnacious and
almost satirical spirit; and then of his verses composed for the use
of the populace in ridicule of the orthodox doctrine \footnote{Socr. i. 5, 6. Theod. Hist. i. 5. Epiphan. Hær. lxix. 7, 8.
Philostorg. ii. 2. Athan. de Decret. 16.}. But
afterwards, when the heresy was arraigned before the Nicene Council, and placed on the defensive, and later still, when its
successes reduced it to the necessity of occupying the chairs of
theology, it suffered the fate of the other dogmatic heresies before
it; split, in spite of court favour, into at least four different
creeds, in less than twenty years \footnote{Petav. Dogm. Theol. t. ii. i. 9 and 10.}; and at length gave way to
the despised but indestructible truth which it had for a time
obscured.

Arianism had in fact a close connexion with the existing
Aristotelic school. This might have been conjectured, even had there
been no proof of the fact, adapted as that philosopher's logical
system confessedly is to baffle an adversary, or at most to detect
error, rather than to establish truth \footnote{``Omnem vim venenorum suorum in dialectica disputatione
constituunt, quæ philosophorum sententia definitur non adstruendi vim
habere, sed studium destruendi. Sed non in dialectica complacuit Deo
salvum facere populum suum.'' Ambros. de Fide, i. 5. [§ 42.]}. But we have actually
reason, in the circumstances of its history, for considering it as the
off-shoot of those schools of inquiry and debate which acknowledged
Aristotle as their principal authority, and were conducted by teachers
who went by the name of Sophists. It was in these schools that the
leaders of the heretical body were educated for the part assigned them
in the troubles of the Church. The oratory of Paulus of Samosata is
characterized by the distinguishing traits of the scholastic eloquence
in the descriptive letter of the Council which condemned him; in
which, moreover, he is stigmatized by the most disgraceful title to
which a Sophist was exposed by the degraded exercise of his
profession \footnote{[\emph{sophist\={e}s kai go\={e}s}], a juggler. Vide
Cressol. Theatr. Rhetor. i. 13. iii. 17.}. The skill of Arius in the art of disputation is
well known. Asterius was a Sophist by profession. Aetius came from the
School of an Aristotelian of Alexandria. Eunomius, his pupil, who
re-constructed the Arian doctrine on its original basis, at the end of
the reign of Constantius, is represented by Ruffinus as
``pre-eminent in dialectic power.'' \footnote{Petav. Theol. prolegom. iii. 3. Baltus, Defense des Peres, ii.
19. Brucker, vol. iii. p. 288. Cave, Hist. Literar. vol. 1.} At a later
period still, the like disputatious spirit and spurious originality
are indirectly ascribed to the heterodox school, in the advice of
Sisinnius to Nectarius of Constantinople, when the Emperor Theodosius
required the latter to renew the controversy with a view to its final
settlement \footnote{Bull, Defens. Fid. Nic. Epilog.}. Well versed in theological learning, and aware
that adroitness in debate was the very life and weapon of heresy,
Sisinnius proposed to the Patriarch, to drop the use of dialectics,
and merely challenge his opponents to utter a general anathema against
all such Ante-Nicene Fathers as had taught what they themselves now
denounced as false doctrine. On the experiment being tried, the
heretics would neither consent to be tried by the opinions of the
ancients, nor yet dared condemn those whom ``all the people
counted as prophets.'' ``Upon this,'' say the historians
who record the story, ``the Emperor perceived that they rested
their cause on their dialectic skill, and not on the testimony of the
early Church.'' \footnote{Socr. Hist. v. 10. Soz. Hist. vii. 12.}

Abundant evidence, were more required, could be added to the
above, in proof of the connexion of the Arians with the schools of
heathen disputation. The two Gregories, Basil, Ambrose, and Cyril,
protest with one voice against the dialectics of their opponents; and
the sum of their declarations is briefly expressed by a writer of the
fourth century, who calls Aristotle the Bishop of the Arians \footnote{Petav. Dogm. Theol. supra. Brucker, vol. iii. pp. 324, 352,
353. Epiph. Hær. lxix. 69. [Vigil. Thaps. contr. Eutych. i. 2.]}.

2.

And while the science of argumentation provided the means, their
practice of disputing for the sake of exercise or amusement supplied
the temptation, of assailing received opinions. This practice, which
had long prevailed in the Schools, was early introduced into the
Eastern Church \footnote{The art was called [\emph{eristik\={e}}] and the actual
discussion, [\emph{gymnasia}]. Cressol. Theatr. Rhet. ii. 3. [Vide
also Athan. Tr. p. 44, e. Also a remarkable instance in Ernesti from
Origen, ap Lumper, t. 10, p. 148. Contrasted with [\emph{gymnastikoi
logoi}] were [\emph{ag\={o}nistikoi}], \emph{in earnest},
according to Sextus Empiricus, vide Hypot. i. 33, p. 57, with
Fabricius's note.]}. It was there employed as a means of
preparing the Christian teacher for the controversy with unbelievers.
The discussion sometimes proceeded in the form of a lecture delivered
by the master of the school to his pupils; sometimes in that of an
inquiry, to be submitted to the criticism of his hearers; sometimes by
way of dialogue, in which opposite sides were taken for argument-sake.
In some cases, it was taken down in notes by the bystanders, at the
time; in others committed to writing by the parties engaged in it \footnote{Dodw. Diss. in Iren. v. 14. Socr. Hist. i. 5.}. Necessary as these exercises would be for the purpose
designed, yet they were obviously open to abuse, though moderated by
ever so orthodox and strictly scriptural a rule, in an age when no
sufficient ecclesiastical symbol existed, as a guide to the memory and
judgment of the eager disputant. It is evident, too, how difficult it
would be to secure opinions or arguments from publicity, which were
but hazarded in the confidence of Christian friendship, and which,
when viewed apart from the circumstances of the case, lent a seemingly
deliberate sanction to heterodox novelties. Athanasius implies \footnote{Athan. de Decret. 25 and 27. [He says the same of Marcellus in
his defence, Apol. contr. Ar. 47.]}, that in the theological works of Origen and Theognostus, while
the orthodox faith was explicitly maintained, nevertheless heretical
tenets were discussed, and in their place more or less defended, by
way of exercise in argument. The countenance thus accidentally given
to the cause of error is evidenced in his eagerness to give the
explanation. But far greater was the evil, when men destitute of
religious seriousness and earnestness engaged in the like theological
discussions, not with any definite ecclesiastical object, but as a
mere trial of skill, or as a literary recreation; regardless of the
mischief thus done to the simplicity of Christian morals, and the evil
encouragement given to fallacious reasonings and sceptical views. The
error of the ancient Sophists had consisted in their indulging without
restraint or discrimination in the discussion of practical topics,
whether religious or political, instead of selecting such as might
exercise, without demoralizing, their minds. The rhetoricians of
Christian times introduced the same error into their treatment of
the highest and most sacred subjects of theology. We are told, that
Julian commenced his opposition to the true faith by defending the
heathen side of religious questions, in disputing with his brother
Gallus \footnote{Greg. Nazianz. Orat. iii. 27, 31. [iv. 30.]}; and probably he would not have been able himself to
assign the point of time at which he ceased merely to take a part, and
became earnest in his unbelief. But it is unnecessary to have recourse
to particular instances, in order to prove the consequences of a
practice so evidently destructive of a reverential and sober spirit.

Moreover, in these theological discussions, the disputants were in
danger of being misled by the unsoundness of the positions which they
assumed, as elementary truths or axioms in the argument. As logic and
rhetoric made them expert in proof and refutation, so there was much
in other sciences, which formed a liberal education, in geometry and
arithmetic, to confine the mind to the contemplation of material
objects, as if these could supply suitable tests and standards for
examining those of a moral and spiritual nature; whereas there are
truths foreign to the province of the most exercised intellect, some
of them the peculiar discoveries of the improved moral sense (or what
Scripture terms ``\emph{the Spirit}''), and others still less
on a level with our reason, and received on the sole authority of
Revelation. Then, however, as now, the minds of speculative men were
impatient of ignorance, and loth to confess that the laws of truth and
falsehood, which their experience of this world furnished, could not
at once be applied to measure and determine the facts of another.
Accordingly, nothing was left for those who would not believe the
incomprehensibility of the Divine Essence, but to conceive of it by
the analogy of sense; and using the figurative terms of theology in
their literal meaning as if landmarks in their inquiries, to suppose
that then, and then only, they steered in a safe course, when they
avoided every contradiction of a mathematical and material nature.
Hence, canons grounded on physics were made the basis of discussions
about possibilities and impossibilities in a spiritual substance, as
confidently and as fallaciously, as those which in modern times have
been derived from the same false analogies against the existence of
moral self-action or free-will. Thus the argument by which Paulus of
Samosata baffled the Antiochene Council, was drawn from a sophistical
use of the very word \emph{substance}, which the orthodox had employed in
expressing the scriptural notion of the unity subsisting between the
Father and the Son \footnote{Bull, Defens. F. N. ii. 1. § 10.}. Such too was the mode of reasoning
adopted at Rome by the Artemas or Artemon, already mentioned, and his
followers, at the end of the second century. A contemporary writer,
after saying that they supported their ``God-denying
apostasy'' by syllogistic forms of argument, proceeds,
``Abandoning the inspired writings, they devote themselves to
geometry, as becomes those who are of the earth, and speak of the
earth, and are ignorant of Him who is from above. Euclid's treatises,
for instance, are zealously studied by some of them; Aristotle and
Theophrastus are objects of their admiration; while Galen may be said even to be adored by others. It is needless to declare that such
perverters of the sciences of unbelievers to the purposes of their own
heresy, such diluters of the simple Scripture faith with heathen
subtleties, have no claim whatever to be called believers.'' \footnote{Euseb. Hist. v. 28.} And such is Epiphanius's description of the Anomœans, the genuine
offspring of the original Arian stock. ``Aiming,'' he says,
``to exhibit the Divine Nature by means of Aristotelic syllogisms
and geometrical data, they are thence led on to declare that Christ
cannot be derived from God.'' \footnote{Epiph. Hær. p. 809.}

3.

Lastly, the absence of an adequate symbol of doctrine increased the
evils thus existing, by affording an excuse and sometimes a reason for
investigations, the necessity of which had not yet been superseded by
the authority of an ecclesiastical decision. The traditionary system,
received from the first age of the Church, had been as yet but
partially set forth in authoritative forms; and by the time of the
Nicene Council, the voices of the Apostles were but faintly heard
throughout Christendom, and might be plausibly disregarded by those
who were unwilling to hear. Even at the beginning of the third
century, the disciples of Artemas boldly pronounced their heresy to be
apostolical, and maintained that all the bishops of Rome had held it
till Victor inclusive \footnote{Euseb. ibid.}, whose episcopate was but a few years
before their own time. The progress of unbelief naturally led them on
to disparage, rather than to appeal to their predecessors; and to
trust their cause to their own ingenuity, instead of defending an
inconvenient fiction concerning the opinions of a former age. It ended
in teaching them to regard the ecclesiastical authorities of former
times as on a level with the uneducated and unenlightened of their own
days. Paulus did not scruple to express contempt for the received
expositors of Scripture at Antioch; and it is one of the first
accusations brought by Alexander against Arius and his party, that
``they put themselves above the ancients, and the teachers of our
youth, and the prelates of the day; considering themselves alone to be
wise, and to have discovered truths, which had never been revealed to
man before them.'' \footnote{Theod. Hist. i. 4. [``Solæ in contemptu sunt divinæ
literæ, quæ nec suam scholam nec magistros habeant, et de quibus
peritissimè disputare se credat, qui nunquam didicit.'' Facund.
p. 581. ed. Sirm.; vide also, p. 565.]}

On the other hand, while the line of tradition, drawn out as it was
to the distance of two centuries from the Apostles, had at length
become of too frail a texture, to resist the touch of subtle and
ill-directed reason, the Church was naturally unwilling to have
recourse to the novel, though necessary measure, of imposing an
authoritative creed upon those whom it invested with the office of
teaching. If I avow my belief, that freedom from symbols and articles
is abstractedly the highest state of Christian communion, and the
peculiar privilege of the primitive Church \footnote{[``Non eguistis literâ, qui spiritu abundabatis, etc. Ubi
sensus conscientiæ periclitatur, illic litera postulator.'' Hilar.
de Syn. 63. Vide the Benedictine note.]}, it is not from
any tenderness towards that proud impatience of control in which many
exult, as in a virtue: but first, because technicality and formalism
are, in their degree, inevitable results of public confessions of
faith; and next, because when confessions do not exist, the mysteries
of divine truth, instead of being exposed to the gaze of the profane
and uninstructed, are kept hidden in the bosom of the Church, far more
faithfully than is otherwise possible; and reserved by a private
teaching, through the channel of her ministers, as rewards in due
measure and season, for those who are prepared to profit by them; for
those, that is, who are diligently passing through the successive
stages of faith and obedience. And thus, while the Church is not
committed to declarations, which, most true as they are, still are
daily wrested by infidels to their ruin; on the other hand, much of
that mischievous fanaticism is avoided, which at present abounds from
the vanity of men, who think that they can explain the sublime
doctrines and exuberant promises of the Gospel, before they have yet
learned to know themselves and to discern the holiness of God, under
the preparatory discipline of the Law and of Natural Religion.
Influenced, as we may suppose, by these various considerations, from
reverence for the free spirit of Christian faith, and still more for
the sacred truths which are the objects of it, and again from
tenderness both for the heathen and the neophyte, who were unequal to
the reception of the strong meat of the full Gospel, the rulers of the
Church were dilatory in applying a remedy, which nevertheless the
circumstances of the times imperatively required. They were loth to
confess, that the Church had grown too old to enjoy the free,
unsuspicious teaching with which her childhood was blest; and that her
disciples must, for the future, calculate and reason before they spoke
and acted. So much was this the case, that in the Council of
Antioch (as has been said), on the objection of Paulus, they actually
withdrew a test which was eventually adopted by the more experienced
Fathers at Nicæa; and which, if then sanctioned, might, as far as the
Church was concerned, have extinguished the heretical spirit in the
very place of its birth.—Meanwhile, the adoption of Christianity, as
the religion of the empire, augmented the evil consequences of this
omission, excommunication becoming more difficult, while entrance into
the Church was less restricted than before.

\xsection{Section 3. The Church of Alexandria}

A\textsc{s} the Church of Antioch was exposed to the influence of
Judaism, so was the Alexandrian Church characterized in primitive
times by its attachment to that comprehensive philosophy, which was
reduced to system about the beginning of the third century, and then
went by the name of the New Platonic, or Eclectic. A supposed
resemblance between the Arian and the Eclectic doctrine concerning the
Holy Trinity, has led to a common notion that the Alexandrian Fathers
were the medium by which a philosophical error was introduced into the
Church; and this hypothetical cause of a disputable resemblance has
been apparently evidenced by the solitary fact, which cannot be
denied, that Arius himself was a presbyter of Alexandria. We have
already seen, however, that Arius was educated at Antioch; and we
shall see hereafter that, so far from being favourably heard at
Alexandria, he was, on the first promulgation of his heresy, expelled
the Church in that city, and obliged to seek refuge among his
Collucianists of Syria. And it is manifestly the opinion of Athanasius,
that he was but the pupil or the tool of deeper men \footnote{Athan. de Decr. Nic. 8, 20; ad Monach. 66; de Synod. 22.}, probably of Eusebius of Nicomedia, who in no sense
belongs to Alexandria. But various motives have led theological
writers to implicate this celebrated Church in the charge of heresy.
Infidels have felt a satisfaction, and heretics have had an interest,
in representing that the most learned Christian community did not
submit implicitly to the theology taught in Scripture and by the
Church; a conclusion, which, even if substantiated, would little
disturb the enlightened defender of Christianity, who may safely admit
that learning, though a powerful instrument of the truth in right
hands, is no unerring guide into it. The Romanists \footnote{[As to the charges made against Petavius,
vide Bull, Defens. N. F. proœm.; Budd Isagog. p. 580; Bayle, Dict. (Petau.);
Brucker, Phil. t. iii. p. 345.]}, on the other hand, have thought by the same line of policy to
exalt the Apostolical purity of their own Church, by the contrast of
unfaithfulness in its early rival; and (what is of greater importance)
to insinuate both the necessity of an infallible authority, by
exaggerating the errors and contrarieties of the Ante-Nicene Fathers,
and the fact of its existence, by throwing us, for exactness of
doctrinal statement, upon the decisions of the subsequent Councils. In
the following pages, I hope to clear the illustrious Church in
question of the grave imputation thus directed against her from
opposite quarters: the imputation of considering the Son of God by
nature inferior to the Father, that is, of platonizing or arianizing.
But I have no need to profess myself her disciple, though, as regards
the doctrine in debate, I might well do so; and, instead of setting
about any formal defence, I will merely place before the reader the
general principles of her teaching, and leave it to him to apply
them, as far as he judges they will go, in explanation of the
language, which has been the ground of the suspicions against her.

1.

St. Mark, the founder of the Alexandrian Church, may be numbered
among the personal friends and associates of that Apostle, who held it
to be his especial office to convert the heathen; an office, which was
impressed upon the community formed by the Evangelist, with a strength
and permanence unknown in the other primitive Churches. The
Alexandrian may peculiarly be called the Missionary and Polemical
Church of Antiquity. Situated in the centre of the accessible world,
and on the extremity of Christendom, in a city which was at once the
chief mart of commerce, and a celebrated seat of both Jewish and Greek
philosophy, it was supplied in especial abundance, both with materials
and instruments prompting to the exercise of Christian zeal. Its
catechetical school, founded (it is said) by the Evangelist himself,
was a pattern to other Churches in its diligent and systematic
preparation of candidates for baptism; while other institutions were
added of a controversial character, for the purpose of carefully
examining into the doctrines revealed in Scripture, and of cultivating
the habit of argument and disputation \footnote{Cave. Hist. Literar. vol. i. p. 80.}. While the internal affairs of the community were administered
by its bishops, on these academical bodies, as subsidiary to the
divinely-sanctioned system, devolved the defence and propagation of
the faith, under the presidency of laymen or inferior
ecclesiastics. Athenagoras, the first recorded master of the
catechetical school, is known by his defence of the Christians, still
extant, addressed to the Emperor Marcus. Pantænus, who succeeded him,
was sent by Demetrius, at that time bishop, as missionary to the
Indians or Arabians. Origen, who was soon after appointed catechist at
the early age of eighteen, had already given the earnest of his future
celebrity, by his persuasive disputations with the unbelievers of
Alexandria. Afterwards he appeared in the character of a Christian
apologist before an Arabian prince, and Mammæa, the mother of
Alexander Severus, and addressed letters on the subject of religion to
the Emperor Philip and his wife Severa; and he was known far and wide
in his day, for his indefatigable zeal and ready services in the
confutation of heretics, for his various controversial and critical
writings, and for the number and dignity of his converts \footnote{Philipp. Sidet. fragm. apud Dodw. in Iren.
Huet. Origen.}.

Proselytism, then, in all its branches, the apologetic, the
polemical, and the didactic, being the peculiar function of the
Alexandrian Church, it is manifest that the writings of its
theologians would partake largely of an exoteric character. I mean,
that such men would write, not with the openness of Christian
familiarity, but with the tenderness or the reserve with which we are
accustomed to address those who do not sympathize with us, or whom we
fear to mislead or to prejudice against the truth, by precipitate
disclosures of its details. The example of the inspired writer of the
Epistle to the Hebrews was their authority for making a broad
distinction between the doctrines suitable to the state of the
weak and ignorant, and those which are the peculiar property of a
baptized and regenerate Christian. The Apostle in that Epistle, when
speaking of the most sacred Christian verities, as hidden under the
allegories of the Old Testament, seems suddenly to check himself, from
the apprehension that he was divulging mysteries beyond the
understanding of his brethren; who, instead of being masters in
Scripture doctrine, were not yet versed even in its elements, needed
the nourishment of children rather than of grown men, nay, perchance,
having quenched the illumination of baptism, had forfeited the
capacity of comprehending even the first elements of the truth. In the
same place he enumerates these elements, or foundation of Christian
teaching \footnote{Hebr. v. 11; vi. 6. [\emph{ta stoicheia t\={e}s
arch\={e}s t\={o}n logi\={o}n tou theou, ho t\={e}s arch\={e}s
tou Christou logos}].}, in contrast
with the esoteric doctrines which the ``long-exercised habit of
moral discernment'' can alone appropriate and enjoy, as follows;—repentance,
faith in God, the doctrinal meaning of the right of baptism,
confirmation as the channel of miraculous gifts, the future
resurrection, and the final separation of good and bad. His first
Epistle to the Corinthians contains the same distinction between the
carnal or imperfect and the established Christian, which is laid down
in that addressed to the Hebrews. While maintaining that in
Christianity is contained a largeness of wisdom, or (to use human
language) a profound philosophy, fulfilling those vague conceptions of
greatness, which had led the aspiring intellect of the heathen sages
to shadow forth their unreal systems, he at the same time insists upon the impossibility of man's arriving at this hidden treasure all
at once, and warns his brethren, instead of attempting to cross by a
short path from the false to the true knowledge, to humble themselves
to the low and narrow portal of the heavenly temple, and to become
fools, that they might at length be really wise. As before, he speaks
of the difference of doctrine suited respectively to neophytes and
confirmed Christians, under the analogy of the difference of food
proper for the old and young; a difference which lies, not in the
arbitrary will of the dispenser, but in the necessity of the case, the
more sublime truths of Revelation affording no nourishment to the
souls of the unbelieving or unstable.

Accordingly, in the system of the early catechetical schools, the \emph{perfect},
or men in Christ, were such as had deliberately taken upon them the
profession of believers; had made the vows, and received the grace of
baptism; and were admitted to all the privileges and the revelations
of which the Church had been constituted the dispenser. But before
reception into this full discipleship, a previous season of
preparation, from two to three years, was enjoined, in order to try
their obedience, and instruct them in the principles of revealed
truth. During this introductory discipline, they were called \emph{Catechumens},
and the teaching itself \emph{Catechetical}, from the careful and
systematic examination by which their grounding in the faith was
effected. The matter of the instruction thus communicated to them,
varied with the time of their discipleship, advancing from the most
simple principle of Natural Religion to the peculiar doctrines of the
Gospel, from moral truths to the Christian mysteries. On their first
admission they were denominated \emph{hearers}, from the leave
granted them to attend the reading of the Scriptures and sermons in
the Church. Afterwards, being allowed to stay during the prayers, and
receiving the imposition of hands as the sign of their progress in
spiritual knowledge, they were called \emph{worshippers}. Lastly, some
short time before their baptism, they were taught the Lord's Prayer
(the peculiar privilege of the regenerate), were entrusted with the
knowledge of the Creed, and, as destined for incorporation into the
body of believers, received the titles of \emph{competent} or \emph{elect}
\footnote{[\emph{teleioi; akro\={o}menoi}], or
audientes; [\emph{gonuklinontes}], or [\emph{euchomenoi}]; competentes,
electi, or [\emph{ph\={o}tizomenoi}]. Bingham, Antiq. book x.
Suicer. Thes. in verb, [\emph{kat\={e}che}\emph{o}].}. Even to the last, they
were granted nothing beyond a formal and general account of the
articles of the Christian faith; the exact and fully developed
doctrines of the Trinity and the Incarnation, and still more, the
doctrine of the Atonement, as once made upon the cross, and
commemorated and appropriated in the Eucharist, being the exclusive
possession of the serious and practised Christian. On the other hand,
the chief subjects of catechisings, as we learn from Cyril \footnote{Bingham, ibid.}, were the doctrines of repentance and pardon, of the necessity
of good works, of the nature and use of baptism, and the immortality
of the soul;—as the Apostle had determined them.

The exoteric teaching, thus observed in the Catechetical Schools,
was still more appropriate, when the Christian teacher addressed
himself, not to the instruction of willing hearers, but to controversy
or public preaching. At the present day, there are very many sincere
Christians, who consider that the evangelical doctrines are the
appointed instruments of conversion, and, as such, exclusively
attended with the Divine blessing. In proof of this position, with an
inconsistency remarkable in those who profess a jealous adherence to
the inspired text, and are not slow to accuse others of ignorance of
its contents, they appeal, not to Scripture, but to the stirring
effects of this (so-called) Gospel preaching, and to the inefficiency,
on the other hand, of mere exhortations respecting the benevolence and
mercy of God, the necessity of repentance, the rights of conscience,
and the obligation of obedience. But it is scarcely the attribute of a
generous faith, to be anxiously inquiring into the consequences of
this or that system, with a view to decide its admissibility, instead
of turning at once to the revealed word, and inquiring into the rule
there exhibited to us. God can defend and vindicate His own command,
whatever it turn out to be; weak though it seem to our vain wisdom,
and unworthy of the Giver; and that His course in this instance is
really that which the hasty religionist condemns as if the theory of
unenlightened formalists, is evident to careful students of Scripture,
and is confirmed by the practice of the Primitive Church.

As to Scripture, I shall but observe, in addition to the remarks
already made on the passages in the Epistles to the Corinthians and
Hebrews, that no one sanction can be adduced thence, whether of
precept or of example, in behalf of the practice of stimulating the
affections, such as gratitude or remorse, by means of the doctrine of
the Atonement, in order to the conversion of the hearers;—that, on
the contrary, it is its uniform method to connect the Gospel with
Natural Religion, and to mark out obedience to the moral law as
the ordinary means of attaining to a Christian faith, the higher
evangelical truths, as well as the Eucharist, which is the visible
emblem of them, being received as the reward and confirmation of
habitual piety;—that, in the preaching of the Apostles and
Evangelists in the Book of Acts, the sacred mysteries are revealed to
individuals in proportion to their actual religious proficiency; that
the first principles of righteousness, temperance, and judgment to
come, are urged upon Felix; while the elders of Ephesus are reminded
of the divinity and vicarious sacrifice of Christ, and the presence
and power of the Holy Spirit in the Church;—lastly, that among those
converts, who were made the chief instruments of the first propagation
of the Gospel, or who are honoured with especial favour in Scripture,
none are found who had not been faithful to the light already given
them, and were not distinguished, previously to their conversion, by a
strictly conscientious deportment. Such are the divine notices given
to those who desire an apostolical rule for dispensing the word of
life; and as such, the ancient Fathers received them. They received
them as the fulfilment of our Lord's command, not to give that which
is holy to dogs, nor to cast pearls before swine; a text cited by
Clement and Tertullian \footnote{Ceillier, Apol. des Pères, ch. ii. Bingh.
Antiq. x. 5.},
among others, in justification of their cautious distribution of
sacred truth. They also considered this caution as the result of the
most truly charitable consideration for those whom they addressed, who
were likely to be perplexed, not converted, by the sudden exhibition
of the whole evangelical scheme. This is the doctrine of Theodoret,
Chrysostom, and others, in their comments upon the passage in the
Epistle to the Hebrews \footnote{Suicer. Thes. in verb. [\emph{stoicheion}].}.
``Should a catechumen ask thee what the teachers have determined,
(says Cyril of Jerusalem) tell nothing to one who is without. For we
impart to thee a secret and a promise of the world to come. Keep safe
the secret for Him who gives the reward. Listen not to one who asks,
'What harm is there in my knowing also?' Even the sick ask for wine,
which, unseasonably given, brings on delirium; and so there come two
ills, the death of the patient and the disrepute of the
physician.'' In another place he says, ``All may hear the
Gospel, but the glory of the Gospel is set apart for the true
disciples of Christ. To all who could hear, the Lord spake, but in
parables; to His disciples He privately explained them. What is the
blaze of Divine glory to the enlightened, is the blinding of
unbelievers. These are the secrets which the Church unfolds to him who
passes on from the catechumens, and not to the heathen. For we do not
unfold to a heathen the truths concerning Father, Son, and Holy
Spirit; nay, not even in the case of catechumens, do we clearly
explain the mysteries, but we frequently say many things indirectly,
so that believers who have been taught may understand, and the others
may not be injured.'' \footnote{Cyril. Hieros. ed. Milles, præf. § 7
catech. vi. 16.}

The work of St. Clement, of Alexandria, called Stromateis, or
Tapestry-work, from the variety of its contents, well illustrates the
Primitive Church's method of instruction, as far as regards the
educated portion of the community. It had the distinct object of
interesting and conciliating the learned heathen who perused it;
but it also exemplifies the peculiar caution then adopted by
Christians in teaching the truth,—their desire to rouse the moral
powers to internal voluntary action, and their dread of loading or
formalizing the mind. In the opening of his work, Clement speaks of
his miscellaneous discussions as mingling truth with philosophy;
``or rather,'' he continues, ``involving and concealing
it, as the shell hides the edible fruit of the nut.'' In another
place he compares them, not to a fancy garden, but to some
thickly-wooded mountain, where vegetation of every sort, growing
promiscuously, by its very abundance conceals from the plunderer the
fruit trees, which are intended for the rightful owner. ``We must
hide,'' he says, ``that wisdom, spoken in mystery, which the
Son of God has taught us. Thus the Prophet Esaias has his tongue
cleansed with fire, that he may be able to declare the vision; and our
ears must be sanctified as well as our tongues, if we aim at being
recipients of the truth. This was a hindrance to my writing; and still
I have anxiety, since Scripture says, 'Cast not your pearls before
swine;' for those pure and bright truths, which are so marvellous and
full of God to goodly natures, do but provoke laughter, when spoken in
the hearing of the many.'' \footnote{Strom. i. 1. 12; v. 3; vi. 1; vii. 18.} The Fathers considered that they had the pattern as well as
the recommendation of this method of teaching in Scripture itself \footnote{``Bonæ sunt in Scripturis sacris
mysteriorum profunditates, quæ ob hoc teguntur, ne vilescant; ob hoc
quæruntur, ut exerceant; ob hoc autem aperiuntur, ut pascant.''
August. in Petav. præf. in Trin. i. 5.}.

2.

This self-restraint and abstinence, practised at least partially, by the Primitive Church in the publication of the most
sacred doctrines of our religion, is termed, in theological language,
the \emph{Disciplina Arcani}; concerning which a few remarks may here
be added, not so much in recommendation of it (which is beside my
purpose), as to prevent misconception of its principle and limits.

Now, first, it may be asked, How was any secrecy practicable,
seeing that the Scriptures were open to every one who chose to consult
them? It may startle those who are but acquainted with the popular
writings of this day, yet, I believe, the most accurate consideration
of the subject will lead us to acquiesce in the statement, as a
general truth, that the doctrines in question have never been learned
merely from Scripture. Surely the Sacred Volume was never intended,
and is not adapted, to teach us our creed; however certain it is that
we can prove our creed from it, when it has once been taught us \footnote{Vide Dr. Hawkins's original and most
conclusive work on Unauthoritative Tradition, which contains in it the
key to a number of difficulties which are apt to perplex the
theological student.}, and in spite of individual producible exceptions to the
general rule. From the very first, that rule has been, as a matter of
fact, that the Church should teach the truth, and then should appeal
to Scripture in vindication of its own teaching. And from the first,
it has been the error of heretics to neglect the information thus
provided for them, and to attempt of themselves a work to which they
are unequal, the eliciting a systematic doctrine from the scattered
notices of the truth which Scripture contains. Such men act, in the
solemn concerns of religion, the part of the self-sufficient natural
philosopher, who should obstinately reject Newton's theory of
gravitation, and endeavour, with talents inadequate to the task, to
strike out some theory of motion by himself. The insufficiency of the
mere private study of Holy Scripture for arriving at the exact and
entire truth which Scripture really contains, is shown by the fact,
that creeds and teachers have ever been divinely provided, and by the
discordance of opinions which exists wherever those aids are thrown
aside; as it is also shown by the very structure of the Bible itself.
And if this be so, it follows that, while inquirers and neophytes in
the first centuries lawfully used the inspired writings for the
purposes of morals and for instruction in the rudiments of the faith,
they still might need the teaching of the Church as a key to the
collection of passages which related to the mysteries of the Gospel,
passages which are obscure from the necessity of combining and
receiving them all.

A more plausible objection to the existence of this rule of secrecy
in the Early Church arises from the circumstance, that the Christian
Apologists openly mention to the whole world the sacred tenets which
have been above represented as the peculiar possession of the
confirmed believer. But it must be observed, that the writers of these
were frequently laymen, and so did not commit the Church as a body,
nor even in its separate authorities, to formal statement or to
theological discussion. The great duty of the Christian teacher was to
unfold the sacred truths in due order, and not prematurely to insist
on the difficulties, or to apply the promises of the Gospel; and if
others erred in this respect, still it remained a duty to him. And further, these disclosures are not so conclusive as they seem to
be at first sight; the approximations of philosophy, and the
corruptions of heresy, being so considerable, as to create a confusion
concerning the precise character of the ecclesiastical doctrine.
Besides, in matter of fact, some of the early apologists themselves,
as Tatian, were tainted with heretical opinions.

But in truth, it is not the actual practice of the Primitive
Church, which I am concerned with, so much as its principle. Men often
break through the rules, which they set themselves for the conduct of
life, with or without good reason. If it was the professed principle
of the early teachers, to speak exoterically to those who were without
the Church, instances of a contrary practice but prove their
inconsistency; whereas the fact of the existence of the principle
answers the purpose which is the ultimate aim of these remarks, viz.
it accounts for those instances in the teaching of the Alexandrians,
whether many or few, and whether extant or not in writing, in which
they were silent as regards the mysterious doctrines of Christianity.
Indeed it is evident, that anyhow the \emph{Disciplina Arcani} could
not be observed for any long time in the Church. Apostates would
reveal its doctrines, even if these escaped in no other way. Perhaps
it was almost abandoned, as far as men of letters were concerned,
after the date of Ammonius; indeed there are various reasons for
limiting its strict enforcement to the end of the second century. And
it is plain, that during the time when the sacred doctrines were
passing into the stock of public knowledge, Christian
controversialists would be in a difficulty how to conduct themselves,
what to deny, explain or complete, in the popular notions of
their creed; and they would consequently be betrayed into
inconsistencies of statement, and vary in their method of disputing.

The \emph{Disciplina Arcani} being supposed, with these
limitations, to have had a real existence, I observe further, in
explanation of its principle, that the elementary information given to
the heathen or catechumen was in no sense undone by the subsequent
secret teaching, which was in fact but the filling up of a bare but
correct outline. The contrary theory was maintained by the Manichees,
who represented the initiatory discipline as founded on a fiction or
hypothesis, which was to be forgotten by the learner as he made
progress in the real doctrine of the gospel \footnote{August. in Advers, Leg. et Proph. lib. ii.}; somewhat after the manner of a school in the present day,
which supposes conversion to be effected by an exhibition of free
promises and threats, and an appeal to our moral capabilities, which
after conversion are discovered to have no foundation in fact. But
``Far be it from so great an Apostle,'' says Augustine,
speaking of St. Paul, ``a vessel elect of God, an organ of the
Holy Ghost, to be one man when he preached, another when he wrote, one
man in private, another in public. He was made all to all men, not by
the craft of a deceiver, but from the affection of a sympathizer,
succouring the diverse diseases of souls with the diverse emotions of
compassion; to the little ones dispensing the lesser doctrines, not
false ones, but the higher mysteries to the perfect, all of them,
however, true, harmonious, and divine.'' \footnote{Mosheim, de Caus. Supp. Libror. § 17. 1
do not find it in this exact form in Augustine's treatise; vide in
Advers. Leg. et Proph. lib. ii. 4, 6, \&c.}

Next, the truths reserved for the baptized Christian were not put
forward as the arbitrary determinations of individuals, as the word of
man, but rather as an apostolical legacy, preserved and dispensed by
the Church. Thus Irenæus when engaged in refuting the heretics of his
age, who appealed from the text of Scripture to a sense independent of
it, as the test between truth and falsehood in its contents, says,
``We know the doctrine of our salvation through none but those who
have transmitted to us the gospel, first proclaiming it, then (by God's
will) delivering it to us in the Scriptures, as a basis and pillar of
our faith. Nor dare we affirm that their announcements were made
previously to their attaining perfect knowledge, as some presume to
say, boasting that they set right the Apostles.'' \footnote{Iren. iii. 1. Vide also Tertull. de
Præscr. Hæret. 22.} He then proceeds to speak of the clearness and cogency of the
traditions preserved in the Church, as containing that true wisdom of
the perfect, of which St. Paul speaks, and to which the Gnostics
pretended. And, indeed, without formal proofs of the existence and the
authority in primitive times of an Apostolical Tradition, it is plain
that there must have been such tradition, granting that the Apostles
conversed, and their friends had memories, like other men. It is quite
inconceivable that they should not have been led to arrange the series
of revealed doctrines more systematically than they record them in
Scripture, as soon as their converts became exposed to the attacks and
misrepresentations of heretics; unless they were forbidden so to do, a
supposition which cannot be maintained. Their statements thus
occasioned would be preserved, as a matter of course; together with those other secret but less important truths, to which St. Paul
seems to allude, and which the early writers more or less acknowledge,
whether concerning the types of the Jewish Church, or the prospective
fortunes of the Christian \footnote{Mosheim, de Reb. ante Const. sæc. ii. §
34.}.
And such recollections of apostolical teaching would evidently be
binding on the faith of those who were instructed in them; unless it
can be supposed, that, though coming from inspired teachers, they were
not of divine origin.

However, it must not be supposed, that this appeal to Tradition in
the slightest degree disparages the sovereign authority and
sufficiency of Holy Scripture, as a record of the truth. In the
passage from Irenæus above cited, Apostolical Tradition is brought
forward, not to supersede Scripture, but in conjunction with
Scripture, to refute the self-authorized, arbitrary doctrines of the
heretics. We must cautiously distinguish, with that Father, between a
tradition supplanting or perverting the inspired records, and a
corroborating, illustrating, and altogether subordinate tradition. It
is of the latter that he speaks, classing the traditionary and the
written doctrine together, as substantially one and the same, and as
each equally opposed to the profane inventions of Valentinus and
Marcion.

Lastly, the secret tradition soon ceased to exist even in theory.
It was authoritatively divulged, and perpetuated in the form of
symbols according as the successive innovations of heretics called for
its publication. In the creeds of the early Councils, it may be
considered as having come to light, and so ended; so that whatever has
not been thus authenticated, whether it was prophetical
information, or comment on the past dispensations \footnote{2 Thes. ii. 5, 15. Heb. v. 11.}, is, from the circumstances of the case, lost to the Church.
What, however, was then (by God's good providence) seasonably
preserved, is in some sense of apostolical authority still; and at
least serves the chief office of the early traditions, viz. that of
interpreting and harmonizing the statements of Scripture.

3.

In the passages lately quoted from Clement and Cyril, mention was
made by those writers of a mode of speaking, which was intelligible to
the well-instructed, but conveyed no definite meaning to ordinary
hearers. This was the Allegorical Method; which well deserves our
attention before we leave the subject of the \emph{Disciplina Arcani},
as being one chief means by which it was observed. The word \emph{allegorizing}
must here be understood in a wide signification; as including in its
meaning, not only the representation of truths, under a foreign,
though analogous exterior, after the manner of our Lord's parables,
but the practice of generalizing facts into principles, of adumbrating
greater truths under the image of lesser, of implying the consequences
or the basis of doctrines in their correlatives, and altogether those
instances of thinking, reasoning, and teaching, which depend upon the
use of propositions which are abstruse, and of connexions which are
obscure, and which, in the case of uninspired authors, we consider
profound, or poetical, or enthusiastic, or illogical, according to our
opinion of those by whom they are exhibited.

This method of writing was the national peculiarity of that
literature in which the Alexandrian Church was educated. The
hieroglyphics of the ancient Egyptians mark the antiquity of a
practice, which, in a later age, being enriched and diversified by the
genius of their Greek conquerors, was applied as a key both to
mythological legends, and to the sacred truths of Scripture. The
Stoics were the first to avail themselves of an expedient which
smoothed the deformities of the Pagan creed. The Jews, and then the
Christians, of Alexandria, employed it in the interpretation of the
inspired writings. Those writings themselves have certainly an
allegorical structure, and seem to countenance and invite an
allegorical interpretation; and in consequence, they have been
referred by some critics to one and the same heathen origin, as if
Moses first, and then St. Paul, borrowed their symbolical system
respectively from the Egyptian and the Alexandrian philosophy.

But it is more natural to consider that the Divine Wisdom used on
the sublimest of all subjects, media, which we spontaneously select
for the expression of solemn thought and elevated emotion; and had no
especial regard to the practice in any particular country, which
afforded but one instance of the operation of a general principle of
our nature. When the mind is occupied by some vast and awful subject
of contemplation, it is prompted to give utterance to its feelings in
a figurative style; for ordinary words will not convey the admiration,
nor literal words the reverence which possesses it; and when, dazzled
at length with the great sight, it turns away for relief, it still
catches in every new object which it encounters, glimpses of its
former vision, and colours its whole range of thought with this one
abiding association. If, however, others have preceded it in the
privilege of such contemplations, a well-disciplined piety will lead
it to adopt the images which they have invented, both from affection
for what is familiar to it, and from a fear of using unsanctioned
language on a sacred subject. Such are the feelings under which a
deeply impressed mind addresses itself to the task of disclosing even
its human thoughts; and this account of it, if we may dare to
conjecture, in its measure applies to the case of a mind under the
immediate influence of inspiration. Certainly, the matter of
Revelation suggests some such hypothetical explanation of the
structure of the books which are its vehicle; in which the
divinely-instructed imagination of the writers is ever glancing to and
fro, connecting past things with future, illuminating God's lower
providences and man's humblest services by allusions to the relations
of the evangelical covenant, and then in turn suddenly leaving the
latter to dwell upon those past dealings of God with man, which must
not be forgotten merely because they have been excelled. No prophet
ends his subject: his brethren after him renew, enlarge, transfigure,
or reconstruct it; so that the Bible, though various in its parts,
forms a whole, grounded on a few distinct doctrinal principles
discernible throughout it; and is in consequence intelligible indeed
in its general drift, but obscure in its text; and even tempts the
student, if I may so speak, to a lax and disrespectful interpretation
of it. History is made the external garb of prophecy, and persons and
facts become the figures of heavenly things. I need only refer, by way
of instance, to the delineation of Abraham as the type of the
accepted worshipper of God; to the history of the brazen serpent; to
the prophetical bearing of the ``call of Israel out of
Egypt;'' to the personification of the Church in the Apostolic
Epistles as the reflected image of Christ; and, further, to the
mystical import, interpreted by our Lord Himself, of the title of God
as the God of the Patriarchs. Above all other subjects, it need
scarcely be said, the likeness of the promised Mediator is conspicuous
throughout the sacred volume as in a picture: moving along the line of
the history, in one or other of His destined offices, the dispenser of
blessings in Joseph, the inspired interpreter of truth in Moses, the
conqueror in Joshua, the active preacher in Samuel, the suffering
combatant in David, and in Solomon the triumphant and glorious king.

Moreover, Scripture assigns the same uses to this allegorical
style, which were contemplated by the Fathers when they made it
subservient to the \emph{Disciplina Arcani}; viz. those of trying the
earnestness and patience of inquirers, discriminating between the
proud and the humble, and conveying instruction to believers, and that
in the most permanently impressive manner, without the world's sharing
in the knowledge. Our Lord's remarks on the design of his own
parables, is a sufficient evidence of this intention.

Thus there seemed every encouragement, from the structure of
Scripture, from the apparent causes which led to that structure, and
from the purposes to which it was actually applied by its Divine
Author, to induce the Alexandrians to consider its text as primarily
and directly the instrument of an allegorical teaching. And since
it sanctions the principle of allegorizing by its own example, they
would not consider themselves confined within the limits of the very
instances which it supplies, because of the evident spiritual drift of
various passages which, nevertheless, it does not interpret
spiritually; thus to the narrative contained in the twenty-second
chapter of Genesis, few people will deny an evangelical import, though
the New Testament itself nowhere assigns it. Yet, on the other hand,
granting that a certain liberty of interpretation, beyond the
precedent, but according to the spirit of Scripture, be allowable in
the Christian teacher, still few people will deny, that some rule is
necessary as a safeguard against its abuse, in order to secure the
sacred text from being explained away by the heretic, and misquoted
and perverted by weak or fanatical minds. Such a safeguard we shall
find in bearing cautiously in mind this consideration: viz. that (as a
general rule), every passage of Scripture has some one definite and
sufficient sense, which was prominently before the mind of the writer,
or in the intention of the Blessed Spirit, and to which all other
ideas, though they might arise, or be implied, still were subordinate.
It is this true meaning of the text, which it is the business of the
expositor to unfold. This it is, which every diligent student will
think it a great gain to discover; and, though he will not shut his
eyes to the indirect and instructive applications of which the text is
capable, he never will so reason as to forget that there is one sense
peculiarly its own. Sometimes it is easily ascertained, sometimes it
can be scarcely conjectured; sometimes it is contained in the literal
sense of the words employed, as in the historical parts;
sometimes it is the allegorical, as in our Lord's parables; or
sometimes the secondary sense may be more important in after ages than
the original, as in the instance of the Jewish ritual; still in all
cases (to speak generally) there is but one main primary sense,
whether literal or figurative; a regard for which must ever keep us
sober and reverent in the employment of those allegorisms, which,
nevertheless, our Christian liberty does not altogether forbid.

The protest of Scripture against all careless expositions of its
meaning, is strikingly implied in the extreme reserve and caution,
with which it unfolds its own typical signification; for instance, in
the Mosaic ritual no hint was given of its undoubted prophetical
character, lest an excuse should be furnished to the Israelitish
worshipper for undervaluing its actual commands. So, again, the
secondary and distinct meaning of prophecy, is commonly hidden from
view by the veil of the literal text, lest its immediate scope should
be overlooked; when that is once fulfilled, the recesses of the sacred
language seem to open, and give up the further truths deposited in
them. Our Lord, probably, in the prophecy recorded in the Gospels, was
not careful (if I may so express myself) that His disciples should
distinguish between His final and immediate coming; thinking it a less
error that they should consider the last day approaching, than that
they should forget their own duties in the contemplation of the future
fortunes of the Church. Nay, even types fulfilled, if they be
historical, seem sometimes purposely to be left without the sanction
of an interpretation, lest we should neglect the instruction still
conveyed by a literal narrative. This accounts for the silence
observed concerning the evangelical import, to which I have already
referred, of the sacrifice of Isaac, which contains a definite and
permanent moral lesson, as a matter of fact, however clear may be its
further meaning as emblematical of our Lord's sufferings on the cross.
In corroboration of this remark, let it be observed, that there seems
to have been in the Church a traditionary explanation of these
historical types, derived from the Apostles, but kept among the secret
doctrines, as being dangerous to the majority of hearers \footnote{Vide Mosheim, de Reb. Ant. Const. sæc.
ii. § 34. Rosenmuller, Hist. Interpr. iii. 2. § 1.}; and certainly St. Paul, in the Epistle to the Hebrews,
affords us an instance of such a tradition, both as existing and as
secret (even though it be shown to be of Jewish origin), when, first
checking himself and questioning his brethren's faith, he
communicates, not without hesitation, the evangelical scope of the
account of Melchisedec, as introduced into the book of Genesis.

As to the Christian writers of Alexandria, if they erred in their
use of the Allegory, their error did not lie in the mere adoption of
an instrument which Philo or the Egyptian hierophants had employed
(though this is sometimes made the ground of objection), for Scripture
itself had taken it out of the hands of such authorities. Nor did
their error lie in the mere circumstance of their allegorizing
Scripture, where Scripture gave no direct countenance; as if we might
not interpret the sacred word for ourselves, as we interpret the
events of life, by the principles which itself supplies. But they
erred, whenever and as far as they carried their favourite rule of
exposition beyond the spirit of the canon above laid down, so as
to obscure the primary meaning of Scripture, and to weaken the force
of historical facts and doctrinal declarations; and much more, if at
any time they degraded the inspired text to the office of conveying
the thoughts of uninspired teachers on subjects not sacred.

And, as it is impossible to draw a precise line between the use and
abuse of allegorizing, so it is impossible also to ascertain the exact
degree of blame incurred by individual teachers who familiarly indulge
in it. They may be faulty as commentators, yet instructive as
devotional writers; and their liberty in interpretation is to be
regulated by the state of mind in which they address themselves to the
work, and by their proficiency in the knowledge and practice of
Christian duty. So far as men use the language of the Bible (as is
often done in poems and works of fiction) as the mere instrument of a
cultivated fancy, to make their style attractive or impressive, so
far, it is needless to say, they are guilty of a great irreverence
towards its Divine Author. On the other hand, it is surely no
extravagance to assert that there are minds so gifted and disciplined
as to approach the position occupied by the inspired writers, and
therefore able to apply their words with a fitness, and entitled to do
so with a freedom, which is unintelligible to the dull or heartless
criticism of inferior understandings. So far then as the Alexandrian
Fathers partook of such a singular gift of grace (and Origen surely
bears on him the tokens of some exalted moral dignity), not incited by
a capricious and presumptuous imagination, but burning with that
vigorous faith, which, seeing God in all things, does and suffers
all for His sake, and, while filled with the contemplation of His
supreme glory, still discharges each command in the exactness of its
real meaning, in the same degree they stand not merely excused, but
are placed immeasurably above the multitude of those who find it so
easy to censure them.—And so much on the Allegory, as the means of
observing the \emph{Disciplina Arcani}.

4.

The same method of interpretation was used for another purpose,
which is more open to censure. When Christian controversialists were
urged by objections to various passages in the history of the Old
Testament, as derogatory to the Divine Perfections or to the Jewish
saints, they had recourse to an allegorical explanation by way of
answer. Thus Origen spiritualizes the account of Abraham's denying his
wife, the polygamy of the Patriarchs, and Noah's intoxication \footnote{Heut. Origen. p. 171, Rosenmuller supra.
[On this subject, vide a striking passage in Facundus, Def. Tr. Cap.
xii. 1, pp. 568-9.]}. It is impossible to defend such a mode of interpretation,
which seems to imply a want of faith in those who had recourse to it.
Doubtless this earnestness to exculpate the saints of the elder
covenant is partly to be attributed to a noble jealousy for the honour
of God, and a reverence for the memory of those who, on the whole,
rise in their moral attainments far above their fellows, and well
deserve the confidence in their virtue which the Alexandrians
manifest. Yet God has given us rules of right and wrong, which we must
not be afraid to apply in estimating the conduct of even the best of
mere men; though errors are thereby detected, the scandal of
which we ourselves have to bear in our own day. So far must be granted
in fairness; but some have gone on to censure the principle itself
which this procedure involved: viz. that of representing religion, for
the purpose of conciliating the heathen, in the form most attractive
to their prejudices: and, as it was generally received in the
Primitive Church, and the considerations which it involves are not
without their bearings upon the doctrinal question in which we shall
be presently engaged, I will devote some space here to the examination
of it.

The mode of arguing and teaching in question which is called \emph{economical}
\footnote{[\emph{kat' oikonomian}].} by the ancients, can
scarcely be disconnected from the \emph{Disciplina Arcani}, as will
appear by some of the instances which follow, though it is convenient
to consider it by itself. If it is necessary to contrast the two with
each other, the one may be considered as withholding the truth, and
the other as setting it out to advantage. The Economy is certainly
sanctioned by St. Paul in his own conduct. To the Jews he became as a
Jew, and as without the Law to the heathen \footnote{[On the economies of St. Peter and St.
Paul, vide Lardner's Heathen Test. ch. xxxvii. 7.]}. His behaviour at Athens is the most remarkable instance in
his history of this method of acting. Instead of uttering any
invective against their Polytheism, he began a discourse upon the
Unity of the Divine Nature; and then proceeded to claim the altar \footnote{[Vide this argument in the mouth of
Dionysius (in Euseb. Hist. vii. 11, [\emph{ou pantes pantas}],
\&c.) as his plea for liberty of worship, with the neat retort of
the Prefect.]}, consecrated in the neighbourhood to the unknown God, as
the property of Him whom he preached to them, and to enforce his
doctrine of the Divine Immateriality, not by miracles, but by
argument, and that founded on the words of a heathen poet. This was
the example which the Alexandrians set before them in their
intercourse with the heathen, as may be shown by the following
instances.

Theonas, Bishop of Alexandria (A.D. 282-300),
has left his directions for the behaviour of Christians who were in
the service of the imperial court. The utmost caution is enjoined
them, not to give offence to the heathen emperor. If a Christian was
appointed librarian, he was to take good care not to show any contempt
for secular knowledge and the ancient writers. He was advised to make
himself familiar with the poets, philosophers, orators, and
historians, of classical literature; and, while discussing their
writings, to take incidental opportunities of recommending the
Scriptures, introducing mention of Christ, and by degrees revealing
the real dignity of His nature \footnote{Rose's Neander, Eccl. Hist. vol. i. p.
145. ``Insurgere poterit Christi mentio, explicabitur paullatim
ejus sola divinitas.'' Tillem. Mem. vol. iv. p. 240, 241.}.

The conversion of Gregory of Neocæsarea, (A.D.
231) affords an exemplification of this procedure in an individual
case. He had originally attached himself to the study of rhetoric and
the law, but was persuaded by Origen, whose lectures he attended, to
exchange these pursuits, first for science, then for philosophy, then
for theology, so far as right notions concerning religion could be
extracted from the promiscuous writings of the various
philosophical sects. Thus, while professedly teaching him Pagan
philosophy, his skilful master insensibly enlightened him in the
knowledge of the Christian faith. Then leading him to Scripture, he
explained to him its difficulties as they arose; till Gregory,
overcome by the force of truth, announced to his instructor his
intention of exchanging the pursuits of this world for the service of
God \footnote{This was Origen's usual method, vide Euseb.
Eccl. Hist. vi. 18. He has signified it himself in these words: [\emph{gumnasion
men phamen einai t\={e}s psuch\={e}s t\={e}n anthr\={o}pin\={e}n
sophian, telos de t\={e}n theian}]. Contr. Cels. vi. 13.}.

Clement's Stromateis (A.D. 200), a work which
has already furnished us with illustrations of the Alexandrian method
of teaching, was written with the design of converting the learned
heathen, and pursues the same plan which Origen adopted towards
Gregory. The author therein professes his wish to blend together
philosophy and religion, refutes those who censure the former, shows
the advantage of it, and how it is to be applied. This leading at once
to an inquiry concerning what particular school of philosophy is to be
held of divine origin, he answers in a celebrated passage, that all
are to be referred thither as far as they respectively inculcate the
principles of piety and morality, and none, except as containing the
portions and foreshadowings of the truth. ``By philosophy,''
he says, ``I do not mean the Stoic, nor the Platonic, nor the
Epicurean and Aristotelic, but all good doctrine in every one of the
schools, all precepts of holiness combined with religious knowledge.
All this, taken together, or the \emph{Eclectic}, I call \emph{philosophy}:
whereas the rest are mere forgeries of the human intellect, and
in no respect to be accounted divine.'' \footnote{Clem. Strom. i. 7.} At the same time, to mark out the peculiar divinity of
Revealed Religion, he traces all the philosophy of the heathen to the
teaching of the Hebrew sages, earnestly maintaining its entire
subserviency to Christianity, as but the love of that truth which the
Scriptures really impart.

The same general purpose of conciliating the heathen, and (as far
as might be,) indulging the existing fashions to which their
literature was subjected, may be traced in the slighter compositions \footnote{[\emph{logoi}]. [Such are those (Pagan) of
Maximus Tyrius. Three sacred narratives of Eusebius Emesenus are to be
found at Vienna. Augusti has published one of them: Bonn, 1820. Vide
Lambec. Bibl. Vind. iv. p. 286.]} which the Christians published in defence of their religion \footnote{Dodwell in Iren. Diss. vi. § 14, 16.}, being what in this day might be called pamphlets, written in
imitation of speeches after the manner of Isocrates, and adorned with
those graces of language which the schools taught, and the inspired
Apostle has exhibited in his Epistle to the Hebrews. Clement's
Exhortation to the Gentiles is a specimen of this style of writing; as
also those of Athenagoras and Tatian, and that ascribed to Justin
Martyr.

Again:—the last-mentioned Father supplies us with an instance of
an economical relinquishment of a sacred doctrine. When Justin Martyr,
in his argument with the Jew Trypho, (A.D. 150.)
finds himself unable to convince him from the Old Testament of the
divinity of Christ, he falls back upon the doctrine of His divine
Mission, as if this were a point indisputable on the one hand,
and on the other, affording a sufficient ground, from which to
advance, when expedient, to the proof of the full evangelical truth \footnote{Vide Bull, Judic. Eccl. vi. 7.}. In the same passage, moreover, as arguing with an unbeliever,
he permits himself to speak without an anathema of those (the
Ebionites) who professed Christianity, and yet denied Christ's
divinity. Athanasius himself fully recognizes the propriety of this
concealment of the doctrine on a fitting occasion, and thus accounts
for the silence of the Apostles concerning it, in their speeches
recorded in the book of Acts, viz. that they were unwilling, by a
disclosure of it, to prejudice the Jews against those miracles, the
acknowledgment of which was a first step towards their receiving it \footnote{Athan. de Sent. Dionys. 8. Theodoret,
Chrysostom, and others say the same. Vide Suicer. Thesaurus, verb [\emph{stoicheion}],
and Whitby on Heb. v. 12.}.

Gregory of Neocæsarea (A.D. 240-270), whose
conversion by Origen has already been adduced in illustration,
furnishes us in his own conduct with a similar but stronger instance
of an economical concealment of the full truth. It seems that certain
heretical teachers, in the time of Basil, ascribed to Gregory, whether
by way of censure or in self-defence, the Sabellian view of the
Trinity; and, moreover, the belief that Christ was a creature. The
occasion of these statements, as imputed to him, was a \emph{vivâ voce}
controversy with a heathen, which had been taken down in writing by
the bystanders. The charge of Sabellianism is refuted by Gregory's
extant writings; both imputations, however, are answered by St. Basil,
and that, on the principle of controversy which I have above
attempted to describe. ``When Gregory,'' he says,
``declared that the Father and Son were two in our conception of
them, one in \emph{hypostasis}, he spoke not as teaching doctrine, but
as arguing with an unbeliever, viz. in his disputation with Ælianus;
but this distinction our heretical opponents could not enter into,
much as they pride themselves on the subtlety of their intellect. Even
granting there were no mistakes in taking the notes (which, please
God, it is my intention to prove from the text as it now stands), it
is to be supposed, that he did not think it necessary to be very exact
in his doctrinal terms, when employed in converting a heathen; but in
some things, even to concede to his feelings, that he might gain him
over to the cardinal points. Accordingly, you may find many
expressions there, of which heretics now take great advantage, such as
'creature,' 'made,' and the like. So again, many statements which he
has made concerning the Incarnation, are referred to the Divine Nature
of the Son by those who do not skilfully enter into his meaning; as,
indeed, is the very expression in question which they have circulated
\footnote{Basil, Epist. ccx. § 5.}.''

I will here again instance a parallel use of the Economy on the
part of Athanasius himself, and will avail myself of the words of the
learned Petavius. ``Even Athanasius,'' he says, ``whose
very gift it was, above all other Fathers, to possess a clear and
accurate knowledge of the Catholic doctrine concerning the Trinity, so
that all succeeding antagonists of Arianism may be truly said to have
derived their powers and their arguments from him, even this keen
and vigilant champion of orthodoxy, in arguing with the Gentiles for
the Divinity and incarnation of the Word, urges them with
considerations drawn from their own philosophical notions concerning
Him. Not that he was ignorant how unlike orthodoxy, and how like
Arianism, such notions were, but he bore in mind the necessity of
favourably disposing the minds of the Gentiles to listen to his
teaching; and he was aware that it was one thing to lay the rudiments
of the faith in an ignorant or heathen mind, and another to defend the
faith against heretics, or to teach it dogmatically. For instance, in
answering their objection to the Divine Word having taken flesh, which
especially offended them, he bids them consider whether they are not
inconsistent in dwelling upon this, while they themselves believe that
there is a Divine Word, the presiding principle and soul of the world,
through the movements of which He is visibly displayed; 'for what (he
asks) does Christianity say more than that the Word has presented
Himself to the inspection of our senses by the instrumentality of a
body?' And yet it is certain that the Father and the pervading Word of
the Platonists, differed materially from the Sacred Persons of the
Trinity, as we hold the doctrine, and Athanasius too, in every page of
his writings.'' \footnote{Petav. de Trin. ii. præf. 3, § 5.
[abridged and re-arranged. Vide ibid. iii. 1, § 6. Vide also Euseb.
contr. Marcel. ii. 22, p. 140; iii. 3. pp. 161, 2].}

There are instances in various ways of the economical method, that
is, of accommodation to the feelings and prejudices of the hearer, in
leading him to the reception of a novel or unacceptable doctrine.
It professes to be founded in the actual necessity of the case;
because those who are strangers to the tone of thought and principles
of the speaker, cannot at once be initiated into his system, and
because they must begin with imperfect views; and therefore, if he is
to teach them at all, he must put before them large propositions,
which he has afterwards to modify, or make assertions which are but
parallel or analogous to the truth, rather than coincident with it.
And it cannot be denied that those who attempt to speak at all times
the naked truth, or rather the commonly-received expression of it, are
certain, more than other men, to convey wrong impressions of their
meaning to those who happen to be below them, or to differ widely from
them, in intelligence and cast of mind. On the other hand, the abuse
of the Economy in the hands of unscrupulous reasoners, is obvious.
Even the honest controversialist or teacher will find it very
difficult to represent without misrepresenting, what it is yet his
duty to present to his hearers with caution or reserve. Here the
obvious rule to guide our practice is, to be careful ever to maintain
substantial truth in our use of the economical method. It is thus we
lead forward children by degrees, influencing and impressing their
minds by means of their own confined conceptions of things, before we
attempt to introduce them to our own; yet at the same time modelling
their thoughts according to the analogy of those to which we mean
ultimately to bring them. Again, the information given to the blind
man, that scarlet was like the sound of a trumpet, is an instance of
an unexceptionable economy, since it was as true as it could be
under the circumstances of the case, conveying a substantially correct
impression as far as it went.

In applying this rule to the instances above given, it is plain
that Justin, Gregory, or Athanasius, were justifiable or not in their
Economy, according as they did or did not practically mislead their
opponents. Merely to leave a man in errors which he had independently
of us, or to abstain from removing them, cannot be blamed as a fault,
and may be a duty; though it is so difficult to hit the mark in these
perplexing cases, that it is not wonderful, should these or other
Fathers have failed at times, and said more or less than was proper.
Again, in the instances of St. Paul, Theonas, Origen, and Clement, the
doctrine which their conduct implies, is the Divinity of Paganism; a
true doctrine, though the heathen whom they addressed would not at
first rightly apprehend it. But I am aware that some persons will
differ from me here, and others will be perplexed about my meaning. So
let this be a reserved point, to be considered when we have finished
the present subject.

The Alexandrian Father who has already been quoted, accurately
describes the rules which should guide the Christian in speaking and
acting economically. ``Being fully persuaded of the omnipresence
of God,'' says Clement, ``and ashamed to come short of the
truth, he is satisfied with the approval of God, and of his own
conscience. Whatever is in his mind, is also on his tongue; towards
those who are fit recipients, both in speaking and living, he
harmonizes his profession with his thoughts. He both thinks and speaks
the truth; except when careful treatment is necessary, and then, as a
physician for the good of his patients, he will lie, or rather
utter a lie, as the Sophists say. For instance, the noble Apostle
circumcised Timothy, while he cried out and wrote down, 'Circumcision
availeth not ... Nothing, however, but his neighbour's good will lead
him to do this ... He gives himself up for the Church, for the friends
whom he hath begotten in the faith for an ensample to those who have
the ability to undertake the high office (\emph{economy}) of a religious and
charitable teacher, for an exhibition of truth in his words, and for
the exercise of love towards the Lord.'' \footnote{Clem. Strom. vii. 8, 9 (abridged). [Vide
Plat. Leg. ii. 8, [\emph{oupote pseudetai, kan pseudos leg\={e}}].
Sext. Empir. adv. Log. p. 378, with notes T and U. On this whole
subject, vide the Author's ``Apologia,'' notes F and G, pp.
343-363.]}

Further light will be thrown upon the doctrine of the Economy, by
considering it as exemplified in the dealings of Providence towards
man. The word occurs in St. Paul's Epistle to the Ephesians, where it
is used for that series of Divine appointments viewed as a whole, by
which the Gospel is introduced and realized among mankind, being
translated in our version ``\emph{dispensation}.'' It will
evidently bear a wider sense, embracing the Jewish and patriarchal
dispensations, or any Divine procedure, greater or less, which
consists of means and an end. Thus it is applied by the Fathers, to
the history of Christ's humiliation, as exhibited in the doctrines of
His incarnation, ministry, atonement, exaltation, and mediatorial
sovereignty, and, as such distinguished from the ``\emph{theologia}''
or the collection of truths relative to His personal indwelling in the
bosom of God. Again, it might with equal fitness be used for the
general system of providence by which the world's course is
carried on; or, again, for the work of creation itself, as opposed to
the absolute perfection of the Eternal God, that internal
concentration of His Attributes in self-contemplation, which took
place on the seventh day, when He rested from all the work which He
had made. And since this everlasting and unchangeable quiescence is
the simplest and truest notion we can obtain of the Deity, it seems to
follow, that strictly speaking, all those so-called Economies or
dispensations, which display His character in action, are but
condescensions to the infirmity and peculiarity of our minds, shadowy
representations of realities which are incomprehensible to creatures
such as ourselves, who estimate everything by the rule of association
and arrangement, by the notion of a purpose and plan, object and
means, parts and whole. What, for instance, is the revelation of
general moral laws, their infringement, their tedious victory, the
endurance of the wicked, and the ``winking at the times of
ignorance,'' but an ``\emph{Economia}'' of greater truths untold,
the best practical communication of them which our minds in their
present state will admit? What are the phenomena of the external
world, but a divine mode of conveying to the mind the realities of
existence, individuality, and the influence of being on being, the
best possible, though beguiling the imagination of most men with a
harmless but unfounded belief in matter as distinct from the
impressions on their senses? This at least is the opinion of some
philosophers, and whether the particular theory be right or wrong, it
serves as an illustration here of the great truth which we are
considering. Or what, again, as others hold, is the popular
argument from final causes but an ``\emph{Economia}'' suited
to the practical wants of the multitude, as teaching them in the
simplest way the active presence of Him, who after all dwells
intelligibly, prior to argument, in their heart and conscience? And
though on the mind's first mastering this general principle, it seems
to itself at the moment to have cut all the ties which bind it to the
universe, and to be floated off upon the ocean of interminable
scepticism; yet a true sense of its own weakness brings it back, the
instinctive persuasion that it must be intended to rely on something,
and therefore that the information given, though philosophically
inaccurate, must be practically certain; a sure confidence in the love
of Him who cannot deceive, and who has impressed the image and thought
of Himself and of His will upon our original nature. Here then we may
lay down with certainty as a consolatory truth, what was but a rule of
duty when we were reviewing the Economies of man; viz. that whatever
is told us from heaven, is true in so full and substantial a sense,
that no possible mistake can arise practically from following it. And
it may be added, on the other hand, that the greatest risk will result
from attempting to be wiser than God has made us, and to outstep in
the least degree the circle which is prescribed as the limit of our
range. This is but the duty of implicit faith in Him who knows what is
good for us, and who has ordained that in our practical concerns
intellectual ability should do no more than enlighten us in the
difficulties of our situation, not in the solutions of them.
Accordingly, we may safely admit the first chapter of the book of Job,
the twenty-second of the first book of Kings, and other passages
of Scripture, to be \emph{Economies}, that is, representations conveying
substantial truth in the form in which we are best able to receive it;
and to be accepted by us and used in their literal sense, as our
highest wisdom, because we have no powers of mind equal to the more
philosophical determination of them. Again, the Mosaic Dispensation
was an Economy, simulating (so to say) unchangeableness, when from the
first it was destined to be abolished. And our Blessed Lord's conduct
on earth abounds with the like gracious and considerate condescension
to the weakness of His creatures, who would have been driven either to
a terrified inaction or to presumption, had they known then as
afterwards the secret of His Divine Nature.

I will add two or three instances, in which this doctrine of the
Divine Economies has been wrongly applied; and I do so from necessity,
lest the foregoing remarks should seem to countenance errors, which I
am most desirous at all times and every where to protest against.

For instance, the Economy has been employed to the disparagement of
the Old Testament Saints; as if the praise bestowed on them by
Almighty God were but economically given, that is, with reference to
their times and circumstances; their real insight into moral truth
being possibly below the average standard of knowledge in matters of
faith and practice received among nations rescued from the rude and
semi-savage state in which they are considered to have lived. And
again, it has been even supposed, that injunctions, as well as praise,
have been thus given them, which an enlightened age is at liberty
to criticize; for instance, the command to slay Isaac has sometimes
been viewed as an economy, based upon certain received ideas in
Abraham's day, concerning the innocence and merit of human sacrifice.
It is enough to have thus disclaimed participation in these theories,
which of course are no objection to the general doctrine of the
Economy, unless indeed it could be shown, that those who hold a
principle are answerable for all the applications arbitrarily made of
it by the licentious ingenuity of others.

Again, the principle of the Economy has sometimes been applied to
the interpretation of the New Testament. It has been said, for
instance, that the Epistle to the Hebrews does not state the simple
truth in the sense in which the Apostles themselves believed it, but
merely as it would be palatable to the Jews. The advocates of this
hypothesis have proceeded to maintain, that the doctrine of the
Atonement is no part of the essential and permanent evangelical
system. To a conscientious reasoner, however, it is evident, that the
structure of the Epistle in question is so intimately connected with
the reality of the expiatory scheme, that to suppose the latter
imaginary, would be to impute to the writer, not an economy (which
always preserves substantial truth), but a gross and audacious deceit.

A parallel theory to this has been put forward by men of piety
among the Predestinarians, with a view of reconciling the
inconsistency between their faith and practice. They have suggested,
that the promises and threats of Scripture are founded on an economy,
which is needful to effect the conversion of the elect, but clears up and vanishes under the light of the true spiritual
perception, to which the converted at length attain. This has been
noticed in another connexion, and will here serve as one among many
illustrations which might be given, of the fallacious application of a
true principle. And so much upon the \emph{Economia}.

5.

A question was just now reserved, as interfering with the subject
then before us. In what sense can it be said, that there is any
connection between Paganism and Christianity so real, as to warrant
the preacher of the latter to conciliate idolaters by allusion to it?
St. Paul evidently connects the true religion with the existing
systems which he laboured to supplant, in his speech to the Athenians
in the Acts, and his example is a sufficient guide to missionaries
now, and a full justification of the line of conduct pursued by the
Alexandrians, in the instances similar to it; but are we able to
account for his conduct, and ascertain the principle by which it was
regulated? I think we can; and the exhibition of it will set before
the reader another doctrine of the Alexandrian school, which it is as
much to our purpose to understand, and which I shall call \emph{the
divinity of Traditionary Religion}.

We know well enough for practical purposes what is meant by
Revealed Religion; viz. that it is the doctrine taught in the Mosaic
and Christian dispensations, and contained in the Holy Scriptures, and
is from God in a sense in which no other doctrine can be said to be
from Him. Yet if we would speak correctly, we must confess, on the
authority of the Bible itself, that all knowledge of religion is from
Him, and not only that which the Bible has transmitted to us.
There never was a time when God had not spoken to man, and told him to
a certain extent his duty. His injunctions to Noah, the common father
of all mankind, is the first recorded fact of the sacred history after
the deluge. Accordingly, we are expressly told in the New Testament,
that at no time He left Himself without witness in the world, and that
in every nation He accepts those who fear and obey Him. It would seem,
then, that there is something true and divinely revealed, in every
religion all over the earth, overloaded, as it may be, and at times
even stifled by the impieties which the corrupt will and understanding
of man have incorporated with it. Such are the doctrines of the power
and presence of an invisible God, of His moral law and governance, of
the obligation of duty, and the certainty of a just judgment, and of
reward and punishment, as eventually dispensed to individuals; so that
Revelation, properly speaking, is an universal, not a local gift; and
the distinction between the state of Israelites formerly and
Christians now, and that of the heathen, is, not that we can, and they
cannot attain to future blessedness, but that the Church of God ever
has had, and the rest of mankind never have had, authoritative
documents of truth, and appointed channels of communication with Him.
The word and the Sacraments are the characteristic of the elect people
of God; but all men have had more or less the guidance of Tradition,
in addition to those internal notions of right and wrong which the
Spirit has put into the heart of each individual.

This vague and uncertain family of religious truths, originally
from God, but sojourning without the sanction of miracle, or a
definite home, as pilgrims up and down the world, and discernible and
separable from the corrupt legends with which they are mixed, by the
spiritual mind alone, may be called the \emph{Dispensation of Paganism},
after the example of the learned Father already quoted \footnote{Clement says, [\emph{T\={e}n philosophian
Hell\={e}sin oion diath\={e}k\={e}n oikeian dedosthai,
hupobathran ousan t\={e}s kata Christon philosophias}]. Strom
vi. p. 648.}. And further, Scripture gives us reason to believe that the
traditions, thus originally delivered to mankind at large, have been
secretly re-animated and enforced by new communications from the
unseen world; though these were not of such a nature as to be produced
as evidence, or used as criteria and tests, and roused the attention
rather than informed the understandings of the heathen. The book of
Genesis contains a record of the Dispensation of Natural Religion, or
Paganism, as well as of the patriarchal. The dreams of Pharaoh and
Abimelech, as of Nebuchadnezzar afterwards, are instances of the
dealings of God with those to whom He did not vouchsafe a written
revelation. Or should it be said, that these particular cases merely
come within the range of the Divine supernatural Governance which was
in their neighbourhood,—an assertion which requires proof,—let the
book of Job be taken as a less suspicious instance of the dealings of
God with the heathen. Job was a pagan in the same sense in which the
Eastern nations are Pagans in the present day. He lived among
idolaters \footnote{Job xxxi. 26-28.}, yet he and
his friends had cleared themselves from the superstitions with which
the true creed was beset; and while one of them was divinely
instructed by dreams \footnote{Ibid. iv. 13, \&c.},
he himself at length heard the voice of God out of the whirlwind, in
recompense for his long trial and his faithfulness under it \footnote{Job xxxviii. 1; xlii. 10, \&c. [Vide
also Gen. xli. 45. Exod. iii. 1. Jon. i. 5-16.]}. Why should not the book of Job be accepted by us, as a
gracious intimation given us, who are God's sons, for our comfort,
when we are anxious about our brethren who are still ``scattered
abroad'' in an evil world; an intimation that the Sacrifice, which
is the hope of Christians, has its power and its success, wherever men
seek God with their whole heart?—If it be objected that Job lived in
a less corrupted age than the times of ignorance which followed,
Scripture, as if for our full satisfaction, draws back the curtain
farther still in the history of Balaam. There a bad man and a heathen
is made the oracle of true divine messages about doing justly, and
loving mercy, and walking humbly; nay, even among the altars of
superstition, the Spirit of God vouchsafes to utter prophecy \footnote{Numb. xxii.-xxiv. Mic. vi. 5, 8.}. And so in the cave of Endor, even a saint was sent from the
dead to join the company of an apostate king, and of the sorceress
whose aid he was seeking \footnote{1 Sam. xxviii. 14.}.
Accordingly, there is nothing unreasonable in the notion, that there
may have been heathen poets and sages, or sibyls again, in a certain
extent divinely illuminated, and organs through whom religious and
moral truth was conveyed to their countrymen; though their knowledge
of the Power from whom the gift came, nay, and their perception of the
gift as existing in themselves, may have been very faint or defective.

This doctrine, thus imperfectly sketched, shall now be presented to
the reader in the words of St. Clement. ``To the Word of
God,'' he says, ``all the host of angels and heavenly powers
is subject, revealing, as He does, His holy office (\emph{economy}),
for Him who has put all things under Him. Wherefore, His are all men;
some actually knowing Him, others not as yet: some as friends''
(Christians), ``others as faithful servants'' (Jews),
``others as simply servants'' (heathen). ``He is the
Teacher, who instructs the enlightened Christian by mysteries, and the
faithful labourer by cheerful hopes, and the hard of heart with His
keen corrective discipline; so that His providence is particular,
public, and universal ... He it is who gives to the Greeks their
philosophy by His ministering Angels ... for He is the Saviour not of
these or those, but of all ... His precepts, both the former and the
latter, are drawn forth from one fount; those who were before the Law,
not suffered to be without law, those who do not hear the Jewish
philosophy, not surrendered to an unbridled course. Dispensing in
former times to some His precepts, to others philosophy, now at
length, by His own personal coming, He has closed the course of
unbelief, which is henceforth inexcusable; Greek and barbarian''
(that is, Jew) ``being led forward by a separate process to that
perfection which is through faith.'' \footnote{Clem. Strom. vii. 2.}

If this doctrine be scriptural, it is not difficult to determine
the line of conduct which is to be observed by the Christian apologist
and missionary. Believing God's hand to be in every system, so far
forth as it is true (though Scripture alone is the depositary of His
unadulterated and complete revelation), he will, after St. Paul's
manner, seek some points in the existing superstitions as the basis of
his own instructions, instead of indiscriminately condemning and
discarding the whole assemblage of heathen opinions and practices; and
he will address his hearers, not as men in a state of actual
perdition, but as being in imminent danger of ``the wrath to
come,'' because they are in bondage and ignorance, and probably
under God's displeasure, that is, the vast majority of them are so in
fact; but not necessarily so, from the very circumstance of their
being heathen. And while he strenuously opposes all that is
idolatrous, immoral, and profane, in their creed, he will profess to
be leading them on to perfection, and to be recovering and purifying,
rather than reversing the essential principles of their belief.

A number of corollaries may be drawn from this view of the relation
of Christianity to Paganism, by way of solving difficulties which
often perplex the mind. For example, we thus perceive the utter
impropriety of ridicule and satire as a means of preparing a heathen
population for the reception of the truth. Of course it is right,
soberly and temperately, to expose the absurdities of idol-worship;
but sometimes it is maintained that a writer, such as the infamous
Lucian, who scoffs at an established religion altogether, is the
suitable preparation for the Christian preacher,—as if infidelity
were a middle state between superstition and truth. This view derives
its plausibility from the circumstance that in drawing out systems in
writing, to erase a false doctrine is the first step towards inserting
the true. Accordingly, the mind is often compared to a tablet or
paper: a state of it is contemplated of absolute freedom from all
prepossessions and likings for one system or another, as a first step
towards arriving at the truth; and infidelity represented as that
candid and dispassionate frame of mind, which is the desideratum. For
instance, at the present day, men are to be found of high religious
profession, who, to the surprise and grief of sober minds, exult in
the overthrow just now of religion in France, as if an unbeliever were
in a more hopeful state than a bigot, for advancement in real
spiritual knowledge. But in truth, the mind never can resemble a blank
paper, in its freedom from impressions and prejudices. Infidelity is a
positive, not a negative state; it is a state of profaneness, pride,
and selfishness; and he who believes a little, but encompasses that
little with the inventions of men, is undeniably in a better condition
than he who blots out from his mind both the human inventions, and
that portion of truth which was concealed in them.

Again: it is plain that the tenderness of dealing, which it is our
duty to adopt towards a heathen unbeliever, is not to be used towards
an apostate. No \emph{economy} can be employed towards those who have been
once enlightened, and have fallen away. I wish to speak explicitly on
this subject, because there is a great deal of that spurious charity
among us which would cultivate the friendship of those who, in a
Christian country, speak against the Church or its creeds. Origen and
others were not unwilling to be on a footing of intercourse with the
heathen philosophers of their day, in order, if it were possible, to
lead them into the truth; but deliberate heretics and apostates, those who had known the truth, and rejected it, were objects of their
abhorrence, and were avoided from the truest charity to them. For what
can be said to those who already know all we have to say? And how can
we show our fear for their souls, nay, and for our own steadfastness,
except by a strong action? Thus Origen, when a youth, could not be
induced to attend the prayers of a heretic of Antioch whom his
patroness had adopted, from a loathing \footnote{[\emph{Bdeluttomenos}]. Eus. Hist. vi. 2
[vii. 7, Eulog. ap. Phot. p. 861]}, as he says, of heresy. And St. Austin himself tells us, that
while he was a Manichee, his own mother would not eat at the same
table with him in her house, from her strong aversion to the
blasphemies which were the characteristic of his sect \footnote{Bingham, Antiq. xvi. 2, § 11.}. And Scripture fully sanctions this mode of acting, by the
severity with which such unhappy men are spoken of, on the different
occasions when mention is made of them \footnote{Rom. xvi. 17. 2 Thess. iii. 14. 2 John 10,
11, \&c.}.

Further: the foregoing remarks may serve to show us, with what view
the early Church cultivated and employed heathen literature in its
missionary labours; viz. not with the notion that the cultivation,
which literature gives, was any substantial improvement of our moral
nature, but as thereby opening the mind, and rendering it susceptible
of an appeal; nor as if the heathen literature itself had any direct
connexion with the matter of Christianity, but because it contained in
it the scattered fragments of those original traditions which might be
made the means of introducing a student to the Christian system, being
the ore in which the true metal was found. The account above given of
the conversion of Gregory is a proof of this.

The only danger to which the Alexandrian doctrine is exposed, is
that of its confusing the Scripture Dispensations with that of Natural
Religion, as if they were of equal authority; as if the Gospel had not
a claim of acceptance on the conscience of all who heard it, nor
became a touchstone of their moral condition; and as if the Bible, as
the Pagan system, were but partially true, and had not been attested
by the discriminating evidence of miracles. This is the heresy of the
Neologians in this day, as it was of the Eclectics in primitive times;
as will be shown in the next section. The foregoing extract from
Clement shows his entire freedom from so grievous an error; but in
order to satisfy any suspicion which may exist of his using language
which may have led to a more decided corruption after his day, I will
quote a passage from the sixth book of his Stromateis, in which he
maintains the supremacy of Revealed Religion, as being in fact the
source and test of all other religions; the extreme imperfection of
the latter; the derivation of whatever is true in these from
Revelation; the secret presence of God in them, by that Word of Life
which is directly and bodily revealed in Christianity; and the
corruption and yet forced imitation of the truth by the evil spirit in
such of them, as he wishes to make pass current among mankind.

``Should it be said that the Greeks discovered philosophy by
human wisdom,'' he says, ``I reply, that I find the Scriptures
declare all wisdom to be a divine gift: for instance, the Psalmist
considers wisdom to be the greatest of gifts, and offers this
petition, 'I am thy servant, make me wise.' And does not David ask for
illumination in its diverse functions, when he says 'Teach me
goodness, discipline, and knowledge, for I have believed Thy
precepts'? Here he confesses that the Covenants of God are of supreme
authority, and vouchsafed to the choice portion of mankind. Again,
there is a Psalm which says of God, 'He hath not acted thus with any
other nation, and His judgments He hath not revealed to them;' where
the words, 'He hath not done \emph{thus},' imply that He hath indeed
done somewhat, but not \emph{thus}. By using \emph{thus} he contrasts
their state with our superiority; else the Prophet might simply have
said, 'He hath not acted with other nations,' without adding \emph{thus}.
The prophetical figure, 'The Lord is over many waters,' refers to the
same truth; that is, a Lord not only of the different covenants, but
also of the various methods of teaching, which lead to righteousness,
whether among the Gentiles or the Jews. David also bears his testimony
to this truth, when he says in the Psalm, 'Let the sinners be turned
into hell, all the nations which \emph{forget} God;' that is, they
forget whom they formerly remembered, they put aside Him whom they
knew before they forgot. It seems then there was some dim knowledge of
God even among the Gentiles ... They who say that philosophy
originates with the devil, would do well to consider what Scripture
says about the devil's being transformed into an Angel of light. For
what will he do then? it is plain he will prophesy. Now if he
prophesies as an Angel of light, of course he will speak what is true.
If he shall prophesy angelic and enlightened doctrine, he will
prophesy what is profitable also; that is, \emph{at the time when} he is thus
changed in his apparent actions, far different as he is at bottom in
his real apostasy. For how would he deceive except by craftily
leading on the inquirer \emph{by means of truth}, to an intimacy with
himself, and so at length seducing him into error? ... Therefore
philosophy is not false, though he who is thief and liar speaks truth
by a change in his manner of acting ... The philosophy of the Greeks,
limited and particular as it is, contains the rudiments of that really
perfect knowledge which is beyond this world, which is engaged in
intellectual objects, and upon those more spiritual, which eye hath
not seen, nor ear heard, nor the heart of man conceived, before they
were made clear to us by our Great Teacher, who reveals the holy of
holies, and still holier truths in an ascending scale, to those who
are genuine heirs of the Lord's adoption.'' \footnote{Strom. vi. 8.}

6.

What I have said about the method of teaching adopted by the
Alexandrian, and more or less by the other primitive Churches, amounts
to this; that they on principle refrained from telling unbelievers all
they believed themselves, and further, that they endeavoured to
connect their own doctrine with theirs, whether Jewish or pagan,
adopting their sentiments and even their language, as far as they
lawfully could. Some instances of this have been given; more will
follow, in the remarks which I shall now make upon the influence of
Platonism on their theological language.

The reasons, which induced the early Fathers to avail themselves of
the language of Platonism, were various. They did so, partly as an \emph{argumentum
ad hominem}; as if the Christian were not professing in the
doctrine of the Trinity a more mysterious tenet, than that which
had been propounded by a great heathen authority; partly to conciliate
their philosophical opponents; partly to save themselves the
arduousness of inventing terms, where the Church had not yet
authoritatively supplied them; and partly with the hope, or even
belief, that the Platonic school had been guided in portions of its
system by a more than human wisdom, of which Moses was the unknown but
real source. As far as these reasons depend upon the rule of the
Economy, they have already been considered; and an instance of their
operation given in the exoteric conduct of Athanasius himself, whose
orthodoxy no one questions. But the last reason given, their suspicion
of the divine origin of the Platonic doctrine, requires some
explanation.

It is unquestionable that, from very early times, traditions have
been afloat through the world, attaching the notion of a Trinity, in
some sense or other, to the First Cause. Not to mention the traces of
this doctrine in the classical and the Indian mythologies, we detect
it in the Magian hypothesis of a supreme and two subordinate
antagonist deities in Plutarch's Trinity of God, matter, and the evil
spirit, and in certain heresies in the first age of the Church, which,
to the Divine Being and the Demiurgus, added a third original
principle, sometimes the evil spirit, and sometimes matter \footnote{Cudworth, Intell. Syst. i. 4, § 13, 16.
Beausobre, Hist. de Manich. iv. 6, § 8, \&c.}. Plato has adopted the same general notion; and with no closer
or more definite approach to the true doctrine. On the whole, it seems
reasonable to infer, that the heathen world possessed traditions too
ancient to be rejected, and too sacred to be used in popular
theology. If Plato's doctrine bears a greater apparent resemblance to
the revealed truth than that of others, this is owing merely to his
reserve in speaking on the subject. His obscurity allows room for an
ingenious fancy to impose a meaning upon him. Whether he includes in
his Trinity the notion of a First Cause, its active energy, and the
influence resulting from it; or again, the divine substance as the
source of all spiritual beings from eternity, the divine power and
wisdom as exerted in time in the formation of the material world, and
thirdly, the innumerable derivative spirits by whom the world is
immediately governed, is altogether doubtful. Nay, even the writers
who revived his philosophy in the third and fourth centuries after
Christ, and embellished the doctrine with additions from Scripture,
discover a like extraordinary variation in their mode of expounding
it. The Maker of the world, the Demiurge, considered by Plato
sometimes as the first, sometimes as the second principle, is by
Julian placed as the second, by Plotinus as the third, and by Proclus
as the fourth, that is, the last of three subordinate powers, all
dependent on a First, or the One Supreme Deity \footnote{Petav. Theol. Dogm. tom. ii. i. 1, § 5.}. In truth, speculations, vague and unpractical as these, made
no impression on the minds of the heathen philosophers, perhaps as
never being considered by them as matters of fact, but as allegories
and metaphysical notions, and accordingly, caused in them no
solicitude or diligence to maintain consistency in their expression of
them.

But very different was the influence of the ancient theory of
Plato, however originated, when it came in contact with believers
in the inspired records, who at once discerned in it that mysterious
Doctrine, brought out as if into bodily shape and almost practical
persuasiveness, which lay hid under the angelic manifestations of the
Law and the visions of the Prophets. Difficult as it is to determine
the precise place in the sacred writings, where the Divine Logos or
Word was first revealed, and how far He is intended in each particular
passage, the idea of Him is doubtless seated very deeply in their
teaching. Appearing first as if a mere created minister of God's will,
He is found to be invested with an ever-brightening glory, till at
length we are bid fall down as before the personal Presence and
consubstantial Representative of the one God. Those then, who were
acquainted with the Sacred Volume, possessed in it a key, more or less
exact according to their degree of knowledge, for that aboriginal
tradition which the heathen ignorantly but piously venerated, and were
prompt in appropriating the language of philosophers, with a changed
meaning, to the rightful service of that spiritual kingdom, of which a
divine personal mediation was the great characteristic. In the books
of Wisdom and Ecclesiasticus, and much more, in the writings of Philo,
the Logos of Plato, which had denoted the divine energy in forming the
world, or the Demiurge, and the previous all-perfect incommunicable
design of it, or the Only-begotten, was arrayed in the attributes of
personality, made the instrument of creation, and the revealed Image
of the incomprehensible God. Amid such bold and impatient
anticipations of the future, it is not wonderful that the Alexandrian
Jews outstepped the truth which they hoped to appropriate; and that
intruding into things not seen as yet, with the confidence of
prophets rather than of disciples of Revelation, they eventually
obscured the doctrine when disclosed, which we may well believe they
loved in prospect and desired to honour. This remark particularly
applies to Philo, who associating it with Platonic notions as well as
words, developed its lineaments with so rude and hasty a hand, as to
separate the idea of the Divine Word from that of the Eternal God; and
so perhaps to prepare the way for Arianism \footnote{This may be illustrated by the theological
language of the Paradise Lost, which, as far as the very words go, is
conformable both to Scripture and the writings of the early Fathers,
but becomes offensive as being dwelt upon as if it were literal, not
figurative. It is scriptural to say that the Son went forth from the
Father to create the worlds; but when this is made the basis of a
scene or pageant, it borders on Arianism. Milton has made Allegory, or
the Economy, \emph{real}. Vide infra. ch. ii. § 4. fin.}.

Even after this Alexandrino-Judaic doctrine had been corrected and
completed by the inspired Apostles St. Paul and St. John, it did not
lose its hold upon the Fathers of the Christian Church, who could not
but discern in the old Scriptures, even more clearly than their
predecessors, those rudiments of the perfect truth which God's former
revelations concealed; and who in consequence called others, (as it
were,) to gaze upon these both as a prophetical witness in confutation
of unbelief, and in gratitude to Him who had wrought so marvellously
with His Church. But it followed from the nature of the case, that,
while they thus traced with watchful eyes, under the veil of the
literal text, the first and gathering tokens of that Divine Agent who
in fulness of time became their Redeemer, they were led to speak of
Him in terms short of that full confession of His divine
greatness, which the Gospel reveals, and which they themselves
elsewhere unequivocally expressed, especially as living in times
before the history of heresy had taught them the necessity of caution
in their phraseology. Thus, for instance, from a text in the book of
Proverbs \footnote{Prov. viii. 22, [\emph{Kurios ektisen}]. \emph{Septuag}.}, which they
understood to refer to Christ, Origen and others speak of Him as
``created by the Lord in the beginning, before His works of
old;'' meaning no more than that it was He, the true Light of man,
who was secretly intended by the Spirit, and mystically (though
incompletely) described, when Solomon spoke of the Divine Wisdom as
the instrument of God's providence and moral governance. In like
manner, when Justin speaks of the Son as the minister of God, it is
with direct reference to those numerous passages of the Old Testament,
in which a ministering angelic presence is more or less characterized
by the titles and attributes of Divine Perfection \footnote{Justin. Apol. i. 63. Tryph. 56, \&c.}. And, in the use of this emblematical diction they were
countenanced (not to mention the Apocalypse) by the almost sacred
authority of the platonizing books of Wisdom and Ecclesiasticus; works
so highly revered by the Alexandrian Church as to be put into the
hands of Catechumens as a preparation for inspired Scripture, contrary
to the discipline observed in the neighbouring Church of Jerusalem \footnote{Bingh. Antiq. x. 1. § 7.}.

The following are additional instances of Platonic language in the
early Fathers; though the reader will scarcely perceive at first sight
what is the fault in them, unless he happens to know the
defective or perverse sense in which philosophy or heresy used them \footnote{Petav. Theol. Dogm. tom. ii. 1. 3, 4.}. For instance, Justin speaks of the Word as ``fulfilling
the Father's will.'' Clement calls Him \footnote{[\emph{enno\={e}ma}].} ``the Thought or Reflection of God;'' and in another
place, ``the Second Principle of all things,'' the Father
Himself being the First. Elsewhere he speaks of the Son as an
``all-perfect, all-holy, all-sovereign, all-authoritative,
supreme, and all-searching nature, reaching close upon the sole
Almighty.'' In like manner Origen speaks of the Son as being
``the immediate Creator, and as it were, Artificer of the
world;'' and the Father, ``the Origin of it, as having
committed to His Son its creation.'' A bolder theology than this
of Origen and Clement is adopted by five early writers connected with
very various schools of Christian teaching; none of whom, however, are
of especial authority in the Church \footnote{Theophilus of Antioch (A.D.
168); Tatian, pupil of Justin Martyr (A.D. 169);
Athenagoras of Alexandria (A.D. 177); Hippolytus,
the disciple of Irenæus and friend of Origen (A.D.
222); and the Author who goes under the name of Novatian (A.D.
250). [This is Bull's view; for that maturely adopted by the author,
vide his ``Theological Tracts.'']}. They explained the Scripture doctrine of the generation of
the Word to mean, His manifestation at the beginning of the world as
distinct from God; a statement, which, by weakening the force of a
dogmatic formula which implies our Lord's Divine Nature, might perhaps
lend some accidental countenance after their day to the Arian denial
of it. These subjects will come before us in the next chapter.

I have now, perhaps, sufficiently accounted for the apparent
liberality of the Alexandrian School; which notwithstanding, was
strict and uncompromising, when its system is fairly viewed as a
whole, and with reference to its objects, and as distinct from that
rival and imitative philosophy, to be mentioned in the next section,
which rose out of it at the beginning of the third century, and with
which it is by some writers improperly confounded. That its principles
were always accurately laid, or the conduct of its masters nicely
adjusted to them, need not be contended; or that they opposed
themselves with an exact impartiality to every form of error which
assailed the Church; or that they duly entered into and soundly
applied the Jewish Scriptures; or that in conducting the Economy they
were altogether free from an ambitious imitation of the Apostles,
nobly conceived indeed, but little becoming uninspired teachers. It
may unreluctantly be confessed, wherever it can be proved, that their
exoteric professions at times affected the purity of their esoteric
doctrine, though this remark scarcely applies to their statements on
the subject of the Trinity; and that they indulged a boldness of
inquiry, such as innocence prompts, rashness and irreverence corrupt,
and experience of its mischievous consequences is alone able to
repress. Still all this, and much more than this, were it to be found,
weighs as nothing against the mass of testimonies producible from
extant documents in favour of the real orthodoxy of their creed.
Against a multitude of the very strongest and most explicit
declarations of the divinity of Christ, some of which will be cited in
their proper place, but a very few apparent exceptions to the
strictest language of technical theology can be gathered from their
writings, and these are sufficiently explained by the above
considerations. And further, such is the high religious temper which
their works exhibit, as to be sufficient of itself to convince the
Christian inquirer, that they would have shrunk from the deliberate
blasphemy with which Arius in the succeeding century assailed and
scoffed at the awful majesty of his Redeemer.

Origen, in particular, that man of strong heart, who has paid for
the unbridled freedom of his speculations on other subjects of
theology, by the multitude of grievous and unfair charges which burden
his name with posterity, protests, by the forcible argument of a life
devoted to God's service, against his alleged connexion with the cold
disputatious spirit, and the unprincipled domineering ambition, which
are the historical badges of the heretical party. Nay, it is a
remarkable fact that it was he who discerned the heresy \footnote{``The Word,'' says Origen, ``being the Image of the Invisible God, must Himself be invisible. Nay,
I will maintain further, that as being the Image He is eternal, as the
God whose Image He is. For when was that God, whom St. John calls the
Light, destitute of the Radiance of His incommunicable glory, so that
a man may dare to ascribe a beginning of existence to the Son? ... Let
a man, who dares to say that the Son is not from eternity, consider
well, that this is all one with saying, Divine Wisdom had a beginning,
or Reason, or Life.'' Athan. de Decr. Nic. § 27. Vide also his [\emph{peri
arch\={o}n}] (if Ruffinus may be trusted), for his denouncement
of the still more characteristic Arianisms of the [\emph{\={e}n hote
ouk \={e}n}] and the [\emph{ex ouk ont\={o}n}]. [On Origen's
disadvantages, vide Lumper Hist. t. x. p. 406, \&c.]} outside the Church on its first rise, and actually gave the
alarm, sixty years before Arius's day. Here let it suffice to set
down in his vindication the following facts, which may be left to the
consideration of the reader;—first, that his habitual hatred of
heresy and concern for heretics were such, as to lead him, even when
left an orphan in a stranger's house, to withdraw from the praying and
teaching of one of them, celebrated for his eloquence, who was in
favour with his patroness and other Christians of Alexandria; that all
through his long life he was known throughout Christendom as the
especial opponent of false doctrine, in its various shapes; and that
his pupils, Gregory, Athenodorus, and Dionysius, were principal actors
in the arraignment of Paulus, the historical forerunner of Arius;—next,
that his speculations, extravagant as they often were, related to
points not yet determined by the Church, and, consequently, were
really, what he frequently professed them to be, inquiries;—further,
that these speculations were for the most part ventured in matters of
inferior importance, certainly not upon the sacred doctrines which
Arius afterwards impugned, and in regard to which even his enemy
Jerome allows him to be orthodox;—that the opinions which brought
him into disrepute in his lifetime concerned the creation of the
world, the nature of the human soul, and the like;—that his
opinions, or rather speculations, on these subjects, were imprudently
made public by his friends;—that his writings were incorrectly
transcribed even in his lifetime, according to his own testimony;—that
after his death, Arian interpolations appear to have been made in some
of his works now lost, upon which the subsequent Catholic testimony of
his heterodoxy is grounded;—that, on the other hand, in his extant
works, the doctrine of the Trinity is cleanly avowed, and in
particular, our Lord's Divinity energetically and variously enforced;—and
lastly, that in matter of fact, the Arian party does not seem to have
claimed him, or appealed to him in self-defence, till thirty years
after the first rise of the heresy, when the originators of it were
already dead, although they had showed their inclination to shelter
themselves behind celebrated names, by the stress they laid on their
connexion with the martyr Lucian \footnote{Huet. Origen. lib. 1. lib. ii. 4. § 1.
Bull, Defens. F. N. ii. 9. Waterland's Works, vol. iii. p. 322. Baltus,
Défense des Ss. Pères, ii. 20. Tillemont, Mem. vol. iii. p. 259.
Socrat. Hist. iv. 26. Athanasius notices the change in the Arian
polemics, from mere disputation to an appeal to authority, in his De
Sent. Dionys. § 1, written about A.D. 354. [\emph{ouden
out' eulogon oute pros apodeixin ek t\={e}s theias graph\={e}s rh\={e}ton
echous\={e}s t\={e}s hairese\={o}s aut\={o}n, aei men
prophaseis anaischuntous eporizonto kai sophismata pithana; nun de kai
diaballein tous pateras tetolm\={e}kasi}].}. But if so much can be adduced in exculpation of Origen from
any grave charge of heterodoxy, what accusation can be successfully
maintained against his less suspected fellow-labourers in the
polemical school? so that, in concluding this part of the subject, we
may with full satisfaction adopt the judgment of Jerome:—``It
may be that they erred in simplicity, or that they wrote in another
sense, or that their writings were gradually corrupted by unskilful
transcribers; or certainly before Arius, like 'the sickness that
destroyeth in the noon-day,' was born in Alexandria, they made
statements innocently and incautiously, which are open to the
misinterpretation of the perverse.'' \footnote{Apolog. adv. Ruffin. ii. Oper. vol. ii. p.
149.}

\xsection{Section 4. The Eclectic Sect}

THE words of St. Jerome, with which the
last section closed, may perhaps suggest the suspicion, that the
Alexandrians, though orthodox themselves, yet incautiously prepared
the way for Arianism by the countenance they gave to the use of the
Platonic theological language. But, before speculating on the medium
of connexion between Platonism and Arianism, it would be well to
ascertain the existence of the connexion itself, which is very
doubtful, whether we look for it in history, or in the respective
characters of the parties professing the two doctrines; though it is
certain that Platonism, and Origenism also, became the excuse and
refuge of the heresy when it was condemned by the Church. I proceed to
give an account of the rise and genius of Eclecticism, with the view
of throwing light upon this question; that is, of showing its relation
both to the Alexandrian Church and to Arianism.

1.

The Eclectic philosophy is so called from its professing to select
the better parts of the systems invented before it, and to
digest these into one consistent doctrine. It is doubtful where the
principle of it originated, but it is probably to be ascribed to the
Alexandrian Jews. Certain it is, that the true faith never could come
into contact with the heathen philosophies, without exercising its
right to arbitrate between them, to protest against their vicious or
erroneous dogmas, and to extend its countenance to whatever bore an
exalted or a practical character. A cultivated taste would be likely
to produce among the heathen the same critical spirit which was
created by real religious knowledge; and accordingly we find in the
philosophers of the Augustan and the succeeding age, an approximation
to an eclectic or syncretistic system, similar to that which is found
in the writings of Philo. Some authors have even supposed, that Potamo,
the original projector of the school based on this principle,
flourished in the reign of Augustus; but this notion is untenable, and
we must refer him to the age of Severus, at the end of the second
century \footnote{Brucker, Hist. Phil. per. ii. part. 1, 2 § 4. [Vide Fabric.
Bibl. Græc. i. v. p. 680, ed. Harles.]}. In the mean
time, the Christians had continued to act upon the discriminative view
of heathen philosophy which the Philonists had opened; and, as we have
already seen, Clement, yet without allusion to particular sect or
theory, which did not exist till after his day, declares himself the
patron of the Eclectic principle. Thus we are introduced to the
history of the School which embodied it.

Ammonius, the contemporary of Potamo, and virtually the founder of
the Eclectic sect, was born of Christian parents, and educated
as a Christian in the catechetical institutions of Alexandria, under
the superintendence of Clement or Pantænus. After a time he
renounced, at least secretly, his belief in Christianity; and opening
a school of morals and theology on the stock of principles, esoteric
and exoteric, which he had learned in the Church, he became the
founder of a system really his own, but which by a dexterous artifice
he attributed to Plato. The philosophy thus introduced into the world
was forthwith patronized by the imperial court, both at Rome and in
the East, and spread itself in the course of years throughout the
empire, with bitter hostility and serious detriment to the interests
of true religion; till at length, obtaining in the person of Julian a
second apostate for its advocate, it became the authorized
interpretation and apology for the state polytheism. It is a
controverted point whether or not Ammonius actually separated from the
Church. His disciples affirm it; Eusebius, though not without some
immaterial confusion of statement, denies it \footnote{Euseb. Hist. Eccl. vi. 19.}. On the whole, it is probable that he began his teaching as a
Christian, and but gradually disclosed the systematic infidelity on
which it was grounded. We are told expressly that he bound his
disciples to secrecy, which was not broken till they in turn became
lecturers in Rome, and were led one by one to divulge the real
doctrines of their master \footnote{Bruckner, ibid.};
nor can we otherwise account for the fact of Origen having attended
him for a time, since he who refused to hear Paulus of Antioch, even
when dependent on the patroness of that heretic, would scarcely
have extended a voluntary countenance to a professed deserter from the
Christian faith and name.

This conclusion is confirmed by a consideration of the nature of
the error substituted by Ammonius for the orthodox belief; which was
in substance what in these times would be called \emph{Neologism}, a
heresy which, even more than others, has shown itself desirous and
able to conceal itself under the garb of sound religion, and to keep
the form, while it destroys the spirit, of Christianity. So close,
indeed, was the outward resemblance between Eclecticism and the Divine
system of which it was the deadly enemy, that St. Austin remarks, in
more than one passage, that the difference between the two professions
lay only in the varied acceptation of a few words and propositions \footnote{Mosheim, Diss. de Turb. per recent Plat.
Eccl. § 12.}. This peculiar character of the Eclectic philosophy must be
carefully noticed, for it exculpates the Catholic Fathers from being
really implicated in proceedings, of which at first they did not
discern the drift; while it explains that apparent connexion which, at
the distance of centuries, exists between them and the real originator
of it.

The essential mark of Neologism is the denial of the exclusive
divine mission and peculiar inspiration of the Scripture Prophets;
accompanied the while with a profession of general respect for them as
benefactors of mankind, as really instruments in God's hand, and as in
some sense the organs of His revelations; nay, in a fuller measure
such, than other religious and moral teachers. In its most specious
form, it holds whatever is good and true in the various
religions in the world, to have actually come from God: in its most
degraded, it accounts them all equally to be the result of mere human
benevolence and skill. In all its shapes, it differs from the orthodox
belief, primarily, in denying the miracles of Scripture to have taken
place, in the peculiar way therein represented, as distinctive marks
of God's presence accrediting the teaching of those who wrought them;
next, as a consequence, in denying this teaching, as preserved in
Scripture, to be in such sense the sole record of religious truth,
that all who hear it are bound to profess themselves disciples of it.
Its apparent connexion with Christianity lies (as St. Austin remarks)
in the ambiguous use of certain terms, such as \emph{divine, revelation,
inspiration}, and the like; which may with equal ease be made to
refer either to ordinary and merely providential, or to miraculous
appointments in the counsels of Almighty Wisdom. And these words would
be even more ambiguous than at the present day, in an age, when
Christians were ready to grant, that the heathen were in some sense
under a supernatural Dispensation, as was explained in the foregoing
section.

The rationalism of the Eclectics, though equally opposed with the
modern to the doctrine of the peculiar divinity of the Scripture
revelations, was circumstantially different from it. The Neologists of
the present day deny that the miracles took place in the manner
related in the sacred record; the Eclectics denied their cogency as an
evidence of the extraordinary presence of God. Instead of viewing them
as events of very rare occurrence, and permitted for important
objects in the course of God's providence, they considered them to be
common to every age and country, beyond the knowledge rather than the
power of ordinary men, attainable by submitting to the discipline of
certain mysterious rules, and the immediate work of beings far
inferior to the Supreme Governor of the world. It followed that, a
display of miraculous agency having no connexion with the truth of the
religious system which it accompanied, at least not more than any gift
merely human was connected with it, such as learning or talent, the
inquirer was at once thrown upon the examination of the doctrines for
the evidence of the divinity of Christianity; and there being no place
left for a claim on his allegiance to it as a whole, and for what is
strictly termed faith, he admitted or rejected as he chose, compared
and combined it with whatever was valuable elsewhere, and was at
liberty to propose to himself that philosopher for a presiding
authority, whom the Christians did but condescend to praise for his
approximation towards some of the truths which Revelation had
unfolded. The chapel of Alexander Severus was a fit emblem of that
system, which placed on a level Abraham, Orpheus, Pythagoras, and the
Sacred Name by which Christians are called. The zeal, the brotherly
love, the beneficence, and the wise discipline of the Church, are
applauded, and held up for imitation in the letters of the Emperor
Julian; who at another time calls the Almighty Guardian of the
Israelites a ``great God,'' \footnote{Gibbon, Hist. ch. xxiii.} while in common with his sect he professed to restore the
Christian doctrine of the Trinity to its ancient and pure
Platonic basis. It followed as a natural consequence, that the claims
of religion being no longer combined, defined, and embodied in a
personal Mediator between God and man, its various precepts were
dissipated back again and confused in the mass of human knowledge, as
before Christ came; and in its stead a mere intellectual literature
arose in the Eclectic School, and usurped the theological chair as an
interpreter of sacred duties, and the instructor of the inquiring
mind. ``In the religion which he (Julian) had adopted,'' says
Gibbon, ``piety and learning were almost synonymous; and a crowd
of poets, of rhetoricians, and of philosophers, hastened to the
Imperial Court, to occupy the vacant places of the bishops, who had
seduced the credulity of Constantius.'' \footnote{Ibid.} Who does not recognize in this old philosophy the chief
features of that recent school of liberalism and false illumination,
political and moral, which is now Satan's instrument in deluding the
nations, but which is worse and more earthly than it, inasmuch as his
former artifice, affecting a religious ceremonial, could not but leave
so much of substantial truth mixed in the system as to impress its
disciples with somewhat of a lofty and serious character, utterly
foreign to the cold, scoffing spirit of modern rationalism?

The freedom of the Alexandrian Christians from the Eclectic error
was shown above, when I was explaining the principles of their
teaching; a passage of Clement being cited, which clearly
distinguished between the ordinary and the miraculous appointments of
Providence. An examination of the dates of the history will show
that they could not do more than bear this indirect testimony against
it by anticipation. Clement himself was prior to the rise of
Eclecticism; Origen, prior to its public establishment as a sect.
Ammonius opened his school at the end of the second century, and
continued to preside in it at least till A.D.
243 \footnote{Fabric. Biblioth. Græc. Harles. iv. 29.}; during which period,
and probably for some years after his death, the real character of his
doctrines was carefully hidden from the world. He committed nothing to
writing, whether of his exoteric or esoteric philosophy, and when
Origen, who was scarcely his junior, attended him in his first years,
probably had not yet decidedly settled the form of his system.
Plotinus, the first promulgator and chief luminary of Eclecticism,
began his public lectures A.D. 244; and for some
time held himself bound by the promise of secrecy made to his master.
Moreover, he selected Rome as the seat of his labours, and there is
even proof that Origen and he never met. In Alexandria, on the
contrary, the infant philosophy languished; no teacher of note
succeeded to Ammonius; and even had it been otherwise, Origen had left
the city for ever, ten years previous to that philosopher's death. It
is clear, then, that he had no means of detecting the secret
infidelity of the Eclectics; and the proof of this is still stronger,
if, as Brucker calculates \footnote{Bruckner, ibid.},
Plotinus did not divulge his master's secret till A.D.
255, since Origen died A.D. 253. Yet, even in
this ignorance of the purpose of the Eclectics, we find Origen, in his
letter to Gregory expressing dissatisfaction at the actual
effects which had resulted to the Church from that literature in which
he himself was so eminently accomplished. ``For my part,'' he
says to Gregory, ``taught by experience, I will own to you, that
rare is the man, who, having accepted the precious things of Egypt,
leaves the country, and uses them in decorating the worship of God.
Most men who descend thither are brothers of Hadad (Jeroboam),
inventing heretical theories with heathen dexterity, and establishing
(so to say) calves of gold in Bethel, the house of God.'' \footnote{Orig. Ep. ad Gregor. § 2.} So much concerning Origen's ignorance of the Eclectic
philosophy. As to his pupils, Gregory and Dionysius, the latter, who
was Bishop of Alexandria, died A.D. 264;
Gregory, on the other hand, pronounced his panegyrical oration upon
Origen, in which his own attachment to heathen literature is avowed,
as early as A.D. 239; and besides, he had no
connexion whatever with Alexandria, having met with Origen at Cæsarea
\footnote{Tillemont, vol. iv. Chronolog.}. Moreover, just at
this time there were heresies actually spreading in the Church of an
opposite theological character, such as Paulianism; which withdrew
their attention from the prospect or actual rise of a Platonic
pseudo-theology; as will hereafter be shown.

Such, then, were the origin and principles of the Eclectic sect. It
was an excrescence of the school of Alexandria, but not attributable
to it, except as other heresies may be ascribed to other Churches,
which give them birth indeed, but cast them out and condemn them when
they become manifest. It went out from the Christians, but it
was not of them:—whether it resembled the Arians, on the other hand,
and what use its tenets were to them, are the next points to consider.

2.

The Arian school has already been attributed to Antioch as its
birth-place, and its character determined to be what we may call
Aristotelico-Judaic. Now, at very first sight, there are striking
points of difference between it and the Eclectics. On its Aristotelic
side, its disputatious temper was altogether uncongenial to the new
Platonists. These philosophers were commonly distinguished by their
melancholy temperament, which disposed them to mysticism, and often
urged them to eccentricities bordering on insanity \footnote{Brucker, supra.}. Far from cultivating the talents requisite for success in
life, they placed the sublimer virtues in an abstraction from sense,
and an indifference to ordinary duties. They believed that an
intercourse with the intelligences of the spiritual world could only
be effected by divesting themselves of their humanity; and that the
acquisition of miraculous gifts would compensate for their neglect of
rules necessary for the well-being of common mortals. In pursuit of
this hidden talent, Plotinus meditated a journey into India, after the
pattern of Apollonius; while bodily privations and magical rites were
methods prescribed in their philosophy for rising in the scale of
being. As might be expected from the professors of such a creed, the
science of argumentation was disdained, as beneath the regard of those
who were walking by an internal vision of the truth, not by the
calculations of a tedious and progressive reason; and was only
employed in condescending regard for such as were unable to rise to
their own level. When Iamblichus was foiled in argument by a
dialectician, he observed that the syllogisms of his sect were not
weapons which could be set before the many, being the energy of those
inward virtues which are the peculiar ornament of the philosopher.
Notions such as these, which have their measure of truth, if we
substitute for the unreal and almost passive illumination of the
mystics, that instinctive moral perception which the practice of
virtue ensures, found no sympathy in the shrewd secular policy and the
intriguing spirit of the Arians; nor again, in their sharp-witted
unimaginative cleverness, their precise and technical disputations,
their verbal distinctions, and their eager appeals to the judgment of
the populace, which is ever destitute of refinement and delicacy, and
has just enough acuteness of apprehension to be susceptible of
sophistical reasonings.

On the other hand, viewing the school of Antioch on its judaical
side, we are met by a different but not less remarkable contrast to
the Eclectics. These philosophers had followed the Alexandrians in
adopting the allegorical rule; both from its evident suitableness to
their mystical turn of mind, and as a means of obliterating the
scandals and reconciling the inconsistencies of the heathen mythology.
Judaism, on the contrary, being carnal in its views, was essentially
literal in its interpretations; and, in consequence, as hostile from
its grossness, as the Sophists from their dryness, to the fanciful
fastidiousness of the Eclectics. It had rejected the Messiah, because
He did not fulfil its hopes of a temporal conqueror and king. It
had clung to its obsolete ritual, as not discerning in it the
anticipation of better promises and commands, then fulfilled in the
Gospel. In the Christian Church, it was perpetuating the obstinacy of
its unbelief in a disparagement of Christ's spiritual authority, a
reliance on the externals of religious worship, and an indulgence in
worldly and sensual pleasures. Moreover, it had adopted in its most
odious form the doctrine of the Chiliasts or Millenarians, respecting
the reign of the saints upon earth, a doctrine which Origen, and
afterwards his pupil Dionysius, opposed on the basis of an allegorical
interpretation of Scripture \footnote{Mosh. de Rebus ante Const. Sæc. iii. c.
38.}.
And in this controversy, Judaism was still in connexion, more or less,
with the school of Antioch; which is celebrated in those times, in
contrast to the Alexandrian, for its adherence to the theory of the
literal sense \footnote{Conybeare, Bampt. Lect. iv. Orig. Opp. ed.
Benedict. vol. ii. præf.}.

It may be added, as drawing an additional distinction between the
Arians and the Eclectics, that while the latter maintained the
doctrine of Emanations, and of the eternity of matter, the hypothesis
of the former required or implied the rejection of both tenets; so
that the philosophy did not even furnish the argumentative foundation
of the heresy, to which its theology outwardly bore a partial
resemblance.

3.

But in seasons of difficulty men look about on all sides for
support; and Eclecticism, which had no attractions for the
Sophists of Antioch while their speculations were unknown to the world
at large, became a seasonable refuge (as we learn from various authors
\footnote{Vide Brucker, Hist. Phil. per. ii. part
ii. i. 2. § 8. Baltus, Défense des Pères, ii. 19.}), in the hands of
ingenious disputants, when pressed by the numbers and authority of the
defenders of orthodoxy. First, there was an agreement between the
Schools of Ammonius and of Paulus, in the cardinal point of an
inveterate opposition to the Catholic doctrine of our Lord's Divinity.
The judaizers admitted at most only His miraculous conception. The
Eclectics, honouring Him as a teacher of wisdom, still, far from
considering Him more than man, were active in preparing from the
heathen sages rival specimens of holiness and power. Next, the two
parties agreed in rejecting from their theology all mystery, in the
ecclesiastical notion of the word. The Trinitarian hypothesis of the
Eclectics was not perplexed by any portion of that difficulty of
statement which, in the true doctrine, results from the very
incomprehensibility of its subject. They declared their belief in a
sublime tenet, which Plato had first propounded and the Christians
corrupted; but their Three Divine Principles were in no sense one,
and, while essentially distinct from each other, there was a
successive subordination of nature in the second and the third \footnote{[\emph{archikai hypostaseis}]. Cudworth,
Intell. Syst. i. 4 § 36.}. In such speculations the judaizing Sophist found the very
desideratum which he in vain demanded of the Church; a scripturally-worded
creed, without its accompanying difficulty of conception. Accordingly, to the doctrine thus put into his hands he might appeal
by way of contrast, as fulfilling his just demands; nay, in proportion
as he out-argued and unsettled the faith of his Catholic opponent, so
did he open a way, as a matter of necessity and without formal effort,
for the perverted creed of that philosophy which had so mischievously
anticipated the labours and usurped the office of an ecclesiastical
Synod.

And, further, it must be observed, that, when the Sophist had
mastered the Eclectic theology, he had in fact a most powerful weapon
to mislead or to embarrass his Catholic antagonist. The doctrine which
Ammonius professed to discover in the Church, and to reclaim from the
Christians, was employed by the Arian as if the testimony of the early
Fathers to the truth of the heretical view which he was maintaining.
What was but incaution, or rather unavoidable liberty, in the
Ante-Nicene theology, was insisted on as apostolic truth. Clement and
Origen, already subjected to a perverse interpretation, were witnesses
provided by the Eclectics by anticipation against orthodoxy. This
express appeal to the Alexandrian writers, seems, in matter of fact,
to have been reserved for a late period of the controversy; but from
the first an advantage would accrue to the Arians, by their agreement
(as far as it went) with received language in the early Church.
Perplexity and doubt were thus necessarily introduced into the minds
of those who only heard the rumour of the discussion, and even of many
who witnessed it, and who, but for this apparent primitive sanction,
would have shrunk from the bold, irreverent inquiries and the idle
subtleties which are the tokens of the genuine Arian temper. Nor
was the allegorical principle of Eclecticism incompatible with the
instruments of the Sophist. This also in the hands of a dexterous
disputant, particularly in attack, would become more serviceable to
the heretical than to the orthodox cause. For, inasmuch as the Arian
controversialist professed to be asking for reasons why he should
believe our Lord's divinity, an answer based on allegorisms did not
silence him, while at the same time, it suggested to him the means of
thereby evading those more argumentative proofs of the Catholic
doctrine, which are built upon the explicit and literal testimonies of
Scripture. It was notoriously the artifice of Arius, which has been
since more boldly adopted by modern heretics, to explain away its
clearest declarations by a forced figurative exposition. Here that
peculiar subtlety in the use of language, in which his school
excelled, supported and extended the application of the allegorical
rule, recommended, as it was, to the unguarded believer, and forced
upon the more wary, by its previous reception on the part of the most
illustrious ornaments and truest champions of the Apostolic faith.

But after all there is no sufficient evidence in history that the
Arians did make this use of Neo-Platonism \footnote{There seems to have been a much earlier
coalition between the Platonic and Ebionitish doctrines, if the works
attributed to the Roman Clement may be taken in evidence of it.
Mosheim (de Turb. Eccl. § 34) says both the Recognitions and
Clementines are infected with the latter, and the Clementines with the
former doctrine. These works were written between A.D.
180 and A.D. 250: are they to be referred to the
school of Theodotus and Artemon, which was humanitarian and Roman,
expressly claimed the Bishops of Rome as countenancing its errors, and
falsified the Scriptures at least? Plotinus came to Rome A.D.
244, and Philostratus commenced his life of Apollonius there as early
as A.D. 217. This would account for the
Platonism of the latter of the two compositions, and its absence from
the earlier.}, considered as a party. I believe they did not, and from
the facts of the history should conclude Eusebius of Cæsarea alone to
be favourable to that philosophy: but some persons may attach
importance to the circumstance, that Syria was one of its chief seats
from its very first appearance. The virtuous and amiable Alexander
Severus openly professed its creed in his Syrian court, and in
consequence of this profession, extended his favour to the Jewish
nation. Zenobia, a Jewess in religion, succeeded Alexander in her
taste for heathen literature, and attachment to the syncretistic
philosophy. Her instructor in the Greek language, the celebrated
Longinus, had been the pupil of Ammonius, and was the early master of
Porphyry, the most bitter opponent of Christianity that issued from
the Eclectic school. Afterwards, Amelius, the friend and successor of
Plotinus, transferred the seat of the philosophy from Rome to Laodicea
in Syria; which became remarkable for the number and fame of its
Eclectics \footnote{Mosheim, Diss. de Turb. Eccl. § 11.}. In the next
century, Iamblicus and Libanius, the friend of Julian, both belonged
to the Syrian branch of the sect. It is remarkable that, in the mean
time, its Alexandrian branch declined in reputation on the death of
Ammonius; probably, in consequence of the hostility it met with from
the Church which had the misfortune to give it birth.

\xsection{Section 5. Sabellianism}

ONE subject more must be discussed in
illustration of the conduct of the Alexandrian school, and the
circumstances under which the Arian heresy rose and extended itself.
The Sabellianism which preceded it has often been considered the
occasion of it;—viz. by a natural reaction from one error into its
opposite; to separate the Father from the Son with the Arians, being
the contrary heresy to that of confusing them together after the
manner of the Sabellians. Here, however, Sabellianism shall be
considered neither as the proximate nor the remote cause, or even
occasion, of Arianism; but first, as drawing off the attention of the
Church from the prospective evil of the philosophical spirit; next, as
suggesting such reasonings, and naturalizing such expressions and
positions in the doctrinal statements of the orthodox, as seemed to
countenance the opposite error; lastly, as providing a sort of
justification of the Arians, when they first showed themselves;—that
is, Sabellianism is here regarded as facilitating rather than
originating the disturbances occasioned by the Arian heresy.

1.

The history of the heresy afterwards called Sabellian is
obscure. Its peculiar tenet is the denial of the distinction of
Persons in the Divine Nature; or the doctrine of the \emph{Monarchia},
as it is called by an assumption of exclusive orthodoxy, like that
which has led to the term ``Unitarianism'' at the present day
\footnote{Burton, Bampt. Lect. note 103. [The word [\emph{Monarchia}] was
adopted in opposition to the three [\emph{archikai hypostaseis}] of
the Eclectics; vide supra p. 112.]}. It was first maintained
as a characteristic of party by a school established (as it appears)
in Proconsular Asia, towards the end of the second century. This
school, of which Noetus was the most noted master, is supposed to be
an offshoot of the Gnostics; and doubtless it is historically
connected with branches of that numerous family. Irenæus is said to
have written against it; which either proves its antiquity, or seems
to imply its origination in those previous Gnostic systems, against
which his extant work is entirely directed \footnote{Dodwell in Iren. Diss. vi. 26.}. It may be added, that Simon Magus, the founder of the
Gnostics, certainly held a doctrine resembling that advocated by the
Sabellians.

At the end of the second century, Praxeas, a presbyter of Ephesus,
passed from the early school already mentioned to Rome. Meeting there
with that determined resistance which honourably distinguishes the
primitive Roman Church in its dealings with heresy, he retired into
Africa, and there, as founding no sect, he was soon forgotten.
However, the doubts and speculations which he had published,
concerning the great doctrine in dispute, remained alive in that part
of the world, though latent \footnote{Tertull. in Prax. [It is not certain Praxeas
was detected at Rome.]},
till they burst into a flame about the middle of the third
century, at the eventful era when the rudiments of Arianism were laid
by the sophistical school at Antioch.

The author of this new disturbance was Sabellius, from whom the
heresy has since taken its name. He was a bishop or presbyter in
Pentapolis, a district of Cyrenaica, included within the territory
afterwards called, and then virtually forming, the Alexandrian
Patriarchate. Other bishops in his neighbourhood adopting his
sentiments, his doctrine became so popular among a clergy already
prepared for it, or hitherto unpractised in the necessity of a close
adherence to the authorized formularies of faith, that in a short time
(to use the words of Athanasius) ``the Son of God was scarcely
preached in the Churches.'' Dionysius of Alexandria, as primate,
gave his judgment in writing; but being misunderstood by some orthodox
but over-zealous brethren, he in turn was accused by them, before the
Roman See, of advocating the opposite error, afterwards the Arian; and
in consequence, instead of checking the heresy, found himself involved
in a controversy in defence of his own opinions \footnote{Vide Athan. de Sent. Dionys.}. Nothing more is known concerning the Sabellians for above a
hundred years; when it is inferred from the fact that the Council of
Constantinople (A.D. 381) rejected their
baptism, that they formed at that time a communion distinct from the
Catholic Church.

Another school of heresy also denominated Sabellian, is obscurely
discernible even earlier than the Ephesian, among the Montanists of
Phrygia. The well-known doctrine of these fanatics, when adopted by minds less heated than its original propagators, evidently tended
to a denial of the Personality of the Holy Spirit. Montanus himself
probably was never capable of soberly reflecting on the meaning of his
own words; but even in his lifetime, Æschines, one of his disciples,
saw their real drift, and openly maintained the unreserved \emph{monarchia}
of the Divine Nature \footnote{Tillemont, Mem. vol. ii. p. 204.}.
Hence it is usual for ancient writers to class the Sabellians and
Montanists together, as if coinciding in their doctrinal views \footnote{Vales, ad Socr, i. 23. Soz. ii. 18.}. The success of Æschines in extending his heresy in Asia Minor
was considerable, if we may judge from the condition of that country
at a later period.—Gregory, the pupil of Origen, appears to have
made a successful stand against it in Pontus. Certainly his writings
were employed in the controversy after his death, and that with such
effect, as completely to banish it from that country, though an
attempt was made to revive it in the time of Basil (A.D.
375 \footnote{Basil, Epist. ccx. § 3.}).—In the
patriarchate of Antioch we first hear of it at the beginning of the
third century, Origen reclaiming from it Beryllus, Bishop of Bostra,
in Arabia. In the next generation the martyr Lucian is said to have
been a vigorous opponent of it; and he was at length betrayed to his
heathen persecutors by a Sabellian presbyter of the Church of Antioch.
At a considerably later date (A.D. 375) we hear
of it in Mesopotamia \footnote{Epiphan. Hær. lxii. 1.}.

At first sight it may seem an assumption to refer these various
exhibitions of heterodoxy in Asia Minor, and the East, to some
one school or system, merely on the ground of their distinguishing
tenet being substantially the same. And certainly, in treating an
obscure subject, on which the opinions of learned men differ, it must
be owned that conjecture is the utmost that I am able to offer. The
following statement will at once supply the grounds on which the above
arrangement has been made, and explain the real nature of the doctrine
itself in which the heresy consisted \footnote{[Vid. Athan. Transl. vol. ii. p. 377.]}.

Let it be considered then, whether there were not two kinds of
Sabellianism; the one taught by Praxeas, the other somewhat
resembling, though less material than, the theology of the Gnostics:—the
latter being a modification of the former, arising from the pressure
of the controversy: for instance, parallel to the change which is said
to have taken place in the doctrine of the Ebionites, and in that of
the followers of Paulus of Samosata. Those who denied the distinction
of Persons in the Divine Nature were met by the obvious inquiry, in
what sense they believed God to be united to the human nature of
Christ. The more orthodox, but the more assailable answer to this
question, was to confess that God was, in such sense, one Person with
Christ, as (on their Monarchistic principle) to be in no sense
distinct from Him. This was the more orthodox answer, as preserving
inviolate what is theologically called the doctrine of the hypostatic
union,—the only safeguard against a gradual declension into the
Ebionite, or modern Socinian heresy. But at the same time such an
answer was repugnant to the plainest suggestions of scripturally-enlightened
reason, which leads us to be sure that, according to the obvious
meaning of the inspired text, there is \emph{some} real sense in which
the Father is not the Son; that the Sender and the Sent cannot be in
all respects the same; nor can the Son be said to make Himself
inferior to the Father, and condescend to become man,—to come from
God, and then again to return to Him,—if, after all, there is no
distinction beyond that of words, between those Blessed and Adorable
Agents in the scheme of our redemption. Besides, without venturing to
intrude into things not as yet seen, it appeared evident to the
primitive Church, that, in matter of fact, the Son of God, though
equal in dignity of nature to the Father, and One with Him in essence,
was described in Scripture as undertaking such offices of ministration
and subjection, as are never ascribed, and therefore may not without
blasphemy be ascribed, to the self-existent Father. Accordingly, the
name of Patripassian was affixed to Praxeas, Noetus, and their
followers, in memorial of the unscriptural tenet which was immediately
involved in their denial of the distinction of Persons in the Godhead.

Such doubtless was the doctrine of Sabellius, if regard be paid to
the express declarations of the Fathers. The discriminating Athanasius
plainly affirms it, in his defence of Dionysius \footnote{De Sent. Dionys. § 5, 9, \&c. [Orat.
iii. 36. Origen, in Ep. ad. Tit. t. iv. p. 695: ``Duos definimus,
ne (ut vestra perversitas infert) Pater ipse credatur natus et passus.''
Tertull. adv. Prax. 13.]}. The Semi-Arian Creed called the Macrostich, published at
Antioch, gives a like testimony \footnote{Athan. de Synod. § 26.}; distinguishing, moreover, between the Sabellian
doctrine, and the doctrines of the Paulianists and Photinians, to
which some modern critics have compared it. Cyprian and Austin, living
in Africa, bear express witness to the existence of the Patripassian
sect \footnote{Cyprian, Epist. lxxiii. Tillemont, Mem.
iv. 100.}. On the other
hand, it cannot be denied, that authorities exist favourable to a view
of the doctrine different from the above, and these accordingly may
lead us, in agreement with certain theological writers \footnote{Beausobre, Hist. de Manich. iii. 6. § 7.
Mosheim, de Reb. ant. Const. sæc. ii. § 68; sæc. iii. § 32.
Lardner, Cred. part ii. ch. 41.}, without interfering with the account of the heresy already
given, to describe a modification of it which commonly succeeded to
its primitive form.

The following apparently inconsistent testimonies, suggest both the
history and the doctrine of this second form of Sabellianism. While
the Montanists and Sabellians are classed together by some authors,
there is separate evidence of the connexion of each of these with the
Gnostics. Again, Ambrosius, the convert and friend of Origen was
originally a Valentinian, or Marcionite, or Sabellian, according to
different writers. Further, the doctrine of Sabellius is compared to
that of Valentinus by Alexander of Alexandria, and (apparently) by a
Roman Council (A.D. 324); and by St. Austin it
is referred indifferently to Praxeas, or to Hermogenes, a Gnostic. On
the other hand, one Leucius is described as a Gnostic and Montanist \footnote{Vide Tillemont, vol. ii. p. 204; iv. p.
100, \&c. Waterland's Works, vol. i. p. 236, 237.}. It would appear then, that it is so repugnant to the plain
word of Scripture, and to the most elementary notions of
doctrine thence derived, to suppose that Almighty God is in every
sense one with the human nature of Christ, that a disputant,
especially an innovator, cannot long maintain such a position. It
removes the mystery of the Trinity, only by leaving the doctrine of
the Incarnation in a form still more strange, than that which it
unavoidably presents to the imagination. Pressed, accordingly, by the
authority of Scripture, the Sabellian, instead of speaking of the
substantial union of God with Christ, would probably begin to obscure
his meaning in the decorum of a figurative language. He would speak of
the presence rather than the existence of God in His chosen servant;
and this presence, if allowed to declaim, he would represent as a
certain power or emanation from the Centre of light and truth; if
forced by his opponent into a definite statement, he would own to be
but an inspiration, the same in kind, though superior in degree, with
that which enlightened and guided the prophets. This is that second
form of the Sabellian tenet, which some learned moderns have
illustrated, though they must be considered to err in pronouncing it
the only true one. That it should have resulted from the difficulties
of the Patripassian creed, is natural and almost necessary; and viewed
merely as a conjecture, the above account of its rise reconciles the
discordant testimonies of ecclesiastical history. But we have almost
certain evidence of the matter of fact in Tertullian's tract against
Praxeas \footnote{In Prax. § 27.}, in which the
latter is apparently represented as holding successively, the two
views of doctrine which have been here described. Parallel instances
meet us in the history of the Gnostics and Montanists. Simon
Magus, for instance, seems to have adopted the Patripassian theory.
But the Gnostic family which branched from him, modified it by means
of their doctrine of emanations or æons, till in the theology of
Valentinus, as in that of Cerinthus and Ebion, the incarnation of the
Word, became scarcely more than the display of Divine Power with a
figurative personality in the life and actions of a mere man. The
Montanists, in like manner, from a virtual assumption of the Divinity
of their founder, were led on, as the only way of extricating
themselves from one blasphemy, into that other of denying the
Personality of the Holy Spirit, and then of the Word. Whether the
school of Noetus maintained its first position, we have no means of
knowing; but the change to the second, or semi-humanitarian, may be
detected in the Sabellians, as in Praxeas before them. In the time of
Dionysius of Alexandria, the majority was Patripassian; but in the
time of Alexander they advocated the Emanative, as it may be called,
or in-dwelling theory \footnote{Theod. Hist. i. 4.}.

2.

What there is further to be said on this subject shall be reserved
for the next chapter. Here, however, it is necessary to examine, how,
under these circumstances, the controversy with the Sabellians would
affect the language of ecclesiastical theology. It will be readily
seen, that the line of argument by which the two errors above
specified are to be met, is nearly the same: viz. that of insisting
upon the personality of the Word as distinct from the Father.
For the Patripassian denied that the Word was in any real respect
distinct from Him; the Emanatist, if he may so be called, denied that
He was a Person, or more than an extraordinary manifestation of Divine
Power. The Catholics, on the other hand, asserted His distinct
personality; and necessarily appealed, in proof of this, to such texts
as speak of His pre-existent relations towards the Father; in other
words, His ministrative office in the revealed Economy of the Godhead.
And thus, being obliged from the course of the controversy, to dwell
on this truly scriptural tenet, and happening to do so without a
protest against a denial, as if involved in it, of His equality with
the Father in the One Indivisible Divine Nature (a protest, which
nothing but the actual experience of that denial among them could
render necessary or natural), they were sometimes forced by the
circumstances of the case into an apparent anticipation of the heresy,
which afterwards arose in the shape of Arianism.

This may be illustrated in the history of the two great pupils of
Origen, who, being respectively opposed to the two varieties of
Sabellianism above described, the Patripassian and the Emanative,
incurred odium in a later age, as if they had been forerunners of
Arius: Gregory of Neocæsarea, and Dionysius of Alexandria.

The controversy in which Dionysius was engaged with the
Patripassians of Pentapolis has already been adverted to. Their tenet
of the incarnation of the Father (that is, of the one God without
distinction of Persons), a tenet most repugnant to every scripturally-informed
mind, was refuted at once, by insisting on the essential
character of the Son as representing and revealing the Father; by
arguing, that on the very face of Scripture, the Christ who is there
set before us, (whatever might be the mystery of His nature,) is
certainly delineated as one absolute and real Person, complete in
Himself, sent by the Father, doing His will, and mediating between Him
and man; and that, this being the case, His Person could not be the
same with that of the Father, who sent Him, by any process of
reasoning, which would not also prove any two individual men to have
one literal personality; that is, if there be any analogy at all
between the ordinary sense of the word ``person'' and that in
which the idea is applied in Scripture to the Father and the Son: for
instance, by what artifice of interpretation can the beginning of St.
John's Gospel, or the second chapter of St. Paul's Epistle to the
Philippians be made to harmonize with the notion, that the one God,
simply became and is man, in every sense in which He can still be
spoken of as God?

Writing zealously and freely on this side of the Catholic doctrine,
Dionysius laid himself open to the animadversion of timid and
narrow-minded men, who were unwilling to receive the truth in that
depth and fulness in which Scripture reveals it, and who thought that
orthodoxy consisted in being at all times careful to comprehend in one
phrase or formula the whole of what is believed on any article of
faith. The Roman Church, even then celebrated for its vigilant,
perhaps its over earnest exactness, in matters of doctrine and
discipline, was made the arbiter of the controversy. A council was
held under the presidency of Dionysius its Bishop (about A.D.
260), in which the Alexandrian prelate was accused by the
Pentapolitans of asserting that the Son of God is made and created,
distinct in nature from the incommunicable essence of the Father,
``as the vine is distinct from the vine-dresser,'' and in
consequence, not eternal. The illustration imputed to Dionysius in
this accusation, being a reference to our Lord's words in the
fifteenth chapter of St. John, is a sufficient explanation by itself
of the real drift of his statement, even if his satisfactory answer
were not extant, to set at rest all doubt concerning his orthodoxy. In
that answer, addressed to his namesake of Rome, he observes first,
that his letter to the Sabellians, being directed against a particular
error, of course contained only so much of the entire Catholic
doctrine as was necessary for the refutation of that error;—that his
use of the words ``Father and Son,'' in itself implied his
belief in a oneness of nature between Them;—that in speaking of the
Son as ``made,'' he had no intention of distinguishing
``made'' from ``begotten,'' but, including all kinds
of origination under the term, he used it to discriminate between the
Son and His underived self-originating Father;—lastly, that in
matter of fact he did confess the Catholic doctrine in its most
unqualified and literal sense, and in its fullest and most accurate
exposition. In this letter he even recognizes the celebrated \emph{Homoüsion}
(\emph{consubstantial}) which was afterwards adopted at Nicæa.
However, in spite of these avowals, later writers, and even Basil
himself, do not scruple to complain of Dionysius as having sown the
first seeds of Arianism; Basil confessing the while that his error was
accidental, occasioned by his vehement opposition to the Sabellian
heresy.

Gregory of Neocæsarea, on the other hand, is so far more hardly
circumstanced than Dionysius, first, inasmuch as the charge against
him was not made till after his death, and next, because he is
strangely accused of a tendency to Sabellian as well as Arian errors.
Without accounting for the former of these charges, which does not now
concern us, I offer to the reader the following explanation of the
latter calumny. Sabellianism, in its second or emanative form, had
considerable success in the East before and at the date of Gregory. In
the generation before him, Hermogenes, who professed it, had been
refuted by Theophilus and Tertullian, as well as by Gregory's master,
Origen, who had also reclaimed from a similar error Ambrosius and
Beryllus \footnote{Euseb. Hist. iv. 24. Theod. Hær. i. 19.
Tertull. in Hermog. Huet. Origen, lib. i.}. Gregory
succeeded him in the controversy with such vigour, that his writings
were sufficient to extinguish the heresy, when it reappeared in Pontus
at a later period. He was, moreover, the principal bishop in the first
Council held against Paulus of Samosata, whose heresy was derived from
the emanative school. The Synodal Letter addressed by the assembled
bishops to the heresiarch, whether we ascribe it to this first
Council, with some critics, or with others to the second, or even with
Basnage reject it as spurious, at least illustrates the line of
argument which it was natural to direct against the heresy, and shows
how easily it might be corrupted into an Arian meaning. To the notion
that the Son was but inhabited by a divine power or presence
impersonal, and therefore had no real existence before He came in the
flesh, it was a sufficient answer to appeal to the great works
ascribed to Him in the beginning of all things, and especially to
those angelic manifestations by which God revealed Himself to the
elder Church, and which were universally admitted to be
representations of the Living and Personal Word. The Synodal Letter
accordingly professes a belief in the Son, as the Image and Power of
God, which was before the worlds, in absolute existence, the living
and intelligent Cause of creation; and cites some of the most striking
texts descriptive of His ministrative office under the Jewish law,
such as His appearance to Abraham and Jacob, and to Moses in the
burning bush \footnote{Routh, Reliq. Sacr. vol. ii. p. 463.}. Such is
the statement, in opposition to Paulus of Samosata, put forth by
Gregory and his associate bishops at Antioch; and, the circumstances
of the controversy being overlooked, it is obvious how easily it may
be brought to favour the hypothesis, that the Son is in all respects
distinct from the Father, and by nature as well as in revealed office
inferior to Him.

Lastly, it so happened, that in the course of the third century,
the word \emph{Homoüsion} became more or less connected with the
Gnostic, Manichæan, and Sabellian theologies. Hence early writers,
who had but opposed these heresies, seemed in a subsequent age to have
opposed what had been by that time received as the characteristic of
orthodoxy; as, on the other hand, the Catholics, on their adopting it
in that later age, were accused of what in an earlier time would have
been the Sabellian error, or again of the introduction of corporeal
notions into their creed. But of this more hereafter.

Here a close may be put to our inquiry into the circumstances under
which Arianism appeared in the early Church. The utmost that has been
proposed has been to classify and arrange phenomena which present
themselves on the surface of the history; and this, with a view of
preparing the reader for the direct discussion of the doctrine which
Arianism denied, and for the proceedings on the part of the Church
which that denial occasioned. Especially has it been my object in this
introduction, following the steps of our great divines, to rescue the
Alexandrian Fathers from the calumnies which, with bad intentions
either to them or to the orthodox cause, have been so freely and so
fearlessly cast upon them. Whether Judaism or whether Platonism had
more or less to do in preparing the way for the Arian heresy, are
points of minor importance, compared with the vindication of those
venerable men, the most learned, most eloquent, and most zealous of
the Ante-Nicene Christians. With this view it has been shown above,
that, though the heresy openly commenced, it but accidentally
commenced in Alexandria; that no Alexandrian of name advocated it, and
that, on its appearance, it was forthwith expelled from the
Alexandrian Church, together with its author \footnote{[Vid. Athan. Apol. adv. Arian. 52, and
Hist. Arian. 78 fin.]};—next, that, even granting Platonism originated it, of which
there is no proof, still there are no grounds for implicating the
Alexandrian Fathers in its formation; that while the old Platonism,
which they did favour, had no part in the origination of the Arian
doctrine, the new Platonism or Eclecticism which may be conceived to
have arianized, received no countenance from them; that Eclecticism must abstractedly be referred to their schools, it arose
out of them in no more exact sense than error ever springs from truth;
that, instead of being welcomed by them, the sight of it, as soon as
it was detected, led them rather to condemn their own older and
innocent philosophy; and that, in Alexandria, there was no Eclectic
successor to Ammonius (who concealed his infidelity to the last), till
after the commencement of the Arian troubles;—further, that granting
(what is undeniable) that the Alexandrian Fathers sometimes use
phrases which are similar to those afterwards adopted by the heretics,
these were accidents, not the characteristics of their creed, and were
employed from a studied verbal imitation of the Jewish and
philosophical systems;—of the philosophical, in order to conceal
their own depth of meaning, and to conciliate the heathen, a duty to
which their peculiar functions in the Christian world especially bound
them, and of the Jewish, from an affectionate reverence for the early
traces, in the Old Testament, of God's long-meditated scheme of mercy
to mankind;—or again, that where they seem to arianize, it is from
incompleteness rather than from unsoundness in their confessions,
occasioned by the necessity of opposing a contrary error then
infecting the Church; that five Fathers, who have more especially
incurred the charge of philosophizing in their creed, belong to the
schools of Rome and Antioch, as well as of Alexandria, and that the
most unguarded speculator in the Alexandrian, Origen, is the very
writer first to detect for us, and to denounce the Arian tenet, at
least sixty years before it openly presented itself to the world.

On the other hand, if, dismissing this side of the question,
we ask whence the heresy actually arose, we find that contemporary
authors ascribe it partially to Judaism and Eclecticism, and more
expressly to the influence of the Sophists; that Alexander, to whose
lot it fell first to withstand it, refers us at once to Antioch as its
original seat, to Judaism as its ultimate source, and to the
subtleties of disputation as the instrument of its exhibition: that
Arius and his principal supporters were pupils of the school of
Antioch; and lastly, that in this school at the date fixed by
Alexander, the above-mentioned elements of the heresy are discovered
in alliance, almost in union, Paulus of Samosata, the judaizing
Sophist, being the favourite of a court which patronized Eclecticism,
when it was neglected at Alexandria.

It is evident that deeper and more interesting questions remain,
than any which have here been examined. The real secret causes of the
heresy; its connexion with the character of the age, with the opinions
then afloat, viewed as active moral influences, not as parts of a
system; its position in the general course of God's providential
dealings with His Church, and in the prophecies of the New Testament;
and its relation towards the subsequently developed corruptions of
Christianity; these are subjects towards which some opening may have
been incidentally made for inquirers, but which are too large to be
imagined in the design of a work such as the present.

\xchapter{Chapter 2. The Teaching of the Ante-Nicene Church in its relation to the Arian Heresy}

\xsection{Section 1. On the Principle of the Formation and Imposition of Creeds}

IT has appeared in the foregoing
Chapter, that the temper of the Ante-Nicene Church was opposed to the
imposition of doctrinal tests upon her members; and on the other hand,
that such a measure became necessary in proportion as the cogency of
Apostolic Tradition was weakened by lapse of time. This is a subject
which will bear some further remarks; and will lead to an
investigation of the principle upon which the formation and imposition
of creeds rests. After this, I shall delineate the Catholic doctrine
itself, as held in the first ages of Christianity; and then, the Arian
substitution for it.

1.

I have already observed, that the knowledge of the Christian
mysteries was, in those times, accounted as a privilege, to be eagerly
coveted. It was not likely, then, that reception of them would
be accounted a test; which implies a concession on the part of the
recipient, not an advantage. The idea of disbelieving, or criticizing
the great doctrines of the faith, from the nature of the case, would
scarcely occur to the primitive Christians. These doctrines were the
subject of an Apostolical Tradition; they were the very truths which
had been lately revealed to mankind. They had been committed to the
Church's keeping, and were dispensed by her to those who sought them,
as a favour. They were facts, not opinions. To come to the Church was
all one with expressing a readiness to receive her teaching; to
hesitate to believe, after coming for the sake of believing, would be
an inconsistency too rare to require a special provision against the
chance of it \footnote{[Hoc penitus absurdum est, ut discipulus, ad magistrum vadens,
ante sit artifex quam doceatur, \&c. Hieron. adv. Lucif. 12.]}. It was
sufficient to meet the evil as it arose: the power of excommunication
and deposition was in the hands of the ecclesiastical authorities,
and, as in the case of Paulus, was used impartially. Yet, in the
matter of fact, such instances of contumacy were comparatively rare;
and the Ante-Nicene heresies were in many instances the innovations of
those who had never been in the Church, or who had already been
expelled from it.

We have some difficulty in putting ourselves into the situation of
Christians in those times, from the circumstance that the Holy
Scriptures are now our sole means of satisfying ourselves on points of
doctrine. Thus, every one who comes to the Church considers himself
entitled to judge and decide individually upon its creed. But in that
primitive age, the Apostolical Tradition, that is, the Creed,
was practically the chief source of instruction, especially
considering the obscurities of Scripture; and being withdrawn from
public view, it could not be subjected to the degradation of a
comparison, on the part of inquirers and half-Christians, with those
written documents which are vouchsafed to us from the same inspired
authorities. As for the baptized and incorporate members of the
Church, they of course had the privilege of comparing the written and
the oral tradition, and might exercise it as profitably as in
comparing and harmonizing Scripture with itself. But before baptism,
the systematic knowledge was withheld; and without it, Scripture,
instead of being the source of instruction on the doctrines of the
Trinity and Incarnation, was scarcely more than a sealed book, needing
an interpretation, amply and powerfully as it served the purpose of
proving those doctrines, when they were once disclosed. And so much on
the reluctance of the primitive Fathers to publish creeds, on the
ground that the knowledge of Christian doctrines was a privilege
reserved for those who were baptized, and in no sense a subject of
hesitation and dispute.—It may be added, that the very love of
power, which in every age will sway the bulk of those who are exposed
to the temptation of it, and ecclesiastics in the number, would
indispose them to innovate upon a principle which made themselves the
especial guardians of revealed truth \footnote{Vide Hawkins on Unauthoritative Tradition.}.

Their backwardness proceeded also from a profound reverence for the
sacred mysteries of which they were the dispensers. Here they present
us with the true exhibition of that pious sensitiveness which
the heathen had conceived, but could not justly execute. The latter
had their mysteries, but their rude attempts were superseded by the
divine discipline of the Gospel, which here acted in the office which
is peculiarly its own, rectifying, combining, and completing the
inventions of uninstructed nature. If the early Church regarded the
very knowledge of the truth as a fearful privilege, much more did it
regard that truth itself as glorious and awful; and scarcely
conversing about it to her children, shrank from the impiety of
subjecting it to the hard gaze of the multitude \footnote{Sozomen gives this reason for not inserting
the Nicene Creed in his history: ``I formerly deemed it necessary
to transcribe the confession of faith drawn up by the unanimous
consent of this Council [the Nicene], in order that posterity might
possess a public record of the truth; but subsequently I was persuaded
to the contrary by some godly and learned friends, who represented
that such matters ought to be kept secret, as being only requisite to
be known by disciples and their instructors [\emph{mystais kai mystag\={o}gois}],
and it is possible that the volume will fall into the hands of the
unlearned [\emph{t\={o}n amu\={e}t\={o}n}]. Hist. i. 20. Bohn's
translation.}. We still pray, in the Confirmation service, for those who are
introduced into the full privileges of the Christian covenant, that
they may be ``filled with the spirit of God's holy fear;'' but
the meaning and practical results of deep-seated religious reverence
were far better understood in the primitive times than now, when the
infidelity of the world has corrupted the Church. Now, we allow
ourselves publicly to canvass the most solemn truths in a careless or
fiercely argumentative way; truths, which it is as useless as it is
unseemly to discuss in public, as being attainable only by the sober
and watchful, by slow degrees, with dependence on the Giver of
wisdom, and with strict obedience to the light which has already been
granted. Then, they would scarcely express in writing, what is now not
only preached to the mixed crowds who frequent our churches, but
circulated in print among all ranks and classes of the unclean and the
profane, and pressed upon all who choose to purchase it. Nay, so
perplexed is the present state of things, that the Church is obliged
to change her course of acting, after the spirit of the alteration
made at Nicæa, and unwillingly to take part in the theological
discussions of the day, as a man crushes venomous creatures of
necessity, powerful to do it, but loathing the employment. This is the
apology which the author of the present work, as far as it is worth
while to introduce himself, offers to all sober-minded and zealous
Christians, for venturing to exhibit publicly the great evangelical
doctrines, not indeed in the medium of controversy or proof (which
would be a still more humiliating office), but in an historical and
explanatory form. And he earnestly trusts, that, while doing so, he
may be betrayed into no familiarity or extravagance of expression,
cautiously lowering the Truth, and (as it were), wrapping it in
reverent language, and so depositing it in its due resting-place,
which is the Christian's heart: guiltless of those unutterable
profanations with which a scrutinizing infidelity wounds and lacerates
it. Here, again, is strikingly instanced the unfitness of books,
compared with private communication, for the purposes of religious
instruction; levelling, as they do, the distinctions of mind and
temper by the formality of the written character, and conveying each
kind of knowledge the less perfectly, in proportion as it is of
a moral nature, and requires to be treated with delicacy and
discrimination.

2.

As to the primitive Fathers, with their reverential feelings
towards the Supreme Being, great must have been their indignation
first, and then their perplexity, when apostates disclosed and
corrupted the sacred truth, or when the heretical or philosophical
sects made guesses approximating to it. Though the heretics also had
their mysteries, yet, it is remarkable, that as regards the high
doctrines of the Gospel, they in great measure dropped that restraint
and reserve by which the Catholics partly signified, and partly
secured a reverence for them. Tertullian sharply exposes the want of a
grave and orderly discipline among them in his day. ``It is
uncertain,'' he says, ``who among them is catechumen, who
believer. They meet alike, they hear alike, they pray alike; nay,
though the heathen should drop in, they will cast holy things to dogs,
and their pearls, false jewels as they are, to swine. This overthrow
of order they call simplicity, and our attention to it they call
meretricious embellishment. They communicate with all men
promiscuously; it being nothing to them in what they differ from them,
provided they join with them for the destruction of the truth. They
are all high-minded; all make pretence of knowledge. Their catechumens
are perfect in the faith before they are fully taught. Even their
women are singularly forward; venturing, that is, to teach, to argue,
to exorcise, to undertake cures, nay, perhaps to baptise.'' \footnote{Tertull. de Præscr. hæret. § 41.}

The heretical spirit is ever one and the same in its various forms:
this description of the Gnostics was exactly paralleled, in all those
points for which we have introduced it here, in the history of
Arianism; historically distinct as is the latter system from
Gnosticism. Arius began by throwing out his questions as a subject of
debate for public consideration; and at once formed crowds of
controversialists from those classes who were the least qualified or
deserving to take part in the discussion. Alexander, his diocesan,
accuses him of siding with the Jews and heathen against the Church;
and certainly we learn from the historians, that the heathen
philosophers were from the first warmly interested in the dispute, so
that some of them attended the Nicene Council, for the chance of
ascertaining the orthodox doctrine. Alexander also charges him with
employing women in his disturbance of the Church, apparently referring
at the same time to the Apostle's prediction of them. He speaks
especially of the younger females as zealous in his cause, and as
traversing Alexandria in their eagerness to promote it;—a fact
confirmed by Epiphanius, who speaks (if he may be credited) of as many
as seven hundred from the religious societies of that city at once
taking part with the heresiarch \footnote{Soc. i. 6. Theod. Hist. i. iv. Soz. i. 18.
Epiph. hær. lxix. 3.}. But Arius carried his agitation lower still. It is on no other
authority than that of the historian Philostorgius, his own partisan,
that we are assured of his composing and setting to music, songs on
the subject of his doctrine for the use of the rudest classes of
society, with a view of familiarizing them to it. Other of his
compositions, of a higher literary excellence, were used at
table as a religious accompaniment to the ordinary meal; one of which,
in part preserved by Athanasius, enters upon the most sacred portions
of the theological question \footnote{Philost. ii. 2. Athan. in Arian. i. 5.; de
Syn. 15.}.
The success of these exertions in drawing public attention to his
doctrine is recorded by Eusebius of Cæsarea, who, though no friend of
the heresiarch himself, is unsuspicious evidence as being one of his
party. ``From a little spark a great fire was kindled. The quarrel
began in the Alexandrian Church, then it spread through the whole of
Egypt, Lybia, and the farther Thebais; then it ravaged the other
provinces and cities, till the war of words enlisted not only the
prelates of the churches, but the people too. At length the exposure
was so extraordinary, that even in the heathen theatres, the divine
doctrine became the subject of the vilest ridicule.'' \footnote{Euseb. Vit. Const. ii. 61. Vid. Greg. Naz.
Orat. i. 142; [ii. 81, 82.]} Such was Arianism at its commencement; and if it was so
indecent in the hands of its originator, who, in spite of his courting
the multitude, was distinguished by a certain reserve and loftiness in
his personal deportment, much more flagrant was its impiety under the
direction of his less refined successors. Valens, the favourite bishop
of Constantius, exposed the solemnities of the Eucharist in a judicial
examination to which Jews and heathen were admitted; Eudoxius, the
Arianizer of the Gothic nations, when installed in the patriarchal
throne of Constantinople, uttered as his first words a profane jest,
which was received with loud laughter in the newly-consecrated Church
of St. Sophia; and Aetius, the founder of the Anomœans, was the
grossest and most despicable of buffoons \footnote{Athan. Apol. contr. Arian. 31. Socr. ii. 43.
Cave, Hist. Literar. vol. i. [Eustathius speaks of the [\emph{paradoxoi t\={e}s
Areiou thumel\={e}s mesochoroi}]. Phot. Bibl. p. 759. 30.]}. Later still, we find the same description of the heretical
party in a discourse of the kind and amiable Gregory of Nazianzus.
With a reference to the Arian troubles he says, ``Now is priest an
empty name; contempt is poured upon the rulers, as Scripture says ...
All fear is banished from our souls, shamelessness has taken its
place. Knowledge is now at the will of him who chooses it, and all the
deep mysteries of the Spirit. We are all pious, because we condemn the
impiety of others. We use the infidels as our arbiters, and cast what
is holy to dogs, and pearls before swine, publishing divine truths to
profane ears and minds; and, wretches as we are, we carefully fulfil
the wishes of our enemies, while, without blushing, we 'pollute
ourselves in our inventions.''' \footnote{Greg. Naz. Orat., i. 135; [ii. 79.]}

Enough has now been said, by way of describing the condition of the
Catholic Church, defenceless from the very sacredness and refinement
of its discipline, when the attack of Arianism was made upon it;
insulting its silence, provoking it to argue, unsettling and seducing
its members \footnote{[``Is it not enough to distract a man,
on mere hearing, though unable to controvert, and to make him stop his
ears, from astonishment at the novelty of what he hears said, which
even to mention is to blaspheme?'' Ath. Orat. i. 35. Hence, as if
feeling the matter to be beyond argument, Athanasius could but call
the innovators ``Ariomaniacs,'' from the fierceness of their
``ipse dixit.'' Vid. Athan. Transl. vol. ii. p. 377.]}, and in
consequence requiring its authoritative judgment on the point in
dispute. And in addition to the instruments of evil which were
internally directed against it, the Eclectics had by this time
extended their creed among the learned, with far greater decorum than
the Arians, but still so as practically to interpret the Scriptures in
the place of the Church, and to state dogmatically the conclusions for
which the Arian controvertists were but indirectly preparing the mind
by their objections and sophisms.

3.

Under these circumstances, it was the duty of the rulers of the
Church, at whatever sacrifice of their feelings, to discuss the
subject in controversy fully and unreservedly, and to state their
decision openly. The only alternative was an unmanly non-interference,
and an arbitrary or treacherous prohibition of the discussion. To
enjoin silence on perplexed inquirers, is not to silence their
thoughts; and in the case of serious minds, it is but natural to turn
to the spiritual ruler for advice and relief, and to feel
disappointment at the timidity, or irritation at the harshness, of
those who refuse to lead a lawful inquiry which they cannot stifle \footnote{[\emph{kindunos gar prodosias, en t\={o}i
m\={e} procheir\={o}s apodidonai tas peri theou apokriseis tois
agap\={o}si ton kurion}]. Basil, Ep. 7. Vide Hil. de Trin.
xii. 20.}. Such a course, then, is most unwise as well as cruel,
inasmuch as it throws the question in dispute upon other arbitrators;
or rather, it is more commonly insincere, the traitorous act of those
who care little for the question in dispute, and are content that
opinions should secretly prevail which they profess to condemn. The
Nicene Fathers might despair of reclaiming the Arian party, but they
were bound to erect a witness for the truth, which might be a
guide and a warning to all Catholics, against the lying spirit which
was abroad in the Church. These remarks apply to a censure which is
sometimes passed on them, as if it was their duty to have shut up the
question in the words of Scripture; for the words of Scripture were
the very subject in controversy, and to have prohibited the
controversy, would, in fact, have been but to insult the perplexed,
and to extend real encouragement to insidious opponents of the truth.—But
it may be expedient here to explain more fully the principle of the
obligation which led to their interposition.

Let it be observed then, that as regards the doctrine of the
Trinity, the mere text of Scripture is not calculated either to
satisfy the intellect or to ascertain the temper of those who profess
to accept it as a rule of faith.

1. Before the mind has been roused to reflection and
inquisitiveness about its own acts and impressions, it acquiesces, if
religiously trained, in that practical devotion to the Blessed
Trinity, and implicit acknowledgment of the divinity of Son and
Spirit, which holy Scripture at once teaches and exemplifies. This is
the faith of uneducated men, which is not the less philosophically
correct, nor less acceptable to God, because it does not happen to be
conceived in those precise statements which presuppose the action of
the mind on its own sentiments and notions. Moral feelings do not
directly contemplate and realize to themselves the objects which
excite them. A heathen in obeying his conscience, implicitly worships
Him of whom he has never distinctly heard. Again, a child feels
not the less affectionate reverence towards his parents, because he
cannot discriminate in words, nay, or in idea, between them and
others. As, however, his reason opens, he might ask himself concerning
the ground of his own emotions and conduct towards them; and might
find that these are the correlatives of their peculiar tenderness
towards him, long and intimate knowledge of him, and unhesitating
assumption of authority over him; all which he continually
experiences. And further, he might trace these characteristics of
their influence on him to the essential relation itself, which
involves his own original debt to them for the gift of life and
reason, the inestimable blessing of an indestructible, never-ending
existence. And now his intellect contemplates the object of those
affections, which acted truly from the first, and are not purer or
stronger merely for this accession of knowledge. This will tend to
illustrate the sacred subject to which we are directing our attention.

As the mind is cultivated and expanded, it cannot refrain from the
attempt to analyze the vision which influences the heart, and the
Object in which that vision centres; nor does it stop till it has, in
some sort, succeeded in expressing in words, what has all along been a
principle both of its affections and of its obedience. But here the
parallel ceases; the Object of religious veneration being unseen, and
dissimilar from all that is seen, reason can but represent it in the
medium of those ideas which the experience of life affords (as we see
in the Scripture account, as far as it is addressed to the intellect);
and unless these ideas, however inadequate, be correctly applied to
it, they react upon the affections, and deprave the religious
principle. This is exemplified in the case of the heathen, who, trying
to make their instinctive notion of the Deity an object of reflection,
pictured to their minds false images, which eventually gave them a
pattern and a sanction for sinning. Thus the systematic doctrine of
the Trinity may be considered as the shadow, projected for the
contemplation of the intellect, of the Object of scripturally-informed
piety: a representation, economical; necessarily imperfect, as being
exhibited in a foreign medium, and therefore involving apparent
inconsistencies or mysteries; given to the Church by tradition
contemporaneously with those apostolic writings, which are addressed
more directly to the heart; kept in the background in the infancy of
Christianity, when faith and obedience were vigorous, and brought
forward at a time when, reason being disproportionately developed, and
aiming at sovereignty in the province of religion, its presence became
necessary to expel an usurping idol from the house of God.

If this account of the connexion between the theological system and
the Scripture implication of it be substantially correct, it will be
seen how ineffectual all attempts ever will be to secure the doctrine
by mere general language. It may be readily granted that the
intellectual representation should ever be subordinate to the
cultivation of the religious affections. And after all, it must be
owned, so reluctant is a well-constituted mind to reflect on its own
motive principles, that the correct intellectual image, from its
hardness of outline, may startle and offend those who have all along
been acting upon it. Doubtless there are portions of the
ecclesiastical doctrine, presently to be exhibited, which may at
first sight seem a refinement, merely because the object and bearings
of them are not understood without reflection and experience. But what
is left to the Church but to speak out, in order to exclude error?
Much as we may wish it, we cannot restrain the rovings of the
intellect, or silence its clamorous demand for a formal statement
concerning the Object of our worship. If, for instance, Scripture bids
us adore God, and adore His Son, our reason at once asks, whether it
does not follow that there are two Gods; and a system of doctrine
becomes unavoidable; being framed, let it be observed, not with a view
of explaining, but of arranging the inspired notices concerning the
Supreme Being, of providing, not a consistent, but a connected
statement. There the inquisitiveness of a pious mind rests, viz., when
it has pursued the subject into the mystery which is its limit. But
this is not all. The intellectual expression of theological truth not
only excludes heresy, but directly assists the acts of religious
worship and obedience; fixing and stimulating the Christian spirit in
the same way as the knowledge of the One God relieves and illuminates
the perplexed conscience of the religious heathen.—And thus much on
the importance of Creeds to tranquillize the mind; the text of
Scripture being addressed principally to the affections, and of a
religious, not a philosophical character.

2. Nor, in the next place, is an assent to the text of Scripture
sufficient for the purposes of Christian fellowship. As the sacred
text was not intended to satisfy the intellect, neither was it given
as a test of the religious temper which it forms, and of which it is
an expression. Doubtless no combination of words will ascertain
an unity of sentiment in those who adopt them; but one form is more
adapted for the purpose than another. Scripture being unsystematic,
and the faith which it propounds being scattered through its
documents, and understood only when they are viewed as a whole, the
Creeds aim at concentrating its general spirit, so as to give security
to the Church, as far as may be, that its members take that definite
view of that faith which alone is the true one. But, if this be the
case, how idle is it to suppose that to demand assent to a form of
words which happens to be scriptural, is on that account sufficient to
effect an unanimity in thought and action! If the Church would be
vigorous and influential, it must be decided and plain-spoken in its
doctrine, and must regard its faith rather as a character of mind than
as a notion. To attempt comprehensions of opinion, amiable as the
motive frequently is, is to mistake arrangements of words, which have
no existence except on paper, for habits which are realities; and
ingenious generalizations of discordant sentiments for that practical
agreement which alone can lead to co-operation. We may indeed
artificially classify light and darkness under one term or formula;
but nature has her own fixed courses, and unites mankind by the
sympathy of moral character, not by those forced resemblances which
the imagination singles out at pleasure even in the most promiscuous
collection of materials. However plausible may be the veil thus thrown
over heterogeneous doctrines, the flimsy artifice is discomposed so
soon as the principles beneath it are called upon to move and act. Nor
are these attempted comprehensions innocent; for, it being the
interest of our enemies to weaken the Church, they have always gained
a point, when they have put upon us words for things, and persuaded us
to fraternize with those who, differing from us in essentials,
nevertheless happen, in the excursive range of opinion, somewhere to
intersect that path of faith, which centres in supreme and zealous
devotion to the service of God.

Let it be granted, then, as indisputable, that there are no two
opinions so contrary to each other, but some form of words may be
found vague enough to comprehend them both. The Pantheist will admit
that there is a God, and the Humanitarian that Christ is God, if they
are suffered to say so without explanation. But if this be so, it
becomes the duty, as well as the evident policy of the Church, to
interrogate them, before admitting them to her fellowship. If the
Church be the pillar and ground of the truth, and bound to contend for
the preservation of the faith once delivered to it; if we are
answerable as ministers of Christ for the formation of one, and one
only, character in the heart of man; and if the Scriptures are given
us, as a means indeed towards that end, but inadequate to the office
of interpreting themselves, except to such as live under the same
Divine Influence which inspired them, and which is expressly sent down
upon us that we may interpret them,—then, it is evidently our duty
piously and cautiously to collect the sense of Scripture, and solemnly
to promulgate it in such a form as is best suited, as far as it goes,
to exclude the pride and unbelief of the world. It will be admitted
that, to deny to individual Christians the use of terms not
found in Scripture, as such, would be a superstition and an
encroachment on their religious liberty; and in like manner,
doubtless, to forbid the authorities of the Church to require an
acceptance of such terms, when necessary, from its members, is to
interfere with the discharge of their peculiar duties, as appointed of
the Holy Ghost to be overseers of the Lord's flock. And, though the
discharge of this office is the most momentous and fearful that can
come upon mortal man, and never to be undertaken except by the
collective illumination of the Heads of the Church, yet, when
innovations arise, they must discharge it to the best of their
ability; and whether they succeed or fail, whether they have judged
rightly or hastily of the necessity of their interposition, whether
they devise their safeguard well or ill, draw the line of Church
fellowship broadly or narrowly, countenance the profane reasoner, or
cause the scrupulous to stumble,—to their Master they stand or fall,
as in all other acts of duty, the obligation itself to protect the
Faith remaining unquestionable.

This is an account of the abstract principle on which
ecclesiastical confessions rest. In its practical adoption it has been
softened in two important respects. First, the Creeds imposed have
been compiled either from Apostolical traditions, or from primitive
writings; so that in fact the Church has never been obliged literally
to collect the sense of Scripture. Secondly, the test has been used,
not as a condition of communion, but of authority. As learning is not
necessary for a private Christian, so neither is the full knowledge of
the theological system. The clergy, and others in station, must
be questioned as to their doctrinal views: but for the mass of the
laity, it is enough if they do not set up such counter-statements of
their own, as imply that they have systematized, and that erroneously.
In the Nicene Council, the test was but imposed on the Rulers of the
Church. Lay communion was not denied to such as refused to take it,
provided they introduced no novelties of their own; the anathemas or
excommunications being directed solely against the Arian innovators.

\xsection{Section 2. The Scripture Doctrine of the Trinity}

I BEGIN by laying out the matter of
evidence for the Catholic Doctrine, as it is found in Scripture; that
is, assuming it to be there contained, let us trace out the form in
which it has been communicated to us,—the disposition of the
phenomena, which imply it, on the face of the Revelation. And here be
it observed, in reference to what has already been admitted concerning
the obscurity of the inspired documents, that it is nothing to the
purpose whether or not we should have been able to draw the following
view of the doctrine from them, had it never been suggested to us in
the Creeds. For it has been (providentially) so suggested to all of
us; and the question is not, what we should have done, had we never
had external assistance, but, taking things as we find them, whether,
the clue to the meaning of Scripture being given, (as it ever has been
given,) we may not deduce the doctrine thence, by as argumentative a
process as that which enables us to verify the received theory of
gravitation, which perhaps we could never have discovered for
ourselves, though possessed of the data from which the inventor drew
his conclusions. Indeed, such a state of the case is analogous to that
in which the evidence for Natural Religion is presented to us. It is
very doubtful, whether the phenomena of the visible world would
in themselves have brought us to a knowledge of the Creator; but the
universal tradition of His existence has been from the beginning His
own comment upon them, graciously preceding the study of the evidence.
With this remark I address myself to an arduous undertaking.

First, let it be assumed as agreeable both to reason and
revelation, that there are Attributes and Operations, or by whatever
more suitable term we designate them, peculiar to the Deity; for
instance, creative and preserving power, absolute prescience, moral
sovereignty, and the like. These are ever included in our notion of
the incommunicable nature of God; and, by a figure of speech, were
there occasion for using it, might be called one with God, present,
actively co-operating, and exerting their own distinguishing
influence, in all His laws, providences, and acts. Thus, if He be
eternal, or omnipresent, we consider His knowledge, goodness, and
holiness, to be co-eternal and co-extensive with Him. Moreover, it
would be an absurdity to form a comparison between these and God
Himself; to regard them as numerically distinct from Him; to
investigate the particular mode of their existence in the Divine Mind;
or to treat them as parts of God, inasmuch as they are all included in
the idea of the one Indivisible Godhead. And, lastly, subtle and
unmeaning questions might be raised about some of these; for instance,
God's power: whether, that is, it did or did not exist from eternity,
on the ground, that bearing a relation to things created, it could not
be said to have existence before the era of creation \footnote{Origen de Principiis, i. 2, § 10.}.

Next, it is to be remarked, that the Jewish Scriptures introduce to
our notice certain peculiar Attributes or Manifestations (as they
would seem) of the Deity, corresponding in some measure to those
already mentioned as conveyed to us by Natural Religion, though of a
more obscure character. Such is what is called ``the Spirit of
God;'' a phrase which denotes sometimes the Divine energy,
sometimes creative or preserving power, sometimes the assemblage of
Divine gifts, moral and intellectual, vouchsafed to mankind; having in
all cases a general connexion with the notion of the vivifying
principle of nature. Such again, is ``the Wisdom of God,'' as
introduced into the book of Proverbs; and such is the
``Name,'' the ``Word,'' the ``Glory,'' of God.

Further, these peculiar Manifestations (to give them a name) are
sometimes in the same elder Scriptures singularly invested with the
properties of personality; and, although the expressions of the sacred
text may in some places be interpreted figuratively, yet there are
passages so strangely worded, as at first sight to be inconsistent
with themselves, and such as would be ascribed, in an uninspired work,
to forgetfulness or inaccuracy in the writer;—as, for instance, when
what is first called the Glory of God is subsequently spoken of as an
intelligent Agent, often with the characteristics, or even the name of
an Angel. On the other hand, it elsewhere occurs, that what is
introduced as an Angel, is afterwards described as God Himself.

Now, when we pass on to the New Testament, we find these peculiar
Manifestations of the Divine Essence concentrated and fixed in two,
called the Word, and the Spirit. At the same time, the apparent
Personality ascribed to Them in the Old Testament, is changed for a
real Personality, so clearly and explicitly marked as to resist all
critical experiments upon the language, all attempts at allegorical
interpretation. Here too the Word is also called the Son of God, and
appears to possess such strict personal attributes, as to be able
voluntarily to descend from heaven, and assume our nature without
ceasing to be identically what He was before; so as to speak of
Himself, though a man, as one and the same with the Divine Word who
existed in the beginning. The Personality of the Spirit in some true
and sufficient sense is as accurately revealed; and that the Son is
not the Spirit, is also evident from the fixed relations which are
described as separating Them from each other in the Divine Essence.

Reviewing this process of revelation, Gregory Nazianzen, somewhat
after the manner of the foregoing account, remarks that, as Almighty
God has in the course of His dispensations changed the ritual of
religion by successive abrogations, so He has changed its theology by
continual additions till it has come to perfection. ``Under the
Old Dispensation,'' he proceeds, ``the Father was openly
revealed, and the Son but obscurely. When the New was given, the Son
was manifested, but the Divinity of the Spirit intimated only. Now the
Spirit dwells with us, affording us clearer evidence about Himself,
… that by gradual additions, and flights, as David says, and by
advancing and progressing from glory to glory, the radiance of the
Trinity might shine out on those who are illuminated.'' \footnote{Greg. Naz. Orat. xxxvii. p. 608; [xxxi. 26.]}

Now from this peculiar method in which the doctrine is unfolded to
us in Scripture, we learn so much as this in our contemplation of it;
viz. the absurdity, as well as the presumption, of inquiring minutely
about the actual relations subsisting between God and His Son and
Spirit, and drawing large inferences from what is told us of Them.
Whether They are equal to Him or unequal, whether posterior to Him in
existence or coeval, such inquiries (though often they must be
answered when once started) are in their origin as superfluous as
similar questions concerning the Almighty's relation to His own
attributes (which still we answer as far as we can, when asked); for
the Son and the Spirit are \emph{one} with Him, the ideas of number
and comparison being excluded. Yet this statement must be qualified
from the evidence of Scripture, by two additional remarks. On the one
hand, the Son and Spirit are represented to us in the Economy of
Revelation, as ministering to God, and as, so far, personally
subordinate to Him; and on the other hand, in spite of this personal
inequality, yet, as being partakers of the fulness of the Father, they
are equal to Him in nature, and in Their claims upon our faith and
obedience, as is sufficiently proved by the form of baptism.

The mysteriousness of the doctrine evidently lies in our inability
to conceive a sense of the word \emph{person}, such, as to be more than a
mere character, yet less than an individual intelligent being; our own
notions, as gathered from our experience of human agents, leading us
to consider \emph{personality} as equivalent, in its very idea, to the
unity and independence of the immaterial substance of which it is
predicated.

\xsection{Section 3. The Ecclesiastical Doctrine of the Trinity}

THIS being the general Scripture view
of the Holy Trinity, it follows to describe the Ecclesiastical
Doctrine, chiefly in relation to our Lord, as contained in the
writings of the Fathers, especially the Ante-Nicene \footnote{The examples cited are principally borrowed from the elaborate
catalogues furnished by Petavius, Bishop Bull, and Suicer, in his
Thesaurus and his Comment on the Nicene Creed.}.

Scripture is express in declaring both the divinity of Him who in
due time became man for us, and also His personal distinction from God
in His pre-existent state. This is sufficiently clear from the opening
of St. John's Gospel, which states the mystery as distinctly as an
ecclesiastical comment can propound it. On these two truths the whole
doctrine turns, viz. that our Lord is one with, yet personally
separate from God. Now there are two appellations given to Him in
Scripture, enforcing respectively these two essentials of the true
doctrine; appellations imperfect and open to misconception by
themselves, but qualifying and completing each other. The title of the
\emph{Son} marks His derivation and distinction from the Father,
that of the \emph{Word} (i.e. Reason) denotes His inseparable
inherence in the Divine Unity; and while the former taken by itself,
might lead the mind to conceive of Him as a second being, and the
latter as no real being at all, both together witness to the mystery,
that He is at once \emph{from}, and yet \emph{in}, the Immaterial,
Incomprehensible God. Whether or not these titles contain the proof of
this statement, (which, it is presumed, they actually do,) at least,
they will enable us to classify our ideas: and we have authority for
so using them. ``The Son,'' says Athanasius, ``is the Word
and Wisdom of the Father: from which titles we infer His impassive and
indivisible derivation from the Father, inasmuch as the word (or
reason) of a man is no mere part of him, nor when exercised, goes
forth from him by a passion; much less, therefore, is it so with the
Word of God. On the other hand, the Father calls Him His Son, lest,
from hearing only that He was the Word, we should consider Him such as
the word of man, impersonal, whereas the title of Son, designates Him
as a Word which exists, and a substantial Wisdom.'' \footnote{Athan. de Syn. 41.

In the same way the Semi-Arian Basil (of Ancyra), speaking of such
heretics as argued that the Son has no existence separate from the
Father, because He is called the Word, says, ``For this reason our
predecessors, in order to signify that the Son has a reality, and is
in being, and not a mere word which comes and goes, were obliged to
call Him a substance ... For a word has no real existence, and cannot
be a Son of God, else were there many Sons.'' Epiph. Hær. lxxiii.
12.}

Availing ourselves of this division, let us first dwell on the
appellation of Son, and then on that of Word or Reason.

1.

Nothing can be plainer to the attentive student of Scripture, than
that our Lord is there called the Son of God, not only in respect of
His human nature, but of His pre-existent state also. And if this be
so, the very fact of the revelation of Him as such, implies that we
are to gather something from it, and attach in consequence of it some
ideas to our notion of Him, which otherwise we should not have
attached; else would it not have been made. Taking then the word in
its most vague sense, so as to admit as little risk as possible of
forcing the analogy, we seem to gain the notion of derivation from
God, and therefore, of the utter dissimilarity and distance existing
between Him and all beings except God His Father, as if He partook of
that unapproachable, incommunicable Divine Nature, which is increate
and imperishable.

But Scripture does not leave us here: in order to fix us in this
view, lest we should be perplexed with another notion of the analogy,
derived from that adopted sonship, which is ascribed therein to
created beings, it attaches a characteristic epithet to His Name, as
descriptive of the peculiar relation of Him who bears it to the
Father. It designates Him as the \emph{Only-begotten} or the \emph{own}
\footnote{[John i. 1, 14, 18; iii. 16; v. 18. Rom.
viii. 32. Heb. i. 1-14.]} Son of God, terms
evidently referring, where they occur, to His heavenly nature, and
thus becoming the inspired comment on the more general title. It is
true that the term \emph{generation} is also applied to certain events
in our Lord's mediatorial history: to His resurrection from the dead \footnote{Ps. ii. 7. Acts xiii. 33. Heb. v. 5. Rev. i.
5. Rom. i. 4.}; and, according to the Fathers \footnote{Bull, Defens. Fid. Nic. iii. 9, § 12.}, to His original mission in the beginning of all things to
create the world; and to His manifestation in the flesh. Still,
granting this, the sense of the word ``only-begotten''
remains, defined by its context to relate to something higher than any
event occurring in time, however great or beneficial to the human
race.

Being taken then, as it needs must be taken, to designate His
original nature, it witnesses most forcibly and impressively to that
which is peculiar in it, viz. His origination from God, and such as to
exclude all resemblance to any being but Him, whom nothing created
resembles. Thus, without irreverently and idly speculating upon the
generation in itself, but considering the doctrine as given us as a
practical direction for our worship and obedience, we may accept it in
token, that whatever the Father is, such is the Son. And there are
some remarkable texts in Scripture corroborative of this view: for
instance, that in the fifth chapter of St. John, ``As the Father
hath life in Himself, so hath He given to the Son to have life in
Himself ... What things soever the Father doeth, these also doeth the
Son likewise. For the Father loveth the Son, and showeth Him all
things that Himself doeth ... As the Father raiseth up the dead and
quickeneth them, even so the Son quickeneth whom He will ... that all
men should honour the Son even as they honour the Father. He that
honoureth not the Son, honoureth not the Father which hath sent
Him.''

This is the principle of interpretation acknowledged by the
primitive Church. Its teachers warn us against resting in the
word ``generation,'' they urge us on to seize and use its
practical meaning. ``Speculate not upon the divine generation (\emph{gennesis}),''
says Gregory Nazianzen, ``for it is not safe ... let the doctrine
be honoured silently; it is a great thing for thee to know the fact;
the mode, we cannot admit that even Angels understand, much less
thou.'' \footnote{Greg. Naz. Orat. xxxv. 29, 30 [xxix. 8].} Basil says,
``Seek not what is undiscoverable, for you will not discover; ...
if you will not comply, but are obstinate, I shall deride you, or
rather I weep at your daring: … believe what is written, seek not
what is not written.'' \footnote{Petav. v. 6, § 2.}
Athanasius and Chrysostom repel the profane inquiry argumentatively.
``Such speculators,'' the former says, ``might as well
investigate, where God is, and how God is, and of what nature the
Father is. But as such questions are irreligious, and argue ignorance
of God, so is it also unlawful to venture such thoughts about the
generation of the Son of God.'' And Chrysostom; ``I know that
He begat the Son: the manner how, I am ignorant of. I know that the
Holy Spirit is from Him; how from Him, I do not understand. I eat
food; but how this is converted into my flesh and blood, I know not.
We know not these things, which we see every day when we eat, yet we
meddle with inquiries concerning the substance of God.'' \footnote{Ibid.}

While they thus prohibited speculation, they boldly used the
doctrine for the purposes for which it was given them in Scripture.
Thus Justin Martyr speaks of Christ as the Son, ``who alone is
literally called by that name:'' and arguing with the heathen, he
says, ``Jesus might well deserve from His wisdom to be
called the Son of God, though He were only a man like others, for all
writers speak of God as the 'Father of both men and gods.' But let it
not be strange to you, if, besides this common generation, we consider
Him, as the Word of God, to have been begotten of God in a special
way.'' \footnote{Bull, Defens. ii. 4, § 2. [The sentence
runs on thus:—[\emph{tois ton Herm\={e} logon ton para theou
angeltikon legousin}]. Apol. i. 22.]} Eusebius of
Cæsarea, unsatisfactory as he is as an authority, has nevertheless
well expressed the general Catholic view in his attack upon Marcellus.
``He who describes the Son as a creature made out of
nothing,'' he says, ``does not observe that he is bestowing on
Him only the name of Son, and denying Him to be really such; for He
who has come out of nothing, cannot truly be the Son of God, more than
other things which are made. But He who is truly the Son, born from
God, as from a Father, He may reasonably be called the singularly
beloved and only-begotten of the Father, and therefore He is Himself
God.'' \footnote{Euseb. de Eccles. Theol. i. 9, 10.} This last
inference, that what is born of God, is God, of course implicitly
appeals to, and is supported by, the numerous texts which expressly
call the Son God, and ascribe to Him the divine attributes \footnote{The following are additional specimens
from primitive theology. Clement calls the Son ``the \emph{perfect}
Word, born of the \emph{perfect} Father.'' Tertullian, after quoting the
text, ``All that the Father hath are Mine,'' adds, ``If
so, why should not the Father's titles be His? When then we read that
God is Almighty, and the Highest, and the God of Hosts, and the King
of Israel, and Jehovah, see to it whether the Son also be not
signified by these passages, as being in His own right the Almighty
God, inasmuch as He is the Word of the Almighty God.'' Bull,
Defens. ii. 6, § 3. 7, § 4.}.

The reverential spirit in which the Fathers held the doctrine of
the \emph{gennesis}, led them to the use of other forms of expression,
partly taken from Scripture, partly not, with a view of signifying the
fact of the Son's full participation in the divinity of Him who is His
Father, without dwelling on the mode of participation or origination,
on which they dared not speculate \footnote{Vid. Athan. ad Serap. i. 20.}. Such were the images of the sun and its radiance, the
fountain and the stream, the root and its shoots, a body and its
exhalation, fire and the fire kindled from it; all which were used as
emblems of the sacred mystery in those points in which it was declared
in Scripture, viz. the mystery of the Son's being from the Father and,
as such, partaker in His Divine perfections. The first of these is
found in the first chapter of the Epistle to the Hebrews, where our
Lord is called, ``the brightness of God's glory.'' These
illustrations had a further use in their very variety, as reminding
the Christian that he must not dwell on any one of them for its own
sake. The following passage from Tertullian will show how they were
applied in the inculcation of the sacred doctrine. ``Even when a
ray is shot forth from the sun, though it be but a part from the
whole, yet the sun is in the ray, inasmuch as it is the ray of the
sun; nor is its substance separated, but drawn out. In like manner
there is Spirit from Spirit, and God from God. As when a light is
kindled from another, the original light remains entire and
undiminished, though you borrow from it many like itself; so That
which proceeds from God, is called at once God, and the Son of God,
and Both are One.'' \footnote{Bull, Defens. ii. 7, § 2.}

So much is evidently deducible from what Scripture tells us
concerning the generation of the Son; that there is, (so to express
it,) a reiteration of the One Infinite Nature of God, a communicated
divinity, in the Person of our Lord; an inference supported by the
force of the word ``only begotten,'' and verified by the
freedom and fulness with which the Apostles ascribe to Christ the high
incommunicable titles of eternal perfection and glory. There is one
other notion conveyed to us in the doctrine, which must be evident as
soon as stated, little as may be the practical usefulness of dwelling
upon it. The very name of Son, and the very idea of derivation, imply
a certain subordination of the Son to the Father, so far forth as we
view Him as distinct from the Father, or in His personality: and
frequent testimony is borne to the correctness of this inference in
Scripture, as in the descriptions of the Divine Angel in the Old
Testament, revived in the closing revelations of the New \footnote{Rev. viii. 5.}; and in such passages as that above cited from St. John's
Gospel \footnote{John v. 19-30.}. This is a
truth which every Christian feels, admits, and acts upon; but from
piety he would not allow himself to reflect on what he does, did not
the attack of heresies oblige him. The direct answer which a true
religious loyalty leads him to make to any question about the
subordination of the Son, is that such comparisons are irreverent,
that the Son is one with the Father, and that unless he honours the
Son in all the fulness of honour which he ascribes to the Father, he
is disobeying His express command. It may serve as a very faint
illustration of the offence given him, to consider the manner in which
he would receive any question concerning the love which he feels
respectively for two intimate friends, or for a brother and sister, or
for his parents: though in such cases the impropriety of the inquiry,
arises from the incommensurableness, not the coincidence, of the
respective feelings. But false doctrine forces us to analyze our own
notions, in order to exclude it. Arius argued that, since our Lord was
a Son, therefore He was not God: and from that time we have been
obliged to determine how much we grant and what we deny, lest, while
praying without watching, we lose all. Accordingly, orthodox theology
has since his time worn a different aspect; first, inasmuch as divines
have measured what they said themselves; secondly, inasmuch as they
have measured the Ante-Nicene language, which by its authors was
spoken from the heart, by the necessities of controversies of a later
date. And thus those early teachers have been made appear technical,
when in fact they have only been reduced to system; just as in
literature what is composed freely, is afterwards subjected to the
rules of grammarians and critics. This must be taken as an apology for
whatever there is that sounds harsh in the observations which I have
now to make, and for the injustice which I may seem incidentally to do
in the course of them to the ancient writers whose words are in
question.

``The Catholic doctors,'' says Bishop Bull, ``both
before and after the Nicene Council, are unanimous in declaring that
the Father is greater than the Son, even as to divinity [paternity?];
i.e. not in nature or any essential perfection, which is in the Father
and not in the Son, but alone in what may be called authority, that is in point of origin, since the Son is from the Father, not the
Father from the Son.'' \footnote{Bull, Defens. iv. 2, § 1. Or, again, to
take the words of Petavius: [``Filius eandem numero cum Patre
divinitatem habet, sed proprietate differt. Proinde Filietas ipsa
Paternitat equodammodo minor est, vel Filius, qua Filius, Patre, ut
Pater est, minor dicitur, quoniam origine est posterior, non autem ut
Deus,'' ii. 2, § 15.] Cudworth, too, observes: ``Petavius
himself, expounding the Athanasian creed, writeth in this manner: 'The
Father is in a right Catholic manner affirmed by most of the ancients,
to be greater than the Son, and He is commonly said also, without
reprehension, to be before Him in respect of original.' Whereupon he
concludeth the true meaning of that Creed to be this, that no Person
of the Trinity is greater or less than other in respect of the essence
of the Godhead common to them all … but that notwithstanding there
may be some inequality in them, as they are Hic Deus et Hæc Persona.
Wherefore when Athanasius, and the other orthodox Fathers, writing
against Arius, do so frequently assert the equality of all the Three
Persons, this is to be understood in way of opposition to Arius only,
who made the Son to be unequal to the Father, as [\emph{heteroousios}]
… one being God, and the other a creature; they affirming on the
contrary, that He was equal to the Father, as [\emph{homoousios}] ...
that is, as God and not a creature.'' Cudw. Intell. Syst. 4, § 36.}
Justin, for instance, speaks of the Son as ``having the second
place after the unchangeable and everlasting God and Father of
all.'' Origen says that ``the Son is not more powerful than
the Father, but subordinate ([\emph{hypodeesteron}]); according to His
own words, 'The Father that sent Me, is greater than I.''' This
text is cited in proof of the same doctrine by the Nicene, and
Post-Nicene Fathers, Alexander, Athanasius, Basil, Gregory Nazianzen,
Chrysostom, Cyril, and others, of whom we may content ourselves with
the words of Basil: ``'My Father is greater than I,' that is, so
far forth as Father, since what else does 'Father' signify, than that
He is cause and origin of Him who was begotten by Him?'' and in
another place, ``The Son is second in order to the Father,
since He is from Him; and in dignity, inasmuch as the Father is the
origin and cause of His existence.'' \footnote{Justin, Apol. i. 13, 60. Bull, Defens. iv.
2, § 6, § 9. Petav. ii. 2, § 2, \&c.}

Accordingly, the primitive writers, with an unsuspicious yet
reverent explicitness, take for granted the ministrative character of
the relation of both Son and Spirit towards the Father; still of
course speaking of Them as included in the Divine Unity, not as
external to it. Thus Irenæus, clear and undeniable as is his
orthodoxy, still declares, that the Father ``is ministered to in
all things by His own Offspring and Likeness, the Son and Holy Ghost,
the Word and Wisdom, of whom all angels are servants and
subjects.'' \footnote{Petav. i. 3, § 7.} In
like manner, a ministry is commonly ascribed to the Son and Spirit,
and a bidding and willing to the Father, by Justin, Irenæus, Clement,
Origen, and Methodius \footnote{[\emph{hyp\={e}resia, boul\={e}sis,
thel\={e}ma}], præceptio. Petav. ibid. et. seqq.},
altogether in the spirit of the Post-Nicene authorities already cited:
and without any risk of misleading the reader, as soon as the second
and third Persons are understood to be internal to the Divine Mind, \emph{connaturalia
instrumenta}, concurrent (at the utmost) in no stronger sense, than
when the human will is said to concur with the reason. Gregory
Nazianzen lays down the same doctrine with an explanation, in the
following sentence: ``It is plain,'' he says, ``that the
things, of which the Father designs in Him the forms, these the Word
executes; not as a servant, nor unskilfully, but with full knowledge
and a master's power, and, to speak more suitably, as if He were
the Father.'' \footnote{Bull, Defens. ii. 13, § 10. [Greg. Orat.
xxx. 11. For the subordination of mediatorship, vid. Athan. Orat. iv.
6.]}

Such is the Scriptural and Catholic sense of the word \emph{Son};
on the other hand, it is easy to see what was the defect of this
image, and the consequent danger in the use of it. First, there was an
appearance of materiality, the more suspiciously to be viewed because
there were heresies at the time which denied or neglected the
spiritual nature of Almighty God. Next, too marked a distinction
seemed to be drawn between the Father and Son, tending to give a
separate individuality to each, and so to introduce a kind of
ditheism; and here too heresy and philosophy had prepared the way for
the introduction of the error. The Valentinians and Manichees are
chargeable with both misconceptions. The Eclectics, with the latter;
being Emanatists, they seem to have considered the Son to be both
individually distinct from the Father, and of an inferior nature.—Against
these errors we have the following among other protests.

Tertullian says, ``We declare that two are revealed as God in
Scripture, two as Lord; but we explain ourselves, lest offence should
be taken. They are not called two, in respect of their both being God,
or Lord, but in respect of their being Father and Son; and this
moreover, not from any division of substance, but from mutual
relation, since we pronounce the Son to be individual with and
inseparable from the Father.'' \footnote{Bull, Defens. ii. 4, § 3. 7, § 5. Petav.
i. 4, § 1.} Origen also, commenting upon the word ``brightness,'' \footnote{[\emph{apaugasma}].}
in the first chapter of the Hebrews, says, ``Holy Scripture
endeavours to give to men a refined perception of its teaching, by
introducing the illustration of breath \footnote{[\emph{atmis}]. Wisd. vii. 25.}. It has selected this material image, in order to our
understanding even in some degree, how Christ, who is Wisdom, issues,
as though Breath, from the perfection of God Himself … In like
manner from the analogy of material objects, He is called a pure and
perfect Emanation of the Almighty glory \footnote{[\emph{aporrhoia}], ibid.}. Both these resemblances most clearly show the fellowship of
nature between the Son and Father. For an emanation seems to be of one
substance with that body of which it is the emanation or breath \footnote{In like manner Justin, after saying that
the Divine Power called the Word is born from the Father, adds,
``but not by separation from Him ([\emph{kat' apotom\={e}n}])
as if the Father lost part of Himself, as corporeal substances are not
the same before and after separation.'' [Tryph. 128.] ``The
Son of God,'' says Clement, ``never relinquishes His place of
watch, not parted or separated off, not passing from place to place,
but always every where, illimitable, all intellect, all the light of
the Father, all eye, all-seeing, all-hearing, all-knowing, searching
the powers with His power.'' [Strom. vii. 2.]}.'' And to guard still more strongly against any
misconception of the real drift of the illustration, he cautions his
readers against ``those absurd fictions which give the notion of
certain literal extensions in the Divine Nature; as if they would
distribute it into parts, and divide God the Father, if they could;
whereas to entertain even the light suspicion of this, is not only an
extreme impiety, but an utter folly also, nay not even intelligible
at all, that an incorporeal nature should be capable of
division.'' \footnote{Bull, Defens. ii. 9, § 19.}

2.

To meet more fully this misconception to which the word \emph{Son}
gave rise, the ancient Fathers availed themselves of the other chief
appellation given to our Lord in Scripture. The Logos or Sophia, the
Word, Reason, or Wisdom of God, is only by St. John distinctly applied
to Christ; but both before his time and by his contemporary Apostles
it is used in that ambiguous sense, half literal, half evangelical,
which, when it is once known to belong to our Lord, guides us to the
right interpretation of the metaphor. For instance, when St. Paul
declares that ``the Word of God is alive and active, and keener
than a two-edged sword, and so piercing as to separate soul and
spirit, joints and nerves, and a judge of our thoughts and designs,
and a witness of every creature,'' it is scarcely possible to
decide whether the revealed law of God be spoken of, or the Eternal
Son. On the whole it would appear that our Lord is called the Word or
Wisdom of God in two respects; first, to denote His essential presence
in the Father, in as full a sense as the attribute of wisdom is
essential to Him; secondly, His mediatorship, as the Interpreter or
Word between God and His creatures. No appellation, surely, could have
been more appositely bestowed, in order to counteract the notions of
materiality and of distinct individuality, and of beginning of
existence, which the title of the Son was likely to introduce into the
Catholic doctrine. Accordingly, after the words lately cited,
Origen uses it (or a metaphor like it) for this very purpose. Having
mentioned the absurd idea, which had prevailed, of parts or extensions
in the Divine Nature, he proceeds: ``Rather, as will proceeds out
of the mind, and neither tears the mind, nor is itself separated or
divided from it, in some such manner must we conceive that the Father
has begotten the Son, who is His Image.'' Elsewhere he says,
``It were impious and perilous, merely because our intellect is
weak, to deprive God, as far as our words go, of His only-begotten
co-eternal Word, viz. the 'wisdom in which He rejoiced.' We might as
well conceive that He was not for ever in joy.'' \footnote{Bull, Defens. iii. 3, § 1.} Hence it was usual to declare that to deny the eternity of our
Lord was all one as saying that Almighty God was once without
intelligence \footnote{[\emph{alogos}].}: for
instance, Athenagoras says, that the Son is ``the first-born of
the Father; not as made, for God being Mind Eternal, had from the
beginning reason in Himself, being eternally intellectual; but as
issuing forth upon the chaotic mass as the Idea and Agent of
Creation.'' \footnote{Bull, Defens. iii. 5, § 2, [\emph{ton logon}...\emph{logikos … proelthon … idea kai}\emph{energeia}].} The
same interpretation of the sacred figure is continued after the Nicene
Council; thus Basil says, ``If Christ be the Power of God, and the
Wisdom, and these be increate and co-eternal with God, (for He never
was without wisdom and power,) then, Christ is increate and co-eternal
with God.'' \footnote{Petav. vi. 9, § 2.}

But here again the metaphor was necessarily imperfect; and,
if pursued, open to misconception. Its obvious tendency was to
obliterate the notion of the Son's Personality, that is, to introduce
Sabellianism. Something resembling this was the error of Paulus of
Samosata and Marcellus: who, from the fleeting and momentary character
of a word spoken, inferred that the Divine Word was but the temporary
manifestation of God's glory in the man Christ. And it was to
counteract this tendency, that is, to witness against it, that the
Fathers speak of Him as the Word in an \emph{hypostasis} \footnote{[\emph{enupostatos Logos}].}, the permanent, real, and living Word.

3.

The above is a sketch of the primitive doctrine concerning our Lord's
divine nature, as contained in the two chief appellations which are
ascribed to Him in Scripture. The opposite ideas they convey may be
further denoted respectively by the symbols ``of God,'' and
``in God;'' \footnote{[\emph{ek theou}] and [\emph{en the\={o}i}].} as
though He were so derived from the simple Unity of God as in no
respect to be divided or extended from it, (to speak metaphorically,)
but to inhere within that ineffable individuality. Of these two
conditions \footnote{[Son and Word, ``\emph{of God},''
and ``\emph{in God}'' however, imply each other. ``If not
Son, neither is He Word: if not Word, neither is He Son.'' Athan.
Orat. iv. 24. ``The Son's Being, because of the Father, is
therefore in the Father.'' Athan. iii. 3. ``Ouia Verbum ideo
Filius.'' August. in Psalm. vii. 14, § 5.]} of the
doctrine, however, the divinity of Christ, and the unity of God, the
latter was much more earnestly insisted on in the early times. The
divinity of our Lord was, on the whole, too plain a truth to
dispute; but in proportion as it was known to the heathen, it would
seem to them to involve this consequence,—that, much as the
Christians spoke against polytheism, still, after all, they did admit
a polytheism of their own instead of the Pagan. Hence the anxiety of
the Apologists, while they assail the heathen creed on this account,
to defend their own against a similar charge. Thus Athenagoras, in the
passage lately referred to, says; ``Let no one ridicule the notion
that God has a Son. For we have not such thoughts either about God the
Father or about the Son as your poets, who, in their mythologies, make
the Gods no better than men. But the Son of God is the Word of the
Father [as Creator] both in idea and in active power \footnote{[\emph{ideai kai energeiai}], as at p.
170.} ... the Father and the Son being one. The Son being in the
Father, and the Father in the Son, in the unity and power of the
Spirit, the Son of God is the Mind and Word of the Father.''
Accordingly, the divinity of the Son being assumed, the early writers
are earnest in protecting the doctrine of the Unity; protecting it
both from the materialism of dividing the Godhead, and the paganism of
separating the Son and Spirit from the Father. And to this purpose
they made both the ``of God,'' and the ``in God,''
subservient, in a manner which shall now be shown.

First, the ``in God.'' It is the clear declaration of
Scripture, which we must receive without questioning, that the Son and
Spirit are in the one God, and He in Them. There is that remarkable
text in the first chapter of St. John which says that the Son is
``in the bosom of the Father.'' In another place it is
said that ``the Son is in the Father and the Father in the
Son.'' (John xiv. 11.) And elsewhere the Spirit of God is compared
to ``the spirit of a man which is in him'' (1 Cor. ii. 11.).
This is, in the language of theology, the doctrine of the \emph{coinherence}
\footnote{[\emph{perich\={o}r\={e}sis}], or
circumincessio.}; which was used from
the earliest times on the authority of Scripture, as a safeguard and
witness of the Divine Unity. A passage from Athenagoras to this
purpose has just been cited. Clement has the following doxology at the
end of his Christian Instructor. ``To the One Only Father and Son,
Son and Father, Son our guide and teacher, with the Holy Spirit also,
to the One in all things, in whom are all things, \&c. … to Him
is the glory, \&c.'' And Gregory of Neocæsarea, if the words
form part of his creed, ``In the Trinity there is nothing created,
nothing subservient, nothing of foreign nature, as if absent from it
once, and afterwards added. The Son never failed the Father, nor the
Spirit the Son, but the Trinity remains evermore unchangeable,
unalterable.'' These authorities belong to the early Alexandrian
School. The Ante-Nicene school of Rome is still more explicit.
Dionysius of Rome says, ``We must neither distribute into three
divinities the awful and divine Unity, nor diminish the dignity and
transcendant majesty of our Lord by the name of creature, but we must
believe in God the Father Almighty, and in Christ Jesus His Son, and
in the Holy Spirit; and believe that the Word is united with the God
of the universe. For He says, I and the Father are One; and, I am in
the Father, and the Father in Me. For thus the Divine Trinity
and the holy preaching of the \emph{monarchia} will be
preserved.'' \footnote{Shortly before he had used the following
still stronger expressions: [\emph{h\={e}n\={o}sthai gar anank\={e}
t\={o}i The\={o}i t\={o}n hol\={o}n ton theion Logon;
emphiloch\={o}rein de t\={o}i The\={o}i kai eudiaitaisthai dei
to Hagion Pneuma}]. The Ante-Nicene African school is as express as
the Roman. Tertullian says, ``Connexus Patris in Filio, et Filii
in Paracleto, tres efficit cohærentes, qui tres unum sint, non unus.''
Bull, Defens. ii. 6, § 4; 12, § 1. 11; iv. 4, 12, § 1. 11; iv. 4,
§ 10.}

This doctrine of the \emph{coinherence}, as protecting the Unity
without intrenching on the perfections of the Son and Spirit, may even
be called the characteristic of Catholic Trinitarianism as opposed to
all counterfeits, whether philosophical, Arian, or Oriental. One
Post-Nicene statement of it shall be added. ``If any one truly
receive the Son,'' says Basil, ``he will find that He brings
with him on one hand His Father, on the other the Holy Spirit. For
neither can He from the Father be severed, who is of and ever in the
Father; nor again from His own Spirit disunited, who in It operates
all things ... For we must not conceive separation or division in any
way; as if either the Son could be supposed without the Father, or the
Spirit disunited from the Son. But there is discovered between them
some ineffable and incomprehensible, both communion and
distinction.'' \footnote{Petav. iv. 16, § 9. The Semi-Arian creed,
called \emph{Macrostichos}, drawn up at Antioch A.D.
345, which is in parts unexceptionable in point of orthodoxy, contains
the following striking exposition of the Catholic notion of the \emph{coinherence}.
``Though we affirm the Son to have a distinct existence and life
as the Father has, yet we do not therefore separate Him from the
Father, inventing place and distance between Their union after a
corporeal manner. For we believe that they are united without medium
or interval, and are inseparable.'' And then follow words to which
our language is unequal: [\emph{holou men tou Patros enesternismenou ton
Huion; holon de tou Huiou ex\={e}rt\={e}menou kai prospephukotos
t\={o}i Patri, kai monou tois patr\={o}ois kolpois anapauomenou
di\={e}nek\={o}s}]. Bull, Defens. iv. 4, § 9.}

Secondly, as the ``in God'' led the Fathers to the doctrine
of the \emph{coinherence}, so did the ``of God'' lead them to
the doctrine of the \emph{monarchia} \footnote{[Vid. Athan. Tr. vol. i. pp. 110-112.]}; still, with the one object of guarding against any
resemblance to Polytheism in their creed. Even the heathen had shown a
disposition, designedly or from a spontaneous feeling, to trace all
their deities up to one Principle or \emph{arche}; as is evident by
their Theogonies \footnote{Cudw. Intell. Syst. 4, § 13.}. Much
more did it become that true religion, which prominently put forth the
Unity of God, jealously to guard its language, lest it should seem to
admit the existence of a variety of original Principles. It is said to
have been the doctrine of the Marcionists and Manichees, that there
were three unconnected independent Beings in the Divine Nature.
Scripture and the Church avoid the appearance of tritheism, by tracing
back, (if we may so say,) the infinite perfections of the Son and
Spirit to Him whose Son and Spirit They are. They are, so to express
it, but the new manifestation and repetition of the Father; there
being no room for numeration or comparison between Them, nor any
resting-place for the contemplating mind, till They are referred to
Him in whom They centre. On the other hand, in naming the Father, we
imply the Son and Spirit, whether They be named or not \footnote{Athan. ad Serap. i. 14.}. Without this key, the language of Scripture is perplexed
in the extreme \footnote{Let John v. 20 be taken as an example; or
again, 1 Cor. xii. 4-6. John xiv. 16-18; xvi. 7-15.}.
Hence it is, that the Father is called ``the only God,'' at a
time when our Lord's name is also mentioned, John xvii. 3, 1 Tim. i.
16, 17, as if the Son was but the reiteration of His Person, who is
the Self-Existent, and therefore not to be contrasted with Him in the
way of number. The Creed, called the Apostles', follows this mode of
stating the doctrine; the title of God standing in the opening against
the Father's name, while the Son and Spirit are introduced as distinct
forms or modes, (so to say,) of and in the One Eternal Being. The
Nicene Creed, commonly so called, directed as it is against the
impugners both of the Son's and of the Spirit's divinity, nevertheless
observes the same rule even in a stricter form, beginning with a
confession of the ``\emph{One} God.'' Whether or not this mode
of speaking was designed in Scripture to guard the doctrine of the
Unity from all verbal infringement (and there seems evidence that it
was so, as in 1 Cor. viii. 5, 6,) it certainly was used for this
purpose in the primitive Church. Thus Tertullian says, that it is a
mistake ``to suppose that the number and arrangement of the
Trinity is a division of its Unity; inasmuch as the Unity drawing out
the Trinity from itself, is not destroyed by it, but is subserved.''
\footnote{Again he says, that ``the Trinity
descending from the Father by closely knit and connected steps, both
is consistent with the \emph{monarchia} (Unity), and protects the \emph{economia}
(revealed dispensation).''} Novatian, in like
manner, says, ``God originating from God, so as to be the Second
Person, yet not interfering with the Father's right to be called the
one God. For, had He not a birth, then indeed when compared with
Him who had no birth, He would seem, from the appearance of equality
in both, to make two who were without birth \footnote{[Or unoriginate; viz, on [\emph{agenn\={e}tos}]
and [\emph{anarchos}] in the next Section.]}, and therefore two Gods.'' \footnote{Petav. Præf. 5, 1. iii.; § 8. Dionysius
of Alexandria implies the same doctrine, when he declares; ``We
extend the indivisible Unity into the Trinity, and again we
concentrate the indestructible Trinity into the Unity.'' And
Hilary, to take a Post-Nicene authority, ``We do not detract from
the Father, His being the one God, when we say also that the Son is
God. For He is God from God, one from one; therefore one God, because
God is from Himself. On the other hand, the Son is not on that account
the less God, because the Father is the one God. For the only-begotten
Son of God is not without birth, so as to detract from the Father His
being the One God, nor is He other than God, but because He is born of
God.'' De Trin. i. Vide also Athan. de Sent. Dionys. 17. Bull,
Defens. iv. 4, § 7.}

Accordingly it is impossible to worship One of the Divine Persons,
without worshipping the Others also. In praying to the Father, we only
arrive at His mysterious presence through His Son and Spirit; and in
praying to the Son and Spirit, we are necessarily carried on beyond
them to the source of Godhead from which They are derived. We see this
in the very form of many of the received addresses to the Blessed
Trinity; in which, without intended reference to the mediatorial
scheme, the Son and Spirit seem, even in the view of the Divine Unity,
to take a place in our thoughts between the Father and His creatures;
as in the ordinary doxologies ``to the Father through the Son and
by the Spirit,'' or ``to the Father and Son in the unity of
the Holy Ghost.''

This gives us an insight into the force of expressions, common with
the primitive Fathers, but bearing, in the eyes of inconsiderate
observers, a refined and curious character. They call the Son,
``God of God, Light of Light,'' \&c., much more frequently
than simply God, in order to anticipate in the very form of words, the
charge or the risk of ditheism. Hence, also, the illustrations of the
sun and his rays, \&c., were in such repute; viz. as containing,
not only a description, but also a defence of the Catholic doctrine.
Thus Hippolytus says, ``When I say that the Son is distinct from
the Father, I do not speak of two Gods; but, as it were, light of
light, and the stream from the fountain, and a ray from the sun.''
\footnote{Bull, Defens. iv. 4, § 5.} It was the same reason
which led the Fathers to insist upon the doctrine of the divine
generation.

\xsection{Section 4. Variations in the Ante-Nicene Theological Statements}

THERE will, of course, be differences
of opinion, in deciding how much of the ecclesiastical doctrine, as
above described, was derived from direct Apostolical Tradition, and
how much was the result of intuitive spiritual perception in
scripturally informed and deeply religious minds. Yet it does not seem
too much to affirm, that copious as it may be in theological terms,
yet hardly one can be pointed out which is not found or strictly
implied in the New Testament itself. And indeed so much perhaps will
be granted by all who have claim to be considered Trinitarians; the
objections, which some among them may be disposed to raise, lying
rather against its alleged over-exactness in systematizing Scripture,
than against the truths themselves which are contained in it. But it
should be remembered, that it is we in after times who systematize the
statements of the Fathers, which, as they occur in their works, are
for the most part as natural and unpremeditated as those of the
inspired volume itself. If the more exact terms and phrases of any
writer be brought together, that is, of a writer who has fixed
principles at all, of course they will appear technical and
severe. We count the words of the Fathers, and measure their
sentences; and so convert doxologies into creeds. That we do so, that
the Church has done so more or less from the Nicene Council downwards,
is the fault of those who have obliged us, of those who, ``while
men slept,'' have ``sowed tares among the wheat.''

This remark applies to the statements brought together in the last
Section, from the early writers: which, even though generally
subservient to certain important ends, as, for instance, the
maintenance of the Unity of God, \&c., are still on the whole
written freely and devotionally. But now the discussion passes on to
that more intentional systematizing on the part of the Ante-Nicene
Fathers, which, unavoidable as it was, yet because it was in part
conventional and individual, was ambiguous, and in consequence
afforded at times an apparent countenance to the Arian heresy. It
often becomes necessary to settle the phraseology of divinity, in
points, where the chief problem is, to select the clearest words to
express notions in which all agree; or to find the proposition which
will best fit in with, and connect, a number of received doctrines.
Thus the Calvinists dispute among themselves whether or not God \emph{wills}
the damnation of the non-elect; both parties agree in doctrine, they
doubt how their own meaning may be best expressed \footnote{Vid. another instance infra, ch. v. § 2, in the controversy
about the use of the word \emph{hypostasis}.}. However clearly we see, and firmly we grasp the truth, we have
a natural fear of the appearance of inconsistency; nay, a becoming
fear of misleading others by our inaccuracy of language; and
especially when our words have been misinterpreted by opponents, are
we anxious to guard against such an inconvenience in future. There are
two characteristics of opinions subjected to this intellectual
scrutiny: first, they are variously expressed during the process;
secondly, they are consigned to arbitrary formulas, at the end of it.
Now, to exemplify this in certain Ante-Nicene statements of the great
Catholic doctrine.

1.

The word [\emph{agenn\={e}tos}], \emph{ingenitus} (\emph{unborn,
ingenerate}), was the philosophical term to denote that which had
existed from eternity. It had accordingly been applied by Aristotle to
the world or to matter, which was according to his system without
beginning; and by Plato to his ideas. Now since the Divine Word was
according to Scripture \emph{generate}, He could not be called \emph{ingenerate}
(or eternal), without a verbal contradiction. In process of time a
distinction was made between [\emph{agen\={e}tos}] and [\emph{agenn\={e}tos}],
(\emph{increate} and \emph{ingenerate},) according as the letter [\emph{n}]
was or was not doubled, so that the Son might be said to be [\emph{agen\={e}t\={o}s
genn\={e}tos}] (\emph{increately generate}). The argument which
arose from this perplexity of language, is urged by Arius himself; who
ridicules the [\emph{agenn\={e}togenes}], \emph{ingenerately-generate},
which he conceives must be ascribed, according to the orthodox creed,
to the Son of God \footnote{Vid. infra, Section 5.}. Some
years afterwards, the same was the palmary, or rather the essential
argument of Eunomius, the champion of the Anomœans.

2.

The [\emph{anarchon}] (\emph{unoriginate}). As Is implied in the
word \emph{monarchia}, as already explained, the Father alone is the \emph{arche},
or \emph{origin}, and the Son and Spirit are not origins. The heresy
of the Tritheists made it necessary to insist upon this. Hence the
condemnation, in the (so-called) Apostolical Canons, of those who
baptized ``into the name of Three Unoriginate.'' \footnote{Bull, Defens. iv. 1, § 6.} And Athanasius says, ``We do not teach three Origins, as
our illustration shows; for we do not speak of three Suns, but of the
Sun and its radiance.'' \footnote{Cudw. Intell. Syst. 4, § 36 [p. 709, ed.
Mosheim. But the Benedictine Ed. in Cyril, Catech. xi., says that
Athanasius maintained the Son's [\emph{anarchon}]. Epiphanius, from 1
Cor. xi. 3, argues that the Father is the [\emph{kephal\={e}}], not
the [\emph{arch\={e}}], of the Son. Hær. 76, fin.]}
For the same reason the early writers spoke of the Father as the Fount
of Divinity. At the same time, lest they should in word dishonour the
Son, they ascribed to Him ``an unoriginate generation'' or
``birth.'' \footnote{Suicer. Symb. Nicen. c. viii.} Thus
Alexander, the first champion of orthodox truth against Arius, in his
letter to his namesake of Byzantium: ``We must reserve to the
unbegotten (or unborn) Father His peculiar prerogative, confessing
that no one is the cause of His existence, and to the Son we must pay
the due honour, attributing to Him the unoriginate generation from the
Father, and as we have said already, paying Him worship, so as ever to
speak of Him piously arid reverently, as 'pre-existent, ever-living,'
and 'before the worlds.''' \footnote{Theod. Hist. i. 4, p. 18.}
This distinction however, as might be expected, was but partially
received among the Catholics. Contrasted with all created
beings, the Son and Spirit are of necessity Unoriginate in the Unity
of the Father. Clement, for instance, calls the Son, ``the
everlasting, unoriginate, origin and commencement of all things.''
\footnote{[\emph{t\={e}n achronon, anarchon, arch\={e}n
te kai aparch\={e}n t\={o}n pant\={o}n}].} It was not till they
became alive to the seeming ditheism of such phrases, which the
Sabellian controversy was sure to charge upon them, that they learned
the accurate discrimination observed by Alexander. On the other hand,
when the Arian contest urged them in the contrary direction to
Sabellius, then they returned more or less to the original language of
Clement, though with a fuller explanation of their own meaning.
Gregory Nyssen gives the following plain account of the variations of
their practice: ``Whereas the word \emph{Origin} has many
significations ... sometimes we say that the appellation of the
Unoriginate is not unsuitable to the Son. For when it is taken to mean
derivation of substance from no cause, this indeed we ascribe to the
Father alone. But according to the other senses of the word, since
creation, time, the order of the world are referred to an origin, in
respect of these we ascribe to the Only-begotten, superiority to any
origin; so as to believe Him to be beyond creation, time, and mundane
order, through whom were made all things. And thus we confess Him, who
is not unoriginate in regard to His subsistence, in all other respects
to be unoriginate, and, while the Father is unoriginate and unborn,
the Son to be unoriginate in the sense explained, but not
unborn.'' \footnote{Gregory Nazianzen says the same more
concisely: [\emph{ho Huios, ean h\={o}s aition ton Patera lamban\={e}s,
ouk anarchos; arch\={e} gar Huiou Pat\={e}r, h\={o}s aitios}].
Bull, Defens. iv. 2, § 8. 1; § 3. Petav. i. 4, § i. Suicer, ibid.}

The word \emph{cause} ([\emph{aitios}]) used in this passage, as a
substitute for that use of Origin which peculiarly applies to the
Father as the Fount of Divinity, is found as early as the time of
Justin Martyr, who in his dialogue with Trypho, declares the Father is
to the Son the [\emph{aitios}], or cause of, His being; and it was
resumed by the Post-Nicene writers, when the Arian controversy was
found to turn in no small degree on the exact application of such
terms. Thus Gregory Nazianzen says, ``There is One God, seeing
that the Son and Spirit are referred to One Cause.'' \footnote{However, here too we have a variation in the
use of the word: [\emph{aitios}] being sometimes applied to the Son in
the sense [\emph{arch\={e}}]. The Latin word answering to [\emph{aitios}]
is sometimes \emph{causa}, more commonly \emph{principium} or \emph{auctor}.
Bull, Defens. iv. 1, § 2; § 4. Petav. v. 5, § 10.}

3.

The Ante-Nicene history of the word homoüsion or \emph{consubstantial},
which the Council of Nicæa adopted as its test, will introduce a more
important discussion.

It is one characteristic of Revelation, that it clears up all
doubts about the existence of God, as separate from, and independent
of nature; and shows us that the course of the world depends not
merely on a system, but on a Being, real, living, and individual. What
we ourselves witness, evidences to us the operation of laws, physical
and moral; but it leaves us unsatisfied, whether or not the principle
of these be a mere nature or fate, whether the life of all things be a
mere Anima Mundi, a spirit connatural with the body in which it
acts, or an Agent powerful to make or unmake, to change or supersede,
according to His will. It is here that Revelation supplies the
deficiency of philosophical religion; miracles are its emblem, as well
as its credentials, forcing on the imagination the existence of an
irresponsible self-dependent Being, as well as recommending a
particular message to the reason. This great truth, conveyed in the
very circumstances under which Revelation was made, is explicitly
recognized in its doctrine. Among other modes of inculcating it, may
be named the appellation under which Almighty God disclosed Himself to
the Israelites; Jehovah (or, as the Septuagint translates it, [\emph{ho \={o}n}])
being an expressive appellation of Him, who is essentially separate
from those variable and perishable beings or substances, which
creation presents to our observation. Accordingly, the description of
Him as [\emph{to on}], or in other words, the doctrine of the [\emph{ousia}]
of God, that is, of God viewed as Being and as the one Being, became
familiar to the minds of the primitive Christians; as embodying the
spirit of the Scriptures, and indirectly witnessing against the
characteristic error of pagan philosophy, which considered the Divine
Mind, not as a reality, but as a mere abstract name, or generalized
law of nature, or at best as a mere mode, principle, or an animating
soul, not a Being external to creation, and possessed of
individuality. Cyril of Alexandria defines the word [\emph{ousia}], (\emph{usia,
being, substance},) to be ``that which has existence in itself,
independent of every thing else to constitute it;'' \footnote{[\emph{pragma authuparkton, m\={e}
deomenon heterou pros t\={e}n heautou sustasin}]. Suicer,
Thesaur. verb. [\emph{ousia}].} that is, an individual. This sense of the word must be
carefully borne in mind, since it was \emph{not} that in which it is used by
philosophers, who by it denoted the genus or species, or the ``ens
unum in multis,''—a sense which of course it could not bear when
applied to the One Incommunicable God. The word, thus appropriated to
the service of the God of Revelation, was from the earliest date used
to express the reality and subsistence of the Son; and no word could
be less metaphorical and more precise for this purpose, although the
Platonists chose to refine, and from an affectation of reverence
refused to speak of God except as \emph{hyperusios} \footnote{[Or [\emph{epekeina ousias}]] Petav. [t. i.
i. 6] t. ii. iv. 5, § 8. [Brucker, t. 2, p. 395. Plot. Enn. v. lib. i.
We find [\emph{hyperousios}] or [\emph{epekeina ousias}] in Orig. c.
Cels. vi. 64. Damasc. F. O. i. 4, 8, and 12.]}. Justin Martyr, for instance, speaks of heretics, who
considered that God put forth and withdrew His Logos when it pleased
Him, as if He were an influence, not a Person \footnote{Justin, Tryph. 128.}, somewhat in the sense afterwards adopted by Paulus of
Samosata and others. To meet this error, he speaks of Him as
inseparable from the substance or being, \emph{usia}, of the Father;
that is, in order to exclude all such evasions of Scripture, as might
represent the man Christ as inhabited by a divine glory, power,
nature, and the like, evasions which in reality lead to the conclusion
that He is not God at all.

For this purpose the word \emph{homoüsion} or \emph{consubstantial}
was brought into use among Christian writers; viz. to express the real
divinity of Christ, and that, as being derived from, and one with the
Father's. Here again, as in the instance of its root, the word was
adopted, from the necessity of the case, in a sense different
from the ordinary philosophical use of it. \emph{Homoüsion} properly
means \emph{of the same nature}, or under the same \emph{general} nature, or
species; that is, it is applied to things, which are but similar to
each other, and are considered as one by an abstraction of our minds;
or, it may mean of the same material. Thus Aristotle speaks of the
stars being consubstantial with each other; and Porphyry of the souls
of brute animals being consubstantial to ours \footnote{Bull, Defens. ii. 1, § 2, \&c.}. When, however, it was used in relation to the incommunicable
Essence of God, there was obviously no abstraction possible in
contemplating Him, who is above all comparison with His works. His
nature is solitary, peculiar to Himself, and one; so that whatever was
accounted to be consubstantial or co-essential with Him, was
necessarily included in His individuality, by all who would avoid
recurring to the vagueness of philosophy, and were cautious to
distinguish between the incommunicable Essence of Jehovah and all
created intelligences. And hence the fitness of the term to denote
without metaphor the relation which the Logos bore in the orthodox
creed to His eternal Father. Its use is explained by Athanasius as
follows. ``Though,'' he says, ``we cannot understand what
is meant by the \emph{usia}, being, or \emph{substance} of God, yet we
know as much as this, that God is, which is the way in which Scripture
speaks of Him; and after this pattern, when we wish to designate Him
distinctly, we say God, Father, Lord. When then He says in Scripture,
'I am [\emph{ho \={o}n}],' the Being, and 'I am Jehovah, God,' or
uses the plain word 'God,' we understand by such statements nothing
but His incomprehensible [\emph{ousia}] (being or substance),
and that He, who is there spoken of, is. Let no one then think it
strange, that the Son of God should be said to be [\emph{ek t\={e}s
ousias}] (from the being or substance) of God; rather, let him
agree to the explanation of the Nicene fathers, who, for the words 'of
God' substituted 'of the divine being or substance.' They considered
the two phrases substantially the same, because, as I have said, the
word 'God' denotes nothing but the [\emph{ousia autou tou ontos}], the
being of Him who is. On the other hand, if the Word be not in such
sense 'of God,' as to be the true Son of the Father according to His
nature, but be said to be 'of God,' merely as all creatures are such
because they are His work, then indeed He is not 'from the being of
the Father,' nor Son 'according to being or substance,' but so called
from His virtue, as we may be, who receive the title from grace.''
\footnote{Athan. de Decr. Nic. 32.}

The term \emph{homoüsios} is first employed for this purpose by
the author of the \emph{Pæmander}, a Christian of the beginning of
the second century. Next it occurs in several writers at the end of
the second and the beginning of the third. In Tertullian, the
equivalent phrase, ``unius substantiæ,'' ``\emph{of one
substance},'' is applied to the Trinity. In Origen's comment on
the Hebrews, the \emph{homoüsion} of the Son is deduced from the
figurative title [\emph{apaugasma}], or \emph{radiance}, there given
to Him. In the same age, it was employed by various writers, bishops
and historians, as we learn from the testimonies of Eusebius and
Athanasius \footnote{[Vide Ath. Tr. vol. ii. p. 438. Also
Archelaus speaks of our Lord as ``de substantiâ Dei.'' Routh,
t. iv. p. 228.]}. But at
this era, the middle of the third century, a change took place
in the use of it and other similar words, which is next to be
explained.

The oriental doctrine of Emanations was at a very early period
combined with the Christian theology. According to the system of
Valentinus, a Gnostic heresiarch, who flourished in the early part of
the second century, the Supreme Intelligence of the world gave
existence to a line of Spirits or Eons, who were all more or less
partakers of His nature, that is, of a nature \emph{specifically} the
same, and included in His glory ([\emph{pl\={e}r\={o}ma}]),
though individually separate from the true and Sovereign Deity. It is
obvious, that such a teaching as this abandons the great revealed
principle above insisted on, the incommunicable character and
individuality of the Divine Essence. It considers all spiritual beings
as like God, in the same sense that one man resembles or has the same
nature as another: and accordingly it was at liberty to apply, and did
actually apply, to the Creator and His creatures the word \emph{homoüsion}
or \emph{consubstantial}, in the philosophical sense which the word
originally bore. We have evidence in the work of Irenæus that the
Valentinians did thus employ it. The Manichees followed, about a
century later; they too were Emanatists, and spoke of the human soul
as being \emph{consubstantial} or \emph{co-essential} with God, of one
substance with God. Their principles evidently allowed of a kind of
Trinitarianism; the Son and Spirit being considered Eons of a superior
order to the rest, \emph{consubstantial} with God because Eons, but
one with God in no sense which was not true also of the soul of man.
It is said, moreover, that they were materialists; and used the word \emph{consubstantial}
as it may be applied to different vessels or instruments,
wrought out from some one mass of metal or wood. However, whether this
was so or not, it is plain that anyhow the word in question would
become unsuitable to express the Catholic doctrine, in proportion as
the ears of Christians were familiarized to the terms employed in the
Gnostic and Manichean theologies; nor is it wonderful that at length
they gave up the use of it.

The history of the word \emph{probole} or \emph{offspring} is
parallel to that of the \emph{consubstantial} \footnote{Beausobre, Hist. Manich. iii. 7, § 6.
[Vide Ath. Tr. vol. ii. p. 458.]}. It properly means any thing which proceeds, or is sent forth
from the substance of another, as the fruit of a tree, or the rays of
the sun; in Latin it is translated by \emph{prolatio}, \emph{emissio},
or \emph{editio}, an offspring or issue. Accordingly Justin employed
it, or rather a cognate phrase \footnote{[\emph{probl\={e}then genn\={e}ma}].
Justin. Tryph. 62.}, to designate what Cyril calls above the self-existence \footnote{[\emph{autogonos}]. [Vide Ath. Tr. art. [\emph{huiopat\={o}r}],
vol. ii. p. 475, ed. 1881.]} of the Son, in opposition to the evasions which were necessary
for the system of Paulus, Sabellius, and the rest. Tertullian does the
same; but by that time, Valentinus had given the word a material
signification. Hence Tertullian is obliged to apologize for using it,
when writing against Praxeas, the forerunner of the Sabellians.
``Can the Word of God,'' he asks, ``be unsubstantial, who
is called the Son, who is even named God? He is said to be in the form
or image of God. Is not God a body [substance], Spirit though He be?
... Whatever then has been the substance of the Word \footnote{[Ibid. p. 340, art. \emph{Word}.]}, that, I call a Person, and claim for it the name of Son, and
being such, He comes next to the Father. Let no one suppose that I am
bringing in the notion of any such \emph{probole} (\emph{offspring})
as Valentinus imagined, drawing out his Eons the one from the other.
Why must I give up the word in a right sense, because heresy uses it
in a wrong? besides, heresy borrowed it from us, and has turned truth
into a lie … This is the difference between the uses of it.
Valentinus separates his \emph{probolæ} from their Father; they know
Him not. But we hold that the Son alone knows the Father, reveals Him,
performs His will, and is within Him. He is ever in the Father, as He
has said; ever with God, as it is written; never separated from Him,
for He and the Father are one. This is the true \emph{probole}, the
safeguard of unity, sent forth, not divided off.'' \footnote{Tertull. in Prax. 7, 8, abridged.} Soon after Tertullian thus defended his use of the word \emph{probole},
Origen in another part of the Church gave it up, or rather assailed
it, in argument with Candidus, a Valentinian. ``If the Son is a \emph{probole}
of the Father,'' he says, ``who begets Him from Himself, like
the birth of animals, then of necessity both offspring and original
are of a bodily nature.'' \footnote{[Periarch. iv. p. 190.]} Here we see two writers, with exactly the same theological
creed before them, taking opposite views as to the propriety of using
a word which heresy had corrupted \footnote{[Vide an apposite note of Coustant. Epp.
Pont. Rom. p. 496, on Damasus's Words: ``nec prolativum, ut
generationem ei demas.'']}.

But to return to the word \emph{consubstantial}: though Origen gave
up the word \emph{probole}, yet he used the word \emph{consubstantial},
as has already been mentioned \footnote{[But he was not consistent. Vide Hieron.
contr. Ruff. ii. 19. Also the dissertation in Jackson's preface to
Novatian, p. xlviii, \&c.]}. But shortly after his death, his pupils abandoned it at the
celebrated Council held at Antioch (A.D.
264) against Paulus of Samosata. When they would have used it as a
test, this heretic craftily objected to it on the very ground on which
Origen had surrendered the \emph{probole}. He urged that, if Father
and Son were of one substance, \emph{consubstantial}, there was some
common substance in which they partook, and which consequently was
distinct from and prior to the Divine Persons Themselves; a wretched
sophism, which of course could not deceive Firmilian and Gregory, but
which, being adapted to perplex weak minds, might decide them on
withdrawing the word. It is remarkable too, that the Council was held
about the time when Manes appeared on the borders of the Antiochene
Patriarchate. The disputative school of Paulus pursued the advantage
thus gained; and from that time used the charge of materialism as a
weapon for attacking all sound expositions of Scripture truth. Having
extorted from the Catholics the condemnation of a word long known in
the Church, almost found in Scripture, and less figurative and
material in its meaning than any which could be selected, and
objectionable only in the mouths of heretics, they employed this
concession as a ground of attacking expressions more directly
metaphorical, taken from visible objects, and sanctioned by less
weighty authority. In a letter which shall afterwards be cited, Arius
charges the Catholics with teaching the errors of Valentinus and
Manes; and in another of the original Arian documents, Eusebius of
Nicomedia, maintains in like manner that their doctrine involves the
materiality of the Divine Nature. Thus they were gradually silencing
the Church by a process which legitimately led to Pantheism,
when the Alexandrians gave the alarm, and nobly stood forward in
defence of the faith \footnote{[Parallel to the above instances is Basil's
objection to [\emph{genn\={e}ma}], when used of the Son, which
Athanasius and others apply to him. Vide Ath. Tr. vol. ii. p. 396.]}.

It is worth observing that, when the Asiatic Churches had given up
the \emph{consubstantial}, they, on the contrary, had preserved it.
Not only Dionysius willingly accepts the challenge of his namesake of
Rome, who reminded him of the value of the symbol; but Theognostus
also, who presided at the Catechetical School at the end of the third
century, recognizes it by implication in the following passage, which
has been preserved by Athanasius. ``The substance \footnote{[It may be questioned, however, whether
the word \emph{substance} in this passage is not equivalent to \emph{hypostasis}
or \emph{subsistence}; vide Appendix, No. 4.]} of the Son,'' he says, ``is not external to the
Father, or created; but it is by natural derivation from that of the
Father, as the radiance comes from light (Heb. i. 3). For the radiance
is not the sun, ... and yet not foreign to it; and in like manner
there is an effluence ([\emph{aporrhoia}], Wisd. vii. 25.) from the
Father's substance, though it be indivisible from Him. For as the sun
remains the same without infringement of its nature, though it pour
forth its radiance, so the Father's substance is unchangeable, though
the Son be its Image.'' \footnote{Athan. de Decr. Nic. 25.}

4.

Some notice of the [\emph{thel\={e}sei genn\={e}then}], or
voluntary generation, will suitably follow the discussion of the \emph{consubstantial}; though the subject does not closely concern
theology. It has been already observed that the tendency of the
heresies of the first age was towards materialism and fatalism. As it
was the object of Revelation to destroy all theories which interfered
with the belief of the Divine Omniscience and active Sovereignty, so
the Church seconded this design by receiving and promulgating the
doctrine of the ``\emph{He that is},'' or the Divine ``\emph{Being}''
or ``\emph{Essence},'' as a symbol of His essential
distinction from the perishable world in which He acts. But when the
word \emph{substance} or \emph{essence} itself was taken by the
Gnostics and Manichees in a material sense, the error was again
introduced by the very term which was intended to witness against it.
According to the Oriental Theory, the emanations from the Deity were
eternal with Himself, and were considered as the result, not of His
will and personal energy, but of the necessary laws to which His
nature was subjected; a doctrine which was but fatalism in another
shape. The Eclectics honourably distinguished themselves in
withstanding this blasphemous, or rather atheistical tenet. Plotinus
declares, that ``God's substance and His will are the same; and if
so, as He willed, so He is; so that it is not a more certain truth
that, as is His substance or nature, so is His will and action, than,
as His will and action, so is His substance.'' Origen had preceded
them in their opposition to the same school. Speaking of the
simplicity and perfection of the Divine Essence, he says, ``God
does not even participate in substance, rather He is partaken; by
those, namely, who have the Spirit of God. And our Saviour does not
share in holiness, but, being holiness itself, is shared by the holy.'' The meaning of this doctrine is clear;—to protest, in
the manner of Athanasius, in a passage lately cited, against the
notion that the substance of God is something distinct from God
Himself, and not God viewed as self-existent, the one immaterial,
intelligent, all-perfect Spirit; but the risk of it lay in its
tendency to destroy the doctrine of His individual and real existence
(which the Catholic use of \emph{substance} symbolized), and to
introduce in its stead the notion that a quality or mode of acting was
the governing principle of nature; in other words, Pantheism. This is
an error of which Origen of course cannot be accused; but it is in its
measure chargeable on the Platonic Masters, and is countenanced even
by their mode of speaking of the Supreme Being, as not substantial,
but above the notion of substance \footnote{[\emph{hyperousios}]. Cudw. Intell. Syst.
iv. § 23. Petav. vi. 8, § 19, ibid. t. i. ii, 6, § 9.}.

The controversy did not terminate in the subject of Theism, but was
pursued by the heretical party into questions of Christian Theology.
The Manichees considered the Son and Spirit as necessary emanations
from the Father; erring, first, in their classing those Divine Persons
with intelligences confessedly imperfect and subservient; next, in
introducing a sort of materialism into their notion of the Deity. The
Eclectics on the other hand, maintained, by a strong figure, that the
Eternal Son originated from the Father at His own will; meaning
thereby, that the everlasting mystery, which constitutes the relation
between Father and Son, has no physical or material conditions, and is
such as becomes Him who is altogether Mind, and bound by no
laws, but those established by His own perfection as a first cause.
Thus Iamblichus calls the Son self-begotten \footnote{[\emph{autogonos}]. [Vide Ath. Tr. vol.
ii. p. 475.]}.

The discussion seems hardly to have entered farther into the
Ante-Nicene Church, than is implied in the above notice of it: though
some suppose that Justin and others referred the divine \emph{gennesis}
or \emph{generation} to the will of God. However, it is easy to see
that the ground was prepared for the introduction of a subtle and
irreverent question, whenever the theologizing Sophists should choose
to raise it. Accordingly, it was one of the first and principal
interrogations put to the Catholics by their Arian opponents, whether
the \emph{generation} of the Son was voluntary or not on the part of
the Father; their dilemma being, that Almighty God was subject to laws
external to Himself, if it were not voluntary, and that, if on the
other hand it was voluntary, the Son was in the number of things
created. But of this more in the next Section.

5.

The Word as internal or external to the Father; [\emph{logos
endiathetos}] and [\emph{prophorikos}] \footnote{[Vide Ath. Tr. vol. ii. pp. 340-342.]}:—One theory there was, adopted by several of the early
Fathers, which led them to speak of the Son's generation or birth as
resulting from the Father's will, and yet did not interfere with His
consubstantiality. Of the two titles ascribed in Scripture to our
Lord, that of the ``\emph{Word}'' expresses with peculiar
force His co-eternity in the One Almighty Father. On the other hand,
the title ``\emph{Son}'' has more distinct reference to
His derivation and ministrative office. A distinction resembling this
had already been applied by the Stoics to the Platonic Logos, which
they represented under two aspects, the [\emph{endiathetos}] and [\emph{prophorikos}],
that is, the internal Thought and Purpose of God, and its external
Manifestation, as if in words spoken. The terms were received among
Catholics; the ``Endiathetic'' standing for the Word, as hid
from everlasting in the bosom of the Father, while the ``Prophoric''
was the Son sent forth into the world, in apparent separation from
God, with His Father's name and attributes upon Him, and His Father's
will to perform \footnote{Burton, Bampt. Lect., note 91. Petav. vi.
1-3.}. This
contrast is acknowledged by Athanasius, Gregory Nyssen, Cyril, and
other Post-Nicene writers; nor can it be confuted, being Scriptural in
its doctrine, and merely expressed in philosophical language, found
ready for the purpose. But further, this change of state in the
Eternal Word, from repose to energetic manifestation, as it took place
at the creation, was called by them a \emph{gennesis}: and here too,
no blame attaches to them, for the expression is used in Scripture in
different senses, one of which appears to be the very signification
which they put on it, the mission of the Word to make and govern all
things. Such is the text in St. Paul, that He is ``the image of
the Invisible God, the First-born of every creature;'' such is His
title in St. John as ``the Beginning of the Creation of God.''
\footnote{Col. i. 15. Rev. iii. 14. Vide also Gen. i.
3. Heb. xi. 3. Eccl. xxiv. 3-9.} This \emph{gennesis}
or \emph{generation} was called also the ``going-forth,''
or ``condescension,'' of the Son, which may Scripturally
be ascribed to the will of the all-bountiful Father \footnote{[\emph{proeleusis, sunkatabasis}], Bull,
Defens. iii. 9. [Other writers support him in this view, as Maranus,
in Just. Tryph. 61, and in his work Divin. Jes. Christi, lib. iv. c.
6. Vide contr. Dissert. 3 and 4 in the Author's ``Theological
Tracts.'']}. However, there were some early writers who seem to interpret
the \emph{gennesis} in this meaning exclusively, ascribing the title
of ``\emph{Son}'' to our Lord only after the date of His
mission or economy, and considering that of the ``\emph{Word}''
as His peculiar appellation during the previous eternity \footnote{[Vide ``Theological Tracts,''
iii.]}. Nay, if we carry off their expressions hastily or perversely,
as some theologians have done, we shall perhaps conclude that they
conceived that God existed in One Person before the ``\emph{going-forth},''
and then, if it may be said, by a change in His nature began to exist
in a Second Person; as if an attribute (the Internal Word, ``\emph{Endiathetic},'')
had come into substantive being, as ``\emph{Prophoric}.'' The
Fathers, who have laid themselves open to this charge, are Athenagoras,
Tatian, Theophilus, Hippolytus, and Novatian, as mentioned in the
first Chapter.

Now that they did not mean what a superficial reader might lay to
their charge, may be argued, first, from the parallel language of the
Post-Nicenes, as mentioned above, whose orthodoxy no one questions.
Next, from the extreme absurdity, not to speak of the impiety, of the
doctrine imputed to them; as if, with a more than Gnostic
extravagance, they conceived that any change or extension could take
place in that Individual Essence, which is without parts or passions,
or that the divine \emph{generation} could be an event in time,
instead of being considered a mere expression of the eternal relation
of the Father towards the Son \footnote{[[\emph{oute arch\={e}n echei h\={e}
akatel\={e}ptos autou genn\={e}sis oute telos, anarch\={o}s,
akatapaust\={o}s}], \&c. Damasc. F. O. p. 8. Vide Ath. Tr.
vol. ii. pp. 350 and 108.]}. Indeed, the very absurdity of the literal sense of the words,
in whatever degree they so expressed themselves, was the mischief to
be apprehended from them. The reader, trying a rhetorical description
by too rigid a rule, would attempt to elicit sense by imputing a
heresy, and would conclude that they meant by the External or \emph{Prophoric}
Word a created being, made in the beginning of all things as the
visible emblem of the Internal or \emph{Endiathetic}, and the
instrument of God's purposes towards His creation. This is in fact the
Arian doctrine, which doubtless availed itself in its defence of the
declarations of incautious piety; or rather we have evidence of the
fact, that it did so avail itself, in the letter of Arius to
Alexander, and from the anathema of the Nicene Creed directed against
such as said that ``the Son was not before His \emph{gennesis}.''

Lastly, the orthodoxy of the five writers in question is
ascertained by a careful examination of the passages, which give
ground for the accusation. Two of these shall here be quoted without
comment. Theophilus then says, ``God having His own Word in His
womb, begat Him together with His Wisdom'' (that is, His Spirit),
``uttering them prior to the universe.'' ``He had this
Word as the Minister of His works, and did all things through Him ...
The prophets were not in existence when the world was made; but the
Wisdom of God, which is in Him, and His Holy Word, who is ever
present with Him.'' \footnote{[\emph{echon … ho theos ton heautou logon
endiatheton en tois idiois splanchnois, egenn\={e}sen outon meta t\={e}s
heautou sophias, exereuxamenos}] (Psalm xlv. 1), [\emph{pro \={o}n
hol\={o}n … ho aei sumpar\={o}n aut\={o}}].}
Elsewhere he speaks of ``the Word, eternally seated in the heart
of God;'' \footnote{[\emph{ton logon diapantos endiatheton en
kardiai theou}].}
``for,'' he presently adds, ``before anything was made, He
possessed this Counseller, as being His mind and providence. And when
He purposed to make all that He had deliberated on, He begat this Word
as external to Him, being the First-born antecedent to the whole
creation; not, however, Himself losing the Word'' (that is, the
Internal), ``but begetting it, and yet everlastingly communing
with it.'' \footnote{[\emph{egenn\={e}se prophorikon}].}

In like manner Hippolytus in his answer to Noetus:—``God was
alone, and there was no being coeval with Him, when He willed to
create the world … Not that He was destitute of reason (the Logos),
wisdom or counsel. They are all in Him, He was all. At the time and in
the manner He willed, He manifested His Word [Logos], ... through whom
He made all things ... Moreover He placed over them His Word, whom He
begat as His Counseller and Instrument; whom He had within Him,
invisible to creation, till He manifested Him, uttering the Word, and
begetting Light from Light ... And so Another stood by Him, not as if
there were two Gods, but as though Light from Light, or a ray from the
Sun.'' \footnote{Vide Bull. Defens. iii. 7, 8.}

And thus closes our survey of Catholic Ante-Nicene theology.

\xsection{Section 5. The Arian Heresy}

IT remains to give some account of the
heretical doctrine, which was first promulgated within the Church by
Arius. There have been attempts to attribute this heresy to Catholic
writers previous to his time; yet its contemporaries are express in
their testimony that he was the author of it, nor can anything be
adduced from the Ante-Nicene theology to countenance such an
imputation. Sozomen expressly says, that Arius was the first to
introduce into the Church the formulæ of the ``out of
nothing,'' and the ``once He was not,'' that is, the
creation and the non-eternity of the Son of God. Alexander and
Athanasius, who had the amplest means of information on the subject,
confirm his testimony \footnote{Soz. i. 15. Theod. His. i. 4. Athan. Decr. Nic. 27. de Sent.
Dionys. 6.}.
That the heresy existed before his time outside the Church, may be
true,—though little is known on the subject; and that there had been
certain speculators, such as Paulus of Samosata, who were simply
humanitarians, is undoubtedly true; but they did not hold the formal
doctrine of Arius, that an Angelic being had been exalted into a God.
However, he and his supporters, though they do not venture to
adduce in their favour the evidence of former Catholics, nevertheless
speak in a general way of their having received their doctrines from
others. Arius too himself appears to be only a partisan of the
Eusebians, and they in turn are referable to Lucian of Antioch, who
for some cause or other was at one time under excommunication. But
here we lose sight of the heresy; except that Origen assails a
doctrine, whose we know not \footnote{The [\emph{\={e}n pote hote ouk \={e}n}];
it might be Tertullian who was aimed at especially as St. Dionysius of
Rome denounces the doctrine also.]},
which bears a resemblance to it; nay, if we may trust Ruffinus, which
was expressed in the very same heterodox formulæ, which Sozomen
declares that Arius was the first to preach within the Church.

1.

Before detailing, however, the separate characteristics of his
heresy, it may be right briefly to confront it with such previous
doctrines, in and out of the Church, as may be considered to bear a
resemblance to it.

The fundamental tenet of Arianism was, that the Son of God was a
creature, not born of the Father, but, in the scientific language of
the times, made ``out of nothing.'' \footnote{[\emph{ex ouk ont\={o}n}]; hence the
Arians were called Exucontii.} It followed that He only possessed a super-angelic nature,
being made at God's good pleasure before the worlds, before time,
after the pattern of the attribute Logos or Wisdom, as existing in the
Divine Mind, gifted with the illumination of it, and in consequence
called after it the Word and the Wisdom, nay inheriting the
title itself of God; and at length united to a human body, in the
place of its soul, in the person of Jesus Christ.

1. This doctrine resembled that of the five philosophizing Fathers,
as described in the foregoing Section, so far as this, that it
identified the Son with the External or Prophoric Logos, spoke of the
Divine Logos Itself as if a mere internal attribute, and yet affected
to maintain a connexion between the Logos and the Son. Their doctrine
differed from it, inasmuch as they believed, that He who was the Son
had ever been in personal existence as the Logos in the Father's
bosom, whereas Arianism dated His personal existence from the time of
His manifestation.

2. It resembled the Eclectic theology, so far as to maintain that
the Son was by nature separate from and inferior to the Father; and
again, formed at the Father's will. It differed from Eclecticism, in
considering the Son to have a beginning of existence, whereas the
Platonists held Him, as they held the universe, to be an eternal
Emanation, and the Father's will to be a concomitant, not an
antecedent, of His \emph{gennesis}.

3. It agreed with the teaching of Gnostics and Manichees, in
maintaining the Son's essential inferiority to the Father: it
vehemently opposed them in their material notions of the Deity.

4. It concurred with the disciples of Paulus, in considering the
Intellectual and Ruling Principle in Christ, the Son of God, to be a
mere creature, by nature subject to a moral probation, as other men,
and exalted on the ground of His obedience, and gifted, moreover, with
a heavenly wisdom, called the Logos, which guided Him. The two
heresies also agreed, as the last words imply, in holding the Logos to
be an attribute or manifestation, not a Person \footnote{[When the Eternal Word, after the Nicene
Council, was defined to have a personal subsistence, then the
Samosatene doctrine would become identical with Nestorianism. Both
heresies came from Antioch.]}. Paulus considered it as if a voice or sound, which comes and
goes; so that God may be said to have \emph{spoken} in Christ. Arius
makes use of the same illustration: ``Many words speaketh
God,'' he says, ``which of them is manifested in the
flesh?'' \footnote{Athan. Decret. Nicen. 16.} He differs
from Paulus, in holding the pre-existence of the spiritual
intelligence in Christ, or the Son, whom he considers to be the first
and only creation of the Father's Hand, superangelic, and the God of
the Christian Economy.

5. Arianism agreed with the heresy of Sabellius, in teaching God to
exist only in one Person, and His true Logos to be an attribute,
manifested in the Son, who was a creature \footnote{Athan. Sent. Dionys. 25.}. It differed from Sabellianism, as regards the sense in which
the Logos was to be accounted as existing in Christ. The Sabellian,
lately a Patripassian, at least insisted much upon the formal and
abiding presence of the Logos in Him. The Arian, only partially
admitting the influence of the Divine Logos on that superangelic
nature, which was the Son, and which in Christ took the place of a
soul, nevertheless gave it the name of Logos, and maintained
accordingly that the incarnate Logos was not the true Wisdom and Word
of God, which was one with Him, but a created semblance of it.

6. Such is Arianism in its relations to the principal errors of its
time; and of these it was most opposed to the Gnostic and Sabellian,
which, as we shall see, it did not scruple to impute to its Catholic
adversaries. Towards the Catholics, on the other hand, it stood thus:
it was willing to ascribe to the Son all that is commonly attributed
to Almighty God, His name, authority, and power; all but the
incommunicable nature or being (\emph{usia}), that is, all but that
which alone could give Him a right to these prerogatives of divinity
in a real and literal sense. Now to turn to the arguments by which the
heresy defended itself, or rather, attacked the Church.

2.

1. Arius commenced his heresy thus, as Socrates informs us:—``(1)
If the Father gave birth to the Son, He who was born has an origin of
existence; (2) therefore once the Son was not; (3) therefore He is
created out of nothing.'' \footnote{Socr. i. 5. That is, the Son, as such, (1)
had [\emph{arch\={e}n hyparxe\={o}s}], (2) [\emph{\={e}n hote ouk
\={e}n}], (3) [\emph{ex ouk ont\={o}n echei t\={e}n
hypostasin}]. The argument thus stated in the history, answers to
the first three propositions anathematized at Nicæa, which are as
follows, the figures prefixed marking the correspondence of each with
Arius's theses, as set down by Socrates:—[\emph{tous legontas}] (2)
[\emph{hoti \={e}n pote hote ouk \={e}n}], (1) [\emph{kai prin genn\={e}th\={e}nai
ouk \={e}n}], (3) [\emph{kai hoti ex ouk ont\={o}n egeneto}],
(4) [\emph{\={e} ex heteras hypostase\={o}s \={e} ousias einai, \={e}
ktiston}], (5) [\emph{\={e} trepton \={e} alloi\={o}ton ton
huion tou theou, anathematizei h\={e} hagia katholik\={e} ekkl\={e}sia}].
[The fourth of these propositions is the denial of the [\emph{homoousion}].]
The last, viz. the mutability of the Son, was probably not one of
Arius's original propositions, but forced from him by his opponents as
a necessary consequence of his doctrine. He retracts it in his letters
to Eusebius and Alexander, who, on the other hand, bear testimony to
his having avowed it.}
It appears, then, that he inferred his doctrine from the very
meaning of the word ``\emph{Son},'' which is the designation
of our Lord in Scripture; and so far he adopted a fair and
unexceptionable mode of reasoning. Human relations, though the merest
shadows of ``heavenly things,'' yet would not of course be
employed by Divine Wisdom without fitness, nor unless with the
intention of instructing us. But what should be the exact instruction
derived by us from the word ``\emph{Son}'' is another question
\footnote{[Non recte faciunt, qui vim adhibent, ut sic
se habeat exemplum, ut prototypum. Non enim esset jam exemplum, nisi
haberet aliquid disisimili.'' Leont. Contr. Nest. i. p. 539, ed.
Canis.]}. The Catholics (not to
speak of their guidance from tradition in determining it) had taken
``\emph{Son}'' in its most obvious meaning; as interpreted
moreover by the title ``\emph{Only-begotten},'' and as
confirmed by the general tenor of Revelation. But the Arians selected
as the sense of the figure, that part of the original import of the
word, which, though undeniably included in it, when referred to us, is
at best what logicians call a \emph{property} deduced from the essence
or nature, not an element of its essential idea, and which was
especially out of place, when the word was used to express a truth
about the Divine Being. That a \emph{father} is prior to his \emph{son},
is not suggested, though it is implied, by the force of the terms, as
ordinarily used; and it is an inference altogether irrelevant, when
the inquiry has reference to that Being, from our notion of whom time
as well as space is necessarily excluded. It is fair, indeed, to
object at the outset to the word ``Father'' being applied at
all in its primary sense to the Supreme Being; but this was not the
Arian ground, which was to argue from, not against, the metaphor
employed. Nor was even this the extent of perverseness which their
argument evidences. Let it be observed, that they admitted the primary
sense of the word, \emph{in order} to introduce a mere secondary
sense, contending that, because our Lord was to be considered really
as a Son, therefore in fact He was no Son at all. In the first
proposition Arius assumes that He is really a Son, and argues as if He
were; in the third he has arrived at the conclusion that He was
created, that is, no Son at all, except in a secondary sense, as
having received from the Father a sort of \emph{adoption}. An attempt
was made by the Arians to smooth over their inconsistency, by adducing
passages of Scripture, in which the works of God are spoken of as
births,—as in the instance from Job, ``He giveth birth to the
drops of dew.'' But this is obviously an entirely new mode of
defending their theory of a divine adoption, and does not relieve
their original fault; which consisted in their arguing from an assumed
analogy, which the result of that argument destroyed. For, if He be
the Son of God, no otherwise than man is, that is, by adoption, what
becomes of the argument from the \emph{anterior} and \emph{posterior}
in existence \footnote{[That is, an adopted son is not necessarily
younger, but might be older, than the person adopting him.]}? as if the
notion of adoption, contained in it any necessary reference to the
nature and circumstances of the two parties between whom it takes
place.

2. Accordingly, the Arians were soon obliged to betake themselves
to a more refined argument. They dropped the consideration of time,
and withdrew the inference involving it, which they had drawn from the
literal sense of the word ``\emph{Son}.'' Instead of this,
they maintained that the relation of Father and Son, as such, in
whatever sense considered, could not but imply the notion of voluntary
originator, and on the other hand, of a free gift conferred; and that
the Son must be essentially inferior to Him, from whose will His
existence resulted. Their argument was conveyed in the form of a
dilemma:—``Whether the Father gave birth to the Son \emph{volens}
or \emph{nolens}?'' The Catholics wisely answered them by a
counter inquiry, which was adapted to silence, without countenancing,
the presumptuous disputant. Gregory of Nazianzus asked them,
``Whether the Father is God, \emph{volens} or \emph{nolens}?''
And Cyril of Alexandria, ``Whether He is good, compassionate,
merciful, and holy, with or against His choice? For, if He is so in
consequence of choosing it, and choice ever precedes what is chosen,
these attributes once did not exist in God.'' Athanasius gives
substantially the same answer, solving, however, rather than
confuting, the objection. ``The Arians,'' he says,
``direct their view to the \emph{contradictory} of willing,
instead of considering the more important and the \emph{previous}
question; for, as \emph{unwillingness} is opposed to \emph{willing},
so is \emph{nature} prior to willing, and leads the way to it.'' \footnote{Petav. ii. 5, § 9; vi. 8. 14. [``Generatio
non potestatis est, sed naturæ.'' Ambros. Incarn. 79. [\emph{H\={e}
genn\={e}sis physe\={o}s ergon, h\={e} de ktisis thel\={e}se\={o}s}].
Damasc. F. O. i. 8. p. 133.]}

3. Further:—the Arians attempted to draw their conclusion as to
the dissimilarity of the Father and the Son, from the divine attribute
of the ``Ingenerate'' (\emph{unborn} or increate), which, as I
have already said, was acknowledged on all hands to be the peculiar
attribute of the Father, while it had been the philosophical as
well as Valentinian appellation of the Supreme God. This was the chief
resource of the Anomœans, who revived the pure Arian heresy, some
years after the death of its first author. Their argument has been
expressed in the following form:—that ``it is the essence of the
Father to be \emph{ingenerate}, and of the Son to be \emph{generate};
but \emph{unborn} and \emph{born} cannot be the same.'' \footnote{Beausobre, Hist. Manich. iii. 7.} The shallowness, as well as the miserable trifling of such
disputations on a serious subject, renders them unworthy of a
refutation.

4. Moreover, they argued against the Catholic sense of the word
``\emph{Son},'' from what they conceived to be its \emph{materiality};
and, unwarrantably contrasting its primary with its figurative
signification, as if both could not be preserved, they contended that,
since the word must be figurative, therefore it could not retain its
primary sense, but must be taken in the secondary sense of \emph{adoption}.

5. Their reasonings (so to call them) had now conducted them thus
far:—to maintain that our Lord was a creature, advanced, after
creation, to be a Son of God. They did not shrink from the inference
which these positions implied, viz. that He had been put on trial as
other moral agents, and adopted on being found worthy; that His
holiness was not essential, but acquired.

6. It was next incumbent on them to explain in what sense our Lord
was the ``\emph{Only-begotten},'' since they refused to
understand that title in the Catholic sense of the \emph{Homoüsion}
or \emph{consubstantial}. Accordingly, while pronouncing the
divine birth to be a kind of creation, or an adoption, they attempted
to hide the offensiveness of the heretical doctrine by the variety and
dignity of the prerogatives, by which they distinguished the Son from
other creatures. They declared that He was, strictly speaking, the
only creature of God, as being alone made immediately by Him; and
hence He was called \emph{Only-begotten}, as ``born alone from Him
alone,'' \footnote{Pearson on the Creed, vol. ii. p. 148.
Suicer. Thes. verb. [\emph{monogen\={e}s}].} whereas
all others were made through Him, as the instrument of Divine Power;
and that in consequence He was ``a creature, \emph{but not} as
being \emph{one} of the creatures, a birth or production, but not as
being one of the produced;'' \footnote{[\emph{ktisma, all' ouch h\={o}s en t\={o}n
ktismat\={o}n; genn\={e}ma, all' ouch h\={o}s en t\={o}n
gegenn\={e}men\={o}n}].} that is, to express their sentiment with something of the same
ambiguity, ``He was not a creature like other creatures.''
Another ambiguity of language followed. The idea of time depending on
that of creation, they were able to grant that He, who was employed in
forming all things, therefore brought time itself into being, and was
``before all time;'' not granting thereby that He was
everlasting, but meaning that He was brought into existence
``timelessly,'' independent of that succession of second
causes (as they are called), that elementary system, seemingly
self-sustained and self-renovating, to the laws of which creation
itself may be considered as subjected.

7. Nor, lastly, had they any difficulty either in allowing or in
explaining away the other attributes of divinity ascribed to Christ in
Scripture. They might safely confess Him to be perfect God, one
with God, the object of worship, the author of good; still with the
reserve, that sacred appellations belonged to Him only in the same
general sense in which they are sometimes accidentally bestowed on the
faithful servants of God, and without interfering with the
prerogatives of the One, Eternal, Self-existing Cause of all things \footnote{It may be added that the chief texts,
which the Arians adduced in controversy were, Prov. viii. 22. Matt.
xix. 17; xx. 23. Mark xiii. 32. John v. 19; xiv. 28. 1 Cor. xv. 28.
Col. i. 15; and others which refer to our Lord's mediatorial office (Petav.
ii. 1, \&c. Theod. Hist. i. 14). But it is obvious, that the
strength of their cause did not lie in the text of Scripture.}.

3.

This account of the Arian theology may be suitably illustrated by
some of the original documents of the controversy. Here, then, shall
follow two letters of Arius himself, an extract from his Thalia, a
letter of Eusebius of Nicomedia, and parts of the encyclical Epistle
of Alexander of Alexandria, in justification of his excommunication of
Arius and his followers \footnote{Theodor. Hist. i. 4-6. Socr. i. 6. Athan.
in Arian. i. 5. Synod 15, 16. Epiphan. Hær. lxix. 6, 7. Hilar. Trin.
iv. 12; vi. 5.}.

1. ``To his most dear Lord, Eusebius, a man of God, faithful
and orthodox, Arius, the man unjustly persecuted by the Pope Alexander
for the all-conquering truth's sake, of which thou too art a champion,
sends health in the Lord. As Ammonius, my father, was going to
Nicomedia, it seemed becoming to address this through him; and withal
to represent to that deep-seated affection which thou bearest towards
the brethren for the sake of God and His Christ, how fiercely
the bishop assaults and drives us, leaving no means untried in his
opposition. At length he has driven us out of the city, as men without
God, for dissenting from his public declarations, that, 'As God is
eternal, so is His Son: where the Father, there the Son; the Son
co-exists in God without a beginning (or birth): ever generate, an
ingenerately-generate; that neither in idea, nor by an instant of
time, does God precede the Son; an eternal God, an eternal Son; the
Son is from God Himself.' Since then, Eusebius, thy brother of
Cæsarea, Theodotus, Paulinus, \&c. ... and all the Bishops of the
East declare that God exists without origin before the Son, they are
made anathema by Alexander's sentence; all but Philogonius, Hellanicus,
and Macarius, heretical, ill-grounded men, who say, one that He is an
utterance, another an offspring, another co-ingenerate. These
blasphemies we cannot bear even to hear; no, not if the heretics
should threaten us with ten thousand deaths. What, on the other hand,
are our statements and opinions, our past and present teaching? that
the Son is not ingenerate, nor in any way a part of the Ingenerate,
nor made of any subject-matter \footnote{The Greek of most of these scientific
expressions has been given; of the rest it is as follows:—men
without God, [\emph{atheous}]; without a beginning or birth, [\emph{agenn\={e}t\={o}s}];
ever-generate, [\emph{aeigen\={e}s}]; ingenerately-generate, [\emph{agenn\={e}togen\={e}s}];
an utterance, [\emph{erug\={e}}] (Psalm xlv. 1); offspring, [\emph{probol\={e}}];
co-ingenerate, [\emph{sunagenn\={e}ton}]; of any subject-matter, [\emph{ex
hypokeimenou tinos}].}; but that, by the will and counsel of God, He subsisted before
times and ages, perfect God, Only-begotten, unchangeable; and that
before this generation, or creation, or determination, or
establishment \footnote{These words are selected by Arius, as
being found in Scripture; [Vide Heb. i. 5. Rom. i. 4. Prov. viii. 22,
23.]}, He was
not, for He is not ingenerate. And we are persecuted for saying, The
Son has an origin, but God is unoriginate; for this we are under
persecution, and for saying that He is out of nothing, inasmuch as He
is neither part of God, nor of any subject-matter. Therefore we are
persecuted; the rest thou knowest. I pray that thou be strong in the
Lord, remembering our afflictions, fellow-Lucianist, truly named
Eusebius.'' \footnote{[i.e. the pious, or rather, the orthodox.]}

2. The second letter is written in the name of himself and his
partisans of the Alexandrian Church; who, finding themselves
excommunicated, had withdrawn to Asia, where they had a field for
propagating their opinions. It was composed under the direction of
Eusebius of Nicomedia, and is far more temperate and cautious than the
former.

``To Alexander, our blessed Pope and Bishop, the Priests and
Deacons send health in the Lord. Our hereditary faith, which thou too,
blessed Pope, hast taught us, is this:—We believe in One God, alone
ingenerate, alone everlasting, alone unoriginate, alone truly God,
alone immortal, alone wise, alone good, alone sovereign, alone judge
of all, ordainer, and dispenser, unchangeable and unalterable, just
and good, of the Law and the Prophets, and of the New Covenant. We
believe that this God gave birth to the Only-begotten Son before
age-long times, through whom He has made those ages themselves, and
all things else; that He generated Him, not in semblance, but in
truth, giving Him a real subsistence (or \emph{hypostasis}), at His
own will, so as to be unchangeable and unalterable, God's perfect
creature, but not as other creatures, His production, but not as other
productions; nor as Valentinus maintained, an offspring (\emph{probole});
nor again, as Manichæus, a consubstantial part; nor, as Sabellius, a
Son-Father, which is to make two out of one; nor, as Hieracas, one
torch from another, or a flame divided into two; nor, as if He were
previously in being, and afterwards generated or created again to be a
Son, a notion condemned by thyself, blessed Pope, in full Church and
among the assembled Clergy; but, as we affirm, created at the will of
God before times and before ages, and having life and being from the
Father, who gave subsistence as to Him, so to His glorious
perfections. For, when the Father gave to Him the inheritance of all
things, He did not thereby deprive Himself of attributes, which are
His ingenerately, who is the Source of all things.

``So there are Three Subsistences (or Persons); and, whereas
God is the Cause of all things, and therefore unoriginate simply by
Himself, the Son on the other hand, born of the Father time-apart, and
created and established before all periods, did not exist before He
was born, but being born of the Father time-apart, was brought into
substantive existence (subsistence), He alone by the Father alone. For
He is not eternal, or co-eternal, or co-ingenerate with the Father;
nor hath an existence together with the Father, as if there were two
ingenerate Origins; but God is before all things, as being a Monad,
and the Origin of all;—and therefore before the Son also, as indeed
we have learned from thee in thy public preaching. Inasmuch then
as it is from God that He hath His being, and His glorious
perfections, and His life, and His charge of all things, for this
reason God is His Origin, as being His God and before Him. As to such
phrases as 'from Him,' and 'from the womb,' and 'issued forth from the
Father, and am come,' if they be understood, as they are by some, to
denote a part of the consubstantial, and a \emph{probole} (offspring),
then the Father will be of a compound nature, and divisible, and
changeable, and corporeal; and thus, as far as their words go, the
incorporeal God will be subjected to the properties of matter. I pray
for thy health in the Lord, blessed Pope.'' \footnote{Before age-long periods, [\emph{pro chron\={o}n
ai\={o}ni\={o}n}]; giving Him a real subsistence, [\emph{hypost\={e}santa}];
Son-Father, [\emph{huiopatora}] [Vide Ath. Tr. p. 97, k and p. 514, o;
also Didym. de Trin. iii. 18]; gave subsistence, as to Him, so to His
glorious perfections, [\emph{tas doxas sunupost\={e}santos aut\={o}i}];
Three Subsistences, [\emph{treis hypostaseis}]; born time-apart [\emph{achron\={o}s
genn\={e}theis}]; of a compound nature, [\emph{sunthetos}]. The texts
to which Arius refers are Ps. cx. 3, and John xvi. 28.}

3. About the same time Arius wrote his Thalia, or song for banquets
and merry-makings, from which the following is extracted. He begins
thus:—``According to the faith of God's elect, who know God,
holy children, sound in their creed, gifted with the Holy Spirit of
God, I have received these things from the partakers of wisdom,
accomplished, taught of God, and altogether wise. Along their track I
have pursued my course with like opinions,—I, the famous among men,
the much-suffering for God's glory; and, taught of God, I have gained
wisdom and knowledge.'' After this exordium, he proceeds to
declare, ``that God made the Son the origin (or beginning) of
creation, being Himself unoriginate, and adopted Him to be His
Son; who, on the other hand, has no property of divinity in His own \emph{Hypostasis},
not being equal, nor consubstantial with Him; that God is invisible,
not only to the creatures created through the Son, but to the Son
Himself; that there is a Trinity, but not with an equal glory, the \emph{Hypostases}
being incommunicable with each other; One infinitely more glorious
than the other; that the Father is foreign in substance to the Son, as
existing unoriginate; that by God's will the Son became Wisdom, Power,
the Spirit, the Truth, the Word, the Glory, and the Image of God; that
the Father, as being Almighty, is able to give existence to a being
equal to the Son, though not superior to Him; that, from the time that
He was made, being a mighty God, He has hymned the praises of His
Superior; that He cannot investigate His Father's nature, it being
plain that the originated cannot comprehend the unoriginate; nay, that
He does not know His own.'' \footnote{Incommunicable, [\emph{anepimiktoi}],
(this is in opposition to the [\emph{perich\={o}r\={e}sis}], or
co-inherence); foreign in substance [\emph{xenos kat' ousian}];
investigate, [\emph{exichniasai}].}

4. On the receipt of the letter from Arius, which was the first
document here exhibited, Eusebius of Nicomedia addressed a letter to
Paulinus of Tyre, of which the following is an extract:—``We
have neither heard of two Ingenerates, nor of One divided into two,
subjected to any material affection; but of One Ingenerate, and one
generated by Him really; not from His substance, not partaking of the
nature of the Ingenerate at all, but made altogether other than He in
nature and in power, though made after the perfect likeness of
the character and excellence of His Maker ... But, if He were of Him
in the sense of 'from Him,' as if a part of Him, or from the effluence
of His substance \footnote{Generated, [\emph{gegonos}]; effluence of
His substance, [\emph{ex aporrhoias t\={e}s ousias}]; being from
the Ingenerate, [\emph{ek tou agenn\={e}tou hyparchon}].}, He
would not be spoken of (in Scripture) as created or established ...
for what exists as being from the Ingenerate ceases to be created or
established, as being from its origin ingenerate. But, if His being
called generate suggests the idea that He is made out of the Father's
substance, and has from Him a sameness of nature, we know that not of
Him alone does Scripture use the word 'generate,' but also of things
altogether unlike the Father in nature. For it says of men, 'I have
begotten sons and exalted them, and they have set Me at nought;' and,
'Thou hast left the God who begat thee;' and in other instances, as
'Who has given birth to the drops of dew?' ... Nothing is of His
substance; but all things are made at His will.''

5. Alexander, in his public accusation of Arius and his party to
Alexander of Constantinople, writes thus:—``They say that once
the Son of God was not, and that He, who before had no existence, was
at length made, made such, when He was made, as any other man is by
nature. Numbering the Son of God among created things, they are but
consistent in adding that He is of an alterable nature, capable of
virtue and vice ... When it is urged on them that the Saviour differs
from others, called sons of God, by the unchangeableness of His
nature, stripping off all reverence, they answer, that God, foreknowing and foreseeing His obedience, chose Him out of all
creatures; chose Him, I say, not as possessing aught by nature and
prerogative above the others (since, as they say, there is no Son of
God by nature), nor bearing any peculiar relation towards God; but, as
being, as well as others, of an alterable nature, and preserved from
falling by the pursuit and exercise of virtuous conduct; so that, if
Paul or Peter had made such strenuous progress, they would have gained
a sonship equal to His.''

In another letter, which was addressed to the Churches, he says,
``It is their doctrine, that 'God was not always a Father', that
'the Word of God has not always existed, but was made out of nothing;
for the self-existing God made Him, who once was not, out of what once
was not ... Neither is He like the Father in substance, nor is He the
true and natural Logos of the Father, nor His true Wisdom, but one of
His works and creatures; and He is catachrestically the Word and
Wisdom, inasmuch as He Himself was made by the proper Logos of God,
and by that Wisdom which is in God, by which God made all things, and
Him in the number. Hence He is mutable and alterable by nature, as
other rational beings; and He is foreign and external to God's
substance, being excluded from it. He was made for our sakes, in order
that God might create us by Him as by an instrument; and He would not
have had subsistence, had not God willed our making.' Some one asked
them, if the Word of God could change, as the devil changed? They
scrupled not to answer, 'Certainly, He can.''' \footnote{Like in substance, [\emph{homoios kat' ousian}]
[This, as we shall see afterwards, in the Homϟsian, the symbol of
the Eusebians or Semi-Arians], mutable and alterable, [\emph{treptos kai
alloi\={o}tos}]; excluded, [\emph{apeschoinis menos}].}

4.

More than enough has now been said in explanation of a controversy,
the very sound of which must be painful to any one who has a loving
faith in the Divinity of the Son. Yet so it has been ordered, that He
who was once lifted up to the gaze of the world, and hid not His face
from contumely, has again been subjected to rude scrutiny and
dishonour in the promulgation of His religion to the world. And His
true followers have been themselves obliged in His defence to raise
and fix their eyes boldly on Him, as if He were one of themselves,
dismissing the natural reverence, which would keep them ever at His
feet. The subject may be dismissed with the following remarks:—

1. First, it is obvious to notice the unscriptural character of the
arguments on which the heresy was founded. It is true that the Arians
did not neglect to support their case from such detached portions of
the Inspired Volume as suited their purpose; but still it can never be
said that they showed that earnest desire of sacred truth, and careful
search into its documents, which alone mark the Christian inquirer.
The question is not merely whether they confined themselves to the
language of Scripture, but whether they began with the study of it.
Doubtless, to forbid in controversy the use of all words but those
which actually occur in Scripture, is a superstition, an encroachment
on Scripture liberty, and an impediment to freedom of thought;
and especially unreasonable, considering that a traditional system of
theology, consistent with, but independent of, Scripture, has existed
in the Church from the Apostolic age. ``Why art thou in that
excessive slavery to the letter,'' says Gregory Nazianzen,
``and employest a Judaical wisdom, dwelling upon syllables, while
letting slip realities? Suppose, on thy saying twice five, or twice
seven, I were to understand thence ten or fourteen; or, if I spoke of
a man, when thou hadst named an animal rational and mortal, should I
in that case appear to thee to trifle? How could I so appear, in
merely expressing your own meaning?'' \footnote{Petav. iv. 5, § 6. [Athanasius ever
exalts the theological sense over the words, whether sacred or
ecclesiastical, which are its vehicle, and this even to the apparent
withholding of the symbol [\emph{homoousion}]. Vide Orat. ii. 3, and
Ath. Tr. vol. i., notes pp. 163, 212, 214, 231, \&c.]} But, inasmuch as this liberty was an evangelical privilege,
which might be allowed to the Arian disputants, on the other hand it
was a dangerous privilege also, ever to be subjected to a profound
respect for the sacred text, a cautious adherence to the whole of the
doctrine therein contained, and a regard also for those received
statements, which, though not given to us as inspired, probably are
derived from inspired teachers. Now the most liberal admission which
can be made in behalf of the Arians, is, to grant that they did not in
controversy throw aside the authority of Scripture altogether; that
is, proclaim themselves unbelievers; for it is evident that they took
only just so much of it as would afford them a basis for erecting
their system of heresy by an abstract logical process. The mere words ``Father and Son,'' ``birth,''
``origin,'' \&c., were all that they postulated of revealed
authority for their argument; they professed to do all the rest for
themselves. The meaning of these terms in their context, the
illustration which they afford to each other, and, much more, the
divine doctrine considered as one undivided message, variously
exhibited and dispersed in the various parts of Scripture, were
excluded from the consideration of controversialists, who thought that
truth was gained by disputing instead of investigating.

2. Next, it will be observed that, throughout their discussions,
they assumed as an axiom, that there could be no mystery in the
Scripture doctrine respecting the nature of God. In this, indeed, they
did but follow the example of the contemporary spurious theologies;
though their abstract mode of reasoning from the mere force of one or
two Scripture terms, necessarily forced them more than other heretics
into the use and avowal of the principle. The Sabellian, to avoid
mystery, denied the distinction of Persons in the Divine Nature.
Paulus, and afterwards Apollinaris, for the same reason, denied the
existence of two Intelligent Principles at once, the Word and the
human soul, in the Person of Christ. The Arians adopted both errors.
Yet what is a mystery in doctrine, but a difficulty or inconsistency
in the intellectual expression of it? And what reason is there for
supposing, that Revelation addresses itself to the intellect, except
so far as intellect is necessary for conveying and fixing its truths
on the heart? Why are we not content to take and use what is given us,
without asking questions? The Catholics, on the other hand,
pursued the intellectual investigation of the doctrine, under the
guidance of Scripture and Tradition, merely as far as some immediate
necessity called for it; and cared little, though one mode of
expression seemed inconsistent with another. Thus, they developed the
notion of ``\emph{substance}'' against the Pantheists, of the
``\emph{Hypostatic Word}'' against the Sabellians, of the
``\emph{Internal Word}'' to meet the imputation of Ditheism;
still they did not use these formulæ for any thing beyond shadows of
sacred truth, symbols witnessing against the speculations into which
the unbridled intellect fell.

Accordingly, they were for a time inconsistent with each other in
the minor particulars of their doctrinal statements, being far more
bent on opposing error, than on forming a theology:—inconsistent,
that is, before the experience of controversy and the voice of
tradition had detached them from less accurate or advisable
expressions, and made them correct, or at least compare and adjust
their several declarations. Thus, some said that there was but one \emph{hypostasis},
meaning \emph{substance}, in God; others three \emph{hypostases},
meaning \emph{Subsistences} or Persons; and some spoke of one \emph{usia},
meaning substance, while others spoke of more than one \emph{usia}.
Some allowed, some rejected, the terms \emph{probole} and \emph{homoüsion},
according as they were guided by the prevailing heresy of the day, and
by their own judgment how best to meet it. Some spoke of the Son as
existing from everlasting in the Divine Mind; others implied that the
Logos was everlasting, and became the Son in time. Some asserted that
He was unoriginate, others denied it. Some, when interrogated by
heretics, taught that He was born of the Father at the Father's
will; others, from His nature, not His will; others, neither with His
willing nor not willing \footnote{Justin, Tryph. 61, 100, etc. Petav. vi. 8,
§ 14, 15, 18.}.
Some declared that God was in number Three; others, that He was
numerically One; while to others it perhaps appeared more
philosophical to exclude the idea of number altogether, in discussions
about that Mysterious Nature, which is beyond comparison with itself,
whether viewed as Three or One, and neither falls under nor involves
any conceivable species \footnote{Petav. iv. 13.}.

In all these various statements, the object is clear and
unexceptionable, being merely that of protesting and practically
guarding against dangerous deductions from the Scripture doctrine; and
the problem implied in all of them is, to determine how this end may
best be effected. There are no signs of an intellectual curiosity in
the tenor of these Catholic expositions, prying into things not seen
as yet; nor of an ambition to account for the representations of the
truth given us in the sacred writings. But such a temper is the very
characteristic of the Arian disputants. They insisted on taking the
terms of Scripture and of the Church for more than they signified, and
expected their opponents to admit inferences altogether foreign to the
theological sense in which they were really used. Hence, they
sometimes accused the orthodox of heresy, sometimes of
self-contradiction. The Fathers of the Church have come down to us
loaded with the imputation of the strangest errors, merely because
they united truths, which heresies only shared among themselves;
nor have writers been wanting in modern times, from malevolence or
carelessness, to aggravate these charges. The mystery of their creed
has been converted into an evidence of concurrent heresies. To believe
in the actual Incarnation of the Eternal Wisdom, has been treated, not
as orthodoxy, but as an Ariano-Sabellianism \footnote{[``Eorum error veritati testimonium
dicit, et in consona perfidorum sententia in unum recte fidei modulum
concinunt.'' Vigil. Thaps. contr. Eut. ii. init.]}. To believe that the Son of God was the Logos, was
Sabellianism; to believe that the pre-existent Logos was the Son of
God, was Valentinianism. Gregory of Neo-Cæsarea was called a
Sabellian, because he spoke of one substance in the Divine Nature; he
was called a forerunner of Arius, because he said that Christ was a
creature. Origen, so frequently accused of Arianism, seemed to be a
Sabellian, when he said that the Son was the Auto-aletheia, the
Archetypal Truth. Athenagoras is charged with Sabellianism by the very
writer (Petau), whose general theory it is that he was one of those
Platonizing Fathers who anticipated Arius \footnote{Bull, Defens. iii. 5. § 4.}. Alexander, who at the opening of the controversy, was accused
by Arius of Sabellianizing, has in these latter times been detected by
the flippant Jortin to be an advocate of Semi-Arianism \footnote{Jortin, Eccles. Hist. vol. ii. pp. 179,
180.}, which was the peculiar enemy and assailant of Sabellianism in
all its forms. The celebrated word, \emph{homoüsion}, has not escaped
a similar contrariety of charges. Arius himself ascribes it to the
Manichees; the Semi-Arians at Ancyra anathematize it, as Sabellian. It
is in the same spirit that Arius, in his letter to Eusebius,
scoffs at the ``eternal birth,'' and the ``ingenerate
generation,'' as ascribed to the Son in the orthodox theology; as
if the inconsistency, which the words involved, when taken in their
full sense, were a sufficient refutation of the heavenly truth, of
which they are, each in its place, the partial and relative
expression.

The Catholics sustained these charges with a prudence, which has
(humanly speaking) secured the success of their cause, though it has
availed little to remove the calumnies heaped upon themselves. The
great Dionysius, who has himself been defamed by the ``accuser of
the brethren,'' declares perspicuously the principle of the
orthodox teaching. ``The particular expressions which I have
used,'' he says, in his defence, ``cannot be taken separate
from each other ... whereas my opponents have taken two bald words of
mine, and sling them at me from a distance; not understanding, that,
in the case of subjects, partially known, illustrations foreign to
them in nature, nay, inconsistent with each other, aid the
inquiry.'' \footnote{Athan. de Sent. Dionys. 18.}

However, the Catholics of course considered it a duty to remove, as
far as they could, their own verbal inconsistencies, and to sanction
one form of expression, as orthodox in each case, among the many which
might be adopted. Hence distinctions were made between the \emph{unborn}
and \emph{unmade}, \emph{origin} and \emph{cause}, as already noticed.
But these, clear and intelligible as they were in themselves, and
valuable, both as facilitating the argument and disabusing the
perplexed inquirer, opened to the heretical party the opportunity
of a new misrepresentation. Whenever the orthodox writers showed
an anxiety to reconcile and discriminate their own expressions, the
charge of Manicheism was urged against them; as if to dwell upon, were
to rest in the material images which were the signs of the unknown
truths. Thus the phrase, ``Light of Light,'' the orthodox and
almost apostolic emblem of the derivation of the Son from the Father,
as symbolizing Their inseparability, mutual relation, and the separate
fulness and exact parallelism and unity of Their perfections, was
interpreted by the gross conceptions of the Manichæan Hieracas \footnote{The [\emph{ek Theou}] was treated thus: [\emph{ei
gar ek Theou esti, kai egenn\={e}sen ex hautou ho Theos, h\={o}s
eipein, ex idias hypostase\={o}s phusei \={e} ek t\={e}s
idias ousias, oukoun \={o}nk\={o}th\={e}, \={e} tom\={e}n
edexato \={e} en t\={o}i gennain eplatunth\={e}, \={e}
sunestal\={e}, \={e} ti t\={o}n kata ta path\={e} ta s\={o}matika
hupest\={e}}]. Epiph. Hær. lxix. 15. Or, to take the objection
made at Nicæa to the [\emph{homoousion}] by Eusebius and some others:
[\emph{epei gar ephasan homoousion einai, o ek tinos estin, \={e} kata
merismon, \={e} kata rheusin, \={e} kata probol\={e}n; kata probol\={e}n
men, h\={o}s
ek rhiz\={o}n blast\={e}ma, kata de rheusin, h\={o}s hoi
patrikoi paides, kata merismon de, h\={o}s b\={o}lou chrusides
duo \={e} treis; kat' ouden de tout\={o}n estin ho Huios, dia
touto ou' sunkatatithesthai t\={e} pistei elogon}]. Socr. i. 8.}.

3. When in answer to such objections the Catholics denied that they
attached other than a figurative meaning to their words, their
opponents suddenly turned round, and professed the figurative meaning
of the terms to be that which they themselves advocated. This
inconsistency in their mode of conducting the argument deserves
notice. It has already been instanced in the original argument of
Arius, who maintained, that, since the word Son in its literal sense
included among other ideas that of a beginning of being, the Son of
God had had a beginning or was created, and therefore was not
really a Son of God at all. It was on account of such unscrupulous
dexterity in the controversy, that Alexander and Athanasius give them
the title of chameleons. ``They are as variable and uncertain in
their opinions,'' (says the latter,) ``as chameleons in their
colour. When refuted, they look confused, and when examined they are
perplexed; however, at length they recover their assurance, and bring
forward some evasion. Then, if this in turn is exposed, they do not
rest till they have devised some new absurdity, and, as Scripture
says, meditate vain things, so that they may secure the privilege of
being profane.'' \footnote{Athan. de Decr. Nic. 1. Socr. i. 6. [Vide
Ath. Tr. vol. ii. p. 71.]}

Let us, however, pursue the Arians on their new ground of allegory.
It has been already observed, that they explain the word \emph{Only-begotten}
in the sense of \emph{only-created}; and considered the oneness of the
Father and Son to consist in an unity of character and will, such as
exists between God and His Saints, not in nature.

Now, surely, the temper of mind, which had recourse to such a
comparison between Christ and us, to defend a heresy, was still more
odious, if possible, than the original impiety of the heresy itself.
Thus, the honours graciously bestowed upon human nature, as well as
the condescending self-abasement of our Lord, were made to subserve
the cause of the blasphemer. It is a known peculiarity of the message
of mercy, that it views the Church of Christ as if clothed with, or
hidden within, the glory of Him who ransomed it; so that there is no
name or title belonging to Him literally, which is not in a secondary
sense applied to the reconciled penitent. As our Lord is the
Priest and King of His redeemed, they, as members of Him, are
accounted kings and priests also. They are said to be Christs, or the
anointed, to partake of the Divine Nature, to be the well-beloved of
God, His sons, one with Him, and heirs of glory; in order to express
the fulness and the transcendent excellence of the blessings gained to
the Saints by Christ. In all these forms of speech, no religious mind
runs the risk of confusing its own privileges with the real
prerogatives of Him who gave them; yet it is obviously difficult in
argument to discriminate between the primary and secondary use of the
words, and to elicit and exhibit the delicate reasons lying in the
context of Scripture for conclusions, which the common sense of a
Christian is impatient as well as shocked to hear disputed. Who would
so trifle with words, to take a parallel case, as to argue that,
because Christians are said by St. John to ``know all
things,'' that therefore God is not omniscient in a sense
infinitely above man's highest intelligence?

It may be observed, moreover, that the Arians were inconsistent in
their application of the allegorical rule, by which they attempted to
interpret Scripture; and showed as great deficiency in their
philosophical conceptions of God, as in their practical devotion to
Him. They seem to have fancied that some of His acts were more
comprehensible than others, and might accordingly be made the basis on
which the rest might be interpreted. They referred the divine \emph{gennesis}
or \emph{generation} to the notion of creation; but creation is in
fact as mysterious as the divine \emph{gennesis}; that is, we are as
little able to understand our own words, when we speak of the
world's being brought out of nothing at God's word, as when we confess
that His Eternal Perfections are reiterated, without being doubled, in
the person of His Son. ``How is it,'' asks Athanasius,
``that the impious men dare to speak flippantly on subjects too
sacred to approach, mortals as they are, and incapable of explaining
even God's works upon earth? Why do I say, His earthly works? Let them
treat of themselves, if so be they can investigate their own nature;
yet venturous and self-confident, they tremble not before the glory of
God, which Angels are fain reverently to look into, though in nature
and rank far more excellent than they.'' \footnote{Athan. on Matt. xi. 22. § 6.} Accordingly, he argues that nothing is gained by resolving one
of the divine operations into another; that to make, when attributed
to God, is essentially distinct from the same act when ascribed to
man, as incomprehensible as to give birth or beget \footnote{Athan. de Decr. Nic. 11; vide also Greg.
Naz. Orat. 35. p. 566. Euseb. Eccl. Theol. i. 12.}; and consequently that it is our highest wisdom to take the
truths of Scripture as we find them there, and use them for the
purposes for which they are vouchsafed, without proceeding accurately
to systematize them or to explain them away. Far from elucidating, we
are evidently enfeebling the revealed doctrine, by substituting \emph{only-created}
for \emph{only-begotten}; for if the words are synonymous, why should
the latter be insisted on in Scripture? Accordingly, it is proper to
make a distinction between the primary and the literal meaning of a
term. All the terms which human language applies to the Supreme Being,
may perhaps be more or less figurative; but their primary and
secondary meaning may still remain as distinct, as when they are
referred to earthly objects. We need not give up the primary meaning
of the word \emph{Son} as opposed to the secondary sense of adoption,
because we forbear to use it in its literal and material sense.

4. This being the general character of the Arian reasonings, it is
natural to inquire what was the object towards which they tended. Now
it will be found, that this audacious and elaborate sophistry could
not escape one of two conclusions:—the establishment either of a
sort of ditheism, or, as the more practical alternative, of a mere
humanitarianism as regards our Lord; either a heresy tending to
paganism, or the virtual atheism of philosophy. If the professions of
the Arians are to be believed, they confessed our Lord to be God, God
in all respects \footnote{[\emph{pl\={e}r\={e}s Theos}].}, full
and perfect, yet at the same time to be infinitely distant from the
perfections of the One Eternal Cause. Here at once they are committed
to a ditheism; but Athanasius drives them on to the extreme of
polytheism. ``If,'' he says, ``the Son were an object of
worship for His transcendent glory, then every subordinate being is
bound to worship his superior.'' \footnote{Cudw. Intell. Syst. 4, § 36. Petav. ii.
12, § 6.} But so repulsive is the notion of a secondary God both to
reason, and much more to Christianity, that the real tendency of
Arianism lay towards the sole remaining alternative, the humanitarian
doctrine.—Its essential agreement with the heresy of Paulus has
already been incidentally shown; it differed from it only when the
pressure of controversy required it. Its history is the proof of this. It started with a boldness not inferior to that of Paulus; but
as soon as it was attacked, it suddenly coiled itself into a defensive
posture, and plunged amid the thickets of verbal controversy. At first
it had not scrupled to admit the peccable nature of the Son; but it
soon learned to disguise such consequences of its doctrine, and avowed
that, in matter of fact, He was indefectible. Next it borrowed the
language of Platonism, which, without committing it to any real
renunciation of its former declarations, admitted of the dress of a
high and almost enthusiastic piety. Then it professed an entire
agreement with the Catholics, except as to the adoption of the single
word \emph{consubstantial}, which they urged upon it, and concerning
which, it affected to entertain conscientious scruples. At this time
it was ready to confess that our Lord was the true God, God of God,
born time-apart, or before all time, and not a creature as other
creatures, but peculiarly the Son of God, and His accurate Image.
Afterwards, changing its ground, it protested, as we shall see,
against non-scriptural expressions, of which itself had been the chief
inventor; and proposed an union of all opinions, on the comprehensive
basis of a creed, in which the Son should be merely declared to be
``\emph{in all things like the Father},'' or simply ``\emph{like
Him}.'' This versatility of profession is an illustration of
the character given of the Arians by Athanasius, some pages back,
which is further exemplified in their conduct at the Council in which
they were condemned; but it is here adduced to show the danger to
which the Church was exposed from a party who had no fixed tenet,
except that of opposition to the true notion of Christ's divinity; and
whose teaching, accordingly, had no firm footing of internal
consistency to rest upon, till it descended to the notion of His
simple humanity, that is, to the doctrine of Artemas and Paulus,
though they too, as well as Arius, had enveloped their impieties in
such admissions and professions, as assimilated it more or less in
appearance to the Faith of the Catholic Church.

The conduct of the Arians at Nicæa, as referred to, was as
follows. ``When the Bishops in Council assembled,'' says
Athanasius, an eye-witness, ``were desirous of ridding the Church
of the impious expressions invented by Arius, '\emph{the Son is out of
nothing},' '\emph{is a creature},' '\emph{once was not},' 'of an \emph{alterable
nature},' and perpetuating those which we receive on the authority
of Scripture, that the Son is the Only-begotten of God by nature, the
Word, Power, the sole Wisdom of the Father, very God, as the Apostle
John says, and as Paul, the Radiance of His glory, and the express
Image of His Person; the Eusebians, influenced by their own
heterodoxy, said one to another, 'Let us agree to this; for we too are
of God, there being one God, \emph{of whom} are all things.' … The
Bishops, however, discerning their cunning, and the artifice adopted
by their impiety, in order to express more clearly the '\emph{of God},'
wrote down 'of God's \emph{substance},' creatures being said to be 'of
God,' as not existing of themselves without cause, but having an
origin of their production; but the Son being peculiarly of the
substance of the Father ... Again, on the Bishops asking the few
advocates of Arianism present, whether they allowed the Son to be, not
a creature, but the sole Power, Wisdom, and Image, eternal and
in all respects \footnote{[\emph{aparallakton}].}, of
the Father, and very God, the followers of Eusebius were detected
making signs to each other, to express that this also could be applied
to ourselves. 'For we too,' they said, 'are called in Scripture the
image and glory of God; we are said to live always ... There are many
powers; the locust is called in Scripture ``a great power.''
Nay, that we are God's own sons, is proved expressly from the text, in
which the Son calls us brethren. Nor does their assertion, that He is
very (true) God, distress us; He is very God, because He was made
such.' This was the unprincipled meaning of the Arians. But here too
the Bishops, seeing through their deceit, brought together from
Scripture, the radiance, source and stream, express Image of Person,
'In Thy Light we shall see light,' 'I and the Father are one,' and
last of all, expressed themselves more clearly and concisely, in the
phrase 'consubstantial with the Father;' for all that was beforesaid
has this meaning. As to their complaint about non-scriptural phrases,
they themselves are evidence of its futility. It was they who began
with their impious expressions; for, after their 'Out of nothing,' and
'Once was not,' going beyond Scripture in order to be impious, now
they make it a grievance, that, in condemning them, we go beyond
Scripture, in order to be pious.'' \footnote{Athan. Ep. ad Afros., 5, 6.} The last remark is important; even those traditional
statements of the Catholic doctrine, which were more explicit than
Scripture, had not as yet, when the controversy began, taken the shape
of formulæ. It was the Arian defined propositions of the ``\emph{out
of nothing},'' and the like, which called for the
imposition of the ``\emph{consubstantial}.''

It has sometimes been said, that the Catholics anxiously searched
for some offensive test, which might operate to the exclusion of the
Arians. This is not correct, inasmuch as they have no need to search;
the ``\emph{from God's substance}'' having been openly denied
by the Arians, five years before the Council, and no practical
distinction between it and the \emph{consubstantial} existing, till
the era of Basil and his Semi-Arians. Yet, had it been necessary,
doubtless it would have been their duty to seek for a test of this
nature; nay, to urge upon the heretical teachers the plain
consequences of their doctrine, and to drive them into the adoption of
them. These consequences are certain of being elicited in the
long-run; and it is but equitable to anticipate them in the persons of
the heresiarchs, rather than to suffer them gradually to unfold and
spread far and wide after their day, sapping the faith of their
deluded and less guilty followers. Many a man would be deterred from
outstepping the truth, could he see the end of his course from the
beginning. The Arians felt this, and therefore resisted a detection,
which would at once expose them to the condemnation of all serious
men. In this lies the difference between the treatment due to an
individual in heresy, and to one who is confident enough to publish
the innovations which he has originated. The former claims from us the
most affectionate sympathy, and the most considerate attention. The
latter should meet with no mercy; he assumes the office of the
Tempter, and, so far forth as his error goes, must be dealt with by
the competent authority, as if he were embodied Evil. To spare
him is a false and dangerous pity. It is to endanger the souls of
thousands, and it is uncharitable towards himself.

\xchapter{Chapter 3. The Ecumenical Council of Nicæa in the Reign of Constantine}

\xsection{Section 1. History of the Nicene Council}

THE authentic account of the
proceedings of the Nicene Council is not extant \footnote{Vide Ittigius, Hist. Conc. Nic. § 1. The rest of this volume is
drawn up from the following authorities: Eusebius, Vit. Const.
Socrates, Sozomen, and Theodoret, Hist. Eccles., the various
historical tracts of Athanasius, Epiphanius Hær. lxix. lxxiii., and
the Acta Conciliorum. Of moderns, especially Tillemont and Petavius;
then, Maimbourg's History of Arianism, the Benedictine Life of
Athanasius, Cave's Life of Athanasius and Literary History, Gibbon's
Roman History and Mr. Bridges' Reign of Constantine.}. It has in consequence been judged expedient to put together in
the foregoing Chapter whatever was necessary for the explanation of
the Catholic and Arian creeds, and the controversy concerning them,
rather than to reserve any portion of the doctrinal discussion for the
present, though in some respects the more appropriate place for its
introduction. Here then the transactions at Nicæa shall be reviewed
in their political or ecclesiastical aspect.

1.

Arius first published his heresy about the year 319. With his
turbulent conduct in 306 and a few years later we are not here
concerned. After this date, in 313, he is said, on the death of
Achillas, to have aspired to the primacy of the Egyptian Church; and,
according to Philostorgius \footnote{Philos. i. 3.},
the historian of his party, a writer of little credit, to have
generously resigned his claims in favour of Alexander, who was
elected. His ambitious character renders it not improbable that he was
a candidate for the vacant dignity; but, if so, the difference of age
between himself and Alexander, which must have been considerable,
would at once account for the elevation of the latter, and be an
evidence of the indecency of Arius in becoming a competitor at all.
His first attack on the Catholic doctrine was conducted with an
openness which, considering the general duplicity of his party, is the
most honourable trait in his character. In a public meeting of the
clergy of Alexandria, he accused his diocesan of Sabellianism; an
insult which Alexander, from deference to the talents and learning of
the objector, sustained with somewhat too little of the dignity
befitting ``the ruler of the people.'' The mischief which
ensued from his misplaced meekness was considerable. Arius was one of
the public preachers of Alexandria; and, as some suppose, Master of
the Catechetical School. Others of the city Presbyters were stimulated
by his example to similar irregularities. Colluthus, Carponas, and
Sarmatas began to form each his own party in a Church which Meletius
had already troubled; and Colluthus went so far as to promulgate
an heretical doctrine, and to found a sect. Still hoping to settle
these disorders without the exercise of his episcopal power, Alexander
summoned a meeting of his clergy, in which Arius was allowed to state
his doctrines freely, and to argue in their defence; and, whether from
a desire not to overbear the discussion, or from distrust in his own
power of accurately expressing the truth, and anxiety about the charge
of heresy brought against himself, the Primate, though in no wise a
man of feeble mind, is said to have refrained from committing himself
on the controverted subject, ``applauding,'' as Sozomen tells
us, ``sometimes the one party, sometimes the other.'' \footnote{Soz. i. 14.} At length the error of Arius appeared to be of so serious and
confirmed a nature, that countenance of it would have been sinful. It
began to spread beyond the Alexandrian Church; the indecision of
Alexander excited the murmurs of the Catholics; till, called
unwillingly to the discharge of a severe duty, he gave public evidence
of his real indignation against the blasphemies which he had so long
endured, by excommunicating Arius with his followers.

This proceeding, obligatory as it was on a Christian Bishop, and
ratified by the concurrence of a provincial Council, and expedient
even for the immediate interests of Christianity, had other Churches
been equally honest in their allegiance to the true faith, had the
effect of increasing the influence of Arius, by throwing him upon his
fellow-Lucianists of the rival dioceses of the East, and giving
notoriety to his name and tenets. In Egypt, indeed, he had already
been supported by the Meletian faction; which, in spite of its
profession of orthodoxy, continued in alliance with him, through
jealousy of the Church, even after he had fallen into heresy. But the
countenance of these schismatics was of small consideration, compared
with the powerful aid frankly tendered him, on his excommunication, by
the leading men in the great Catholic communities of Asia Minor and
the East. Cæsarea was the first place to afford him a retreat from
Alexandrian orthodoxy, where he received a cordial reception from the
learned Eusebius, Metropolitan of Palestine; while Athanasius, Bishop
of Anazarbus in Cilicia, and others, did not hesitate, by letters on
his behalf, to declare their concurrence with him in the full extent
of his heresy. Eusebius even declared that Christ was not very or true
God; and his associate Athanasius asserted, that He was in the number
of the hundred sheep of the parable, that is, one of the creatures of
God.

Yet, in spite of the countenance of these and other eminent men,
Arius found it difficult to maintain his ground against the general
indignation which his heresy excited. He was resolutely opposed by
Philogonius, Patriarch of Antioch, and Macarius of Jerusalem; who
promptly answered the call made upon them by Alexander, in his
circulars addressed to the Syrian Churches. In the meanwhile Eusebius
of Nicomedia, the early friend of Arius, and the ecclesiastical
adviser of Constantia, the Emperor's sister, declared in his favour;
and offered him a refuge, which he readily accepted, from the growing
unpopularity which attended him in Palestine. Supported by the
patronage of so powerful a prelate, Arius was now scarcely to be
considered in the position of a schismatic or an outcast. He assumed
in consequence a more calm and respectful demeanour towards Alexander;
imitated the courteous language of his friend; and in his Epistle,
which was introduced into the foregoing Chapter, addresses his
diocesan with studious humility, and defers or appeals to previous
statements made by Alexander himself on the doctrine in dispute \footnote{[Alexander's siding with Arius, was nothing
more than his disclaiming the views of the Five Fathers, vide supr.
pp. 202, 220; also Appendix, No. 2, [\emph{genn\={e}sis}]. As to
the Arian evasions which follow, vide supr. pp. 193, 216, 223, 238,
\&c.]}. At this time also he seems to have corrected and completed his
system. George, afterwards Bishop of Laodicea, taught him an evasion
for the orthodox test ``\emph{of God},'' by a reference to 1
Cor. xi. 12. Asterius, a sophist of Cappadocia, advocated the
secondary sense of the word Logos as applied to Christ, with a
reference to such passages as Joel ii. 25; and, in order to explain
away the force of the word ``\emph{Only-begotten},'' ([\emph{monogen\={e}s}],)
maintained, that to Christ alone out of all creatures it had been
given, to be fashioned under the immediate presence and perilous
weight of the Divine Hand. Now too, as it appears, the title of
``True God'' was ascribed to Him by the heretical party; the
``\emph{of an alterable nature}'' was withdrawn; and an
admission of His actual indefectibility substituted for it. The heresy
being thus placed on a less exceptionable basis, the influence of
Eusebius was exerted in Councils both in Bithynia and Palestine; in
which Arius was acknowledged, and more urgent solicitations addressed
to Alexander, with the view of effecting his re-admission into the
Church.

This was the history of the controversy for the first four or five
years of its existence; that is, till the era of the battle of
Hadrianople (A.D. 323), by the issue of which
Constantine, becoming master of the Roman world, was at liberty to
turn his thoughts to the state of Christianity in the Eastern
Provinces of the Empire. From this date it is connected with civil
history; a result natural, and indeed necessary under the existing
circumstances, though it was the occasion of subjecting Christianity
to fresh persecutions, in place of those which its nominal triumph had
terminated. When a heresy, condemned and excommunicated by one Church,
was taken up by another, and independent Christian bodies thus stood
in open opposition, nothing was left to those who desired peace, to
say nothing of orthodoxy, but to bring the question under the notice
of a General Council. But as a previous step, the leave of the civil
power was plainly necessary for so public a display of that
wide-spreading Association, of which the faith of the Gospel was the
uniting and animating principle. Thus the Church could not meet
together in one, without entering into a sort of negotiation with the
powers that be; whose jealousy it is the duty of Christians, both as
individuals and as a body, if possible, to dispel. On the other hand,
the Roman Emperor, as a professed disciple of the truth, was of course
bound to protect its interests, and to afford every facility for its
establishment in purity and efficacy. It was under these circumstances
that the Nicene Council was convoked.

2.

Now we must direct our view for a while to the character and
history of Constantine. It is an ungrateful task to discuss the
private opinions and motives of an Emperor who was the first to
profess himself the Protector of the Church, and to relieve it from
the abject and suffering condition in which it had lain for three
centuries. Constantine is our benefactor; inasmuch as we, who now
live, may be considered to have received the gift of Christianity by
means of the increased influence which he gave to the Church. And,
were it not that in conferring his benefaction he burdened it with the
bequest of an heresy, which outlived his age by many centuries, and
still exists in its effects in the divisions of the East, nothing
would here be said, from mere grateful recollection of him, by way of
analyzing the state of mind in which he viewed the benefit which he
has conveyed to us. But his conduct, as it discovers itself in the
subsequent history, natural as it was in his case, still has somewhat
of a warning in it, which must not be neglected in after times.

It is of course impossible accurately to describe the various
feelings with which one in Constantine's peculiar situation was likely
to regard Christianity; yet the joint effect of them all may be
gathered from his actual conduct, and the state of the civilized world
at the time. He found his empire distracted with civil and religious
dissensions, which tended to the dissolution of society; at a time
too, when the barbarians without were pressing upon it with a vigour,
formidable in itself, but far more menacing in consequence of the
decay of the ancient spirit of Rome. He perceived the powers of its
old polytheism, from whatever cause, exhausted; and a newly-risen
philosophy vainly endeavouring to resuscitate a mythology which
had done its work, and now, like all things of earth, was fast
returning to the dust from which it was taken. He heard the same
philosophy inculcating the principles of that more exalted and refined
religion, which a civilized age will always require; and he witnessed
the same substantial teaching, as he would consider it, embodied in
the precepts, and enforced by the energetic discipline, the union, and
the example of the Christian Church. Here his thoughts would rest, as
in a natural solution of the investigation to which the state of his
empire gave rise; and, without knowing enough of the internal
characters of Christianity to care to instruct himself in them, he
would discern, on the face of it, a doctrine more real than that of
philosophy, and a rule of life more severe and energetic even than
that of the old Republic. The Gospel seemed to be the fit instrument
of a civil reformation \footnote{Gibbon, Hist. ch. xx.},
being but a new form of the old wisdom, which had existed in the world
at large from the beginning. Revering, nay, in one sense, honestly
submitting to its faith, still he acknowledged it rather as a school
than joined it as a polity; and by refraining from the sacrament of
baptism till his last illness, he acted in the spirit of men of the
world in every age, who dislike to pledge themselves to engagements
which they still intend to fulfil, and to descend from the position of
judges to that of disciples of the truth \footnote{Vide his speech, Euseb. Vit. Const. iv. 62.}.

Concord is so eminently the perfection of the Christian temper,
conduct, and discipline, and it had been so wonderfully exemplified in
the previous history of the Church, that it was almost
unavoidable in a heathen soldier and statesman to regard it as the
sole precept of the Gospel. It required a far more refined moral
perception, to detect and to approve the principle on which this
internal peace is grounded in Scripture; to submit to the dictation of
truth, as such, as a primary authority in matters of political and
private conduct; to understand how belief in a certain creed was a
condition of Divine favour, how the social union was intended to
result from an unity of opinions, the love of man to spring from the
love of God, and zeal to be prior in the succession of Christian
graces to benevolence. It had been predicted by Him, who came to offer
peace to the world, that, in matter of fact, that gift would be
changed into the sword of discord; mankind being offended by the
doctrine, more than they were won over by the amiableness of
Christianity. But He alone was able thus to discern through what a
succession of difficulties Divine truth advances to its final victory;
shallow minds anticipate the end apart from the course which leads to
it. Especially they who receive scarcely more of His teaching than the
instinct of civilization recognizes (and Constantine must, on the
whole, be classed among such), view the religious dissensions of the
Church as simply evil, and (as they would fain prove) contrary to His
own precepts; whereas in fact they are but the history of truth in its
first stage of trial, when it aims at being ``pure,'' before
it is ``peaceable;'' and are reprehensible only so far as
baser passions mix themselves with that true loyalty towards God,
which desires His glory in the first place, and only in the second
place, the tranquillity and good order of society.

The Edict of Milan (A.D. 313) was among the
first effects of Constantine's anxiety to restore fellowship of
feeling to the members of his distracted empire. In it an absolute
toleration was given by him and his colleague Licinius, to the
Christians and all other persuasions, to follow the form of worship
which each had adopted for himself; and it was granted with the
professed view of consulting for the peace of their people.

A year did not elapse from the date of this Edict, when Constantine
found it necessary to support it by severe repressive measures against
the Donatists of Africa, though their offences were scarcely of a
civil nature. Their schism had originated in the disappointed ambition
of two presbyters; who fomented an opposition to Cæcilian, illegally
elevated, as they pretended, to the episcopate of Carthage. Growing
into a sect, they appealed to Constantine, who referred their cause to
the arbitration of successive Councils. These pronounced in favour of
Cæcilian; and, on Constantine's reviewing and confirming their
sentence, the defeated party assailed him with intemperate complaints,
accused Hosius, his adviser, of partiality in the decision, stirred up
the magistrates against the Catholic Church, and endeavoured to
deprive it of its places of worship. Constantine in consequence took
possession of their churches, banished their seditious bishops, and
put some of them to death. A love of truth is not irreconcilable
either with an unlimited toleration, or an exclusive patronage of a
selected religion; but to endure or discountenance error, according as
it is, or is not, represented in an independent system and existing
authority, to spare the pagans and to tyrannize over the
schismatics, is the conduct of one who subjected religious principle
to expediency, and aimed at peace, as a supreme good, by forcible
measures where it was possible, otherwise by conciliation.

It must be observed, moreover, that subsequently to the celebrated
vision of the Labarum (A.D. 312), he publicly
invoked the Deity as one and the same in all forms of worship; and at
a later period (A.D. 321), he promulgated
simultaneous edicts for the observance of Sunday, and the due
consultation of the aruspices \footnote{Gibbon, Hist. ibid.}.
On the other hand, as in the Edict of Milan, so in his Letters and
Edicts connected with the Arian controversy, the same reference is
made to external peace and good order, as the chief object towards
which his thoughts were directed. The same desire of tranquillity led
him to summon to the Nicene Council the Novatian Bishop Acesius, as
well as the orthodox prelates. At a later period still when he
extended a more open countenance to the Church as an institution, the
same principle discovers itself in his conduct as actuated him in his
measures against the Donatists. In proportion as he recognizes the
Catholic body, he drops his toleration of the sectaries. He prohibited
the conventicles of the Valentinians, Montanists, and other heretics;
who, at his bidding, joined the Church in such numbers (many of them,
says Eusebius, ``through fear of the Imperial threat, with
hypocritical minds'' \footnote{Euseb. Vit. Const. iii. 66. [\emph{nun pepl\={e}r\={o}tai
h\={e} ekkl\={e}sia kekrummen\={o}n hairetik\={o}n}].
Cyril. Catech. xv. 4.]}),
that at length both heresy and schism might be said to disappear from
the face of society. Now let us observe his conduct in the Arian
controversy.

Doubtless it was a grievous disappointment to a generous and
large-minded prince, to discover that the Church itself, from which he
had looked for the consolidation of his empire, was convulsed by
dissensions such as were unknown amid the heartless wranglings of
Pagan philosophy. The disturbances caused by the Donatists, which his
acquisition of Italy (A.D. 312) had opened upon
his view, extended from the borders of the Alexandrian patriarchate to
the ocean. The conquest of the East (A.D. 323)
did but enlarge his prospect of the distractions of Christendom. The
patriarchate just mentioned had lately been visited by a deplorable
heresy, which having run its course through the chief parts of Egypt,
Lybia, and Cyrenaica, had attacked Palestine and Syria, and spread
thence into the dioceses of Asia Minor and the Lydian Proconsulate.

Constantine was informed of the growing schism at Nicomedia, and at
once addressed a letter to Alexander and Arius jointly \footnote{Euseb. Vit. Const. ii. 64-72.}; a reference to which will enable the reader to verify for
himself the account above given of the nature of the Emperor's
Christianity. He professes therein two motives as impelling him in his
public conduct; first, the desire of effecting the reception,
throughout his dominions, of some one definite and complete form of
religious worship; next, that of settling and invigorating the civil
institutions of the empire. Desirous of securing an unity of sentiment
among all the believers in the Deity, he first directed his attention
to the religious dissensions of Africa, which he had hoped, with
the aid of the Oriental Christians, to terminate. ``But,'' he
continues, ``glorious and Divine Providence! how fatally were my
ears, or rather my heart, wounded, by the report of a rising schism
among you, far more acrimonious than the African dissensions ... On
investigation, I find that the reason for this quarrel is
insignificant and worthless ... As I understand it, you, Alexander,
were asking the separate opinions of your clergy on some passage of
your law, or rather were inquiring about some idle question, when you,
Arius, inconsiderately committed yourself to statements which should
either never have come into your mind, or have been at once repressed.
On this a difference ensued, Christian intercourse was suspended, the
sacred flock was divided into two, breaking the harmonious unity of
the common body ... Listen to the advice of me, your fellow-servant:—neither
ask nor answer questions which are not upon any injunction of your
law, but from the altercation of barren leisure; at best keep them to
yourselves, and do not publish them ... Your contention is not about
any capital commandment of your law; neither of you is introducing any
novel scheme of divine worship; you are of one and the same way of
thinking, so that it is in your power to unite in one communion. Even
the philosophers can agree together, one and all, in one dogma, though
differing in particulars ... Is it right for brothers to oppose
brothers, for the sake of trifles? ... Such conduct might be expected
from the multitude, or from the recklessness of boyhood; but is little
in keeping with your sacred profession, and with your personal
wisdom.'' Such is the substance of his letter, which, written on
an imperfect knowledge of the facts of the case, and with somewhat of
the prejudices of Eclectic liberalism, was inapplicable, even where
abstractedly true; his fault lying in his supposing, that an
individual like himself, who had not even received the grace of
baptism, could discriminate between great and little questions in
theology. He concludes with the following words, which show the
amiableness and sincerity of a mind in a measure awakened from the
darkness of heathenism, though they betray the affectation of the
rhetorician. ``Give me back my days of calm, my nights of
security; that I may experience henceforth the comfort of the clear
light, and the cheerfulness of tranquillity. Otherwise, I shall sigh
and be dissolved in tears ... So great is my grief, that I put off my
journey to the East on the news of your dissension ... Open for me
that path towards you, which your contentions have closed up. Let me
see you and all other cities in happiness; that I may offer due
thanksgivings to God above, for the unanimity and free intercourse
which is seen among you.''

This letter was conveyed to the Alexandrian Church by Hosius, who
was appointed by the Emperor to mediate between the contending
parties. A Council was called, in which some minor irregularities were
arranged, but nothing settled on the main question in dispute. Hosius
returned to his master to report an unsuccessful mission, and to
advise, as the sole measure which remained to be adopted, the calling
of a General Council, in which the Catholic doctrine might be formally
declared, and a judgment promulgated as to the basis upon which
communion with the Church was henceforth to be determined. Constantine
assented; and, discovering that the ecclesiastical authorities were
earnest in condemning the tenets of Arius, as being an audacious
innovation on the received creed, he suddenly adopted a new line of
conduct towards the heresy; and in a Letter which he addressed to
Arius, professes himself a zealous advocate of Christian truth,
ventures to expound it, and attacks Arius with a vehemence which can
only be imputed to his impatience in finding that any individual had
presumed to disturb the peace of the community. It is remarkable, as
showing his utter ignorance of doctrines, which were never intended
for discussion among the unbaptized heathen, or the secularized
Christian, that, in spite of this bold avowal of the orthodox faith in
detail, yet shortly after he explained to Eusebius one of the Nicene
declarations in a sense which even Arius would scarcely have allowed,
expressed as it is almost after the manner of Paulus \footnote{Theod. Hist. i. 12.}.

3.

The first Ecumenical Council met at Nicæa in Bithynia, in the
summer of A.D. 325. It was attended by about 300
Bishops, chiefly from the eastern provinces of the empire, besides a
multitude of priests, deacons, and other functionaries of the Church.
Hosius, one of the most eminent men of an age of saints, was
president. The Fathers who took the principal share in its proceedings
were Alexander of Alexandria, attended by his deacon Athanasius, then
about 27 years of age, and soon afterwards his successor in the
see; Eustathius, patriarch of Antioch, Macarius of Jerusalem,
Cæcilian of Carthage, the object of the hostility of the Donatists,
Leontius of Cæsarea in Cappadocia, and Marcellus of Ancyra, whose
name was afterwards unhappily notorious in the Church. The number of
Arian Bishops is variously stated at 13, 17, or 22; the most
conspicuous of these being the well-known prelates of Nicomedia and
Cæsarea, both of whom bore the name of Eusebius.

The discussions of the Council commenced in the middle of June, and
were at first private. Arius was introduced and examined; and
confessed his impieties with a plainness and vehemence far more
respectable than the hypocrisy which was the characteristic of his
party, and ultimately was adopted by himself. Then followed his
disputation with Athanasius \footnote{[``It is difficult,'' say the
Notes, Ath. Tr. vol. ii. p. 17, ``to gain a clear idea of the
character of Arius. Athanasius speaks as if his Thalia was but in
keeping with his life, calling him the 'Sotadean Arius,' while
Constantine, Alexander, and Epiphanius give us a contrary view of him,
still differing one from the other. Constantine, indeed, is not
consistent with himself; first he cries out to him (as if with
Athanasius), 'Arius, Arius, at least let the society of Venus keep you
back,' then 'Look, look all men ... how his veins and flesh are
possessed with poison, and are in a ferment of severe pain; how his
whole body is wasted, and is all withered and sad and pale and
shaking, and all that is miserable and fearfully emaciated. How
hateful to see, and how filthy is his mass of hair, how he is half
dead all over, with failing eyes and bloodless countenance, and woe-begone;
so that, all these things combining in him at once, frenzy, madness,
and folly, from the continuance of the complaint, have made thee wild
and savage. But, not having any sense of the bad plight he is in, he
cries out, ``I am transported with delight, and I leap and skip
for joy, and I fly;'' and again, with boyish impetuosity, ``Be
it so,'' he says, ``we are lost.''''' Harduin. Conc.
t. i. p. 457. St. Alexander speaks of Arius's melancholy temperament.
Epiphanius's account of him is as follows: ``From elation of mind
this old man swerved from the truth. He was in stature very tall,
downcast in visage, with manners like a wily serpent, captivating to
every guileless heart by that same crafty bearing. For, ever habited
in cloke and vest, he was pleasant of address, ever persuading souls
and flattering,'' \&c. Hær. 69, 3. Arius is here said to be
tall; Athanasius, unless Julian's description of him is but
declamation, was short, [\emph{m\={e}de an\={e}r, all' anthr\={o}piskos
entel\={e}s}] (``not even a man, but a common little
fellow''). Ep. 51. However, Gregory Nazianzen, who had never seen
him, speaks of him, as ``high in prowess, and humble in spirit,
mild, meek, full of sympathy, pleasant in speech, more pleasant in
manners, \emph{angelical} in \emph{person}, more angelical in mind,
serene in his rebukes, instructive in his praises,'' \&c. Orat,
21. 8.]},
who afterwards engaged the Arian Eusebius of Nicomedia, Maris,
and Theognis. The unfortunate Marcellus also distinguished himself in
the defence of the Catholic doctrine.

Reference has been already made to Gibbon's representation \footnote{[Supr. p. 234.]}, that the Fathers of the Council were in doubt for a time, how
to discriminate between their own doctrine and the heresy; but the
discussions of the foregoing Chapter contain sufficient evidence, that
they had rather to reconcile themselves to the adoption of a formula
which expedience suggested, and to the use of it as a test, than to
discover a means of ejecting or subduing their opponents. In the very
beginning of the controversy, Eusebius of Nicomedia had declared, that
he would not admit the ``\emph{from the substance}'' as an
attribute of our Lord \footnote{Theod. Hist. i. 6. [Vide Ath. Tr. vol. ii.
p. 438.]}.
A letter containing a similar avowal was read in the Council, and made
clear to its members the objects for which they had met; viz. to
ascertain the character and tendency of the heresy; to raise a protest
and defence against it; lastly, for that purpose, to overcome
their own reluctance to the formal and unauthoritative adoption of a
word, in explanation of the true doctrine, which was not found in
Scripture, had actually been perverted in the previous century to an
heretical meaning, and was in consequence forbidden by the Antiochene
Council which condemned Paulus.

The Arian party, on the other hand, anxious to avoid a test, which
they themselves had suggested, presented a Creed of their own, drawn
up by Eusebius of Cæsarea. In it, though the expression ``\emph{of
the substance}'' or ``\emph{consubstantial}'' was
omitted, every term of honour and dignity, short of this, was bestowed
therein upon the Son of God; who was designated as the Logos of God,
God of God, Light of Light, Life of Life, the Only-begotten Son, the
First-born of the whole creation, of the Father before all worlds, and
the Instrument of creating them. The Three Persons were confessed to
be in real \emph{hypostasis} or \emph{subsistence} (in opposition to
Sabellianism), and to be truly Father, Son, and Holy Ghost. The
Catholics saw very clearly, that concessions of this kind on the part
of the Arians did not conceal the real question in dispute. Orthodox
as were the terms employed by them, naturally and satisfactorily as
they would have answered the purposes of a test, had the existing
questions never been agitated, and consistent as they were with
certain producible statements of the Ante-Nicene writers, they were
irrelevant at a time when evasions had been found for them all, and
triumphantly proclaimed. The plain question was, whether our Lord was
God in as full a sense as the Father, though not to be viewed as
separable from Him; or whether, as the sole alternative, He was
a creature; that is, whether He was literally of, and in, the one
Indivisible Essence which we adore as God, ``consubstantial with
God,'' or of a substance which had a beginning. The Arians said
that He was a creature, the Catholics that He was very God; and all
the subtleties of the most fertile ingenuity could not alter, and
could but hide, this fundamental difference. A specimen of the Arian
argumentation at the Council has already been given on the testimony
of Athanasius; happily it was not successful. A form of creed was
drawn up by Hosius, containing the discriminating terms of orthodoxy \footnote{[Justice has not been done here to the
ground of tradition, on which the Fathers specially took their stand.
For example, ``Whoever heard such doctrine?'' says Athanasius;
``whence, from whom did they gain it? Who thus expounded to them
when they were at school?'' Orat. i. § 8. ``Is it not enough
to distract a man, and to make him stop his ears?'' § 35. Vide
Ath. Tr. vol. ii. pp. 247-253, 311.]}; and anathemas were added against all who maintained the
heretical formulæ, Arius and his immediate followers being mentioned
by name. In order to prevent misapprehension of the sense in which the
test was used, explanations accompanied it. Thus carefully defined, it
was offered for subscription to the members of the Council; who in
consequence bound themselves to excommunicate from their respective
bodies all who actually obtruded upon the Church the unscriptural and
novel positions of Arius. As to the laity, they were not required to
subscribe any test as the condition of communion; though they were of
course exposed to the operation of the anathema, in case they ventured
on positive innovations on the rule of faith.

While the Council took this clear and temperate view of its
duties, Constantine acted a part altogether consistent with his own
previous sentiments, and praiseworthy under the circumstances of his
defective knowledge. He had followed the proceedings of the assembled
prelates with interest, and had neglected no opportunity of impressing
upon them the supreme importance of securing the peace of the Church.
On the opening of the Council, he had set the example of conciliation,
by burning publicly, without reading, certain charges which had been
presented to him against some of its members; a noble act, as
conveying a lesson to all present to repress every private feeling,
and to deliberate for the well-being of the Church Catholic to the end
of time. Such was his behaviour, while the question in controversy was
still pending; but when the decision was once announced, his tone
altered, and what had been a recommendation of caution, at once became
an injunction to conform. Opposition to the sentence of the Church was
considered as disobedience to the civil authority; the prospect of
banishment was proposed as the alternative of subscription; and it was
not long before seven of the thirteen dissentient Bishops submitted to
the pressure of the occasion, and accepted the creed with its
anathemas as articles of peace.

Indeed the position in which Eusebius of Nicomedia had placed their
cause, rendered it difficult for them consistently to refuse
subscription. The violence, with which Arius originally assailed the
Catholics, had been succeeded by an affected earnestness for unity and
concord, so soon as his favour at Court allowed him to dispense with
the low popularity by which he first rose into notice. The
insignificancy of the points in dispute which had lately been
the very ground of complaint with him and his party against the
particular Church which condemned him, became an argument for their
yielding, when the other Churches of Christendom confirmed the
sentence of the Alexandrian. It is said, that some of them substituted
the ``\emph{homœusion}'' (``\emph{like in substance}''),
for the ``\emph{homoüsion}'' (``\emph{one in substance}'')
in the confessions which they presented to the Council; but it is
unsafe to trust the Anomœan Philostorgius, on whose authority the
report rests \footnote{Philost. i. 9.}, in a
charge against the Eusebian party, and perhaps after all he merely
means, that they explained the latter by the former as an excuse for
their own recantation. The six, who remained unpersuaded, had founded
an objection, which the explanations set forth by the Council had gone
to obviate, on the alleged materialism of the word which had been
selected as the test. At length four of them gave way; and the other
two, Eusebius of Nicomedia and another, withdrawing their opposition
to the ``\emph{homoüsion},'' only refused to sign the
condemnation of Arius. These, however, were at length released from
their difficulty, by the submission of the heresiarch himself; who was
pardoned on the understanding, that he never returned to the Church,
which had suffered so much from his intrigues. There is, however, some
difficulty in this part of the history. Eusebius shortly afterwards
suffered a temporary exile, on a detection of his former practices
with Licinius to the injury of Constantine; and Arius, apparently
involved in his ruin, was banished with his followers into Illyria.

\xsection{Section 2. Consequences of the Nicene Council}

FROM the time that the Eusebians
consented to subscribe the Homoüsion in accordance with the wishes of
a heathen prince, they became nothing better than a political party.
They soon learned, indeed, to call themselves Homϟsians, or
believers in the ``like'' substance (\emph{homϟsion},) as
if they still held the peculiarities of a religious creed; but in
truth it is an abuse of language to say that they had any definite
belief at all. For this reason, the account of the Homϟsian or
Semi-Arian doctrine shall be postponed, till such time as we fall in
with individuals whom we may believe to be serious in their
professions, and to act under the influence of religious convictions
however erroneous. Here the Eusebians must be described as a secular
faction, which is the true character of them in the history in which
they bear a part.

Strictly speaking, the Christian Church, as being a visible
society, is necessarily a political power or party. It may be a party
triumphant, or a party under persecution; but a party it always must
be, prior in existence to the civil institutions with which it is
surrounded, and from its latent divinity formidable and
influential, even to the end of time. The grant of permanency was made
in the beginning, not to the mere doctrine of the Gospel, but to the
Association itself built upon the doctrine \footnote{Matt. xvi. 18.}; in prediction, not only of the indestructibility of
Christianity, but of the medium also through which it was to be
manifested to the world. Thus the Ecclesiastical Body is a
divinely-appointed means, towards realizing the great evangelical
blessings. Christians depart from their duty, or become in an
offensive sense political, not when they act as members of one
community, but when they do so for temporal ends or in an illegal
manner; not when they assume the attitude of a party, but when they
split into many. If the primitive believers did not interfere with the
acts of the civil government, it was merely because they had no civil
rights enabling them legally to do so. But where they have rights, the
case is different \footnote{Acts xvi. 37-39.}; and
the existence of a secular spirit is to be ascertained, not by their
using these, but their using them for ends short of the ends for which
they were given. Doubtless in criticizing the mode of their exercising
them in a particular case, differences of opinion may fairly exist;
but the principle itself, the duty of using their civil rights in the
service of religion, is clear; and since there is a popular
misconception, that Christians, and especially the Clergy, as such,
have no concern in temporal affairs, it is expedient to take every
opportunity of formally denying the position, and demanding proof of
it. In truth, the Church was framed for the express purpose of
interfering, or (as irreligious men will say) meddling with the world.
It is the plain duty of its members, not only to associate
internally, but also to develope that internal union in an external
warfare with the spirit of evil, whether in Kings' courts or among the
mixed multitude; and, if they can do nothing else, at least they can
suffer for the truth, and remind men of it, by inflicting on them the
task of persecution.

1.

These principles being assumed, it is easy to enter into the
relative positions of the Catholics and Arians at the era under
consideration. As to the Arians, it is a matter of fact, that Arius
and his friends commenced their career with the deliberate commission
of disorderly and schismatical acts; and it is a clear inference from
their subsequent proceedings, that they did so for private ends. For
both reasons, then, they were a mere political faction, usurping the
name of religion; and, as such, essentially anti-christian. The
question here is not whether their doctrine was right or wrong; but,
whether they did not make it a secondary object of their exertions, an
instrument towards attaining ends which they valued above it. Now it
will be found, that the party was prior to the creed. They grafted
their heresy on the schism of the Meletians, who continued to support
them after they had published it; and they readily abandoned it, when
their secular interests required the sacrifice. At the Council of
Nicæa, they began by maintaining an erroneous doctrine; they ended by
concessions which implied the further heresy that points of faith are
of no importance; and, if they were odious when they blasphemed the
truth, they were still more odious when they confessed it. It
was the very principle of Eclecticism to make light of differences in
belief; while it was involved in the primary notion of a Revelation
that these differences were of importance, and it was taught with
plainness in the Gospel, that to join with those who denied the right
faith was a sin.

This adoption, however, on the part of the Eusebians, of the dreams
of Pagan philosophy, served in some sort as a recommendation of them
to a prince who, both from education and from knowledge of the world,
was especially tempted to consider all truth as a theory which was not
realized in a present tangible form. Accordingly, when once they had
rid themselves of the mortification caused by their forced
subscription, they had the satisfaction of finding themselves the most
powerful party in the Church, as being the representative and organ of
the Emperor's sentiments. They then at once changed places with the
Catholics; who sustained a double defeat, both in the continued power
of those whom they had hoped to exclude from the Church, and again, in
the invidiousness of their own unrelenting suspicion and dislike of
men, who had seemed by subscription to satisfy all reasonable doubt
respecting their orthodoxy.

The Arian party was fortunate, moreover, in its leaders; one the
most dexterous politician, the other the most accomplished theologian
of the age. Eusebius of Nicomedia was a Lucianist, the fellow-disciple
of Arius. He was originally Bishop of Berytus, in Phœnicia; but,
having gained the confidence of Constantia, sister to Constantine, and
wife to Licinius, he was by her influence translated to
Nicomedia, where the Eastern Court then resided. Here he secretly
engaged in the cause of Licinius against his rival, and is even
reported to have been indifferent to the security of the Christians
during the persecution which followed; a charge which certainly
derives some confirmation from Alexander's circular epistle, in which
the Arians are accused of directing the violence of the civil power
against the orthodox of Alexandria. On the ruin of Licinius, he was
screened by Constantia from the resentment of the conqueror; and,
being recommended by his polished manners and shrewd and persuasive
talent, he soon contrived to gain an influence over the mind of
Constantine himself. From the time that Arius had recourse to him on
his flight from Palestine, he is to be accounted the real head of the
heretical party; and his influence is quickly discernible in the
change which ensued in its language and conduct. While a courteous
tone was assumed towards the defenders of the orthodox doctrine, the
subtleties of dialectics, in which the sect excelled, were used, not
in attacking, but in deceiving its opponents, in making unbelief
plausible, and obliterating the distinctive marks of the true creed.
It must not be forgotten that it was from Nicomedia, the see of
Eusebius, that Constantine wrote his epistle to Alexander and Arius.

In supporting Arianism in its new direction, the other Eusebius,
Bishop of Cæsarea, was of singular service. This distinguished
writer, to whom the Christian world has so great a debt at the present
day, though not characterized by the unprincipled ambition of his
namesake, is unhappily connected in history with the Arian
party. He seems to have had the faults and the virtues of the mere man
of letters: strongly excited neither to good nor to evil, and careless
at once of the cause of truth and the prizes of secular greatness, in
comparison of the comforts and decencies of literary ease. His first
master was Dorotheus of Antioch \footnote{Danz. de Eus. Cæsar, 22.}; afterwards he became a pupil of the School of Cæsarea, which
seems to have been his birth-place, and where Origen had taught. Here
he studied the works of that great master, and the other writers of
the Alexandrian school. It does not appear when he first began to
arianize. At Cæsarea he is celebrated as the friend of the Orthodox
Pamphilus, afterwards martyred, whom he assisted in his defence of
Origen, in answer to the charges of heterodoxy then in circulation
against him. The first book of this work is still extant in the Latin
translation of Ruffinus, and its statements of the Catholic doctrines
are altogether explicit and accurate. In his own writings, numerous as
they are, there is very little which fixes on Eusebius any charge,
beyond that of an attachment to the Platonic phraseology. Had he not
connected himself with the Arian party, it would have been unjust to
have suspected him of heresy. But his acts are his confession. He
openly sided with those whose blasphemies a true Christian would have
abhorred; and he sanctioned and shared their deeds of violence and
injustice perpetrated on the Catholics.

But it is a different reason which has led to the mention of
Eusebius in this connection. The grave accusation under which he lies,
is not that of arianizing, but of corrupting the simplicity of
the Gospel with an Eclectic spirit. While he held out the ambiguous
language of the schools as a refuge, and the Alexandrian imitation of
it as an argument, against the pursuit of the orthodox, his conduct
gave countenance to the secular maxim, that difference in creeds is a
matter of inferior moment, and that, provided we confess as far as the
very terms of Scripture, we may speculate as philosophers, and live as
the world \footnote{In this association of the Eusebian with the
Eclectic temper, it must not be forgotten, that Julian the Apostate
was the pupil of Eusebius of Nicomedia, his kinsman; that he took part
with the Arians against the Catholics; and that, in one of his extant
epistles, he speaks in praise of the writings of an Arian Bishop,
George of Laodicea. Vide Weisman, sec. iv. 35. § 12.}. A more
dangerous adviser Constantine could hardly have selected, than a man
thus variously gifted, thus exalted in the Church, thus disposed
towards the very errors against which he required especially to be
guarded. The remark has been made that, throughout his Ecclesiastical
History no instance occurs of his expressing abhorrence of the
superstitions of paganism, and that his custom is either to praise, or
not to blame, such heretical writers as fall under his notice \footnote{Kestner de Euseb. Auctor. prolegom. § 17.
Yet it must be confessed, he is strongly opposed to [\emph{go\={e}teia}]
in all its forms; i.e. as being unworthy a philosopher.}.

Nor must the influence of the Court pass unnoticed, in recounting
the means by which Arianism secured a hold over the mind of the
Emperor. Constantia, his favourite sister, was the original patroness
of Eusebius of Nicomedia; and thus a princess, whose name would
otherwise be dignified by her misfortunes, is known to
Christians of later times only as a principal instrument of the
success of heresy. Wrought upon by a presbyter, a creature of the
bishop's, who was in her confidence, she summoned Constantine to her
bed-side in her last illness, begged him, as her parting request, to
extend his favour to the Arians, and especially commended to his
regard the presbyter himself, who had stimulated her to this
experiment on the feelings of a brother. The hangers-on of the
Imperial Court imitated her in her preference for the polite and
smooth demeanour of the Eusebian prelates, which was advantageously
contrasted to the stern simplicity of the Catholics. The eunuchs and
slaves of the palace (strange to say) embraced the tenets of Arianism;
and all the most light-minded and frivolous of mankind allowed
themselves to pervert the solemn subject in controversy into matter
for fashionable conversation or literary amusement.

The arts of flattery completed the triumph of the heretical party.
So many are the temptations to which monarchs are exposed of
forgetting that they are men, that it is obviously the duty of the
Episcopal Order to remind them that there is a visible Power in the
world, divinely founded and protected, superior to their own. But
Eusebius places himself at the feet of a heathen; and forgetful of his
own ordination-grace, allows the Emperor to style himself ``the
bishop of Paganism,'' and ``the predestined Apostle of virtue
to all men.'' \footnote{Euseb. Vit. Const. iii. 58, iv. 24. Vide
also, i. 4, 24.} The
shrine of the Church was thrown open to his inspection; and, contrary
to the spirit of Christianity, its mysteries were officiously
explained to one who was not yet even a candidate for baptism. The restoration and erection of Churches, which is the honourable
distinction of his reign, assimilated him, in the minds of his
courtiers, to the Divine Founder and Priest of the invisible temple;
and the magnificence, which soothed the vanity of a monarch, seemed in
its charitable uses almost a substitute for personal religion \footnote{Ibid. iv. 22, and alibi. Vide Gibbon, ch. xx.}.

2.

While events thus gradually worked for the secular advancement of
the heretical party, the Catholics were allotted gratifications and
anxieties of a higher character. The proceedings of the Council had
detected the paucity of the Arians among the Rulers of the Church;
which had been the more clearly ascertained, inasmuch as no temporal
interests had operated to gain for the orthodox cause that vast
preponderance of advocates which it had actually obtained. Moreover,
it had confirmed by the combined evidence of the universal Church, the
argument from Scripture and local tradition, which each separate
Christian community already possessed. And there was a satisfaction in
having found a formula adequate to the preservation of the
all-important article in controversy in all its purity. On the other
hand, in spite of these immediate causes of congratulation, the
fortunes of the Church were clouded in prospect, by the Emperor's
adoption of its Creed as a formula of peace, not of belief, and by the
ready subscription of the unprincipled faction, which had previously
objected to it. This immediate failure, which not unfrequently attends
beneficial measures in their commencement, issued, as has been
said, in the temporary triumph of the Arians. The disease, which had
called for the Council, instead of being expelled from the system, was
thrown back upon the Church, and for a time afflicted it \footnote{Theod. Hist. i. 6, fin.}; nor was it cast out, except by the persevering fasting and
prayer, the labours and sufferings, of the oppressed believers.
Meanwhile, the Catholic prelates could but retire from the Court
party, and carefully watch its movements; and, in consequence,
incurred the reproach and the penalty of being ``troublers of
Israel.'' This may be illustrated from the subsequent history of
Arius himself, with which this Chapter shall close.

It is doubtful, whether or not Arius was persuaded to sign the
symbol at the Nicene Council; but at least he professed to receive it
about five years afterwards. At this time Eusebius of Nicomedia had
been restored to the favour of Constantine; who, on the other hand,
influenced by his sister, had become less zealous in his adherence to
the orthodox side of the controversy. An attempt was made by the
friends of Arius to effect his re-admission into the Church at
Alexandria. The great Athanasius was at this time Primate of Egypt;
and in his instance the question was tried, whether or not the Church
would adopt the secular principles, to which the Arians were willing
to subject it, and would abandon its faith, as the condition of
present peace and prosperity. He was already known as the counsellor
of Alexander in the previous controversy; yet, Eusebius did not at
once give up the hope of gaining him over, a hope which was
strengthened by his recent triumph over the orthodox prelates of
Antioch, Gaza, and Hadrianople, whom he had found means to deprive of
their sees to make way for Arians. Failing in his attempt at
conciliation, he pursued the policy which might have been anticipated,
and accused the Bishop of Alexandria of a youthful rashness, and an
obstinate contentious spirit, incompatible with the good understanding
which ought to subsist among Christians. Arius was summoned to Court,
presented an ambiguous confession, and was favourably received by
Constantine. Thence he was despatched to Alexandria, and was quickly
followed by an imperial injunction addressed to Athanasius, in order
to secure the restoration of the heresiarch to the Church to which he
had belonged. ``On being informed of my pleasure,'' says
Constantine, in the fragment of the Epistle preserved by Athanasius,
``give free admission to all, who are desirous of entering into
communion with the Church. For if I learn of your standing in the way
of any who were seeking it, or interdicting them, I will send at once
those who shall depose you instead, by my authority, and banish you
from your see.'' \footnote{Athan. Apol. contr. Arian 59.} It
was not to be supposed, that Athanasius would yield to an order,
though from his sovereign, which was conceived in such ignorance of
the principles of Church communion, and of the powers of its Rulers;
and, on his explanation, the Emperor professed himself well satisfied,
that he should use his own discretion in the matter. The intrigues of
the Eusebians, which followed, shall elsewhere be related; they ended
in effecting the banishment of Athanasius into Gaul, the restoration
of Arius at a Council held at Jerusalem, his return to
Alexandria, and, when the anger of the intractable populace against
him broke out into a tumult, his recall to Constantinople to give
further explanations respecting his real opinions.

There the last and memorable scene of his history took place, and
furnishes a fresh illustration of the clearness and integrity, with
which the Catholics maintained the true principles of Church union,
against those who would have sacrificed truth to peace. The aged
Alexander, bishop of the see, underwent a persecution of entreaties
and threats, such as had already been employed against Athanasius. The
Eusebians urged upon him, by way of warning, their fresh successes
over the Bishops of Ancyra and Alexandria; and appointed a day, by
which he was to admit Arius to communion, or to be ejected from his
see. Constantine confirmed this alternative. At first, indeed, he had
been struck with doubts respecting the sincerity of Arius; but, on the
latter professing with an oath that his tenets were orthodox, and
presenting a confession, in which the terms of Scripture were made the
vehicle of his characteristic impieties, the Emperor dismissed his
scruples, observing with an anxiety and seriousness which rise above
his ordinary character, that Arius had well sworn if his words had no
double meaning; otherwise, God would avenge. The miserable man did not
hesitate to swear, that he professed the Creed of the Catholic Church
without reservation, and that he had never said nor thought otherwise,
than according to the statements which he now made.

For seven days previous to that appointed for his re-admission, the
Church of Constantinople, Bishop and people, were given up to
fasting and prayer. Alexander, after a vain endeavour to move the
Emperor, had recourse to the most solemn and extraordinary form of
anathema allowed in the Church \footnote{Bingham, Antiq. xvi. 2. § 17.}; and with tears besought its Divine Guardian, either to take
himself out of the world, or to remove thence the instrument of those
extended and increasing spiritual evils, with which Christendom was
darkening. On the evening before the day of his proposed triumph,
Arius passed through the streets of the city with his party, in an
ostentatious manner; when the stroke of death suddenly overtook him,
and he expired before his danger was discovered.

Under the circumstances, a thoughtful mind cannot but account this
as one of those remarkable interpositions of power, by which Divine
Providence urges on the consciences of men in the natural course of
things, what their reason from the first acknowledges, that He is not
indifferent to human conduct. To say that these do not fall within the
ordinary course of His governance, is merely to say that they \emph{are}
judgments; which, in the common meaning of the word, stand for events
extraordinary and unexpected. That such do take place under the
Christian Dispensation, is sufficiently proved by the history of
Ananias and Sapphira. It is remarkable too, that the similar
occurrences, which happen at the present day, are generally connected
with some unusual perjury or extreme blasphemy; and, though we may not
infer the sin from the circumstance of the temporal infliction, yet,
the commission of the sin being ascertained, we may well account, that
its guilt is providentially impressed on the minds and enlarged
in the estimation of the multitude, by the visible penalty by which it
is followed. Nor do we in such cases necessarily pass any absolute
sentence upon the person, who appears to be the object of Divine
Visitation; but merely upon the particular act which provoked it, and
which has its fearful character of evil stamped upon it, independent
of the punishment which draws our attention to it. The man of God, who
prophesied against the altar in Bethel, is not to be regarded by the
light of his last act, though a judgment followed it, but according to
the general tenor of his life. Arius also must thus be viewed; though,
unhappily, his closing deed is but the seal of a prevaricating and
presumptuous career.

Athanasius, who is one of the authorities from whom the foregoing
account is taken, received it from Macarius, a presbyter of the Church
of Constantinople, who was in that city at the time. He adds,
``while the Church was rejoicing at the deliverance, Alexander
administered the communion in pious and orthodox form, praying with
all the brethren and glorifying God greatly; not as if rejoicing over
his death, (God forbid! for to all men it is appointed once to die,)
but because in this event there was displayed somewhat more than a
human judgment. For the Lord Himself, judging between the threats of
the Eusebians and the prayer of Alexander, has in this event given
sentence against the heresy of the Arians; showing it to be unworthy
of ecclesiastical fellowship, and manifesting to all, that though it
have the patronage of Emperor and of all men yet that by the Church
itself it is condemned.'' \footnote{Epist. ad Scrap. 4.}

\xchapter{Chapter 4. Councils in the Reign of Constantius}

\xsection{Section 1. The Eusebians}

\footnote{[In this Chapter a change in the structure of the sentences has
been made here and there, with the view of relieving the intricacies
of the narrative.]} \footnote{[Vid. Appendix, No. 6.]} THE death of Arius was productive of no
important consequences in the history of his party. They had never
deferred to him as their leader, and since the Nicene Council had even
abandoned his creed. The theology of the Eclectics had opened to
Eusebius of Cæsarea a language less obnoxious to the Catholics and to
Constantine, than that into which he had been betrayed in Palestine;
while his namesake, possessing the confidence of the Emperor, was
enabled to wield weapons more decisive in the controversy than those
which Arius had used. From that time Semi-Arianism was their
profession, and calumny their weapon, for the deposition, by legal
process, of their Catholic opponents. This is the character of their
proceedings from A.D. 328 to A.D.
350; when circumstances led them to adopt a third creed, and enabled
them to support it by open force.

1.

It may at first sight excite our surprise, that men who were so
little careful to be consistent in their professions of faith, should
be at the pains to find evasions for a test, which they might have
subscribed as a matter of course, and then dismissed from their
thoughts. But, not to mention the natural desire of continuing an
opposition to which they had once committed themselves, and especially
after a defeat, there is, moreover, that in religious mysteries which
is ever distasteful to secular minds. The marvellous, which is sure to
excite the impatience and resentment of the baffled reason, becomes
insupportable when found in those solemn topics, which it would fain
look upon, as necessary indeed for the uneducated, but irrelevant when
addressed to those who are already skilled in the knowledge and the
superficial decencies of virtue. The difficulties of science may be
dismissed from the mind, and virtually forgotten; the precepts of
morality, imperative as they are, may be received with the
condescension, and applied with the modifications, of a
self-applauding refinement. But what at once demands attention, yet
refuses to satisfy curiosity, places itself above the human mind,
imprints on it the thought of Him who is eternal, and enforces the
necessity of obedience for its own sake. And thus it becomes to the
proud and irreverent, what the consciousness of guilt is to the
sinner; a spectre haunting the field, and disturbing the complacency,
of their intellectual investigations. In this at least, throughout
their changes, the Eusebians are consistent,—in their hatred of the
Sacred Mystery.

It has sometimes been scornfully said, on the other hand, that the
zeal of Christians, in the discussion of theological subjects, has
increased with the mysteriousness of the doctrine in dispute. There is
no reason why we should shrink from the avowal. Doubtless, a subject
that is dear to us, does become more deeply fixed in our affections by
its very peculiarities and incidental obscurities. We desire to revere
what we already love; and we seek for the materials of reverence in
such parts of it, as exceed our intelligence or imagination. It should
therefore excite our devout gratitude, to reflect how the truth has
been revealed to us in Scripture in the most practical manner; so as
both to humble and to win over, while it consoles, those who really
love it. Moreover, with reference to the particular mystery under
consideration, since a belief in our Lord's Divinity is closely
connected (how, it matters not) with deep religious feeling generally,—involving
a sense both of our need and of the value of the blessings which He
has procured for us, and an emancipation from the tyranny of the
visible world,—it is not wonderful, that those, who would confine
our knowledge of God to things seen, should dislike to hear of His
true and only Image. If the unbeliever has attempted to account for
the rise of the doctrine, by the alleged natural growth of a
veneration for the Person and acts of the Redeemer, let it at least be
allowed to Christians to reverse the process of argument, and to
maintain rather, that a low estimation of the evangelical blessings
leads to unworthy conceptions of the Author of them. In the case of
laymen it will show itself in a sceptical neglect of the subject of
religion altogether; while ecclesiastics, on whose minds
religion is forced, are tempted either to an undue exaltation of their
order, or to a creed dishonourable to their Lord. The Eusebians
adopted the latter alternative, and so merged the supremacy of Divine
Truth amid the multifarious religions and philosophies of the world.

Their skilfulness in reasoning and love of disputation afford us an
additional explanation of their pertinacious opposition to the Nicene
Creed. Though, in possessing the favour of the Imperial Court, they
had already the substantial advantages of victory, they disdained
success without a battle. They loved the excitement of suspense, and
the triumph of victory. And this sophistical turn of mind accounts,
not only for their incessant wranglings, but for their frequent
changes of view, as regards the doctrine in dispute. It may be doubted
whether men, so practised in the gymnastics of the Aristotelic school,
could carefully develope and consistently maintain a definite view of
doctrine; especially in a case, where the difficulties of an unsound
cause combined with their own habitual restlessness and levity to
defeat the attempt. Accordingly, in their conduct of the argument,
they seem to be aiming at nothing beyond ``living from hand to
mouth,'' as the saying is; availing themselves of some or other
expedient, which would suffice to carry them through existing
difficulties; admissions, whether to satisfy the timid conscience of
Constantius, or to deceive the Western Church; or statements so
faintly precise and so decently ambiguous, as to embrace the greatest
number of opinions possible, and to deprive religion, in consequence,
of its austere and commanding aspect.

That I may not seem to be indulging in vague accusation, I here
present the reader with a sketch of the lives of the chief of them;
from which he will be able to decide, whether the above explanation of
their conduct is unnecessary or gratuitous.

The most distinguished of the party, after Eusebius himself, for
ability, learning, and unscrupulousness, was Acacius, the successor of
the other Eusebius in the see of Cæsarea. He had been his pupil, and
on his death inherited his library. Jerome ranks him among the most
learned commentators on Scripture. The Arian historian, Philostorgius,
praises his boldness, penetration, and perspicuity in unfolding his
views: and Sozomen speaks of his talents and influence as equal to the
execution of the most difficult designs \footnote{Tillemont, Mem. des Ariens, vol. vi. c. 28.}. He began at first with professing himself a Semi-Arian after
the example of Eusebius, his master; next, he became the founder of
the party, which will presently be described as the \emph{Homœan} or \emph{Scriptural};
thirdly, he joined himself to the Anomœans or pure Arians, so as even
to be the intimate associate of the wretched Aetius; fourthly, at the
command of Constantius, he deserted and excommunicated him; fifthly,
in the reign of the Catholic Jovian, he signed the \emph{Homoüsion}
or symbol of Nicæa.

George, of Laodicæa, another of the leading members of the
Eusebian party, was originally a presbyter of the Alexandrian Church,
and deposed by Alexander for the assistance afforded by him to Arius
at Nicomedia. At the end of the reign of Constantius, he professed for
a while the sentiments of the Semi-Arians; whether seriously or not,
we have not the means of deciding, although the character given
of him by Athanasius, who is generally candid in his judgments, is
unfavourable to his sincerity. Certainly he deserted the Semi-Arians
in no long time, and died an Anomœan. He is also accused of open and
habitual irregularities of life.

Leontius, the most crafty of his party, was promoted by the Arians
to the see of Antioch \footnote{A strange and scandalous transaction in
early life, gave him the appellation of [\emph{ho apokopos}]. Athan.
ad Monach. 4.};
and though a pupil of the school of Lucian, and consistently attached
to the opinions of Arius throughout his life, he seems to have
conducted himself in his high position with moderation and good
temper. The Catholic party was at that time still strong in the city,
particularly among the laity; the crimes of Stephen and Placillus, his
immediate Arian predecessors, had brought discredit on the heretical
cause; and the theological opinions of Constantius, who was attached
to the Semi-Arian doctrine, rendered it dangerous to avow the plain
blasphemies of the first founder of their creed. Accordingly, with the
view of seducing the Catholics to his own communion, he was anxious to
profess an agreement with the Church, even where he held an opposite
opinion; and we are told that in the public doxology, which was
practically the test of faith, not even the nearest to him in the
congregation could hear from him more than the words ``for ever
and ever,'' with which it concludes. It was apparently with the
same design, that he converted the almshouses of the city, destined
for the reception of strangers, into seminaries for propagating the
Christian faith; and published a panegyrical account of St.
Babylas, when his body was to be removed to Daphne, by way of
consecrating a place which had been before devoted to sensual
excesses. In the meanwhile, he gradually weakened the Church, by a
systematic promotion of heretical, and a discountenance of the
orthodox Clergy; one of his most scandalous acts being his ordination
of Aetius, the founder of the Anomœans, who was afterwards promoted
to the episcopacy in the reign of Julian.

Eudoxius, the successor of Leontius, in the see of Antioch, was his
fellow-pupil in the school of Lucian. He is said to have been
converted to Semi-Arianism by the writings of the Sophist Asterius;
but he afterwards joined the Anomœans, and got possession of the
patriarchate of Constantinople. It was there, at the dedication of the
cathedral of St. Sophia, that he uttered the wanton impiety, which has
characterized him with a distinctness, which supersedes all historical
notice of his conduct, or discussion of his religious opinions.
``When Eudoxius,'' says Socrates, ``had taken his seat on
the episcopal throne, his first words were these celebrated ones, 'the
Father is [\emph{aseb\={e}s}], irreligious; the Son [\emph{euseb\={e}s}],
religious.' When a noise and confusion ensued, he added, 'Be not
distressed at what I say; for the Father is irreligious, as
worshipping none; but the Son is religious towards the Father.' On
this the tumult ceased, and in its place an intemperate laughter
seized the congregation; and it remains as a good saying even to this
time.'' \footnote{Socr. Hist. ii. 43. [[\emph{Eusebeia, asebeia,
dussebeia}], and their derivatives, in the language of Athanasius
or his age, means orthodoxy, heterodoxy, orthodox, \&c. This
circumstance gives its point to the jest. This sense is traceable to
St. Paul's words, ``Great is the mystery of godliness ([\emph{eusebeias}]),''
orthodoxy. Vide Athan. Opp. passim. Thus Arius also ends his letter to
Eusebius with ``[\emph{al\={e}th\={o}s eusebie}]'' And
St. Basil, defending his own freedom from Arian error, says that St.
Macrina, his grandmother, ``moulded him from his infancy in the
dogmas of religion ([\emph{eusebeias}]),'' and that, when he grew
up, and travelled, he ever chose those for his fathers and guides,
whom he found walking according to ``the rule of religion ([\emph{eusebeias}])
handed down.'' Ep. 204, 6. Vide also, Basil. Opp. t. 2, p. 599.
Greg. Naz. Orat. ii. 80. Euseb. cont. Marc. i. 7. Joan. Antioch. apud
Facund. i. 1. Sozomen, i. 20, as supr. note p. 140.]}

Valens, Bishop of Mursa, in Pannonia, shall close this list of
Eusebian Prelates. He was one of the immediate disciples of Arius;
and, from an early age, the champion of his heresy in the Latin
Church. In the conduct of the controversy, he inherited more of the
plain dealing as well as of the principles of his master, than his
associates; he was an open advocate of the Anomœan doctrine, and by
his personal influence with Constantius balanced the power of the
Semi-Arian party, derived from the Emperor's private attachment to
their doctrine. The favour of Constantius was gained by a fortunate
artifice, at the time the latter was directing his arms against the
tyrant Magnentius. ``While the two armies were engaged in the
plains of Mursa,'' says Gibbon, ``and the fate of the two
rivals depended on the chance of war, the son of Constantine passed
the anxious moments in a church of the martyrs, under the walls of the
city. His spiritual comforter Valens, the Arian Bishop of the diocese,
employed the most artful precautions to obtain such early
intelligence, as might secure either his favour or his escape. A
secret chain of swift and trusty messengers informed him of the
vicissitudes of the battle; and while the courtiers stood trembling
around their affrighted master, Valens assured him that the
Gallic legions gave way; and insinuated, with some presence of mind,
that the glorious event had been revealed to him by an Angel. The
grateful Emperor ascribed his success to the merits and intercession
of the Bishop of Mursa, whose faith had deserved the public and
miraculous approbation of Heaven.'' \footnote{Gibbon, Hist. ch. xxi.}

Such were the leaders of the Eusebian or Court faction; and on the
review of them, do we not seem to see in each a fresh exhibition of
their great type and forerunner, Paulus, on one side or other of his
character, though surpassing him in extravagance of conduct, as
possessing a wider field, and more powerful incentives for ambitious
and energetic exertion? We see the same accommodation of the Christian
Creed to the humour of an earthly Sovereign, the same fertility of
disputation in support of their version of it, the same reckless
profanation of things sacred, the same patient dissemination of error
for the services of the age after them; and, if they are free from the
personal immoralities of their master, they balance this favourable
trait of character by the cruel and hard-hearted temper, which
discovers itself in their persecution of the Catholics.

2.

This persecution was conducted till the middle of the century
according to the outward forms of ecclesiastical law. Charges of
various kinds were preferred in Council against the orthodox prelates
of the principal sees, with a profession at least of regularity,
whatever unfairness there might be in the details of the
proceedings. By this means all the most powerful Churches of Eastern
Christendom, by the commencement of the reign of Constantius (A.D.
337), had been brought under the influence of the Arians;
Constantinople, Heraclea, Hadrianople, Ephesus, Ancyra, both Cæsareas,
Antioch, Laodicæa, and Alexandria. Eustathius of Antioch, for
instance, had incurred their hatred, by his strenuous resistance to
the heresy in the seat of its first origin. After the example of his
immediate predecessor Philogonius, he refused communion to Stephen,
Leontius, Eudoxius, George, and others; and accused Eusebius of
Cæsarea openly of having violated the faith of Nicæa. The heads of
the party assembled in Council at Antioch; and, on charges of heresy
and immorality, which they professed to be satisfactorily maintained,
pronounced sentence of deposition against him. Constantine banished
him to Philippi, together with a considerable number of the priests
and deacons of his Church. So again, Marcellus of Ancyra, another of
their inveterate opponents, was deposed, anathematized, and banished
by them, with greater appearance of justice, on the ground of his
leaning to the errors of Sabellius. But their most rancorous enmity
and most persevering efforts were directed against the high-minded
Patriarch of Alexandria; and, in illustration of their principles and
conduct, the circumstances of his first persecution shall here be
briefly related.

When Eusebius of Nicomedia failed to effect the restoration of
Arius into the Alexandrian Church by persuasion, he had threatened to
gain his end by harsher means. Calumnies were easily invented against
the man who had withstood his purpose; and it so happened, that
willing tools were found on the spot for conducting the attack. The
Meletian sectaries have already been noticed, as being the original
associates of Arius; who had troubled the Church by taking part in
their schism, before he promulgated his peculiar heresy. They were
called after Meletius, Bishop of Lycopolis in the Thebaid; who, being
deposed for lapsing in the Dioclesian persecution, separated from the
Catholics, and, propagating a spurious succession of clergy by his
episcopal prerogative, formed a powerful body in the heart of the
Egyptian Church. The Council of Nicæa, desirous of terminating the
disorder in the most temperate manner, instead of deposing the
Meletian bishops, had arranged, that they should retain a nominal rank
in the sees, in which they had respectively placed themselves; while,
by forbidding them to exercise their episcopal functions, it provided
for the termination of the schism at their death. But, with the bad
fortune which commonly attends conciliatory measures, unless
accompanied by such a display of vigour as shows that concession is
but condescension, the clemency was forgotten in the restriction,
which irritated, without repressing them; and, being bent on the
overthrow of the dominant Church, they made a sacrifice of their
principles, which had hitherto been orthodox, and joined the Eusebians.
By this intrigue, the latter gained an entrance into the Egyptian
Church, as effectual as that which had already been opened to them, by
means of their heresy itself, in Syria and Asia Minor \footnote{The Meletians, on the other hand, were not
in the event equally advantaged by the coalition; for after the
success of their attack upon Athanasius, Constantine, true to his
object of restoring tranquillity to the Church, while he banished
Athanasius to Treves, banished also John, the leader of the Meletians,
who had been forward in procuring his condemnation.}.

Charges against Athanasius were produced and examined in Councils
successively held at Cæsarea and Tyre (A.D.
333-335); the Meletians being the accusers, and the Eusebians the
judges in the trial. At an earlier date, it had been attempted to
convict him of political offences; but, on examination, Constantine
became satisfied of his innocence. It had been represented, that, of
his own authority, he had imposed and rigorously exacted a duty upon
the Egyptian linen cloth; the pretended tribute being in fact nothing
beyond the offerings, which pious persons had made to the Church, in
the shape of vestments for the service of the sanctuary. It had
moreover been alleged, that he had sent pecuniary aid to one
Philumenus, who was in rebellion against the Emperor; as at a later
period they accused him of a design of distressing Constantinople, by
stopping the corn vessels of Alexandria, destined for the supply of
the metropolis.

The charges brought against him before these Councils were both of
a civil and of an ecclesiastical character; that he, or Macarius, one
of his deacons, had broken a consecrated chalice, and the holy table
itself, and had thrown the sacred books into the fire; next, that he
had killed Arsenius, a Meletian bishop, whose hand, amputated and
preserved for magical purposes, had been found in Athanasius's house.
The latter of these strange accusations was refuted at the Council of
Cæsarea by Arsenius himself, whom Athanasius had gained, and who, on
the production of a human hand at the trial, presented himself
before the judges, thus destroying the circumstantial evidence by
which it was to be identified as his. The former charge was refuted at
Tyre by the testimony of the Egyptian bishops; who, after exposing the
equivocating evidence of the accuser, went on to prove that at the
place where their Metropolitan was said to have broken the chalice,
there was neither church, nor altar, nor chalice, existing. These were
the principal allegations brought against him; and their extraordinary
absurdity, (certain as the charges are as matters of history, from
evidence of various kinds,) can only be accounted for by supposing,
that the Eusebians were even then too powerful and too bold, to care
for much more than the bare forms of law, or to scruple at any
evidence, which the unskilfulness of their Egyptian coadjutors might
set before them. A charge of violence in his conduct towards certain
Meletians was added to the above; and, as some say, a still more
frivolous accusation of incontinence, but whether this was ever
brought, is more than doubtful.

Cæsarea and Tyre were places too public even for the audacity of
the Eusebians, when the facts of the case were so plainly in favour of
the accused. It was now proposed that a commission of inquiry should
be sent to the Mareotis, which was in the neighbourhood, and formed
part of the diocese, of Alexandria, and was the scene of the alleged
profanation of the sacred chalice. The leading members of this
commission were Valens and Ursacius, Theognis, Maris, and two others,
all Eusebians; they took with them the chief accuser of Athanasius as
their guide and host, leaving Athanasius and Macarius at Tyre, and
refusing admittance into the court of inquiry to such of the
clergy of the Mareotis, as were desirous of defending their Bishop's
interests in his absence. The issue of such proceedings may be
anticipated. On the return of the commission to Tyre, Athanasius was
formally condemned of rebellion, sedition, and a tyrannical use of his
episcopal power, of murder, sacrilege, and magic; was deposed from the
see of Alexandria, and prohibited from ever returning to that city.
Constantine confirmed the sentence of the Council, and Athanasius was
banished to Gaul.

3.

It has often been remarked that persecutions of Christians, as in
St. Paul's case, ``fall out rather unto the furtherance of the
Gospel.'' [Phil. i. 12.] The dispersion of the disciples, after
the martyrdom of St. Stephen, introduced the word of truth together
with themselves among the Samaritans; and in the case before us, the
exile of Athanasius led to his introduction to the younger
Constantine, son of the great Emperor of that name, who warmly
embraced his cause, and gave him the opportunity of rousing the zeal,
and gaining the personal friendship of the Catholics of the West.
Constans also, another son of Constantine, declared in his favour; and
thus, on the death of their father, which took place two years after
the Council of Tyre, one third alone of his power, in the person of
the Semi-Arian Constantius, Emperor of the East, remained with that
party, which, while Constantine lived, was able to wield the whole
strength of the State against the orthodox Bishops. The support
of the Roman See was a still more important advantage gained by
Athanasius. Rome was the natural mediator between Alexandria and
Antioch, and at that time possessed extensive influence among the
Churches of the West. Accordingly, when Constantius re-commenced the
persecution, to which his father had been persuaded, the exiles betook
themselves to Rome; and about the year 340 or 341 we read of Bishops
from Thrace, Syria, Phœnicia, and Palestine, collected there, besides
a multitude of Presbyters, and among the former, Athanasius himself,
Marcellus, Asciepas of Gaza, and Luke of Hadrianople. The first act of
the Roman See in their favour was the holding a provincial Council, in
which the charges against Athanasius and Marcellus were examined, and
pronounced to be untenable. And its next act was to advocate the
summoning of a Council of the whole Church with the same purpose,
referring it to Athanasius to select a place of meeting, where his
cause might be secure of a more impartial hearing, than it had met
with at Cæsarea and Tyre.

The Eusebians, on the other hand, perceiving the danger which their
interests would sustain, should a Council be held at any distance from
their own peculiar territory, determined on anticipating the projected
Council by one of their own, in which they might both confirm the
sentence of deposition against Athanasius, and, if possible, contrive
a confession of faith, to allay the suspicions which the Occidentals
entertained of their orthodoxy \footnote{[``After the Nicene Council, the
Eusebians did not dare avow their heresy in Constantine's time, but
merely attempted the banishment of Athanasius, and the restoration of
Arius. Their first Council was A.D. 341, four
years after Constantine's death and Constantius's accession.''—Ath.
Tr. vol. i. pp. 92, 93.]}.
This was the occasion of the Council of the Dedication, as it is
called, held by them at Antioch, in the year 341, and which is one of
the most celebrated Councils of the century. It was usual to solemnize
the consecration of places of worship, by an attendance of the
principal prelates of the neighbouring districts; and the great Church
of the Metropolis of Syria, called the Dominicum Aureum, which had
just been built, afforded both the pretext and the name to their
assembly. Between ninety and a hundred bishops came together on this
occasion, all Arians or Arianizers, and agreed without difficulty upon
the immediate object of the Council, the ratification of the Synods of
Cæsarea and Tyre in condemnation of Athanasius.

So far their undertaking was in their own hands; but a more
difficult task remained behind, viz., to gain the approval and consent
of the Western Church, by an exposition of the articles of their
faith. Not intending to bind themselves by the decision at Nicæa,
they had to find some substitute for the \emph{Homoüsion}. With this
view four, or even five creeds, more or less resembling the Nicene in
language, were successively adopted. The first was that ascribed to
the martyr Lucian, though doubts are entertained concerning its
genuineness. It is in itself almost unexceptionable; and, had there
been no controversies on the subjects contained in it, would have been
a satisfactory evidence of the orthodoxy of its promulgators. The Son
is therein styled the exact Image of the substance, will, power, and
glory of the Father; and the Three Persons of the Holy Trinity
are said to be three in substance, one in will \footnote{Exact image, [\emph{aparallaktos eik\={o}n}];
substance, [\emph{ousia}]; subsistence, or person, [\emph{hypostasis}].}. An evasive condemnation was added of the Arian tenets;
sufficient, as it might seem, to delude the Latins, who were unskilled
in the subtleties of the question. For instance, it was denied that
our Lord was born ``in time,'' but in the heretical school, as
was shown above, time was supposed to commence with the creation of
the world; and it was denied that He was ``in the number of the
creatures,'' it being their doctrine, that He was the sole
immediate work of God, and, as such, not like others, but separate
from the whole creation, of which indeed He was the author. Next, for
some or other reason, two new creeds were proposed, and partially
adopted by the Council; the same in character of doctrine, but
shorter. These three were all circulated, and more or less received in
the neighbouring Churches; but, on consideration, none of them seemed
adequate to the object in view, that of recommending the Eusebians to
the distant Churches of the West. Accordingly, a fourth formulary was
drawn up after a few months' delay, among others by Mark, Bishop of
Arethusa, a Semi-Arian Bishop of religious character, afterwards to be
mentioned; its composers were deputed to present it to Constans; and,
this creed proving unsatisfactory, a fifth confession was drawn up
with considerable care and ability; though it too failed to quiet the
suspicions of the Latins. This last is called the Macrostich, from the
number of its paragraphs, and did not make its appearance till three
years after the former.

In truth, no such exposition of the Catholic faith could satisfy
the Western Christians, while they were witnesses to the exile of its
great champion on account of his fidelity to it. Here the Eusebians
were wanting in their usual practical shrewdness. Words, however
orthodox, could not weigh against so plain a fact. The Occidentals,
however unskilled in the niceties of the Greek language, were able to
ascertain the heresy of the Eusebians in their malevolence towards
Athanasius. Nay, the anxious attempts of his enemies, to please them
by means of a confession of faith, were a refutation of their
pretences. For, inasmuch as the sense of the Catholic world, had
already been recorded in the \emph{Homoüsion}, why should they devise
a new formulary, if after all they agreed with the Church? or, why
should they themselves be so fertile in confessions, if they had all
of them but one faith? It is brought against them by Athanasius, that
in their creeds they date their exposition of the Catholic doctrine,
as if it were something new, instead simply of its being declared,
which was the sole design of the Nicene Fathers; while at other times,
they affected to acknowledge the authority of former Councils, which
nevertheless they were indirectly opposing \footnote{Athan. de Syn. 3. 37.}. Under these circumstances the Roman Church, as the
representative of the Latins, only became more bent upon the
convocation of a General Council in which the Nicene Creed might be
ratified, and any innovation upon it reprobated; and the innocence of
Athanasius, which it had already ascertained in its provincial Synod,
might be formally proved, and proclaimed to the whole of Christendom.
This object was at length accomplished. Constans, whom
Athanasius had visited and gained, successfully exerted his influence
with his brother Constantius, the Emperor of the East; and a Council
of the whole Christian world was summoned at Sardica for the above
purposes, the exculpation of Marcellus and others being included with
that of Athanasius.

Sardica was chosen as the place of meeting, as lying on the
confines of the two divisions of the Empire. It is on the borders of Mœsia,
Thrace, and Illyricum, and at the foot of Mount Hæmus, which
separates it from Philippopolis. There the heads of the Christian
world assembled in the year 347, twenty-two years after the Nicene
Council, in number above 380 bishops, of whom seventy-six were Arian.
The President of the Council was the venerable Hosius; whose name was
in itself a pledge, that the decision of Nicæa was simply to be
preserved, and no fresh question raised on a subject already exhausted
by controversy. But, almost before the opening of the Council, matters
were brought to a crisis; a schism took place in its members; the
Arians retreated to Philippopolis, and there excommunicated the
leaders of the orthodox, Julius of Rome, Hosius, and Protogenes of
Sardica, issued a sixth confession of faith, and confirmed the
proceedings of the Antiochene Council against Athanasius and the other
exiles.

This secession of the Arians arose in consequence of their finding,
that Athanasius was allowed a seat in the Council; the discussions of
which they refused to attend, while a Bishop took part in them, who
had already been deposed by Synods of the East. The orthodox replied,
that a later Council, held at Rome, had fully acquitted and
restored him; moreover, that to maintain his guilt was but to assume
the principal point, which they were then assembled to debate; and,
though very consistent with their absenting themselves from the
Council altogether, could not be permitted to those, who had by their
coming recognized the object, for which it was called. Accordingly,
without being moved by their retreat, the Council proceeded to the
condemnation of some of the more notorious opponents among them of the
Creed of Nicæa, examined the charges against Athanasius and the rest,
reviewed the acts of the investigations at Tyre and the Mareotis,
which the Eusebians had sent to Rome in their defence, and confirmed
the decree of the Council of Rome, in favour of the accused. Constans
enforced this decision on his brother by the arguments peculiar to a
monarch; and the timid Constantius, yielding to fear what he denied to
justice, consented to restore to Alexandria a champion of the truth,
who had been condemned on the wildest of charges, by the most hostile
and unprincipled of judges.

The journey of Athanasius to Alexandria elicited the fullest and
most satisfactory testimonies of the real orthodoxy of the Eastern
Christians; in spite of the existing cowardice or misapprehension,
which surrendered them to the tyrannical rule of a few determined and
energetic heretics. The Bishops of Palestine, one of the chief holds
of the Arian spirit, welcomed, with the solemnity of a Council, a
restoration, which, under the circumstances of the case, was almost a
triumph over their own sovereign; and so excited was the Catholic
feeling even at Antioch, that Constantius feared to grant to the
Athanasians a single Church in that city, lest it should have been the
ruin of the Arian cause.

One of the more important consequences of the Council of Sardica,
was the public recantation of Valens, and his accomplice Ursacius,
Bishop of Singidon, in Pannonia, two of the most inveterate enemies
and calumniators of Athanasius. It was addressed to the Bishop of
Rome, and was conceived in the following terms: ``Whereas we are
known heretofore to have preferred many grievous charges against
Athanasius the Bishop, and, on being put on our defence by your
excellency, have failed to make good our charges, we declare to your
excellency, in the presence of all the presbyters, our brethren, that
all which we have heretofore heard against the aforesaid, is false,
and altogether foreign to his character; and therefore, that we
heartily embrace the communion of the aforesaid Athanasius, especially
considering your Holiness, according to your habitual clemency, has
condescended to pardon our mistake. Further we declare, that, should
the Orientals at any time, or Athanasius, from resentful feelings, be
desirous to bring us to account, that we will not act in the matter
without your sanction. As for the heretic Arius, and his partisans,
who say that ``\emph{Once the Son was not},'' that ``\emph{He
is of created Substance},'' and that ``\emph{He is not the Son
of God before all time},'' we anathematize them now, and once
for all, according to our former statement which we presented at
Milan. Witness our hand, that we condemn once for all the Arian
heresy, as we have already said, and its advocates. Witness also the
hand of Ursacius.—I, Ursacius the Bishop, have set my name to this
statement.'' \footnote{Athan. Apol. cont. Arian. 58.}

The Council of Milan, referred to in the conclusion of this letter,
seems to have been held A.D. 347; two years
after the Arian creed, called Macrostich, was sent into the West, and
shortly after the declaration of Constans in favour of the restoration
of the Athanasians.

\xsection{Section 2. The Semi-Arians}

THE events recorded in the last Section
were attended by important consequences in the history of Arianism.
The Council of Sardica led to a separation between the Eastern and
Western Churches; which seemed to be there represented respectively by
the rival Synods of Sardica and Philippopolis, and which had before
this time hidden their differences from each other, and communicated
together from a fear of increasing the existing evil \footnote{Soz. iii. 13.}. Not that really there was any discordance of doctrine between
them. The historian, from whom this statement is taken, gives it at
the same time as his own opinion, that the majority of the Asiatics
were Homoüsians, though tyrannized over by the court influence, the
sophistry, the importunity, and the daring, of the Eusebian party.
This mere handful of divines, unscrupulously pressing forward into the
highest ecclesiastical stations, set about them to change the
condition of the Churches thus put into their power; and, as has been
remarked in the case of Leontius of Antioch, filled the inferior
offices with their own creatures, and sowed the seeds of future
discords and disorders, which they could not hope to have themselves
the satisfaction of beholding. The orthodox majority of Bishops and
divines, on the other hand, timorously or indolently, kept in the
background; and allowed themselves to be represented at Sardica by
men, whose tenets they knew to be unchristian, and professed to
abominate. And in such circumstances, the blame of the open
dissensions, which ensued between the Eastern and Western divisions of
Christendom, was certain to be attributed to those who urged the
summoning of the Council, not to those who neglected their duty by
staying away. In qualification of this censure, however, the
intriguing spirit of the Eusebians must be borne in mind; who might
have means, of which we are not told, of keeping away their orthodox
brethren from Sardica. Certainly the expense of the journey was
considerable, whatever might be the imperial or the ecclesiastical
allowances for it \footnote{[On the \emph{cursus publicus}, vid.
Gothofred. in Cod. Theod. viii. tit. 5. It was provided for the
journeys of the Emperor, for persons whom he summoned, for
magistrates, ambassadors, and for such private persons as the Emperor
indulged in the use of it, which was gratis. The use was granted by
Constantine to the Bishops who were summoned to Nicæa, as far as it
went, in addition to other means of travelling. Euseb. V. Const. iii.
6. (though aliter Valesius in loc.) The \emph{cursus publicus} brought
the Bishops to the Council of Tyre. Ibid. iv. 43. In the conference
between Liberius and Constantius (Theod. Hist. ii. 13), it is objected
that the \emph{cursus publicus} is not sufficient to convey Bishops to
the Council, as Liberius proposes; he answers that the Churches are
rich enough to convey their Bishops as far as the seas. Thus St.
Hilary was compelled (datâ evectionis copiâ, Sulp. Sev. Hist. ii.
57) to attend at Seleucia, as Athanasius at Tyre. Julian complains of
the abuse of the \emph{cursus publicus}, perhaps with an allusion to
these Councils of Constantius. Vide Cod. Theod. viii. tit. 5, l. 12;
where Gothofred quotes Liban. Epitaph. in Julian. (vol. i. p. 569, ed.
Reiske). Vide the well-known passage of Ammianus, who speaks of the
Councils as being the ruin of the \emph{res vehicularia}, Hist. xxi.
16. The Eusebians at Philippopolis say the same, Hilar. Fragm. iii.
25. The Emperor provided board and perhaps lodging for the Bishops at
Ariminum; which the Bishops of Aquitaine, Gaul, and Britain declined,
except three British from poverty. Sulp. Hist. ii. 56. Hunmeric in
Africa, after assembling 466 Bishops at Carthage, dismissed them
without mode of conveyance, provision, or baggage. Victor. Utic. Hist.
iii. init. In the Emperor's letter previous to the assembling of the
sixth Ecumenical Council, A.D. 678 (Harduin.
Conc. t. 3, p. 1043, fin.), he says he has given orders for the
conveyance and maintenance of its members. Pope John VIII. reminds
Ursus, Duke of Venice (A.D. 876), of the same
duty of providing for the members of a Council, ``secundum pios
principes, qui in talibus munificè semper erant intenti.'' Colet.
Concil. (Ven. 1730) t. xi. p. 14.]}, and
their absence from their flocks, especially in an age fertile in
Councils, was an evil. Still there is enough in the history of the
times, to evidence a culpable negligence on the part of the orthodox
of Asia.

However, this rupture between the East and West has here been
noticed, not to censure the Asiatic Churches, but for the sake of its
influence on the fortunes of Arianism. It had the effect of pushing
forward the Semi-Arians, as they are called, into a party distinct
from the Eusebian or Court party, among whom they had hitherto been
concealed. This party, as its name implies, professed a doctrine
approximating to the orthodox; and thus served as a means of deceiving
the Western Churches, which were unskilled in the evasions, by which
the Eusebians extricated themselves from even the most explicit
confessions of the Catholic doctrine. Accordingly, the six heretical
confessions hitherto recounted were all Semi-Arian in character, as
being intended more or less to justify the heretical party in the eyes
of the Latins. But when this object ceased to be feasible, by
the event of the Sardican Council, the Semi-Arians ceased to be of
service to the Eusebians, and a separation between the parties
gradually took place.

1.

The Semi-Arians, whose history shall here be introduced,
originated, as far as their doctrine is concerned, in the change of
profession which the Nicene anathema was the occasion of imposing upon
the Eusebians; and had for their founders Eusebius of Cæsarea, and
the Sophist Asterius. But viewed as a party, they are of a later date
\footnote{[Vide Ath. Tr. vol. ii. pp. 282-286.]}. The genuine Eusebians
were never in earnest in the modified creeds, which they so
ostentatiously put forward for the approbation of the West. However,
while they clamoured in defence of the inconsistent doctrine contained
in them, which, resembling the orthodox in word, might in fact subvert
it, and at once confessed and denied our Lord, it so happened, that
they actually recommended that doctrine to the judgment of some of
their followers, and succeeded in creating a direct belief in an
hypothesis, which in their own case was but the cloke for their own
indifference to the truth. This at least seems the true explanation of
an intricate subject in the history. There are always men of sensitive
and subtle minds, the natural victims of the bold disputant; men, who,
unable to take a broad and common-sense view of an important subject,
try to satisfy their intellect and conscience by refined distinctions
and perverse reservations. Men of this stamp were especially to be
found among a people possessed of the language and acuteness of
the Greeks. Accordingly, the Eusebians at length perceived, doubtless
to their surprise and disgust, that a party had arisen from among
themselves, with all the positiveness (as they would consider it), and
nothing of the straightforward simplicity of the Catholic
controversialists, more willing to dogmatize than to argue, and
binding down their associates to the real import of the words, which
they had themselves chosen as mere evasions of orthodoxy; and to their
dismay they discovered, that in this party the new Emperor himself was
to be numbered. Constantius, indeed, may be taken as a type of a
genuine Semi-Arian; resisting, as he did, the orthodox doctrine from
over-subtlety, timidity, pride, restlessness, or other weakness of
mind, yet paradoxical enough to combat at the same time and condemn
all, who ventured to teach anything short of that orthodoxy. Balanced
on this imperceptible centre between truth and error, he alternately
banished every party in the controversy, not even sparing his own; and
had recourse in turn to every creed for relief, except that in which
the truth was actually to be found.

The symbol of the Semi-Arians was the \emph{Homϟsion}, ``\emph{like
in substance},'' which they substituted for the orthodox \emph{Homoüsion},
``\emph{one in substance},'' or ``\emph{consubstantial}.''
Their objections to the latter formula took the following form. If the
word \emph{usia}, ``\emph{substance},'' denoted the
``first substance,'' or an individual being, then \emph{Homoüsios}
seemed to bear a Sabellian meaning, and to involve a denial of the
separate personality of the Son \footnote{Epiph. Hær. lxxiii. 11, fin.}. On the other hand, if the word was understood as including two
distinct Persons (or \emph{Hypostases}), this was to use it, as
it is used of created things; as if by substance were meant some
common nature, either divided in fact, or one merely by abstraction \footnote{Soz. iii. 18.}. They were strengthened in this view by the decree of the
Council, held at Antioch between the years 260 and 270, in
condemnation of Paulus, in which the word \emph{Homoüsion} was
proscribed. They preferred, accordingly, to name the Son ``like in
substance,'' \footnote{[\emph{homoios kat' ousian}].} or \emph{Homϟsios},
with the Father, that is, of a substance like in all things, except in
not being the Father's substance; maintaining at the same time, that,
though the Son and Spirit were separate in substance from the Father,
still they were so included in His glory that there was but one God.

Instead of admitting the evasion of the Arians, that the word \emph{Son}
had but a secondary sense, and that our Lord was in reality a
creature, though ``not like other creatures,'' they plainly
declared that He was not a creature, but truly the Son, born of the
substance (\emph{usia}) of the Father, as if an Emanation from Him at
His will; yet they would not allow Him simply to be God, as the Father
was; but, asserting that there were various energies in the Divine
Being, they considered creation to be one, and the \emph{gennesis} or \emph{generation}
to be another, so that the Son, though distinct in substance from God,
was at the same time essentially distinct from every created nature.
Or they suggested that He was the offspring of the Person (\emph{hypostasis}),
not of the \emph{substance} or \emph{usia} of the Father; or, so to
say, of the Divine Will, as if the force of the word ``\emph{Son}''
consisted in this point. Further, instead of the ``\emph{once}
\emph{He was not},'' they adopted the ``\emph{generated
time-apart},'' for which even Arius had changed it. That is, as
holding that the question of the beginning of the Son's existence was
beyond our comprehension, they only asserted that there was such a
beginning, but that it was before time and independent of it; as if it
were possible to draw a distinction between the Catholic doctrine of
the derivation or order of succession in the Holy Trinity (the ``\emph{unoriginately
generated}'') and this notion of a beginning simplified of the
condition of time.

Such was the Semi-Arian Creed, really involving contradictions in
terms, parallel to those of which the orthodox were accused;—that
the Son was born before all times, yet not eternal; not a creature,
yet not God; of His substance, yet not the same in substance; and His
exact and perfect resemblance in all things, yet not a second Deity.

2.

Yet the men were better than their creed; and it is satisfactory to
be able to detect amid the impiety and worldliness of the heretical
party any elements of a purer spirit, which gradually exerted itself
and worked out from the corrupt mass, in which it was embedded. Even
thus viewed as distinct from their political associates, the
Semi-Arians are a motley party at best; yet they may be considered as
Saints and Martyrs, when compared with the Eusebians, and in fact some
of them have actually been acknowledged as such by the Catholics of
subsequent times. Their zeal in detecting the humanitarianism of
Marcellus and Photinus, and their good service in withstanding the
Anomœans, who arrived at the same humanitarianism by a bolder
course of thought, will presently be mentioned. On the whole they were
men of correct and exemplary life, and earnest according to their
views; and they even made pretensions to sanctity in their outward
deportment, in which they differed from the true Eusebians, who, as
far as the times allowed it, affected the manners and principles of
the world. It may be added, that both Athanasius and Hilary, two of
the most uncompromising supporters of the Catholic doctrine, speak
favourably of them. Athanasius does not hesitate to call them brothers
\footnote{[However, he is severe upon Eustathius and
Basil (ad Ep. Æg. 7.), as St. Basil is on the former, who had been
his friend.]}; considering that,
however necessary it was for the edification of the Church at large,
that the Homoüsion should be enforced on the clergy, yet that the
privileges of private Christian fellowship were not to be denied to
those, who from one cause or other stumbled at the use of it \footnote{Athan. de Syn. 41.}. It is remarkable, that the Semi-Arians, on the contrary, in
their most celebrated Synod (at Ancyra, A.D.
358) anathematized the holders of the Homoüsion, as if crypto-Sabellians
\footnote{Epiph. supra.}.

Basil, the successor of Marcellus, in the see of Ancyra, united in
his person the most varied learning with the most blameless life, of
all the Semi-Arians \footnote{Theod. Hist. ii. 25.}.
This praise of rectitude in conduct was shared with him by Eustathius
of Sebaste, and Eleusius of Cyzicus. These three Bishops especially
attracted the regard of Hilary, on his banishment to Phrygia by the
intrigues of the Arians (A.D. 356). The zealous
confessor feelingly laments the condition, in which he found the
Churches in those parts. ``I do not speak of things strange to
me:'' he says, ``I write not without knowledge; I have heard
and seen in my own person the faults, not of laics merely, but of
bishops. For, excepting Eleusius and a few with him, the ten provinces
of Asia, in which I am, are for the most part truly ignorant of
God.'' \footnote{Hilar. de Syn. 63.} His
testimony in favour of the Semi-Arians of Asia Minor, must in fairness
be considered as delivered with the same force of assertion, which
marks his protest against all but them; and he elsewhere addresses
Basil, Eustathius, and Eleusius, by the title of ``Sanctissimi
viri.'' \footnote{Ibid. 90. Vid. also the Life of St. Basil
of Cæsarea, who was intimate for a time with Eustathius and others.}

Mark, Bishop of Arethusa, in Syria, has obtained from the Greek
Church the honours of a Saint and Martyr. He indulged, indeed, a
violence of spirit, which assimilates him to the pure Arians, who were
the first among Christians to employ force in the cause of religion.
But violence, which endures as freely as it assails, obtains our
respect, if it is denied our praise. His exertions in the cause of
Christianity were attended with considerable success. In the reign of
Constantius, availing himself of his power as a Christian Bishop, he
demolished a heathen temple, and built a church on its site. When
Julian succeeded, it was Mark's turn to suffer. The Emperor had been
saved by him, when a child, on the massacre of the other princes of
his house; but on this occasion he considered that the claims at once
of justice and of paganism outweighed the recollection of ancient
services. Mark was condemned to rebuild the temple, or to pay
the price of it; and, on his flight from his bishoprick, many of his
flock were arrested as his hostages. Upon this, he surrendered himself
to his persecutors, who immediately subjected him to the most
revolting, as well as the most cruel indignities. ``They
apprehended the aged prelate,'' says Gibbon, selecting some out of
the number, ``they inhumanly scourged him; they tore his beard;
and his naked body, anointed with honey, was suspended, in a net,
between heaven and earth, and exposed to the stings of insects and the
rays of a Syrian sun.'' \footnote{Gibbon, Hist. ch. xxiii.}
The payment of one piece of gold towards the rebuilding of the temple,
would have rescued him from these torments; but, resolute in his
refusal to contribute to the service of idolatry, he allowed himself,
with a generous insensibility, even to jest at his own sufferings \footnote{Soz. v. 10.}, till he wore out the fury, or even, it is said, effected the
conversion of his persecutors. Gregory Nazianzen, and Theodoret,
besides celebrating his activity in making converts, make mention of
his wisdom and piety, his cultivated understanding, his love of
virtue, and the honourable consistency of his life \footnote{Tillem. Mem. vol. vii. p. 340.}.

Cyril of Jerusalem, and Eusebius of Samosata, are both Saints in
the Roman Calendar, though connected in history with the Semi-Arian
party. Eusebius was the friend of St. Basil, surnamed the Great; and
Cyril is still known to us in his perspicuous and eloquent discourses
addressed to the Catechumens.

Others might be named of a like respectability, though deficient,
with those above-mentioned, either in moral or in intellectual
judgment. With these were mingled a few of a darker character. George
of Laodicea, one of the genuine Eusebians, joined them for a time, and
took a chief share together with Basil in the management of the
Council of Ancyra. Macedonius, who was originally an Anomœan, passed
through Semi-Arianism to the heresy of the Pneumatomachists, that is,
the denial of the Divinity of the Holy Ghost, of which he is
theologically the founder.

3.

The Semi-Arians, being such as above described, were at first both
in faith and conduct an ornament and recommendation of the Eusebians.
But, when once the latter stood at variance with the Latin Church by
the event of the Sardican Council, they ceased to be of service to
them as a blind, which was no longer available, or rather were an
incumbrance to them, and formidable rivals in the favour of
Constantius. The separation between the two parties was probably
retarded for a while by the forced submission and recantation of the
Eusebian Valens and Ursacius; but an event soon happened, which
altogether released those two Bishops and the rest of the Eusebians
from the embarrassments, in which the influence of the West and the
timidity of Constantius had for the moment involved them. This was the
assassination of the Catholic Constans which took place A.D.
350; in consequence of which (Constantine, the eldest of the brothers,
being already dead) Constantius succeeded to the undivided empire.
Thus the Eusebians had the whole of the West opened to their ambition
\footnote{[The Eusebians, or political party, were
renewed in the Acacians, immediately to be mentioned, Athanasius
calling the latter the heirs of the former, (Hist. Arian. §§ 19 and
28; vid. also Ath. Tr. vol. ii. p. 28.) He ever distinguishes the
Arians proper from the Eusebians (in his Ep. Enc. and Apol. Contr.
Arian.), as afterwards the Anomœans were to be distinguished from the
Acacians.]}; and were bound
by no impediment, except such as the ill-instructed Semi-Arianism of
the Emperor might impose upon them. Their proceedings under these
fortunate circumstances will come before us presently; here I will
confine myself to the mention of the artifice, by which they succeeded
in recommending themselves to Constantius, while they opposed and
triumphed over the Semi-Arian Creed.

This artifice, which, obvious as it is, is curious, from the place
which it holds in the history of Arianism, was that of affecting on
principle to limit confessions of faith to Scripture terms; and was
adopted by Acacius, Bishop of Cæsarea, in Palestine, the successor of
the learned Eusebius, one of the very men, who had advocated the
Semi-Arian non-scriptural formularies of the Dedication and of
Philippopolis \footnote{Athan. de Syn. 36-38.}. From
the earliest date, the Arians had taken refuge from the difficulties
of their own unscriptural dogmas in the letter of the sacred writers;
but they had scarcely ventured on the inconsistency of objecting to
the terms of theology, as such. But here Eusebius of Cæsarea
anticipated the proceedings of his party and, as he opened upon his
contemporaries the evasion of Semi-Arianism, so did he also anticipate
his pupil Acacius, in the more specious artifice now under
consideration. It is suggested in the apology which he put forth for
signing the Nicene anathema of the Arian formulæ; which anathema he
defends on the principle, that these formulæ were not conceived in the language of Scripture \footnote{Vid. also Theod. Hist. ii. 3. [who tells
us that the objection of ``unscripturalness'' had been
suggested to Constantius by the Arian priest, the favourite of
Constantia, to whom Constantine had entrusted his will. Eusebius, in
his Letter about the Nicene Creed, does scarcely more than glance at
this objection.]}. Allusion is made to the same principle from time to time in
the subsequent Arian Councils, as if even then the laxer Eusebians
were struggling against the dogmatism of the Semi-Arians. Though the
Creed of Lucian introduces the ``usia,'' the three other
Creeds of the Dedication omit it; and this hypothesis of differences
of opinion in the heretical body at these Councils partly accounts for
that hesitation and ambiguity in declaring their faith, which has been
noticed in its place. Again, the Macrostich omits the ``usia,''
professes generally that the Son is ``\emph{like in all things to the
Father},'' and enforces the propriety of keeping to the
language of Scripture \footnote{Vid. Athan. de Synod.}.

About the time which is at present more particularly before us,
that is, after the death of Constans, this modification of Arianism
becomes distinct, and collects around it the Eastern Eusebians, under
the skilful management of Acacius. It is not easy to fix the date of
his openly adopting it; the immediate cause of which was his quarrel
with the Semi-Arian Cyril, which lies between A.D.
349-357. The distinguishing principle of his new doctrine was
adherence to the Scripture phraseology, in opposition to the
inconvenient precision of the Semi-Arians; its distinguishing tenet is
the vague confession that the Son is generally ``\emph{like},''
or at most ``\emph{in all things like}'' the Father,—``\emph{like}''
as opposed to the ``\emph{one in substance},'' ``\emph{like
in substance},'' and ``\emph{unlike},'' \footnote{[\emph{homoion}] or [\emph{kata panta homoion}]
is the tenet of the Acacians or Homœans, as opposed to Catholic [\emph{homoousion}],
the Semi-Arian [\emph{homoiousion}], and the [\emph{anomoion}] of the
Eunomians or Aetians. [St. Cyril, however adopts the [\emph{kata panta
homoion}], as does Damascene.]}—that is, the vague confession that the Son is \emph{generally
like}, or \emph{altogether like}, the Father. Of these two
expressions, the ``\emph{in all things like}'' was allowed by
the Semi-Arians, who included ``\emph{in substance}'' under
it; whereas the Acacians (for so they may now be called), or Homœans
(as holding the \emph{Homœon} or \emph{like}), covertly intended to
exclude the ``\emph{in substance}'' by that very expression,
mere similarity always implying difference, and ``\emph{substance}''
being, as they would argue, necessarily excluded from the ``\emph{in
all things},'' if the ``\emph{like}'' were intended to
stand for any thing short of identity. It is plain then that, in the
meaning of its authors, and in the practical effect of it, this new
hypothesis was neither more nor less than the pure Arian, or, as it
was afterwards called, Anomœan, though the phrase, in which it was
conveyed, bore in its letter the reverse sense.

Such was the state of the heresy about the year 350; before
reviewing its history, as carried on between the two rival parties
into which its advocates, the Eusebians, were dividing, the Semi-Arian
and Homœan, I shall turn to the sufferings of the Catholic Church at
that period.

\xsection{Section 3. The Athanasians}

THE second Arian Persecution is spread
over the space of about twelve years, being the interval between the
death of Constans, and that of Constantius (A.D.
350-361). Various local violences, particularly at Alexandria and
Constantinople, had occurred with the countenance of the Eusebians at
an earlier date; but they were rather acts of revenge, than intended
as means of bringing over the Catholics, and were conducted on no
plan. The chief sees, too, had been seized, and their occupants
banished. But now the alternative of subscription or suffering was
generally introduced; and, though Arianism was more sanguinary in its
later persecutions, it could not be more audacious and abandoned than
it showed itself in this.

The artifice of the Homœon, of which Acacius had undertaken the
management, was adapted to promote the success of his party, among the
orthodox of the West, as well as to delude or embarrass the Oriental
Semi-Arians, for whom it was particularly provided. The Latin
Churches, who had not been exposed to those trials of heretical
subtlety of which the Homoüsion was reluctantly made the
remedy, had adhered with a noble simplicity to the decision of Nicæa;
being satisfied (as it would seem), that, whether or not they had need
of the test of orthodoxy at present, in it lay the security of the
great doctrine in debate, whenever the need should come. At the same
time, they were naturally jealous of the introduction of such terms
into their theology, as chiefly served to remind them of the
dissensions of foreigners; and, as influenced by this feeling, even
after their leaders had declared against the Eusebians at Sardica,
they were exposed to the temptation of listening favourably to the
artifice of the ``\emph{Homœon}'' or ``\emph{like}.''
To shut up the subject in Scripture terms, and to say that our Lord
was like His Father, no explanation being added, seemed to be a
peaceful doctrine, and certainly was in itself unexceptionable; and,
of course would wear a still more favourable aspect, when contrasted
with the threat of exile and poverty, by which its acceptance was
enforced. On the other hand, the proposed measure veiled the grossness
of that threat itself, and fixed the attention of the solicited
Churches rather upon the argument, than upon the Imperial command.
Minds that are proof against the mere menaces of power, are overcome
by the artifices of an importunate casuistry. Those, who would rather
have suffered death than have sanctioned the impieties of Arius,
hardly saw how to defend themselves in refusing creeds, which were
abstractedly true, though incomplete, and intolerable only because the
badges of a prevaricating party. Thus Arianism gained its first
footing in the West. And, when one concession was made, another was
demanded; or, at other times, the first concession was
converted, not without speciousness, into a principle, as allowing
change altogether in theological language, as if to depart from the
Homoüsion were in fact to acquiesce in the open impieties of Arius
and the Anomœans. This is the character of the history as more or
less illustrated in this and the subsequent Section; the Catholics
being harassed by sophistry and persecution, and the Semi-Arians first
acquiescing in the Homœon, then retracting, and becoming more
distinct upon the scene, as the Eusebians or Acacians ventured to
speak of our Lord in less honourable terms.

But there was another subscription, required of the Catholics
during the same period and from an earlier date, as painful, and to
all but the most honest minds as embarrassing, as that to the creed of
the Homœon; and that was the condemnation of Athanasius. The
Eusebians were incited against him by resentment and jealousy; they
perceived that the success of their schemes was impossible, while a
Bishop was on the scene, so popular at home, so respected abroad, the
bond of connexion between the orthodox of Europe and Asia, the organ
of their sentiments, and the guide and vigorous agent of their
counsels. Moreover, the circumstances of the times had attached an
adventitious importance to his fortunes; as if the cause of the Homoüsion
were providentially committed to his custody, and in his safety or
overthrow, the triumph or loss of the truth were actually involved.
And, in the eyes of the Emperor, the Catholic champion appeared as a
rival of his own sovereignty; type, as he really was, and instrument
of that Apostolic Order, which, whether or not united to the civil
power, must, to the end of time, divide the rule with Cæsar as
the minister of God. Considering then Athanasius too great for a
subject, Constantius, as if for the peace of his empire, desired his
destruction at any rate \footnote{Gibbon, Hist. ch. xxi.}.
Whether he was unfortunate or culpable it mattered not; whether
implicated in legal guilt, or forced by circumstances into his present
position; still he was the fit victim of a sort of ecclesiastical
ostracism, which, accordingly, he called upon the Church to inflict.
He demanded it of the Church, for the very eminence of Athanasius
rendered it unsafe, even for the Emperor, to approach him in any other
way. The Patriarch of Alexandria could not be deposed, except after a
series of successes over less powerful Catholics, and with the forced
acquiescence or countenance of the principal Christian communities.
And thus the history of the first few years of the persecution,
presents to us the curious spectacle of a party warfare raging
everywhere, except in the neighbourhood of the person who was the real
object of it, and who was left for a time to continue the work of God
at Alexandria, unmolested by the Councils, conferences, and
usurpations, which perplexed the other capitals of Christendom.

As regards the majority of Bishops who were called upon to condemn
him, there was, it would appear, little room for error of judgment, if
they dealt honestly with their consciences. Yet, in the West, there
were those, doubtless, who hardly knew enough of him to give him their
confidence, or who had no means of forming a true opinion of the fresh
charges to which he was subjected. Those, which were originally urged against him, have already been stated; the new allegations were
as follows: that he had excited differences between Constantius and
his brother; that he had corresponded with Magnentius, the usurper of
the West; that he had dedicated, or used, a new Church in Alexandria
without the Emperor's leave; and lastly, that he had not obeyed his
mandate summoning him to Italy.—Now to review some of the prominent
passages in the persecution:—

1.

Paul had succeeded Alexander in the See of Constantinople, A.D.
336. At the date before us (A.D. 350), he had
already been thrice driven from his Church by the intrigues of the
Arians; Pontus, Gaul, and Mesopotamia, being successively the places
of his exile. He had now been two years restored, when he was called a
fourth time, not merely to exile, but to martyrdom. By authority of
the Emperor, he was conveyed from Constantinople to Cucusus in
Cappadocia, a dreary town amid the deserts of the Taurus, afterwards
the place of banishment of his successor St. Chrysostom. Here he was
left for six days without food; when his conductors impatiently
anticipated the termination of his sufferings by strangling him in
prison. Macedonius, the Semi-Arian, took possession of the vacant see,
and maintained his power by the most savage excesses. The confiscation
of property, banishment, brandings, torture, and death, were the means
of his accomplishing in the Church of Constantinople, a conformity
with the tenets of heresy. The Novatians, as maintaining the
Homoüsion, were included in the persecution. On their refusing to
communicate with him, they were seized and scourged, and the
sacred elements violently thrust into their mouths. Women and children
were forcibly baptized; and, on the former resisting, they were
subjected to cruelties too miserable to be described.

2.

The sufferings of the Church of Hadrianople occurred about the same
time, or even earlier. Under the superintendence of a civil officer,
who had already acted as the tool of the Eusebians in the Mareotis,
several of the clergy were beheaded; Lucius, their Bishop, for the
second time loaded with chains and sent into exile, where he died; and
three other Bishops of the neighbourhood visited by an Imperial edict,
which banished them, at the peril of their lives, from all parts of
the Empire.

3.

Continuing their operations westward, the Arians next possessed
themselves of the province of Sirmium in Pannonia, in which the
dioceses of Valens and Ursacius were situated. These Bishops, on the
death of Constans, had relapsed into the heresy of his brother, who
was now master of the whole Roman world; and from that time they may
be accounted as the leaders of the Eusebian party, especially in the
West. The Church of Sirmium was opened to their assaults under the
following circumstances. It had always been the policy of the Arians
to maintain that the Homoüsion involved some or other heresy by
necessary consequence. A Valentinian or a Manichean materialism was
sometimes ascribed to the orthodox doctrine; and at another
time, Sabellianism, which was especially hateful to the Semi-Arians.
And it happened, most unhappily for the Church, that one of the most
strenuous of her champions at Nicæa, had since fallen into a heresy
of a Sabellian character; and had thus confirmed the prejudice against
the true doctrine, by what might be taken to stand as an instance of
its dangerous tendency. In the course of a work in refutation of the
Sophist Asterius, one of the first professed Semi-Arians, Marcellus,
Bishop of Ancyra, was led to simplify (as he conceived) the creed of
the Church, by statements which savoured of Sabellianism; that is, he
maintained the unity of the Son with the Father, at the expense of the
doctrine of the personal distinction between the Two. He was answered,
not only by Asterius himself, but by Eusebius of Cæsarea and Acacius;
and, A.D. 335, he was deposed from his see by
the Eusebians, in order to make way for the Semi-Arian Basil. In spite
of the suspicions against him, the orthodox party defended him for a
considerable time, and the Council of Sardica (A.D.
347) acquitted him and restored him to his see; but at length, perhaps
on account of the increasing definiteness of his heretical views, he
was abandoned by his friends as hopeless, even by Athanasius, who
quietly put him aside with the acquiescence of Marcellus himself. But
the evil did not end there; his disciple Photinus, Bishop of Sirmium,
increased the scandal, by advocating, and with greater boldness, an
almost Unitarian doctrine. The Eusebians did not neglect the
opportunity thus offered them, both to calumniate the Catholic
teaching, and to seize on so considerable a see, which its present
occupier had disgraced by his heresy. They held a Council at
Sirmium (A.D. 351), to inquire into his
opinions; and at his request a formal disputation was held. Basil, the
rival of Marcellus, was selected to be the antagonist of his pupil;
and having the easier position to defend, gained the victory in the
judgment of impartial arbiters, who had been selected. The deposition
of Photinus followed, and an Arian, Germinius, placed in his see. Also
a new creed was promulgated of a structure between Homœusian and
Homœan, being the first of three which are dated from Sirmium.
Germinius some years afterwards adopted a Semi-Arianism bordering upon
the Catholic doctrine, and that at a time when it may be hoped that
secular views did not influence his change.

4.

The first open attack upon Athanasius and the independence of the
West, was made two years later at Aries, at that time the residence of
the Court. There the Emperor held a Council, with the intention of
committing the Bishops of the West to an overt act against the
Alexandrian prelate. It was attended by the deputies of Liberius, the
new Bishop of Rome, whom the Eusebian party had already addressed,
hoping to find him more tractable than his predecessor Julius.
Liberius, however, had been decided in Athanasius's favour by the
Letter of an Egyptian Council; and, in order to evade the Emperor's
overtures, he addressed to him a submissive message, petitioning him
for a general and final Council at Aquileia, a measure which
Constantius had already led the Catholics to expect. The Western
Bishops at Arles, on their part, demanded that, as a previous
step to the condemnation of Athanasius, the orthodox Creed should be
acknowledged by the Council, and Arius anathematized. However, the
Eusebians carried their point; Valens followed up with characteristic
violence the imperiousness of Constantius; ill treatment was added,
till the Fathers of the Council, worn out by sufferings, consented to
depose and even excommunicate Athanasius. Upon this, an edict was
published, denouncing punishment on all Bishops who refused to
subscribe the decree thus obtained. Among the instances of cowardice,
which were exhibited at Arles, none was more lamentable than that of
Vincent of Capua, one of the deputies from Liberius to the Emperor.
Vincent had on former occasions shown himself a zealous supporter of
orthodoxy. He is supposed to be the presbyter of the same name who was
one of the representatives of the Roman Bishop at Nicæa; he had acted
with the orthodox at Sardica, and had afterwards been sent by Constans
to Constantius, to effect the restoration of the Athanasians in A.D.
348. It was on this occasion, that he and his companion had been
exposed to the malice of Stephen, the Arian Bishop of Antioch; who,
anxious to destroy their influence, caused a woman of light character
to be introduced into their chamber, with the intention of founding a
calumny against them; and who, on the artifice being discovered, was
deposed by order of Constantius. On the present occasion, Vincent was
entirely in the confidence of Liberius; who, having entrusted him with
his delicate commission from a sense of his vigour and experience, was
deeply afflicted at his fall. It is satisfactory to know, that
Vincent retrieved himself afterwards at Ariminum; where he boldly
resisted the tyrannical attempt of the Eusebians, to force their creed
on the Western Church.

5.

Times of trial bring forward men of zeal and boldness, who thus are
enabled to transmit their names to posterity. Liberius, downcast at
the disgrace of his representative, and liable himself to fluctuations
of mind, was unexpectedly cheered by the arrival of the famous
Lucifer, Bishop of Cagliari, in Sardinia, and Eusebius of Vercellæ.
These, joined by a few others, proceeded as his deputies and advocates
to the great Council of Milan, which was held by Constantius (A.D.
355), two years later than that in which Vincent fell. The Fathers
collected there were in number above 300, almost all of the Western
Church. Constantius was present, and Valens conducted the Arian manœuvres;
and so secure of success were he and his party, that they did not
scruple to insult the Council with the proposal of a pure Arian, or
Anomœan, creed.

Whether this creed was generally subscribed, does not appear; but
the condemnation of Athanasius was universally agreed upon, scarcely
one or two of the whole number refusing to sign it. This is
remarkable; inasmuch as, at first, the Occidentals demanded of the
Eusebians an avowal of the orthodox faith, as the condition of
entering upon the consideration of the charges against him. But herein
is the strength of audacious men; who gain what is unjust, by asking
what is extravagant. Sozomen attributes the concession of the
Council to fear, surprise, and ignorance \footnote{Soz. iv. 9.}. In truth, a collection of men, who were strangers to each
other, and without organization or recognized leaders, without
definite objects or policy, was open to every variety of influence,
which the cleverness of the usurping faction might direct against
them. The simplicity of honesty, the weakness of an amiable temper,
the inexperience of a secluded life, and the slowness of the
unpractised intellect, all combined with their alarm at the Emperor's
manifested displeasure, to impel them to take part with his heresy.
When some of them ventured to object the rule of the Church against
his command, that they should condemn Athanasius, and communicate with
the Arians, ``My will must be its rule,'' he replied; ``so
the Syrian Bishops have decided; and so must yourselves, would you
escape exile.''

Several of the more noble-minded prelates of the principal Churches
submitted to the alternative, and left their sees. Dionysius, Exarch
of Milan, was banished to Cappadocia or Armenia, where he died before
the end of the persecution; Auxentius being placed in his see, a
bitter Arian, brought for the purpose from Cappadocia, and from his
ignorance of Latin, singularly ill-fitted to preside over a Western
province. Lucifer was sent off into Syria, and Eusebius of Vercellæ
into Palestine. A fresh and more violent edict was published against
Athanasius; orders were given to arrest him as an impious person, and
to put the Arians in possession of his churches, and of the
benefactions, which Constantine had left for ecclesiastical and
charitable uses. All Bishops were prohibited from communion with
him, under pain of losing their sees; and the laity were to be
compelled by the magistrates to join themselves to the heretical
party. Hilary of Poictiers was the next victim of the persecution. He
had taken part in a petition, presented to Constantius, in behalf of
the exiled bishops. In consequence a Gallic Council was called, under
the presidency of Saturninus, Bishop of Arles; and Hilary was banished
into Phrygia.

6.

The history of Liberius, the occupier of the most powerful see in
the West, possesses an interest, which deserves our careful attention.
In 356, the year after the Council of Milan, the principal eunuch of
the Imperial Court had been sent, to urge on him by threats and
promises the condemnation of Athanasius; and, on his insisting on a
fair trial for the accused, and a disavowal of Arianism on the part of
his accusers, as preliminary conditions, had caused him to be forced
away to Milan. There the same arguments were addressed to him in the
more impressive words of the Emperor himself; who urged upon him
``the notoriously wicked life of Athanasius, his vexatious
opposition to the peace of the Church, his intrigues to effect a
quarrel between the imperial brothers, and his frequent condemnation
in the Councils of Eastern and Western Christendom;'' and further
exhorted him, as being by his pastoral office especially a man of
peace, to be cautious of appearing the sole obstacle to the happy
settlement of a question, which could not otherwise be arranged.
Liberius replied by demanding of Constantius even more than his own
deputies had proposed to the Milanese Council;—first, that
there should be a general subscription to the Nicene faith throughout
the Church; next, that the banished bishops should be restored to
their sees; and lastly, should the trial of Athanasius be still
thought advisable, that a Council should be held at Alexandria, where
justice might be fairly dealt between him and his accusers. The
conference between them ended in Liberius being allowed three days to
choose between making the required subscription, and going into exile;
at the end of which time he manfully departed for Berœa, in Thrace.
Constantius and the empress, struck with the nobleness of his conduct,
sent after him a thousand pieces of gold; but he refused a gift, which
must have laid him under restraint towards heretical benefactors. Much
more promptly did he reject the offer of assistance, which Eusebius,
the eunuch before-mentioned, from whatever feeling, made him.
``You have desolated the Churches of Christendom,'' he said to
the powerful favourite, ``and then you offer me alms as a convict.
Go, first learn to be a Christian.'' \footnote{Soz. iv. 11. Theod. Hist. ii. 16.}

There are men, in whose mouths sentiments, such as these, are
becoming and admirable, as being the result of Christian magnanimity,
and imposed upon them by their station in the Church. But the sequel
of the history shows, that in the conduct of Liberius there was more
of personal feeling and intemperate indignation, than of deep-seated
fortitude of soul. His fall, which followed, scandalous as it is in
itself, may yet be taken to illustrate the silent firmness of those
others his fellow-sufferers, of whom we hear less, because they
bore themselves more consistently. Two years of exile, among the
dreary solitudes of Thrace, broke his spirit; and the triumph of his
deacon Felix, who had succeeded to his power, painfully forced upon
his imagination his own listless condition, which brought him no work
to perform, and no witness of his sufferings for the truth's sake.
Demophilus, one of the foremost of the Eusebian party, was bishop of
Berœa, the place of Liberius's banishment; and gave intelligence of
his growing melancholy to his own associates. Wise in their
generation, they had an instrument ready prepared for the tempter's
office. Fortunatian, Bishop of Aquileia, who stood high in the opinion
of Liberius for disinterestedness and courage, had conformed to the
court-religion in the Arian Council of Milan; and he was now employed
by the Eusebians, to gain over the wavering prelate. The arguments of
Fortunatian and Demophilus shall be given in the words of Maimbourg.
``They told him, that they could not conceive, how a man of his
worth and spirit could so long obstinately resolve to be miserable
upon a chimerical notion, which subsisted only in the imagination of
people of weak or no understanding: that, indeed, if he suffered for
the cause of God and the Church, of which God had given him the
government, they should not only look upon his sufferings as glorious,
but, being willing to partake of his glory, they should also become
his companions in banishment themselves. But that this matter related
neither to God nor religion; that it concerned merely a private
person, named Athanasius, whose cause had nothing in common with that
of the Church, whom the public voice had long since accused of
numberless crimes, whom Councils had condemned, and who had been
turned out of his see by the great Constantine, whose judgment alone
was sufficient to justify all that the East and West had so often
pronounced against him. That, even if he were not so guilty as men
made him, yet it was necessary to sacrifice him to the peace of the
Church, and to throw him into the sea to appease the storm, which he
was the occasion of raising; but that, the greater part of the Bishops
having condemned him, the defending him would be causing a schism, and
that it was a very uncommon sight to see the Roman prelate abandon the
care of the Church, and banish himself into Thrace, to become the
martyr of one, whom both divine and human justice had so often
declared guilty. That it was high time to undeceive himself, and to
open his eyes at last; to see, whether it was not passion in
Athanasius, which gave a false alarm, and opposed an imaginary heresy,
to make the world believe that they had a mind to establish
error.'' \footnote{Webster's translation is used: one or two
irrelevant phrases, introduced by Maimbourg on the subject of Roman
supremacy, being omitted.}

The arguments, diffusively but instructively reported in the above
extract, were enforced by the threat of death as the consequence of
obstinacy; while, on the other hand, a temptation of a peculiar nature
presented itself to the exiled bishop in his very popularity with the
Roman people, which was such, that Constantius had already been
obliged to promise them his restoration. Moreover, as if to give a
reality to the inducements by which he was assailed, a specific plan
of mutual concession and concord had been projected, in which Liberius
was required to take part. The Western Catholics were, as we
have seen, on all occasions requiring evidence of the orthodoxy of the
Eusebians, before they consented to take part with them against
Athanasius. Constantius then, desirous of ingratiating himself with
the people of Rome, and himself a Semi-Arian, and at that time alarmed
at the increasing boldness of the Anomœans, or pure Arians, presently
to be mentioned, perceived his opportunity for effecting a general
acceptance of a Semi-Arian creed; and thus, while sacrificing the
Anomœans, whom he feared, to the Catholics, and claiming from the
Catholics in turn what were scarcely concessions, in the imperfect
language of the West, for realizing that religious peace, which he
held to be incompatible with the inflexible orthodoxy of Athanasius.
Moreover, the heresies of Marcellus and Photinus were in favour of
this scheme; for, by dwelling upon them, he withdrew the eyes of
Catholics from the contrary errors of Semi-Arianism. A creed was
compiled from three former confessions, that of the orthodox Council
against Paulus (A.D. 264), that of the
Dedication (A.D. 341), and one of the three
published at Sirmium. Thus carefully composed, it was signed by all
parties, by Liberius \footnote{[Vide supr. pp. 131, 294, 323. There is much
difference of opinion, however, among writers, which was the creed
which Liberius signed; vide Appendix, No. 3.]}, by
the Semi-Arians, and by the Eusebians; the Eusebians being compelled
by the Emperor to submit for the time to the dogmatic formulæ, which
they had gradually abandoned. Were it desirable to enlarge on this
miserable apostasy, there are abundant materials in the letters, which
Liberius wrote in renunciation of Athanasius, to his clergy, and to
the Arian bishops. To Valens he protests, that nothing but his
love of peace, greater than his desire of martyrdom itself, would have
led him to the step which he had taken; in another he declares, that
he has but followed his conscience in God's sight \footnote{Hilar. Fragm. iv. and vi. [The authority for
Fragm. is very doubtful.]}. To add to his misery, Constantius suffered him for a while to
linger in exile, after he had given way. At length he was restored;
and at Ariminum in a measure retrieved his error, together with
Vincent of Capua.

7.

The sufferings and trials of Hosius, which took place about the
same time, are calculated to impress the mind with the most sorrowful
feelings, and still more with a lively indignation against his inhuman
persecutors. Shortly before the conference at Sirmium, at which
Liberius gave his allegiance to the supremacy of Semi-Arianism, a
creed had been drawn up in the same city by Valens and the other more
daring members of the Eusebian body. It would seem, that at this date
Constantius had not taken the alarm against the Anomœans, to the
extent in which he felt it soon afterwards, on the news probably of
their proceedings in the East. Accordingly, the creed in question is
of a mixed character. Not venturing on the \emph{Anomœon}, as at
Milan, it nevertheless condemns the use of the \emph{usia} (\emph{substance}),
\emph{Homoüsion}, and \emph{Homœüsion}, on somewhat of the
equivocal plan, of which Acacius, as I have said above, was the most
conspicuous patron; and being such, it was presented for signature to
the aged Bishop of Corduba. The cruelty which they exercised to
accomplish their purpose, was worthy of that singularly wicked faction
which Eusebius had organized. Hosius was at this time 101 years old;
and had passed a life, prolonged beyond the age of man, in services
and sufferings in the cause of Christ. He had assisted in the
celebrated Council of Elvira, in Spain (about the year 300), and had
been distinguished as a confessor in the Maximinian persecution. He
presided at the General Councils of Nicæa and Sardica, and was
perhaps the only Bishop, besides Athanasius, who was known and
reverenced at once in the East and West. When Constantius became
possessed of the Western world, far from relaxing his zeal in a cause
discountenanced at the Court, Hosius had exerted himself in his own
diocese for the orthodox faith; and, when the persecution began,
endeavoured by letter to rouse other bishops to a sense of the
connexion between the acquittal of Athanasius, and the maintenance of
divine truth. The Eusebians were irritated by his opposition; he was
summoned to the Court at Milan, and, after a vain attempt to shake his
constancy, dismissed back to his see. The importunities of Constantius
being shortly after renewed, both in the way of threats and of
promises, Hosius addressed him an admirable letter, which Athanasius
has preserved. After declaring his willingness to repeat, should it be
necessary, the good confession which he had made in the heathen
persecution, he exhorts the Emperor to abandon his unscriptural creed,
and to turn his ear from Arian advisers. He states his conviction,
that the condemnation of Athanasius was urged merely for the
establishment of the heresy; declares, that at Sardica his accusers
had been challenged publicly to produce the proof of their
allegations, and had failed, and that he himself had conversed with
them in private, and could gain nothing satisfactory from them; and he
further reminds Constantius, that Valens and Ursacius had before now
retracted the charges, which they once urged against him. ``Change
your course of action, I beseech you,'' continues the earnest
Prelate; ``remember that you are a man. Fear the day of judgment;
keep your hands clean against it; meddle not with Church matters; far
from advising us about them, rather seek instruction from us. God has
put dominion into your hands; to us He has entrusted the management of
the Church; and, as a traitor to you is a rebel to the God who
ordained you, so be afraid on your part, lest, usurping ecclesiastical
power, you become guilty of a great sin. It is written, 'Render unto
Cæsar, Cæsar's, and what is God's, to God.' We may not bear rule;
you, O Emperor, may not burn incense. I write this from a care for
your soul. As to your message, I remain in the same mind. I do not
join the Arians. I anathematize them. I do not subscribe the
condemnation of Athanasius.'' \footnote{Athan. Hist. Arian. ad Monach. 44.} Hosius did not address such language with impunity to a Court,
which affected the majesty of oriental despotism. He was summoned to
Sirmium, and thrown into prison. There he remained for a whole year.
Tortures were added to force the old man from his resolution. He was
scourged, and afterwards placed upon the rack. Mysterious it was, that
so honoured a life should be preserved to an extremity of age, to
become the sport and triumph of the Enemy of mankind. At length broken
in spirit, the contemporary of Gregory and Dionysius \footnote{Vide supr. p. 125.} was induced to countenance the impieties of the generation,
into which he had lived; not indeed signing the condemnation of
Athanasius, for he spurned that baseness to the last, but yielding
subscription to a formulary, which forbad the mention of the \emph{Homoüsion},
and thus virtually condemned the creed of Nicæa, and countenanced the
Arian proceedings. Hosius lived about two years after this tragical
event: and, on his deathbed, he protested against the compulsion which
had been used towards him, and, with his last breath, abjured the
heresy which dishonoured his Divine Lord and Saviour.

8.

Meanwhile, the great Egyptian prelate, seated on his patriarchal
throne, had calmly prosecuted the work, for which he was raised up, as
if his name had not been mentioned in the Arian Councils, and the
troubles, which agitated the Western Church, were not the prelude to
the blow, which was to fall on himself. Untutored in concession to
impiety, by the experience or the prospect of suffering, yet,
sensitively alive to the difference between misbelief and
misapprehension, while he punished he spared, and restored in the
spirit of meekness, while he rebuked and rejected with power. On his
return to Alexandria, seven years previous to the events last
recorded, congratulations and professions of attachment poured in upon
him from the provinces of the whole Roman world, near and distant.
From Africa to Illyricum, and from England to Palestine, 400
episcopal letters solicited his communion or patronage; and apologies,
and the officiousness of personal service were liberally tendered by
those, who, through cowardice, dulness, or self-interest, had joined
themselves to the heretical party. Nor did Athanasius fail to improve
the season of prosperity, for the true moral strength and substantial
holiness of the people committed to him. The sacred services were
diligently attended; alms and benefactions supplied the wants of the
friendless and infirm; and the young turned their thoughts to that
generous consecration of themselves to God, recommended by St. Paul in
times of trouble and persecution.

In truth the sufferings, which the Church of Alexandria had lately
undergone from the hands of the Eusebians, were sufficient to
indispose serious minds towards secular engagements, or vows of duty
to a fellow-mortal; to quench those anticipations of quietness and
peace, which the overthrow of paganism had at first excited; and to
remind them, that the girdle of celibacy and the lamp of watchers best
became those, on whom God's judgments might fall suddenly. Not more
than ten years were gone by, since Gregory, appointed to the see of
Athanasius by the Council of the Dedication \footnote{Vid. supra, p. 286.}, had been thrust upon them by the Imperial Governor, with the
most frightful and revolting outrages. Philagrius, an apostate from
the Christian faith, and Arsacius, an eunuch of the Court, introduced
the Eusebian Bishop into his episcopal city. A Church besieged and
spoiled, the massacre of the assembled worshippers, the clergy trodden
under foot, the women subjected to the most infamous profanations, these were the first benedictory greetings scattered by
the Arian among his people. Next, bishops were robbed, beaten,
imprisoned, banished; the sacred elements of the Eucharist were
scornfully cast about by the heathen rabble, which seconded the
usurping party; birds and fruits were offered in sacrifice on the holy
table; hymns chanted in honour of the idols of paganism; and the
Scriptures given to the flames.

Such had already been the trial of a much-enduring Church; and it
might suddenly be renewed in spite of its present prosperity. The
Council of Sardica, convoked principally to remedy these miserable
disorders, had in its Synodal Letter warned the Alexandrian Catholics
against relaxing in the brave testimony they were giving to the faith
of the Gospel. ``We exhort you, beloved brethren, before all
things, that ye hold the right faith of the Catholic Church. Many and
grievous have been your sufferings, and many are the insults and
injuries inflicted on the Catholic Church, but 'he, who endureth unto
the end, the same shall be saved.' Wherefore, should they essay
further enormities against you, let affliction be your rejoicing. For
such sufferings are a kind of martyrdom, and such confessions and
tortures have their reward. Ye shall receive from God the combatant's
prize. Wherefore struggle with all might for the sound faith, and for
the exculpation of our brother Athanasius, your bishop. We on our part
have not been silent about you, nor neglected to provide for your security;
but have been mindful, and done all that Christian love requires of
us, suffering with our suffering brethren, and accounting their trials
as our own.'' \footnote{Athan. Apol. cont. Arian. 38.}

The time was now at hand, which was anticipated by the prophetic
solicitude of the Sardican Fathers. The same year in which Hosius was
thrown into prison, the furies of heretical malice were let loose upon
the Catholics of Alexandria. George of Cappadocia, a man of illiterate
mind and savage manners, was selected by the Eusebians as their new
substitute for Athanasius in the see of that city; and the charge of
executing this extraordinary determination was committed to Syrianus,
Duke of Egypt. The scenes which followed are but the repetition, with
more aggravated horrors, of the atrocities perpetrated by the intruder
Gregory. Syrianus entered Alexandria at night; and straightway
proceeded with his soldiers to one of the churches, where the
Alexandrians were engaged in the services of religion. We have the
account of the irruption from Athanasius himself; who, being accused
by the Arians of cowardice, on occasion of his subsequent flight,
after defending his conduct from Scripture, describes the
circumstances, under which he was driven from his Church. ``It was
now night,'' he says, ``and some of our people were keeping
vigil, as communion was in prospect; when the Duke Syrianus suddenly
came upon us, with a force of above 5000 men, prepared for attack,
with drawn swords, bows, darts, and clubs, ... and surrounded the
church with close parties of the soldiery, that none might escape from
within. There seemed an impropriety in my deserting my congregation in
such a riot, instead of hazarding the danger in their stead; so I
placed myself in my bishop's chair, and bade the deacon read the Psalm
(Ps. cxxxvi.), and the congregation alternate 'for His mercy
endureth for ever,' and then all retire and go home. But the General
bursting at length into the church, and his soldiers blocking up the
chancel, with a view of arresting me, the clergy and some of my people
present began in their turn clamorously to urge me to withdraw myself.
However, I refused to do so, before one and all in the church were
gone. Accordingly I stood up, and directed prayer to be said; and then
I urged them all to depart first, for that it was better that I should
run the risk, than any of them suffer. But by the time that most of
them were gone out, and the rest were following, the Religious
Brethren and some of the clergy, who were immediately about me, ran up
the steps, and dragged me down. And so, be truth my witness, though
the soldiers blockaded the chancel, and were in motion round about the
church, the Lord leading, I made my way through them, and by His
protection got away unperceived; glorifying God mightily, that I had
been enabled to stand by my people, and even to send them out before
me, and yet had escaped in safety from the hands of those who sought
me.'' \footnote{Athan. Apol. de Fug. 24.}

The formal protest of the Alexandrian Christians against this
outrage, which is still extant, gives a stronger and fuller statement
of the violences attending it. ``While we were watching in
prayer,'' they say, ``suddenly about midnight, the most noble
Duke Syrianus came upon us with a large force of legionaries, with
arms, drawn swords, and other military weapons, and their helmets on.
The prayers and sacred reading were proceeding, when they assaulted
the doors, and, on these being laid open by the force of numbers,
he gave the word of command. Upon which, some began to let fly
their arrows, and others to sound a charge; and there was a clashing
of weapons, and swords glared against the lamplight. Presently, the
sacred virgins were slaughtered, numbers trampled down one over
another by the rush of the soldiers, and others killed by arrows. Some
of the soldiers betook themselves to pillage, and began to strip the
females, to whom the very touch of strangers was more terrible than
death. Meanwhile, the Bishop sat on his throne, exhorting all to pray
... He was dragged down, and almost torn to pieces. He swooned away,
and became as dead; we do not know how he got away from them, for they
were bent upon killing him.'' \footnote{Athan. Hist. Arian. ad Monach. 81.}

The first purpose of Athanasius on his escape was at once to betake
himself to Constantius; and he had begun his journey to him, when news
of the fury, with which the persecution raged throughout the West,
changed his intention. A price was set on his head, and every place
was diligently searched in the attempt to find him. He retired into
the wilderness of the Thebaid, then inhabited by the followers of Paul
and Anthony, the first hermits. Driven at length thence by the
activity of his persecutors, he went through a variety of strange
adventures, which lasted for the space of six years, till the death of
Constantius allowed him to return to Alexandria.

His suffragan bishops did not escape a persecution, which was
directed, not against an individual, but against the Christian faith.
Thirty of them were banished, ninety were deprived of their churches;
and many of the inferior clergy suffered with them. Sickness and
death were the ordinary result of such hardships as exile involved;
but direct violence in good measure superseded a lingering and
uncertain vengeance. George, the representative of the Arians, led the
way in a course of horrors, which he carried through all ranks and
professions of the Catholic people; and the Jews and heathen of
Alexandria, sympathizing in his brutality, submitted themselves to his
guidance, and enabled him to extend the range of his crimes in every
direction. Houses were pillaged, churches were burned, or subjected to
the most loathsome profanations, and cemeteries were ransacked. On the
week after Whitsuntide, George himself surprised a congregation, which
had refused to communicate with him. He brought out some of the
consecrated virgins, and threatened them with death by burning, unless
they forthwith turned Arians. On perceiving their constancy of
purpose, he stripped them of their garments, and beat them so
barbarously on the face, that for some time afterwards their features
could not be distinguished. Of the men, forty were scourged; some died
of their wounds, the rest were banished. This is one out of many
notorious facts, publicly declared at the time, and uncontradicted;
and which were not merely the unauthorized excesses of an uneducated
Cappadocian, but recognized by the Arian body as their own acts, in a
state paper from the Imperial Court, and perpetrated for the
maintenance of the peace of the Church, and of a good understanding
among all who agreed in the authority of the sacred Scriptures.

In the manifesto, issued for the benefit of the people of
Alexandria (A.D. 356), the infatuated Emperor
applauds their conduct in turning from a cheat and impostor, and
siding with those who were venerable men, and above all praise.
``The majority of the citizens,'' he continues, ``were
blinded by the influence of one who rose from the abyss, darkly
misleading those who seek the truth; who had at no time any fruitful
exhortation to communicate, but abused the souls of his hearers with
frivolous and superficial discussions ... That noble personage has not
ventured to stand a trial, but has adjudged himself to banishment;
whom it is the interest even of the barbarians to get rid of, lest by
pouring out his griefs as in a play to the first comer, he persuade
some of them to be profane. So we will wish him a fair journey. But
for yourselves, only the select few are your equals, or rather, none
are worthy of your honours; who are allotted excellence and sense,
such as your actions proclaim, celebrated as they are almost in every
place ... You have roused yourselves from the grovelling things of
earth to those of heaven, the most reverend George undertaking to be
your leader, a man of all others the most accomplished in such
matters; under whose care you will enjoy in days to come honourable
hope, and tranquillity at the present time. May all of you hang upon
his words as upon a holy anchor, that any cutting and burning may be
needless on our part against men of depraved souls, whom we seriously
advise to abstain from paying respect to Athanasius, and to dismiss
from their minds his troublesome garrulity; or such factious men will
find themselves involved in extreme peril, which perhaps no skill will
be able to avert from them. For it were absurd indeed, to drive
about the pestilent Athanasius from country to country, aiming at his
death, though he had ten lives, and not to put a stop to the
extravagances of his flatterers and juggling attendants, such as it is
a disgrace to name, and whose death has long been determined by the
judges. Yet there is a hope of pardon, if they will desist from their
former offences. As to their profligate leader Athanasius, he
distracted the harmony of the state, and laid on the most holy men
impious and sacrilegious hands.'' \footnote{Athan. Apol. ad Constant. 30. [Aug. 10,
1886. There is great reason for concluding that the documentary
fragments used above and ascribed to St. Hilary and Liberius, pp. 322,
323, are not genuine. It is safer to confine ourselves to the
following judgment of Bishop Hefele in his ``Councils,'' vol.
ii. pp. 245, 246, ed. 1875:—

``We therefore conclude without
doubt that Liberius, yielding to force and sinking under many years of
confinement and exile, signed the so-called third Sirmian formula,
that is, the collection of older formulas of faith accepted at the
third Sirmian Synod of 358. He did not do this without scruples, for
the Semi-Arian character and origin of these formulas were not unknown
to him; but, as they contained no direct or express rejection of the
orthodox faith, and as it was represented to him, on the other side,
that the Nicene [\emph{homoousios}] formed a cloak for Sabellianism
and Photinism, he allowed himself to be persuaded to accept the third
Sirmian confession. But by so doing he only renounced the letter of
the Nicene faith, not the orthodox faith itself.'']}

The ignorance and folly of this remarkable document are at first
sight incredible; but to an observant mind the common experience of
life brings sufficient proof, that there is nothing too audacious for
party spirit to assert, nothing too gross for monarch or inflamed
populace to receive.

\xsection{Section 4. The Anomœans}

IT remains to relate the circumstances
of the open disunion and schism between the Semi-Arians and the
Anomœans. In order to set this clearly before the reader, a brief
recapitulation must first be made of the history of the heresy, which
has been thrown into the shade in the last Section, by the narrative
of the ecclesiastical events to which it gave occasion.

The Semi-Arian school was the offspring of the ingenious
refinements, under which the Eusebians concealed impieties, which the
temper of the faithful made it inexpedient for them to avow \footnote{[Plato made Semi-Arians, and Aristotle Arians.]}. Its creed preceded the party; that is, those subtleties, which
were too feeble to entangle the clear intellects of the school of
Lucian, produced after a time their due effect upon the natural
subjects of them, viz. men who, with more devotional feeling than the
Arians, had less plain sense, and a like deficiency of humility. A
Platonic fancifulness made them the victims of an Aristotelic
subtlety; and in the philosophising Eusebius and the sophist Asterius,
we recognize the appropriate inventors, though hardly the sincere
disciples, of the new creed. For a time, the distinction between the
Semi-Arians and the Eusebians did not openly appear; the creeds
put forth by the whole party being all, more or less, of a Semi-Arian
cast, down to the Council of Sirmium inclusive (A.D.
351), in which Photinus was condemned. In the meanwhile the Eusebians,
little pleased with the growing dogmatism of members of their own
body, fell upon the expedient of confining their confessions to
Scripture terms; which, when separated from their context, were of
course inadequate to concentrate and ascertain the true doctrine.
Hence the formula of the \emph{Homœon}; which was introduced by
Acacius with the express purpose of deceiving or baffling the
Semi-Arian members of his party. This measure was the more necessary
for Eusebian interests, inasmuch as a new variety of the heresy arose
in the East at the same time, advocated by Aetius and Eunomius; who,
by professing boldly the pure Arian tenet, alarmed Constantius, and
threw him back upon Basil, and the other Semi-Arians. This new
doctrine, called Anomœan, because it maintained that the \emph{usia}
or \emph{substance} of the Son was unlike ([\emph{anomoios}]) the
Divine \emph{usia}, was actually adopted by one portion of the
Eusebians, Valens and his rude Occidentals; whose language and temper,
not admitting the refinements of Grecian genius, led them to rush from
orthodoxy into the most hard and undisguised impiety. And thus the
parties stand at the date now before us (A.D.
356-361); Constantius being alternately swayed by Basil, Acacius, and
Valens, that is, by the Homœüsian, the Homœan, and the Anomœan,—the
Semi-Arian, the Scripturalist, and the Arian pure; by his respect for
Basil and the Semi-Arians, the talent of Acacius, and his personal
attachment to Valens.

1.

Aetius, the founder of the Anomœans, is a remarkable instance of
the struggles and success of a restless and aspiring mind under the
pressure of difficulties. He was a native of Antioch; his father, who
had an office under the governor of the province, dying when he was a
child, he was made the servant or slave of a vine-dresser. He was
first promoted to the trade of a goldsmith or travelling tinker,
according to the conflicting testimony of his friends and enemies.
Falling in with an itinerant practitioner in medicine, he acquired so
much knowledge of the art, as to profess it himself; and, the further
study of his new profession introducing him to the disputations of his
more learned brethren, he manifested such acuteness and boldness in
argument, that he was soon engaged, after the manner of the Sophists,
as a paid advocate for such physicians as wished their own theories
exhibited in the most advantageous form. The schools of Medicine were
at that time infected with Arianism, and thus introduced him to the
science of theology, as well as that of disputation; giving him a bias
towards heresy, which was soon after confirmed by the tuition of
Paulinus, Bishop of Tyre. At Tyre he so boldly conducted the
principles of Arianism to their legitimate results, as to scandalize
the Eusebian successor of Paulinus; who forced him to retire to
Anazarbus, and to resume his former trade of a goldsmith. The energy
of Aetius, however, could not be restrained by the obstacles which
birth, education, and decency threw in his way. He made acquaintance
with a teacher of grammar; and, readily acquiring a smattering of
polite literature, he was soon enabled to criticise his master's
expositions of sacred Scripture before his pupils. A quarrel, as might
be expected, ensued; and Aetius was received into the house of the
Bishop of Anazarbus, who had been one of the Arian prelates at Nicæa.
This man was formerly mentioned as one of the rudest and most daring
among the first assailants of our Lord's divinity \footnote{[Vide supra, p. 239.]}. It is probable, however, that, after signing the \emph{Homoüsion},
he had surrendered himself to the characteristic duplicity and
worldliness of the Eusebian party; for Aetius is said to have
complained, that he was deficient in depth, and, in spite of his
hospitality, looked out for another instructor. Such an one he found
in the person of a priest of Tarsus, who had been from the first a
consistent Arian; and with him he read the Epistles of St. Paul.
Returning to Antioch, he became the pupil of Leontius, in the
prophetical Scriptures; and, after a while, put himself under the
instruction of an Aristotelic sophist of Alexandria. Thus
accomplished, he was ordained deacon by Leontius (A.D.
350), who had been lately raised to the patriarchal See of Antioch.
Thus the rise of the Anomœan sect coincides in point of time with the
death of Constans, an event already noticed in the history of the
Eusebians, as transferring the Empire of the West to Constantius, and,
thereby furthering their division into the Homœan and Homœusian
factions. Scarcely had Aetius been ordained, when the same notorious
irregularities in his carriage, whatever they were, which had more
than once led to his expulsion from the lay communion of the Arians,
caused his deposition from the diaconate, by the very prelate
who had promoted him to it. After this, little is known of him for
several years; excepting a dispute, which he held with the Semi-Arian
Basil, which marks his rising importance. During the interval, he
ingratiated himself with Gallus, the brother of Julian; and was
implicated in his political offences. Escaping, however, the anger of
Constantius, by his comparative insignificance, he retired to
Alexandria, and lived for some time in the train of George of
Cappadocia, who allowed him to officiate as deacon. Such was at this
time the character of the clergy, whom the Arians had introduced into
the Syrian Churches, that this despicable adventurer, whose manners
were as odious, as his life was eccentric, and his creed blasphemous,
had sufficient influence to found a sect, which engaged the attention
of the learned Semi-Arians at Ancyra (A.D. 358),
and has employed the polemical powers of the orthodox Fathers, Basil
and Gregory Nyssen.

Eunomius, his disciple, was the principal disputant in the
controversy. With more learning than Aetius, he was enabled to
complete and fortify the Anomœan system, inheriting from his master
the two peculiarities of character which belong to his school; the
first, a faculty of subtle disputation and hard mathematical
reasoning, the second, a fierce, and in one sense an honest, disdain
of compromise and dissimulation. These had been the two marks of
Arianism at its first rise; and the first associates of Arius, who,
after his submission to Constantine, had kept aloof from the Court
party in disgust, now joyfully welcomed and joined the Anomœans. The
new sect justified their anticipations of its boldness. The same
impatience, with which Aetius had received the ambiguous
explanations of the Eusebian Bishop of Anazarbus, was expressed by
Eunomius for the Acacianism of Eudoxius of Antioch, who in vain
endeavoured to tutor him into a less real and systematic profession of
the Arian tenets. So far did his party carry their vehemence, as even
to re-baptize their Christian converts, as though they had been
heathen; and that, not in the case of Catholics only, but, to the
great offence of the Eusebians, even of those, whom they converted
from the other forms of Arianism \footnote{Epiph. Hær. lxxvi. fin. Bingham, xi. 1. §
10. [Thus, bold as were the original Arians, the Anomœans were bolder
and more consistent. Athanasius challenges the former, if they dare,
to speak out. Basil says ``Aetius was the first to teach openly
that the Father's substance was unlike the Son's.'' Vide Ath. Tr.
vol. ii. pp. 34, 287-292. However, Athanasius interprets Arius's
Thalia to say that the Persons of the Holy Trinity are utterly unlike
([\emph{anomoioi}]) each other in substance and glory without
limit.'' Orat. § 6. De Syn. § 15. Again, Arius held that the
Divine Being was incomprehensible (Athan. de Syn. § 15), but the Anomœans denied it. Socr. iv. 7.]}. Earnestness is always respectable; and, if it be allowable to
speak with a sort of moral catachresis, the Anomœans merited on this
account, as well as ensured, a success, which a false conciliation
must not hope to obtain.

2.

The progress of events rapidly carried them forward upon the scene
of ecclesiastical politics. Valens, who by this time had gained the
lead of the Western Bishops, was seconded in his patronage of them by
the eunuchs of the Court; of whom Eusebius, the Grand Chamberlain, had
unlimited sway over the weak mind of the Emperor. The concessions,
made by Liberius and Hosius to the Eusebian party, furnished an
additional countenance to Arianism, being misrepresented as actual
advances towards the heretical doctrine. The inartificial cast of the
Western theology, which scarcely recognized any middle hypothesis
between that of the Homoüsion and pure Arianism, strengthened the
opinion that those, who had abandoned the one, must in fact have
embraced the other. And, as if this were not enough, it appears that
an Anomœan creed was circulated in the East, with the pretence that
it was the very formula which Hosius and Liberius had subscribed.
Under these circumstances, the Anomœans were soon strong enough to
aid the Eusebians of the East in their contest with the Semi-Arians \footnote{Petav. tom. ii. i. 9, § 6. [Tillemont, t.
6, p. 429.]}. Events in the Churches of Antioch and Jerusalem favoured their
enterprise. It happening that Acacius of Cæsarea and Cyril of
Jerusalem were rivals for the primacy of Palestine, the reputed
connexion of Cyril with the Semi-Arian party had the effect of
throwing Acacius, though the author of the Homœon, on the side of its
Anomœan assailants; accordingly, with the aid of the neighbouring
Bishops, he succeeded in deposing Cyril, and sending him out of the
country. At Antioch, the cautious Leontius, Arian Bishop, dying (A.D.
357), the eunuchs of the Court contrived to place Eudoxius in his see,
a man of restless and intriguing temper, and opposed to the
Semi-Arians. One of his first acts was to hold a Council, at which
Acacius was present, as well as Aetius and Eunomius, the chiefs of the
Anomœans. There the assembled Bishops did not venture beyond the
language of the second creed of Sirmium, which Hosius had
signed, and which kept clear of Anomœan doctrine; but they had no
difficulty in addressing a letter of thanks and congratulations to the
party of the Anomœan Valens, for having at Sirmium brought the
troubles of the West to so satisfactory a termination.

The election, however, of Eudoxius, and this Council which followed
it were not to pass unchallenged by the Semi-Arians. Mention has
already been made of one George \footnote{Vide supr. p. 240.}, a presbyter of Alexandria; who, being among the earliest
supporters of Arius, was degraded by Alexander, but, being received by
the Eusebians into the Church of Antioch, became at length Bishop of
Laodicea. George was justly offended at the promotion of Eudoxius,
without the consent of himself and Mark of Arethusa, the most
considerable Bishops of Syria; and, at this juncture, took part
against the combination of Homœans and Anomœans, at Antioch, who had
just published their assent to the second creed of Sirmium. Falling in
with some clergy whom Eudoxius had excommunicated, he sent letters by
them to Macedonius, Basil of Ancyra, and other leaders of the
Semi-Arians, intreating them to raise a protest against the
proceedings of the Council of Antioch, and so to oblige Eudoxius to
separate himself from Aetius and the Anomœans. This remonstrance
produced its effect; and, under pretence of the dedication of a
Church, a Council was immediately held by the Semi-Arian party at
Ancyra (A.D. 358), in which the Anomœan heresy
was condemned. The Synodal letter, which they published, professed to
be grounded on the Semi-Anian creeds of the Dedication (A.D.
341), of Philippopolis (A.D. 347), and of
Sirmium (A.D. 351), when Photinus was
condemned and deposed. It is a valuable document, even as a defence of
orthodoxy; its error consisting in its obstinate rejection of the
Nicene \emph{Homoüsion}, the sole practical bulwark of the Catholic
faith against the misrepresentations of heresy,—against a sort of
tritheism on the one hand, and a degraded conception of the Son and
Spirit on the other.

The two parties thus at issue, appealed to Constantius at Sirmium.
That weak Prince had lately sanctioned the almost Acacian creed of
Valens, which Hosius had been compelled to subscribe, when the
deputation from Antioch arrived at the Imperial Court; and he readily
gave his assent to the new edition of it which Eudoxius had
promulgated. Scarcely had he done so, when the Semi-Arians made their
appearance from Ancyra, with Basil at their head; and succeeded so
well in representing the dangerous character of the creed passed at
Antioch, that, recalling the messenger who had been sent off to that
city, he forthwith held the Conference, mentioned in the foregoing
Section, in which he imposed a Semi-Arian creed on all parties,
Eudoxius and Valens, the representatives of the Eusebians, being
compelled, as well as the orthodox Liberius, to sign a formulary,
which Basil compiled from the creeds against Paulus of Samosata, and
Photinus (A.D. 264, 351), and the creed of
Lucian, published by the Council of the Dedication (A.D.
341). Yet in spite of the learning, and personal respectability of the
Semi-Arians, which at the moment exerted this strong influence over
the mind of Constantius, the dexterity of the Eusebians in disputation
and intrigue was ultimately successful. Though seventy Bishops
of their party were immediately banished, these were in a few months
reinstated by the capricious Emperor, who from that time inclined
first to the Acacian or Homœan, and then to the open Anomœan or pure
Arian doctrine; and who before his death (A.D.
361) received baptism from the hands of Euzoius, one of the original
associates of Arius, then recently placed in the see of Antioch.—The
history of this change, with the Councils attending it, will bring us
to the close of this chapter.

3.

The Semi-Arians, elated by their success with the Emperor, followed
it up by obtaining his consent for an Ecumenical Council, in which the
faith of the Christian Church should definitely be declared for good.
A meeting of the whole of Christendom had not been attempted, except
in the instance of the Council of Sardica, since the Nicene; and the
Sardican itself had been convoked principally to decide upon the
charges urged against Athanasius, and not to open the doctrinal
question. Indeed it is evident, that none but the heterodox party, now
dominant, could consistently debate an article of belief, which the
united testimony of the Churches of the East and West had once for all
settled at Nicæa. This, then, was the project of the Semi-Arians.
They aimed at a renewal on an Ecumenical scale of the Council of the
Dedication at Antioch in A.D. 341. The Eusebian
party, however, had no intention of tamely submitting to defeat.
Perceiving that it would be more for their own interest that the
prelates of the East and West should not meet in the same place (two
bodies being more manageable than one), they exerted themselves
so strenuously with the assistance of the eunuchs of the palace, that
at last it was determined, that, while the Orientals met at Seleucia
in Isauria, the Occidental Council should be held at Ariminum, in
Italy. Next, a previous Conference was held at Sirmium, in order to
determine on the creed to be presented to the bipartite Council; and
here again the Eusebians gained an advantage, though not at once to
the extent of their wishes. Warned by the late indignation of
Constantius against the Anomœan tenet, they did not attempt to rescue
it from his displeasure; but they struggled for the adoption of the
Acacian \emph{Homœon}, which the Emperor had already both received
and abandoned, and they actually effected the adoption of the ``\emph{like
in all things according to the Scriptures}''—a phrase in
which the Semi-Arians indeed included their ``\emph{like in substance}''
or \emph{Homϟsion}, but which did not necessarily refer to \emph{substance}
or \emph{nature} at all. Under these circumstances the two Councils
met in the autumn of A.D. 359, under the nominal
superintendence of the Semi-Arians; but on the Eusebian side, the
sharp-witted Acacius undertaking to deal with the disputatious Greeks,
the overbearing and cruel Vatens with the plainer Latins.

About 160 Bishops of the Eastern Church assembled at Seleucia \footnote{[Vide Ath. Tr. vol. i. p. 78, notes 8, 9.]}, of whom not above forty were Eusebians. Far the greater number
were professed Semi-Arians; the Egyptian prelates alone, of whom but
twelve or thirteen were present, displaying themselves, as at the
first, the bold and faithful adherents of the \emph{Homoüsion}. It
was soon evident that the forced reconciliation which
Constantius had imposed on the two parties at Sirmium, was of no avail
in their actual deliberations. On each side an alteration of the
proposed formula was demanded. In spite of the sanction given by Basil
and Mark to the ``\emph{like in all things},'' the majority of
their partisans would be contented with nothing short of the definite
``\emph{like in substance},'' or \emph{Homϟsion}, which
left no opening (as they considered) to evasion; and in consequence
proposed to return to Lucian's creed, adopted by the Council of the
Dedication. Acacius, on the other hand, not satisfied with the
advantage he had just gained in the preliminary meeting at Sirmium,
where the mention of the \emph{usia} or \emph{substance} was dropped
(although but lately imposed by Constantius on all parties, in the
formulary which Liberius signed), proposed a creed in which the \emph{Homoüsion}
and \emph{Homœüsion}, were condemned, the \emph{Anomœon}
anathematized, as the source of confusion and schism, and his own \emph{Homœon}
adopted (that is, ``\emph{like},'' without the addition of
``\emph{in all things}''); and when he found himself unable to
accomplish his purpose, not waiting for the formal sentence of
deposition, which the Semi-Arians proceeded to pronounce upon himself
and eight others, he set off to Constantinople, where the Emperor then
was, hoping there, in the absence of Basil and his party, to gain what
had been denied him in the preliminary meeting at Sirmium. It so
happened, however, that his object had been effected even before his
arrival; for, a similar quarrel having resulted from the meeting at
Ariminum, and deputies from the rival parties having thence similarly
been despatched to Constantius, a Conference had already taken place
at a city called Nice or Nicæa, in the neighbourhood of
Hadrianople, and an emendated creed adopted, in which, not only the
safeguard of the ``\emph{in all things}'' was omitted, and the
\emph{usia} condemned, but even the word \emph{Hypostasis} (\emph{Subsistence}
or \emph{Person}) also, on the ground of its being a refinement on
Scripture. So much had been already gained by the influence of Valens,
when the arrival of Acacius at Constantinople gave fresh activity to
the Eusebian party.

Thereupon a Council was summoned in the Imperial city of the
neighbouring Bishops, principally of those of Bithynia, and the
Acacian formula of Ariminum confirmed. Constantius was easily
persuaded to believe of Basil, what had before been asserted of
Athanasius, that he was the impediment to the settlement of the
question, and to the tranquillity of the Church. Various charges of a
civil and ecclesiastical nature were alleged against him and other
Semi-Arians, as formerly against Athanasius, with what degree of truth
it is impossible at this day to determine; and a sentence of
deposition was issued against them. Cyril of Jerusalem, Eleusius of
Cyzicus, Eustathius of Sebaste, and Macedonius of Constantinople, were
in the number of those who suffered with Basil; Macedonius beng
succeeded by Eudoxius, who, thus seated in the first see of the East,
became subsequently the principal stay of Arianism under the Emperor
Valens.

This triumph of the Eusebian party in the East, took place in the
beginning of A.D. 360; by which time the Council
of Ariminum in the West, had been brought to a conclusion. To it we
must now turn our attention.

The Latin Council had commenced its deliberations, before the
Orientals had assembled at Seleucia; yet it did not bring them to a
close till the end of the year. The struggle between the Eusebians and
their opponents had been so much the more stubborn in the West, in
proportion as the latter were more numerous there, and further removed
from Arian doctrine, and Valens on the other hand more unscrupulous,
and armed with fuller powers. Four hundred Bishops were collected at
Ariminum, of whom but eighty were Arians; and the civil officer, to
whom Constantius had committed the superintendence of their
proceedings, had orders not to let them stir out of the city, till
they should agree upon a confession of faith. At the opening of the
Council, Valens, Ursacius, Germinius, Auxentius, Calus, and Demophilus,
the Imperial Commissioners, had presented to the assembly the formula
of the ``\emph{like in all things},'' agreed upon in the
preliminary conference at Sirmium; and demanded, that, putting aside
all strange and mysterious terms of theology, it should be at once
adopted by the assembled Fathers. They had received for answer, that
the Latins determined to adhere to the formulary of Nicæa; and that,
as a first step in their present deliberations, it was necessary that
all present should forthwith anathematize all heresies and
innovations, beginning with that of Arius. The Commissioners had
refused to do so, and had been promptly condemned and deposed, a
deputation of ten being sent from the Council to Constantius, to
acquaint him with the result of its deliberations. The issue of this
mission to the Court, to which Valens opposed one from his own party,
has been already related. Constantius, with a view of wearing
out the Latin Fathers, pretended that the barbarian war required his
immediate attention, and delayed the consideration of the question
till the beginning of October, several months after the opening of the
Council; and then, frightening the Catholic deputation into
compliance, he effected at Nice the adoption of the Homœan creed
(that is, the ``\emph{like}'' without the ``\emph{in all
things}'') and sent it back to Ariminum.

The termination of the Council there assembled was disgraceful to
its members, but more so to the Emperor himself. Distressed by their
long confinement, impatient at their absence from their respective
dioceses, and apprehensive of the approaching winter, they began to
waver. At first, indeed, they refused to communicate with their own
apostate deputies; but these, almost in self-defence, were active and
successful in bringing over others to their new opinions. A threat was
held out by Taurus, the Prætonian Prefect, who superintended the
discussions, that fifteen of the most obstinate should be sent into
banishment; and Valens was importunate in the use of such theological
arguments and explanations, as were likely to effect his object. The
Prefect conjured them with tears to abandon an unfruitful obstinacy,
to reflect on the length of their past confinement, the discomfort of
their situation, the rigours of the winter, and to consider, that
there was but one possible termination of the difficulty, which lay
with themselves, not with him. Valens, on the other hand, affirmed
that the Eastern bishops at Seleucia had abandoned the \emph{usia};
and he demanded of those who still stood their ground, what objection
they could make to the Scriptural creed proposed to them, and
whether, for the sake of a word, they would be the authors of a schism
between Eastern and Western Christendom. He affirmed, that the danger
apprehended by the Catholics was but chimerical; that he and his party
condemned Arius and Arianism, as strongly as themselves, and were only
desirous of avoiding a word, which confessedly is not in Scripture,
and had in past time been productive of much scandal. Then, to put his
sincerity to the proof, he began with a loud voice to anathematize the
maintainers of the Arian blasphemies in succession; and he concluded
by declaring, that he believed the Word to be God, begotten of God
before all time, and not in the number of the creatures, and that
whoever should say that He was a creature as other creatures, was
anathema. The foregoing history of the heresy has sufficiently
explained how the Arians evaded the force of these strong
declarations; but the inexperienced Latins did not detect their
insincerity. Satisfied, and glad to be released, they gave up the \emph{Homoüsion},
and signed the formula of the \emph{Homœon}; and scarcely had they
separated, when Valens, as might be expected, boasted of his victory,
arguing that the faith of Nicæa had been condemned by the very
circumstance of his being allowed to confess, that the Son was
``not a creature as other creatures,'' and so to imply, that,
though not like \emph{other} creatures, still He was created. Thus
ended this celebrated Council; the result of which is well
characterized in the lively statement of Jerome: ``The whole world
groaned in astonishment to find itself Arian.'' \footnote{[``Ingemuit totus orbis, et Arianum se
esse miratus est.'']}

In the proceedings attendant on the Councils of Seleucia and
Ariminum, the Eusebians had skilfully gained two important objects, by
means of unimportant concessions on their part. They had sacrificed
Aetius and his \emph{Anomœon}; and effected in exchange the disgrace
of the Semi-Arians as well as of the Catholics, and the establishment
of the \emph{Homœon}, the truly characteristic symbol of a party,
who, as caring little for the sense of Scripture, found an excuse and
an indulgence of their unconcern, in a pretended maintenance of the
letter. As to the wretched mountebank just mentioned, whose
profaneness was so abominable, as to obtain for him the title of the
``Atheist,'' he was formally condemned in the Council at
Constantinople (A.D. 360) already mentioned, in
which the Semi-Arian Basil, Macedonius, and their associates had been
deposed. During the discussions which attended it, Eleusius, one of
the latter party, laid before the Emperor an Anomœan creed, which he
ascribed to Eudoxius. The latter, when questioned, disowned it; and
named Aetius as its author, who was immediately summoned. Introduced
into the Imperial presence, he was unable to divine, in spite of his
natural acuteness, whether the Emperor was pleased or displeased with
the composition; and, hazarding an acknowledgement of it, he drew down
on himself the full indignation of Constantius, who banished him into
Cilicia, and obliged his patron Eudoxius to anathematize both the
confession in question, and all the positions of the pure Arian
heresy. Such was the fall of Aetius, at the time of the triumph of the
Eusebians; but soon afterwards he was promoted to the episcopate
(under what circumstances is unknown), and was favourably
noticed, as a former friend of Gallus, by the Emperor Julian, who gave
him a territory in the Island of Mitelene.

Eunomius, his disciple, escaped the jealousy of Constantius through
the good offices of Eudoxius, and was advanced to the Bishoprick of
Cyzicus; but, being impatient of dissimulation, he soon fell into
disgrace, and was banished. The death of the Emperor took place at the
end of A.D. 361; his last acts evincing a
further approximation to the unmitigated heresy of Arius. At a Council
held at Antioch in the course of that year, he sanctioned the Anomœan
doctrine in its most revolting form; and shortly before his decease,
received the sacrament of baptism, as has been stated above, from
Euzoius, the personal friend and original associate of Arius himself \footnote{[``At this critical moment Constantius
died, when the cause of truth was only not in the lowest state of
degradation, because a party was in authority and vigour who could
reduce it to a lower still; the Latins committed to an Anti-Catholic
creed, the Pope a renegade, Hosius fallen and dead, Athanasius
wandering in the deserts, Arians in the sees of Christendom, and their
doctrine growing in blasphemy, and their profession of it in boldness,
every day. The Emperor had come to the throne when almost a boy, and
at this time was but forty-four years old. In the ordinary course of
things, he might have reigned till orthodoxy, humanly speaking, was
extinct.'' Ath. Tr. vol. i. p. 121.]}.

\xchapter{Chapter 5. Councils after the Reign of Constantius}

\xsection{Section 1. The Council of Alexandria in the Reign of Julian}

THE accession of Julian was followed by
a general restoration of the banished Bishops; and all eyes throughout
Christendom were at once turned towards Alexandria, as the Church,
which, by its sufferings and its indomitable spirit, had claim to be
the arbiter of doctrine, and the guarantee of peace to the Catholic
world. Athanasius, as the story goes, was, on the death of his
persecutor, suddenly found on his episcopal throne in one of the
Churches of Alexandria \footnote{Cave, Life of Athan. x. 9.}; a
legend, happily expressive of the unwearied activity and almost
ubiquity of that extraordinary man, who, while a price was set on his
head, mingled unperceived in the proceedings at Seleucia and Ariminum
\footnote{[This is doubtful; vide Montfaucon, Athan.,
though Tillemont and Gibbon seem to admit it.]}, and directed the
movements of his fellow-labourers by his writings, when he was
debarred the exercise of his dexterity in debate, and his persuasive
energy in private conversation. He was soon joined by his
fellow-exile, Eusebius of Vercellæ; Lucifer, who had journeyed with
the latter from the Upper Thebaid, on his return to the West, having
gone forward to Antioch on business which will presently be explained.
Meanwhile, no time was lost in holding a Council at Alexandria (A.D.
362) on the general state of the Church.

The object of Julian in recalling the banished Bishops, was the
renewal of those dissensions, by means of toleration, which
Constantius had endeavoured to terminate by force. He knew these
prelates to be of various opinions, Semi-Arians, Macedonians, Anomœans,
as well as orthodox; and, determining to be neuter himself, he waited
with the satisfaction of an Eclectic for the event; being persuaded,
that Christianity could not withstand the shock of parties, not less
discordant, and far more zealous, than the sects of philosophy. It is
even said that he ``invited to his palace the leaders of the
hostile sects, that he might enjoy the agreeable spectacle of their
furious encounters.'' \footnote{Gibbon, ch. xxiii.}
But, in indulging such anticipations of overthrowing Christianity, he
but displayed his own ignorance of the foundation on which it was
built. It could scarcely be conceived that an unbeliever, educated
among heretics, would understand the vigour and indestructibility of
the true Christian spirit; and Julian fell into the error, to which in
all ages men of the world are exposed, of mistaking whatever shows
itself on the surface of the Apostolic Community, its
prominences and irregularities, all that is extravagant, and all that
is transitory, for the real moving principle and life of the system.
It is trying times alone that manifest the saints of God; but they
live notwithstanding, and support the Church in their generation,
though they remain in their obscurity. In the days of Arianism,
indeed, they were in their measure, revealed to the world; still to
such as Julian, they were unavoidably unknown, both in respect to
their numbers and their divine gifts. The thousand of silent
believers, who worshipped in spirit and in truth, were obscured by the
tens and twenties of the various heretical factions, whose clamorous
addresses besieged the Imperial Court; and Athanasius would be
portrayed to Julian's imagination after the picture of his own
preceptor, the time-serving and unscrupulous Eusebius. The event of
his experiment refuted the opinion which led to it. The impartial
toleration of all religious persuasions, malicious as was its intent,
did but contribute to the ascendancy of the right faith; that faith,
which is the only true aliment of the human mind, which can be held as
a principle as well as an opinion, and which influences the heart to
suffer and to labour for its sake.

Of the subjects which engaged the notice of the Alexandrian
Council, two only need here be mentioned; the treatment to be pursued
towards the bishops, who had arianized in the reign of Constantius,
and the settlement of the theological sense of the word \emph{Hypostasis}.
And here, of the former of these.

1.

Instances have already occurred, of the line of conduct pursued by
Athanasius in ecclesiastical matters. Deliberate apostasy and
systematic heresy were the objects of his implacable opposition; but
in his behaviour towards individuals, and in his judgment of the
inconsistent, whether in conduct or creed, he evinces an admirable
tenderness and forbearance. Not only did he reluctantly abandon his
associate, the unfortunate Marcellus, on his sabellianizing, but he
even makes favourable notice of the Semi-Arians, hostile to him both
in word and deed, who rejected the orthodox test, and had confirmed
against him personally at Philippopolis, the verdict of the commission
at the Mareotis. When bishops of his own party, as Liberius of Rome,
were induced to excommunicate him, far from resenting it, he speaks of
them with a temper and candour, which, as displayed in the heat of
controversy, evidences an enlarged prudence, to say nothing of
Christian charity \footnote{Athan. de Syn. 41. Apol. contr. Arian. 89.
Hist. Arian ad Monach. 41, 42.}. It is
this union of opposite excellences, firmness with discrimination and
discretion, which is the characteristic praise of Athanasius: as well
as of several of his predecessors in the See of Alexandria. The
hundred years, preceding his episcopate, had given scope to the
enlightened zeal of Dionysius, and the patient resoluteness of
Alexander. On the other hand, when we look around at the other more
conspicuous champions of orthodoxy of his time, much as we must revere
and bless their memory, yet as regards this maturity and
completeness of character, they are far inferior to Athanasius. The
noble-minded Hilary was intemperate in his language, and assailed
Constantius with an asperity unbecoming a dutiful subject. The fiery
Bishop of Cagliari, exemplary as is his self-devotion, so openly
showed his desire for martyrdom, as to lead the Emperor to exercise
towards him a contemptuous forbearance. Eusebius of Vercellæ
negotiated in the Councils, with a subtlety bordering on Arian
insincerity. From these deficiencies of character Athanasius was
exempt; and on the occasion which has given rise to these remarks, he
had especial need of the combination of gifts, which has made his name
immortal in the Church.

The question of the arianizing bishops was one of much difficulty.
They were in possession of the Churches; and could not be deposed, if
at all, without the risk of a permanent schism. It is evident,
moreover, from the foregoing narrative, how many had been betrayed
into an approval of the Arian opinions, without understanding or
acting upon them. This was particularly the case in the West, where
threats and ill-usage, had been more or less substituted for those
fallacies, which the Latin language scarcely admitted. And even in the
remote Greek Churches, there was much of that devout and unsuspecting
simplicity, which was the easy sport of the supercilious sophistry of
the Eusebians. This was the case with the father of Gregory Nazianzen;
who, being persuaded to receive the Acacian confession of
Constantinople (A.D. 359, 360), on the ground of
its unmixed scripturalness, found himself suddenly deserted by a large
portion of his flock, and was extricated from the charge of
heresy, only by the dexterity of his learned son. Indeed, to many of
the Arianizing bishops, may be applied the remarks, which Hilary makes
upon the laity subjected to Arian teaching; that their own piety
enabled them to interpret expressions religiously, which were
originally invented as evasions of the orthodox doctrine \footnote{``Sanctiores sunt aures plebis,'' he
says, ``quam corda sacerdotum.'' Bull, Defens. epilog. [Vide
infr. Appendix, No. 5.]}.

And even in parts of the East, where a much clearer perception of
the difference between truth and error existed, it must have been an
extreme difficulty to such of the orthodox as lived among Arians, to
determine, in what way best to accomplish duties, which were in
opposition to each other. The same obligation of Christian unity,
which was the apology for the laity who remained, as at Antioch, in
communion with an Arian bishop, would lead to a similar recognition of
his authority by clergy or bishops who were ecclesiastically
subordinate to him. Thus Cyril of Jerusalem, who was in no sense
either Anomœan or Eusebian, received consecration from the hands of
his metropolitan Acacius; and St. Basil, surnamed the Great, the
vigorous champion of orthodoxy against the Emperor Valens, attended
the Council of Constantinople (A.D. 359, 360),
as a deacon, in the train of his namesake Basil, the leader of the
Semi-Arians.

On the other hand, it was scarcely safe to leave the deliberate
heretic in possession of his spiritual power. Many bishops too were
but the creatures of the times, raised up from the lowest of the
people, and deficient in the elementary qualifications of learning and
sobriety. Even those, who had but conceded to the violence of
others, were the objects of a just suspicion; since, frankly as they
now joined the Athanasians, they had already shown as much interest
and reliance in the opposite party.

Swayed by these latter considerations, some of the assembled
prelates advocated the adoption of harsh measures towards the
Arianizers, considering that their deposition was due both to the
injured dignity and to the safety of the Catholic Church. Athanasius,
however, proposed a more temperate policy; and his influence was
sufficient to triumph over the excitement of mind which commonly
accompanies a deliverance from persecution. A decree was passed, that
such bishops as had communicated with the Arians through weakness or
surprise, should be recognized in their respective sees, on their
signing the Nicene formulary; but that those, who had publicly
defended the heresy, should only be admitted to lay-communion. No act
could evince more clearly than this, that it was no party interest,
but the ascendancy of the orthodox doctrine itself, which was the aim
of the Athanasians. They allowed the power of the Church to remain in
the hands of men indifferent to the interests of themselves, on their
return to that faith, which they had denied through fear; and their
ability to force on the Arianizers this condition, evidences what they
might have done, had they chosen to make an appeal against the more
culpable of them to the clergy and laity of their respective churches,
and to create and send out bishops to supply their places. But they
desired peace, as soon as the interests of truth were secured; and
their magnanimous decision was forthwith adopted by Councils
held at Rome, in Spain, Gaul, and Achaia. The state of Asia was less
satisfactory. As to Antioch, its fortunes will immediately engage our
attention. Phrygia and the Proconsulate were in the hands of the
Semi-Arians and Macedonians; Thrace and Bithynia, controlled by the
Imperial Metropolis, were the stronghold of the Eusebian or Court
faction.

2.

The history of the Church of Antioch affords an illustration of the
general disorders of the East at this period, and of the intention of
the sanative measure passed at Alexandria respecting them. Eustathius,
its Bishop, one of the principal Nicene champions, had been an early
victim of Eusebian malice, being deposed on calumnious charges, A.D.
331. A series of Arian prelates succeeded; some of whom, Stephen,
Leontius, and Eudoxius, have been commemorated in the foregoing pages
\footnote{Vide supra, p. 280.}. The Catholics of Antioch
had disagreed among themselves, how to act under these circumstances.
Some, both clergy and laity, refusing the communion of heretical
teachers, had holden together for the time, as a distinct body, till
the cause of truth should regain its natural supremacy; while others
had admitted the usurping succession, which the Imperial will forced
upon the Church. When Athanasius passed through Antioch on his return
from his second exile (A.D. 348), he had
acknowledged the seceders, from a respect for their orthodoxy, and for
the rights of clergy and laity in the election of a bishop. Yet it
cannot be denied, that men of zeal and boldness were found among those
who remained in the heretical communion. Two laymen, Flavian and
Diodorus, protested with spirit against the heterodoxy of the crafty
Leontius, and kept alive an orthodox party in the midst of the
Eusebians.

On the translation of Eudoxius to Constantinople, the year before
the death of Constantius, an accident occurred, which, skilfully
improved, might have healed the incipient schism among the
Trinitarians. Scarcely had Meletius, the new Bishop of the Eusebian
party, taken possession of his see, when he conformed to the Catholic
faith. History describes him as gifted with remarkable sweetness and
benevolence of disposition. Men thus characterized are often deficient
in sensibility, in their practical judgment of heresy; which they
abhor indeed in the abstract, yet countenance in the case of their
friends, from a false charitableness; which leads them, not merely to
hope the best, but to overlook the guilt of opposing the truth, where
the fact is undeniable. Meletius had been brought up in the communion
of the Eusebians; a misfortune, in which nearly all the Oriental
Christians of his day were involved. Being considered as one of their
party, he had been promoted by them to the see of Sebaste, in Armenia;
but, taking offence at the conduct of his flock, he had retired to
Berœa, in Syria. During the residence of the Court at Antioch, A.D.
361, the election of the new prelate of that see came on; and the
choice of both Arians and Arianizing orthodox fell on Meletius.
Acacius was the chief mover in this business. He had lately \footnote{Vide supra, pp. 347, 350.} succeeded in establishing the principle of liberalism at
Constantinople, where a condemnation had been passed on the use of
words not found in Scripture, in confessions of faith; and he
could scarcely have selected a more suitable instrument, as it
appeared, of extending its influence, than a prelate, who united
purity of life and amiableness of temper, to a seeming indifference to
the distinctions between doctrinal truth and error.

On the new Patriarch's arrival at Antioch, he was escorted by the
court bishops, and his own clergy and laity, to the cathedral.
Desirous of solemnizing the occasion, the Emperor himself had
condescended to give the text, on which the assembled prelates were to
comment. It was the celebrated passage from the Proverbs, in which
Origen has piously detected, and the Arians perversely stifled, the
great article of our faith; ``the Lord hath created [possessed] Me
in the beginning of His ways, before His works of old.'' George of
Laodicea, who, on the departure of Eudoxius from Antioch, had left the
Semi-Arians and rejoined the Eusebians, opened the discussion with a
dogmatic explanation of the words. Acacius followed with that
ambiguity of language, which was the characteristic of his school. At
length the new Patriarch arose, and to the surprise of the assembly,
with a subdued manner, and in measured words, avoiding indeed the
Nicene \emph{Homoüsion}, but accurately fixing the meaning of his
expressions, confessed the true Catholic tenet, so long exiled from
the throne and altars of Antioch. A scene followed, such as might be
expected from the excitable temper of the Orientals. The congregation
received his discourse with shouts of joy; while the Arian archdeacon
of the church running up, placed his hand before his mouth to prevent
his speaking; on which Meletius thrust out his hand in sight of the
people, and raising first three fingers, and then one,
symbolized the great truth which he was unable to utter \footnote{Soz. iv. 28.}. The consequences of this bold confession might be expected.
Meletius was banished, and a fresh Bishop appointed, Euzoius, the
friend of Arius. But an important advantage resulted to the orthodox
cause by this occurrence; Catholics and heretics were no longer united
in one communion, the latter being thrown into the position of
schismatics, who had rejected their own bishop. Such was the state of
things, when the death of Constantius occasioned the return of
Meletius, and the convocation of the Council of Alexandria, in which
his case was considered.

The course to be pursued in this matter by the general Church was
evident. There were now in Antioch, besides the heretical party, two
communions professing orthodoxy, of which what may be called the
Protestant body was without a head, Eustathius having died some years
before. It was the obvious duty of the Council, to recommend the
Eustathians to recognize Meletius, and to join in his communion,
whatever original intrusion there might be in the episcopal succession
from which he received his Orders, and whatever might have been his
own previous errors of doctrine. The general principle of restoration,
which they had made the rule of their conduct towards the Arianizers,
led them to this. Accordingly, a commission was appointed to proceed
to Antioch, and to exert their endeavours to bring the dissension to a
happy termination.

Their charitable intentions, however, had been already frustrated
by the unfortunate interference of Lucifer. This Latin Bishop,
strenuous in contending for the faith, had little of the knowledge of
human nature, or of the dexterity in negotiation, necessary for the
management of so delicate a point as that which he had taken upon
himself to settle. He had gone straight to Antioch, when Eusebius of
Vercellæ proceeded to Alexandria; and, on the Alexandrian commission
arriving at the former city, the mischief was done, and the mediation
ineffectual. Indulging, instead of overcoming, the natural reluctance
of the Eustathians to submit to Meletius, Lucifer had been induced,
with the assistance of two others, to consecrate a separate head for
their communion, and by so doing re-animate a dissension, which had
run its course and was dying of itself. The result of this
indiscretion was the rise of an additional, instead of the termination
of the existing schism. Eusebius, who was at the head of the
commission, retired from Antioch in disgust. Lucifer, offended at
becoming the object of censure, separated first from Eusebius, and at
length from all who acknowledged the conforming Arianizers. He founded
a sect, which was called after his name, and lasted about fifty years.

As to the schism at Antioch, it was not terminated till the time of
Chrysostom about the end of the century. Athanasius and the Egyptian
Churches continued in communion with the Eustathians. Much as they had
desired and exerted themselves for a reconciliation between the
parties, they could not but recognize, while it existed, that body
which had all along suffered and laboured with themselves. And
certainly the intercourse, which Meletius held with the unprincipled
Acacius, in the Antiochene Council the following year, and his
refusal to communicate with Athanasius, were not adapted to make them
repent their determination \footnote{Vit. S. Basil, p. cix, ed. Benedict. [Basil
at length succeeded in reconciling Meletius to Athanasius. Vitt.
Benedictt. S. Athanasii, p. lxxxvii, and S. Basilii, p. cix.]}.
The Occidentals and the Churches of Cyprus followed their example. The
Eastern Christians, on the contrary, having for the most part
themselves arianized, took part with the Meletians. At length St.
Chrysostom successfully exerted his influence with the Egyptian and
Western Catholics in behalf of Flavian, the successor of Meletius; a
prelate, it must be admitted, not blameless in the ecclesiastical
quarrel, though he had acted a bold part with Diodorus, afterwards
Bishop of Tarsus, in resisting the insidious attempts of Leontius to
secularize the Church.

3.

The Council of Alexandria was also concerned in determining a
doctrinal question; and here too it exercised a virtual mediation
between the rival parties in the Antiochene Church.

The word \emph{Person} which we venture to use in speaking of those
three distinct and real modes in which it has pleased Almighty God to
reveal to us His being, is in its philosophical sense too wide for our
meaning. Its essential signification, as applied to ourselves, is that
of an individual intelligent agent, answering to the Greek \emph{hypostasis},
or \emph{reality}. On the other hand, if we restrict it to its
etymological sense of \emph{persona} or \emph{prosopon}, that is \emph{character},
it evidently means less than the Scripture doctrine, which we wish to
define by means of it, as denoting merely certain outward
manifestations of the Supreme Being relatively to ourselves, which are
of an accidental and variable nature. The statements of Revelation
then lie between these antagonistic senses in which the doctrine of
the Holy Trinity may be erroneously conceived, between Tritheism, and
what is popularly called Unitarianism.

In the choice of difficulties, then, between words which say too
much and too little, the Latins, looking at the popular and practical
side of the doctrine, selected the term which properly belonged to the
external and defective notion of the Son and Spirit, and called Them
Personæ, or Characters; with no intention, however, of infringing on
the doctrine of their completeness and reality, as distinct from the
Father, but aiming at the whole truth, as nearly as their language
would permit. The Greeks, on the other hand, with their instinctive
anxiety for philosophical accuracy of expression, secured the notion
of Their existence in Themselves, by calling them Hypostases or
Realities; for which they considered, with some reason, that they had
the sanction of the Apostle in his Epistle to the Hebrews. Moreover,
they were led to insist upon this internal view of the doctrine, by
the prevalence of Sabellianism in the East in the third century; a
heresy, which professed to resolve the distinction of the Three
Persons, into a mere distinction of character. Hence the prominence
given to the Three Hypostases or Realities, in the creeds of the
Semi-Arians (for instance, Lucian's and Basil's, A.D.
341-358), who were the especial antagonists of Sabellius, Marcellus,
Photinus, and kindred heretics. It was this praiseworthy jealousy of
Sabellianism, which led the Greeks to lay stress upon the
doctrine of the \emph{Hypostatic Word} \footnote{[[\emph{logos enupostatos}]. Vide supr. p.
171.]} (the Word in real existence), lest the bare use of the terms,
Word, Voice, Power, Wisdom, and Radiance, in designating our Lord,
should lead to a forgetfulness of His Personality. At the same time,
the word \emph{usia} (\emph{substance}) was adopted by them, to
express the simple individuality of the Divine Nature, to which the
Greeks, as scrupulously as the Latins, referred the separate
Personalities of the Son and Spirit.

Thus the two great divisions of Christendom rested satisfied each
with its own theology, agreeing in doctrine, though differing in the
expression of it. But, when the course of the detestable controversy,
which Arius had raised, introduced the Latins to the phraseology of
the Greeks, accustomed to the word Persona, they were startled at the
doctrine of the three Hypostases; a term which they could not
translate except by the word \emph{substance}, and therefore
considered synonymous with the Greek \emph{usia}, and which, in matter
of fact, had led to Arianism on the one hand, and Tritheism on the
other. And the Orientals, on their part, were suspicious of the Latin
maintenance of the One Hypostasis, and Three Personæ; as if such a
formula tended to Sabellianism \footnote{[For the meaning of \emph{Usia} and \emph{Hypostasis},
Vide Appendix, No. 4.]}.

This is but a general account of the difference between the Eastern
and Western theology; for it is difficult to ascertain, when the
language of the Greeks first became fixed and consistent. Some eminent
critics have considered, that \emph{usia} was not discriminated from \emph{hypostasis},
till the Council which has given rise to these remarks. Others
maintain, that the distinction between them is recognized in the
``substance \emph{or} hypostasis'' \footnote{[\emph{ex ousias \={e} hypostase\={o}s}].} of the Nicene Anathema; and these certainly have the authority
of St. Basil on their side \footnote{Vid. Petav. Theol. Dogm. tom. ii. lib. iv.
Bull, Defens. Fid. Nic. ii. 9, § 11.}.
Without attempting an opinion on a point, obscure in itself, and not
of chief importance in the controversy, the existing difference
between the Greeks and Latins, at the times of the Alexandrian
Council, shall be here stated.

At this date, the formula of the Three Hypostases seems, as a
matter of fact, to have been more or less a characteristic of the
Arians. At the same time, it was held by the orthodox of Asia, who had
communicated with them; that is, interpreted by them, of course, in
the orthodox sense which it now bears. This will account for St. Basil's
explanation of the Nicene Anathema; it being natural in an Asiatic
Christian, who seems (unavoidably) to have arianized \footnote{i.e. Semi-Arianized.} for the first thirty years of his life, to imagine (whether
rightly or not) that he perceived in it the distinction between \emph{Usia}
and \emph{Hypostasis}, which he himself had been accustomed to
recognize. Again, in the schism at Antioch, which has been above
narrated, the party of Meletius, which had so long arianized,
maintained the Three Hypostases, in opposition to the Eustathians,
who, as a body, agreed with the Latins, and had in consequence been
accused by the Arians of Sabellianism. Moreover, this connexion of the
Oriental orthodox with the Semi-Arians, partly accounts for some
apparent tritheisms of the former; a heresy into which the latter
certainly did fall \footnote{Petav. i. fin. iv. 13, § 3. The
illustration of three men, as being under the same nature (which is
the ground of the accusation which some writers have brought against
Gregory Nyssen and others, vid. Cudw. iv. 36. p. 597, 601, \&c.
Petav. iv. 7. and 10. Gibbon, ch. xxi.), was but an illustration of a
particular point in the doctrine, and directed against the [\emph{heterousiot\={e}s}]
of the Arians. It is no evidence of tritheism. Vid. Petav. tom. i. iv.
13, § 6-16; and tom. i. ii. 4.}.

Athanasius, on the other hand, without caring to be uniform in his
use of terms, about which the orthodox differed, favours the Latin
usage, speaking of the Supreme Being as one Hypostasis, i.e.
substance. And in this he differed from the previous writers of his
own Church; who, not having experience of the Latin theology, nor of
the perversions of Arianism, adopt, not only the word \emph{hypostasis}
but (what is stronger) the words ``\emph{nature}'' and ``\emph{substance},''
to denote the separate Personalities of the Son and Spirit.

As to the Latins, it is said that, when Hosius came to Alexandria
before the Nicene Council, he was desirous that some explanation
should be made about the \emph{Hypostasis}; though nothing was settled
in consequence. But, soon after the Council of Sardica, an addition
was made to its confession, which in Theodoret runs as follows:
``Whereas the heretics maintain that the \emph{Hypostases} of
Father, Son, and Holy Ghost, are distinct and separate, we declare
that according to the Catholic faith there is but one \emph{Hypostasis}
(which they call \emph{Usia}) of the Three; and the \emph{Hypostasis}
of the Son is the same as the Father's.'' \footnote{Theod. Hist. ii. 8.}

Such was the state of the controversy, if it may so be
called, at the time of the Alexandrian Council; the Church of Antioch
being, as it were, the stage, upon which the two parties in dispute
were represented, the Meletians siding with the orthodox of the East,
and the Eustathians with those of the West. The Council, however,
instead of taking part with either, determined, in accordance with the
writings of Athanasius himself, that, since the question merely
related to the usage of words, it was expedient to allow Christians to
understand the ``\emph{hypostasis}'' in one or other sense
indifferently. The document which conveys its decision, informs us of
the grounds of it. ``If any propose to make additions to the Creed
of Nicæa, (says the Synodal letter,) stop such persons and rather
persuade them to pursue peace; for we ascribe such conduct to nothing
short of a love of controversy. Offence having been given by a
declaration on the part of certain persons, that there are Three \emph{Hypostases},
and it having been urged that this language is not scriptural, and for
that reason suspicious, we desired that the inquiry might not be
pushed beyond the Nicene Confession. At the same time, because of this
spirit of controversy, we questioned them, whether they spoke, as the
Arians, of Hypostases foreign and dissimilar to each other, and
diverse in substance, each independent and separate in itself, as in
the case of individual creatures, or the offspring of man, or, as
different substances, gold, silver, or brass; or, again, as other
heretics hold, of Three Origins, and Three Gods. In answer, they
solemnly assured us, that they neither said nor had imagined any such
thing. On our inquiring, 'In what sense then do you say this, or why
do you use such expressions at all?' they answered, 'Because we
believe in the Holy Trinity, not as a Trinity in name only, but in
truth and reality \footnote{[\emph{huphest\={o}san}].}. We
acknowledge the Father truly and in real subsistence, and the Son
truly in substance, and subsistent, and the Holy Ghost subsisting and
existing.' \footnote{[\emph{Huion al\={e}th\={o}s enousion
onta kai huphest\={o}ta, kai Pneuma Hagion huphestos kai hyparchon}].} They said
too, that they had not spoken of Three Gods, or Three Origins, nor
would tolerate that statement or notion; but acknowledged a Holy
Trinity indeed, but only One Godhead, and One Origin, and the Son
consubstantial with the Father, as the Council declared, and the Holy
Spirit, not a creature, nor foreign, but proper to and indivisible
from, the substance of the Son and the Father.

``Satisfied with this explanation of the expressions in
question, and the reasons for their use, we next examined the other
party, who were accused by the above-mentioned as holding but One
Hypostasis, whether their teaching coincided with that of the
Sabellians, in destroying the substance of the Son and the subsistence
of the Holy Spirit. They were as earnest as the others could be, in
denying both the statement and thought of such a doctrine; 'but we use
\emph{Hypostasis'} (\emph{subsistence}), they said, 'considering it
means the same as \emph{Usia} (\emph{substance}), and we hold that
there is but one, because the Son is from the \emph{Usia} (\emph{substance})
of the Father, and because of the identity of Their nature; for we
believe, as in One Godhead, so in One Divine Nature, and not that the
Father's is one, and that the Son's is foreign, and the Holy Ghost's
also.' It appeared then, that both those, who were accused of
holding three \emph{Hypostases}, agreed with the other party, and
those, who spoke of one Substance, professed the doctrine of the
former in the sense of their interpretation; by both was Arius
anathematized as an enemy of Christ, Sabellius and Paulus of Samosata
as impious, Valentinus and Basilides as strangers to the truth,
Manichæus, as an originator of evil doctrines. And, after these
explanations, all, by God's grace, unanimously agree, that such
expressions were not so desirable or accurate as the Nicene Creed, the
words of which they promised for the future to acquiesce in and to
use.'' \footnote{Athan. Tom. ad Antioch, 5 and 6.}

Plain as was this statement, and natural as the decision resulting
from it, yet it could scarcely be expected to find acceptance in a
city, where recent events had increased dissensions of long standing.
In providing the injured and zealous Eustathians with an
ecclesiastical head, Lucifer had, under existing circumstances,
administered a stimulant to the throbbings and festerings of the baser
passions of human nature—passions, which it requires the strong
exertion of Christian magnanimity and charity to overcome. The
Meletians, on the other hand, recognized as they were by the Oriental
Church as a legitimate branch of itself, were in the position of an
establishment, and so exposed to the temptation of disdaining those
whom the surrounding Churches considered as schismatics. How far each
party was in fault, we are not able to determine; but blame lay
somewhere, for the controversy about the \emph{Hypostasis}, verbal as
it was, became the watchword of the quarrel between the two parties,
and only ended, when the Eustathians were finally absorbed by the
larger and more powerful body.

\xsection{Section 2. The Ecumenical Council of Constantinople in the reign of Theodosius}

THE second Ecumenical Council was held
at Constantinople, A.D. 381-383. It is
celebrated in the history of theology for its condemnation of the
Macedonians, who, separating the Holy Spirit from the unity of the
Father and Son, implied or inferred that He was a creature. A brief
account of it is here added in its ecclesiastical aspect; the doctrine
itself, to which it formally bore witness, having been incidentally
discussed in the second Chapter of this Volume.

Eight years before the date of this Council, Athanasius had been
taken to his rest. After a life of contest, prolonged, in spite of the
hardships he encountered, beyond the age of seventy years, he fell
asleep in peaceable possession of the Churches, for which he had
suffered. The Council of Alexandria was scarcely concluded, when he
was denounced by Julian, and saved his life by flight or concealment.
Returning on Jovian's accession, he was, for a fifth and last time,
forced to retreat before the ministers of his Arian successor Valens;
and for four months lay hid in his father's sepulchre. On a
representation being made to the new Emperor, even with the consent of
the Arians themselves, he was finally restored; and so it
happened, through the good Providence of God, that the fury of
persecution, heavily as it threatened in his last years, yet was
suspended till his death, when it at once burst forth upon the Church
with renewed vigour. Thus he was permitted to muse over his past
trials and his prospects for the future; to collect his mind to meet
his God, gathering himself up with Jacob on his bed of age, and
yielding up the ghost peaceably among his children. Yet, amid the
decay of nature, and the visions of coming dissolution, the attention
of Athanasius was in no wise turned from the active duties of his
station. The vigour of his obedience to those duties remained
unabated; one of his last acts being the excommunication of one of the
Dukes of Lybia, for irregularity of life.

At length, when the great Confessor was removed, the Church
sustained a loss, from which it never recovered. His resolute
resistance of heresy had been but one portion of his services; a more
excellent praise is due to him, for his charitable skill in binding
together his brethren in unity. The Church of Alexandria was the
natural mediator between the East and West; and Athanasius had well
improved the advantages thus committed to him. His judicious
interposition in the troubles at Antioch has lately been described;
and the dissensions between his own Church and Constantinople, which
ensued upon his death, may be taken to show how much the combination
of the Catholics depended on his silent authority. Theological
subtleties were for ever starting into existence among the Greek
Christians; and the Arian controversy had corrupted their spirit,
where it had failed to impair their orthodoxy. Disputation was
the rule of belief, and ambition of conduct, in the Eusebian school;
and these evil introductions out-lived its day. Patronized by the
secular power, the great Churches of Christendom conceived a jealousy
of each other, and gradually fortified themselves in their own
resources. As Athanasius drew towards his end, the task of mediation
became more difficult. In spite of his desire to keep aloof from
party, circumstances threw him against his will into one of the two
divisions, which were beginning to discover themselves in the
Christian world. Even before his time, traces appear of a rivalry
between the Asiatic and Egyptian Churches. The events of his own day,
developing their differences of character, at the same time connected
the Egyptians with the Latins. The mistakes of his own friends obliged
him to side with a seeming faction in the body of the Antiochene
Church; and, in the schism which followed, he found himself in
opposition to the Catholic communities of Asia Minor and the East.
Still, though the course of events tended to ultimate disruptions in
the Catholic Church, his personal influence remained unimpaired to the
last, and enabled him to interpose with good effect in the affairs of
the East. This is well illustrated by a letter addressed to him
shortly before his death, by St. Basil, who belonged to the contrary
party, and had then recently been elevated to the exarchate of
Cæsarea. It shall be here inserted, and may serve as a sort of
valediction in parting with one, who, after the Apostles, has been a
principal instrument, by which the sacred truths of Christianity have
been conveyed and secured to the world.

``To Athanasius, Bishop of Alexandria. The more the sicknesses
of the Church increase, so much the more earnestly do we all turn
towards thy Perfection, persuaded that for thee to lead us is our sole
remaining comfort in our difficulties. By the power of thy prayers, by
the wisdom of thy counsels, thou art able to carry us through this
fearful storm; as all are sure, who have heard or made trial of that
perfection ever so little. Wherefore cease not both to pray for our
souls, and to stir us up by thy letters; didst thou know the profit of
these to us, thou wouldst never let pass an opportunity of writing to
us. For me, were it vouchsafed to me, by the co-operation of thy
prayers, once to see thee, and to profit by the gift lodged in thee,
and to add to the history of my life a meeting with so great and
apostolical a soul, surely I should consider myself to have received
from the loving mercy of God a compensation for all the ills, with
which my life has ever been afflicted.'' \footnote{Basil, Ep. 80.}

1.

The trials of the Church, spoken of by Basil in this letter, were
the beginnings of the persecution directed against it by the Emperor
Valens. This prince, who succeeded Jovian in the East, had been
baptized by Eudoxius; who, from the time he became possessed of the
See of Constantinople, was the chief, and soon became the sole, though
a powerful, support of the Eusebian faction. He is said to have bound
Valens by oath, at the time of his baptism, that he would establish
Arianism as the state religion of the East; and thus to have prolonged
its ascendancy for an additional sixteen years after the death
of Constantius (A.D. 361-378). At the beginning
of this period, the heretical party had been weakened by the secession
of the Semi-Arians, who had not merely left it, but had joined the
Catholics. This part of the history affords a striking illustration,
not only of the gradual influence of truth over error, but of the
remarkable manner in which Divine Providence makes use of error itself
as a preparation for truth; that is, employing the lighter forms of it
in sweeping away those of a more offensive nature. Thus Semi-Arianism
became the bulwark and forerunner of the orthodoxy which it opposed.
From A.D. 357, the date of the second and
virtually Homœan formulary of Sirmium \footnote{[Vide supra, pp. 322, 323.]}, it had protested against the impiety of the genuine Arians. In
the successive Councils of Ancyra and Seleucia, in the two following
years, it had condemned and deposed them; and had established the
scarcely objectionable creed of Lucian. On its own subsequent disgrace
at Court, it had concentrated itself on the Asiatic side of the
Hellespont; while the high character of its leading bishops for
gravity and strictness of life, and its influence over the monastic
institutions, gave it a formidable popularity among the lower classes
on the opposite coast of Thrace.

Six years after the Council of Seleucia (A.D.
365), in the reign of Valens, the Semi-Arians held a Council at
Lampsacus, in which they condemned the Homoean formulary of Ariminum,
confirmed the creed of the Dedication (A.D.
341), and, after citing the Eudoxians to answer the accusations
brought against them, proceeded to ratify that deposition of them,
which had already been pronounced at Seleucia. At this time they
seem to have entertained hopes of gaining the Emperor; but, on finding
the influence of Eudoxius paramount at Court, their horror or jealousy
of his party led them to a bolder step. They resolved on putting
themselves under the protection of Valentinian, the orthodox Emperor
of the West; and, finding it necessary for this purpose to stand well
with the Latin Church, they at length overcame their repugnance to the
\emph{Homoüsion}, and subscribed a formula, of which (at least till
the Council of Constantinople, A.D. 360) they
had been among the most eager and obstinate opposers. Fifty-nine
Semi-Arian Bishops gave in their assent to orthodoxy on this memorable
occasion, which took place A.D. 366. Their
deputies were received into communion by Liberius, who had recovered
himself at Ariminum, and who wrote letters in favour of these new
converts to the Churches of the East. On their return, they presented
themselves before an orthodox Council then sitting at Tyana, exhibited
the commendatory letters which they had received from Italy, Gaul,
Africa, and Sicily, as well as Rome, and were joyfully acknowledged by
the assembled Fathers as members of the Catholic body. A final Council
was appointed at Tarsus; whither it was hoped all the Churches of the
East would send representatives, in order to complete the
reconciliation between the two parties. But enough had been done, as
it would seem, in the external course of events, to unite the
scattered portions of the Church; and, when that end was on the point
of accomplishment, the usual law of Divine Providence intervened, and
left the sequel of the union as a task and a trial for
Christians individually. The project of the Council failed;
thirty-four Semi-Arian bishops suddenly opposed themselves to the
purpose of their brethren, and protested against the \emph{Homoüsion}.
The Emperor, on the other hand, recently baptized by Eudoxius,
interfered; forbade the proposed Council, and proceeded to issue an
edict, in which all bishops were deposed from their Sees who had been
banished under Constantius, and restored by Julian. It was at this
time, that the fifth exile of Athanasius took place, which was lately
mentioned. A more cruel persecution followed in A.D.
371, and lasted for several years. The death of Valens, A.D.
378, was followed by the final downfall of Arianism in the Eastern
Church.

As to Semi-Arianism, it disappears from ecclesiastical history at
the date of the proposed Council of Tarsus (A.D.
367); from which time the portion of the party, which remained
non-conformist, is more properly designated Macedonian, or
Pneumatomachist, from the chief article of their heresy.

2.

During the reign of Valens, much had been done in furtherance of
evangelical truth, in the still remaining territory of Arianism, by
the proceedings of the Semi-Arians; but at the same period symptoms of
returning orthodoxy, even in its purest form, had appeared in
Constantinople itself. On the death of Eudoxius (A.D.
370), the Catholics elected an orthodox successor, by name Evagrius.
He was instantly banished by the Emperor's command; and the population
of Constantinople seconded the act of Valens, by the most unprovoked excesses towards the Catholics. Eighty of their clergy, who
were in consequence deputed to lay their grievances before the
Emperor, lost their lives, under circumstances of extreme treachery
and barbarity. Faith, which was able to stand its ground in such a
season of persecution, was naturally prompted to more strenuous acts,
when prosperous times succeeded. On the death of Valens, the Catholics
of Constantinople looked beyond their own community for assistance, in
combating the dominant heresy. Evagrius, whom they had elected to the
See, seems to have died in exile; and they invited to his place the
celebrated Gregory Nazianzen, a man of diversified accomplishments,
distinguished for his eloquence, and still more for his orthodoxy, his
integrity, and the innocence, amiableness, and refinement of his
character.

Gregory was a native of Cappadocia, and an intimate friend of the
great Basil, with whom he had studied at Athens. On Basil's elevation
to the exarchate of Cæsarea, Gregory had been placed by him in the
bishoprick of Sasime; but, the appointment being contested by Anthimus,
who claimed the primacy of the lower Cappadocia, he retired to
Nazianzus, his father's diocese, where he took on himself those
duties, to which the elder Gregory had become unequal. After the death
of the latter, he remained for several years without pastoral
employment, till the call of the Catholics brought him to
Constantinople. His election was approved by Meletius, patriarch of
Antioch and by Peter, the successor of Athanasius, who by letter
recognized his accession to the metropolitan see.

On his first arrival there, he had no more suitable place of
worship than his own lodgings, where he preached the Catholic doctrine
to the dwindled communion over which he presided. But the result which
Constantius had anticipated, when he denied to Athanasius a Church in
Antioch, soon showed itself at Constantinople. His congregation
increased; the house, in which they assembled, was converted into a
church by the pious liberality of its owner, with the name of
Anastasia, in hope of that resurrection which now awaited the
long-buried truths of the Gospel. The contempt, with which the Arians
had first regarded him, was succeeded by a persecution on the part of
the populace. An attempt was made to stone him; his church was
attacked, and he himself brought before a magistrate, under pretence
of having caused the riot. Violence so unjust did but increase the
influence, which a disdainful toleration had allowed him to establish;
and the accession of the orthodox Theodosius secured it.

On his arrival at Constantinople, the new Emperor resolved on
executing in his capital the determination, which he had already
prescribed by edict to the Eastern Empire. The Arian Bishops were
required to subscribe the Nicene formulary, or to quit their sees.
Demophilus, the Eusebian successor of Eudoxius, who has already been
introduced to our notice as an accomplice in the seduction of Liberius,
was first presented with this alternative; and, with an honesty of
which his party affords few instances, he refused at once to assent to
opinions, which he had all through his life been opposing, and retired
from the city. Many bishops, however, of the Arian party conformed;
and the Church was unhappily inundated by the very evil, which
in the reign of Constantine the Athanasians had strenuously and
successfully withstood.

The unfortunate policy, which led to this measure, might seem at
first sight to be sanctioned by the decree of the Alexandrian Council,
which made subscription the test of orthodoxy; but, on a closer
inspection, the cases will be found to be altogether dissimilar. When
Athanasius acted upon that principle, in the reign of Julian, there
was no secular object to be gained by conformity; or rather, the
malevolence of the Emperor was peculiarly directed against those,
whether orthodox or Semi-Arians, who evinced any earnestness about
Christian truth. Even then, the recognition was not extended to those
who had taken an active part on the side of heresy. On the other hand,
the example of Athanasius himself, and of Alexander of Constantinople,
in the reign of Constantine, sufficiently marked their judgment in the
matter; both of them having resisted the attempt of the Court to force
Arius upon the Church, even though he professed his assent to the \emph{Homoüsion}.

Whether or not it was in Gregory's power to hinder the recognition
of the Arianizers, or whether his firmness was not equal to his
humility and zeal, the consequences of the measure are visible in the
conduct of the General Council, which followed it. He himself may be
considered as the victim of it; and he has left us in poetry and in
oratory his testimony to the deterioration of religious principle,
which the chronic vicissitudes of controversy had brought about in the
Eastern Church.

The following passage, from one of his orations, illustrates
both the state of the times, and his own beautiful character, though
unequal to struggle against them. ``Who is there,'' he says,
``but will find, on measuring himself by St. Paul's rules for the
conduct of Bishops and Priests,—that they should be sober, chaste,
not fond of wine, not strikers, apt to teach, unblamable in all
things, unassailable by the wicked,—that he falls far short of its
perfection? ... I am alarmed to think of our Lord's censure of the
Pharisees, and his reproof of the Scribes; disgraceful indeed would it
be, should we, who are bid be so far above them in virtue, in order to
enter the kingdom of heaven, appear even worse than they ... These
thoughts haunt me night and day; they consume my bones, and feed on my
flesh; they keep me from boldness, or from walking with erect
countenance. They so humble me and cramp my mind, and place a chain on
my tongue, that I cannot think of a Ruler's office, nor of correcting
and guiding others, which is a talent above me; but only, how I myself
may flee from the wrath to come, and scrape myself some little from
the poison of my sin. First, I must be cleansed, and then cleanse
others; learn wisdom, and then impart it; draw near to God, and then
bring others to Him; be sanctified, and then sanctify. 'When will you
ever get to the end of this?' say the all-hasty and unsafe, who are
quick to build up and to pull down. 'When will you place your light on
a candlestick? Where is your talent?' So say friends of mine, who have
more zeal for me than religious seriousness. Ah, my brave men, why ask
my season for acting, and my plan? Surely the last day of payment is
soon enough, old age in its extreme term. Grey hairs have
prudence, and youth is untaught. Best be slow and sure, not quick and
thoughtless; a kingdom for a day, not a tyranny for a life; a little
gold, not a weight of lead. It was the shallow earth shot forth the
early blade. Truly there is cause of fear, lest I be bound hand and
foot, and cast without the marriage-chamber, as an audacious intruder
without fitting garment among the assembled guests. And yet I was
called thither from my youth (to confess a matter which few know), and
on God was I thrown from the womb; made over to Him by my mother's
promise, confirmed in His service by dangers afterwards. Yea, and my
own wish grew up beside her purpose, and my reason ran along with it;
and all I had to give, wealth, name, health, literature, I brought and
offered them to Him, who called and saved me; my sole enjoyment of
them being to despise them, and to have something which I could resign
for Christ. To undertake the direction and government of souls is
above me, who have not yet well learnt to be guided, nor to be
sanctified as far as is fitting. Much more is this so in a time like
the present, when it is a great thing to flee away to some place of
shelter, while others are whirled to and fro, and so to escape the
storm and darkness of the evil one; for this is a time when the
members of the Christian body war with each other, and whatever there
was left of love is come to nought. Moabites and Ammonites, who were
forbidden even to enter the Church of Christ, now tread our holiest
places. We have opened to all, not gates of righteousness, but of
mutual reviling and injury. We think those the best of men, not who
keep from every idle word through fear of God, but such as have openly
or covertly slandered their neighbour most. And we mark the sins
of others, not to lament, but to blame them; not to cure, but to
second the blow; and to make the wounds of others an excuse for our
own. Men are judged good and bad, not by their course of life, but by
their enmities and friendships. We praise today, we call names
tomorrow. All things are readily pardoned to impiety. So magnanimously
are we forgiving in wicked ways!'' \footnote{Greg. Orat. i. 119-137. [ii. 69-73, 77-80,
abridged.]}

The first disturbance in the reviving Church of Constantinople had
arisen from the ambition of Maximus, a Cynic philosopher, who aimed at
supplanting Gregory in his see. He was a friend and countryman of
Peter, the new Patriarch of Alexandria; and had suffered banishment in
the Oasis, on the persecution which followed the death of Athanasius.
His reputation was considerable among learned men of the day, as is
shown by the letters addressed to him by Basil. Gregory fell in with
him at Constantinople; and pleased at the apparent strictness and
manliness of his conduct, he received him into his house, baptized
him, and at length admitted him into inferior orders. The return made
by Maximus to his benefactor, was to conduct an intrigue with one of
his principal Presbyters; to gain over Peter of Alexandria, who had
already recognized Gregory; to obtain from him the presence of three
of his bishops; and, entering the metropolitan church during the
night, to instal himself, with their aid, in the episcopal throne. A
tumult ensued, and he was obliged to leave the city; but, far from
being daunted at the immediate failure of his plot, he laid his case
before a Council of the West, his plea consisting on the one
hand, in the allegation that Gregory, as being Bishop of another
Church, held the See contrary to the Canons, and on the other hand, in
the recognition which he had obtained from the Patriarch of
Alexandria. The Council, deceived by his representations, approved of
his consecration; but Theodosius, to whom he next addressed himself,
saw through his artifices, and banished him.

Fresh mortifications awaited the eloquent preacher, to whom the
Church of Constantinople owed its resurrection. While the Arians
censured his retiring habits, and his abstinence from the innocent
pleasures of life, his own flock began to complain of his neglecting
to use his influence at Court for their advantage. Overwhelmed with
the disquietudes, to which these occurrences gave birth, Gregory
resolved to bid adieu to a post which required a less sensitive or a
more vigorous mind than his own. In a farewell oration, he recounted
his labours and sufferings during the time he had been among them,
commemorated his successes, and exhorted them to persevere in the
truth, which they had learned from him. His congregation were affected
by this address; and, a reaction of feeling taking place, they
passionately entreated him to abandon a resolve, which would involve
the ruin of orthodoxy in Constantinople, and they declared that they
would not quit the church till he acceded to their importunities. At
their entreaties, he consented to suspend the execution of his purpose
for a while; that is, until the Eastern prelates who were expected at
the General Council, which had by that time been convoked, should
appoint a Bishop in his room.

The circumstances attending the arrival of Theodosius at
Constantinople, connected as they were with the establishment of the
true religion, still were calculated to inflict an additional wound on
his feelings, and to increase his indisposition to continue in his
post, endeared though it was to him by its first associations. The
inhabitants of an opulent and luxurious metropolis, familiarized to
Arianism by its forty years' ascendancy among them, and disgusted at
the apparent severity of the orthodox school, prepared to resist the
installation of Gregory in the cathedral of St. Sophia. A strong
military force was appointed to escort him thither; and the Emperor
gave countenance to the proceedings by his own presence. Allowing
himself to be put in possession of the church, Gregory was
nevertheless firm to his purpose of not seating himself upon the
Archiepiscopal throne; and when the light-minded multitude clamorously
required it, he was unequal to the task of addressing them, and
deputed one of his Presbyters to speak in his stead.

Nor were the manners of the Court more congenial to his
well-regulated mind, than the lawless spirit of the people. Offended
at the disorders which he witnessed there, he shunned the
condescending advances of the Emperor; and was with difficulty
withdrawn from the duties of his station, the solitude of his own
thoughts, and the activity of pious ministrations, prayer and fasting,
the punishment of offenders and the visitation of the sick. Careless
of personal splendour, he allowed the revenues of his see to be
expended in supporting its dignity, by inferior ecclesiastics, who
were in his confidence; and, while he defended the principle, on which
Arianism had been dispossessed of its power, he exerted himself
with earnestness to protect the heretics from all intemperate
execution of the Imperial decree.

Nor was the elevated refinement of Gregory better adapted to sway
the minds of the corrupt hierarchy which Arianism had engendered, than
to rule the Court and the people. ``If I must speak the
truth,'' he says in one of his letters, ``I feel disposed to
shun every conference of Bishops; because I never saw Synod brought to
a happy issue, nor remedying, but rather increasing, existing evils.
For ever is there rivalry and ambition, and these have the mastery of
reason;—do not think me extravagant for saying so;—and a mediator
is more likely to be attacked himself, than to succeed in his
pacification. Accordingly, I have fallen back upon myself, and
consider quiet the only security of life.'' \footnote{Greg. Naz. Ep. 55. [Ep. 130.]}

3.

Such was the state of things, under which the second Œcumenical
Council, as it has since been considered, was convoked. It met in May,
A.D. 381; being designed to put an end, as far
as might be, to those very disorders, which unhappily found their
principal exercise in the assemblies which were to remove them. The
Western Church enjoyed at this time an almost perfect peace, and sent
no deputies to Constantinople. But in the Oriental provinces, besides
the distractions caused by the various heretical off-shoots of
Arianism, its indirect effects existed in the dissensions of the
Catholics themselves; in the schism at Antioch; in the claims of
Maximus to the see of Constantinople; and in recent disturbances
at Alexandria, where the loss of Athanasius was already painfully
visible. Added to these, was the ambiguous position of the
Macedonians; who resisted the orthodox doctrine, yet were only by
implication heretical, or at least some of them far less than others.
Thirty-six of their Bishops attended the Council, principally from the
neighbourhood of the Hellespont; of the orthodox there were 150,
Meletius, of Antioch, being the president. Other eminent prelates
present were Gregory Nyssen, brother of St. Basil, who had died some
years before; Amphilochius of Iconium, Diodorus of Tarsus, Cyril of
Jerusalem, and Gelasius of Cæsarea, in Palestine.

The Council had scarcely accomplished its first act, the
establishment of Gregory in the see of Constantinople, to the
exclusion of Maximus, when Meletius, the President, died; an unhappy
event, as not only removing a check from its more turbulent members,
but in itself supplying the materials of immediate discord. An
arrangement had been effected between the two orthodox communions at
Antioch, by which it was provided, that the survivor of the rival
Bishops should be acknowledged by the opposite party, and a
termination thus put to the schism. This was in accordance with the
principle acted upon by the Alexandrian Council, on the separation of
the Meletians from the Arians. At that time the Eustathian party was
called on to concede, by acknowledging Meletius; and now, on the death
of Meletius, it became the duty of the Meletians in turn to submit to
Paulinus, whom Lucifer had consecrated as Bishop of the Eustathians.
Schism, however, admits not of these simple remedies. The
self-will of a Latin Bishop had defeated the plan of conciliation in
the former instance; and now the pride and jealousy of the Orientals
revolted from communion with a prelate of Latin creation. The attempt
of Gregory, who had succeeded to the presidency of the Council, to
calm their angry feelings, and to persuade them to deal fairly with
the Eustathians, as well as to restore peace to the Church, only
directed their violence against himself. It was in vain that his own
connection with the Meletian party evidenced the moderation and
candour of his advice; in vain that the age of Paulinus gave
assurance, that the nominal triumph of the Latins could be of no long
continuance. Flavian, who, together with others, had solemnly sworn,
that he would not accept the bishoprick in case of the death of
Meletius, permitted himself to be elevated to the vacant see; and
Gregory, driven from the Council, took refuge from its clamours in a
remote part of Constantinople.

About this time the arrival of the Egyptian bishops increased the
dissension. By some inexplicable omission they had not been summoned
to the Council; and they came, inflamed with resentment against the
Orientals. They had throughout taken the side of Paulinus, and now
their earnestness in his favour was increased by their jealousy of his
opponents. Another cause of offence was given to them, in the
recognition of Gregory before their arrival; nor did his siding with
them in behalf of Paulinus, avail to avert from him the consequences
of their indignation. Maximus was their countryman, and the deposition
of Gregory was necessary to appease their insulted patriotism.
Accordingly, the former charge was revived of the illegality of
his promotion. A Canon of the Nicene Council prohibited the
translation of bishops, priests, or deacons, from Church to Church;
and, while it was calumniously pretended, that Gregory had held in
succession three bishopricks, Sasime, Nazianzus, and Constantinople,
it could not be denied, that, at least, he had passed from Nazianzus,
the place of his original ordination, to the Imperial city. Urged by
this fresh attack, Gregory once more resolved to retire from an
eminence, which he had from the first been reluctant to occupy, except
for the sake of the remembrances, with which it was connected. The
Emperor with difficulty accepted his resignation; but at length
allowed him to depart from Constantinople, Nectarius being placed on
the patriarchal throne in his stead.

In the mean while, a Council had been held at Aquileia of the
bishops of the north of Italy, with a view of inquiring into the faith
of two Bishops of Dacia, accused of Arianism. During its session, news
was brought of the determination of the Constantinopolitan Fathers to
appoint a successor to Meletius; and, surprised both by the unexpected
continuation of the schism, and by the slight put on themselves, they
petitioned Theodosius to permit a general Council to be convoked at
Alexandria, which the delegates of the Latin Church might attend. Some
dissatisfaction, moreover, was felt for a time at the appointment of
Nectarius, in the place of Maximus, whom they had originally
recognized. They changed their petition shortly after, and expressed a
wish that a Council should be held at Rome.

These letters from the West were submitted to the Council of
Constantinople, at its second, or, (as some say,) third sitting, A.D.
382 or 383, at which Nectarius presided. An answer was returned to the
Latins, declining to repair to Rome, on the ground of the
inconvenience, which would arise from the absence of the Eastern
bishops from their dioceses; the Creed and other doctrinal statements
of the Council were sent them, and the promotion of Nectarius and
Flavian was maintained to be agreeable to the Nicene Canons, which
determined, that the Bishops of a province had the right of
consecrating such of their brethren, as were chosen by the people and
clergy, without the interposition of foreign Churches; an exhortation
to follow peace was added, and to prefer the edification of the whole
body of Christians, to personal attachments and the interests of
individuals.

Thus ended the second General Council. As to the addition made by
it to the Nicene Creed, it is conceived in the temperate spirit, which
might be expected from those men, who took the more active share in
its doctrinal discussions. The ambitious and tumultuous part of the
assembly seems to have been weary of the controversy, and to have left
its settlement to the more experienced and serious-minded of their
body. The Creed of Constantinople is said to be the composition of
Gregory Nyssen \footnote{Whether or not the Macedonians explicitly
denied the divinity of the Holy Spirit, is uncertain; but they viewed
Him as essentially separate from, and external to, the One Indivisible
Godhead. Accordingly, the Creed (which is that since incorporated into
the public services of the Church), without declaring more than the
occasion required, closes all speculations concerning the
incomprehensible subject, by simply confessing his \emph{unity with}
the Father and Son. It declares, moreover, that He is the \emph{Lord}
([\emph{kurios}]) or Sovereign Spirit, because the heretics considered
Him to be but a minister of God; and the supreme \emph{Giver} of life,
because they considered Him a mere instrument, by whom we received the
gift. The last clause of the second paragraph in the Creed, is
directed against the heresy of Marcellus of Ancyra.}.

———————

From the date of this Council, Arianism was formed into a sect
exterior to the Catholic Church; and, taking refuge among the
Barbarian Invaders of the Empire, is merged among those external
enemies of Christianity, whose history cannot be regarded as strictly
ecclesiastical. Such is the general course of religious error; which
rises within the sacred precincts, but in vain endeavours to take root
in a soil uncongenial to it. The domination of heresy, however
prolonged, is but one stage in its existence; it ever hastens to an
end, and that end is the triumph of the Truth. ``I myself have
seen the ungodly in great power,'' says the Psalmist, ``and
flourishing like a green bay tree; I went by, and lo, he was gone; I
sought him, but his place could nowhere be found.'' And so of the
present perils, with which our branch of the Church is beset, as they
bear a marked resemblance to those of the fourth century, so are the
lessons, which we gain from that ancient time, especially cheering and
edifying to Christians of the present day. Then as now, there was the
prospect, and partly the presence in the Church, of an Heretical Power
enthralling it, exerting a varied influence and a usurped claim in the
appointment of her functionaries, and interfering with the management
of her internal affairs. Now as then, ``whosoever shall fall upon
this stone shall be broken, but on whomsoever it shall fall, it
will grind him to powder.'' Meanwhile, we may take comfort in
reflecting, that, though the present tyranny has more of insult, it
has hitherto had less of scandal, than attended the ascendancy of
Arianism; we may rejoice in the piety, prudence, and varied graces of
our Spiritual Rulers; and may rest in the confidence, that, should the
hand of Satan press us sore, our Athanasius and Basil will be given us
in their destined season, to break the bonds of the Oppressor, and let
the captives go free.

\xsection{The original Creed of Nicæa, as contained in Socr. Hist. i. 8.}

[\emph{Pisteuomen eis hena theon, patera pantokratora, pant\={o}n
horat\={o}n te kai aorat\={o}n poi\={e}t\={e}n. Kai eis hena kurion I\={e}soun Christon, ton huion tou theou;
genn\={e}thenta ek tou patros monogen\={e}; tout' estin ek t\={e}s
ousias tou patros, theon ek theou, kai ph\={o}s ek ph\={o}tos,
theon al\={e}thinon ek theou al\={e}thinou; genn\={e}thenta
ou poi\={e}thenta, homoousion t\={o}i patri; di' ou ta panta
egeneto, ta te en t\={o}i ouran\={o}i kai ta en t\={e}i g\={e}i.
Di' h\={e}mas tous anthr\={o}pous kai dia t\={e}n h\={e}meteran
s\={o}t\={e}rian katelthonta, kai sark\={o}thenta, kai
enanthr\={o}p\={e}santa; pathonta, kai anastanta t\={e}i trit\={e}i
h\={e}merai, anelthonta eis tous ouranous, erchomenon krinai z\={o}ntas
kai nekrous. Kai eis to hagion pneuma. Tous de legontas, hoti \={e}n pote hote ouk \={e}n; kai prin
genn\={e}th\={e}nai ouk \={e}n; kai hoti ex ouk ont\={o}n
egeneto; \={e} ex heteras hypostase\={o}s \={e} ousias
phaskontas einai; \={e} ktiston, \={e} trepton, \={e} alloi\={o}ton
ton huion tou theou; anathematizei h\={e} hagia katholik\={e}
kai apostolik\={e} ekkl\={e}sia}.]

\xsection{Chronological Table}

\xsubsection{Persons and Events introduced into the foregoing History}

(The dates are, for the most part, according to
Tillemont.)

\textbf{A.D.}

Philo Judæus, pp. 62, 101
40

St. Mark in Egypt, p. 41
49

Cerinthus and Ebion, heretics, pp. 20, 21
90

Nazarenes, heretics, p. 21
137

Valentine, heretic, p. 55
140

Marcion, heretic, p. 55
144

Justin, p. 68, martyred
167

Theophilus, Bishop of Antioch, p. 95
168

Tatian, heretic, pp. 52, 95
169

Montanus, heresiarch, p. 16
171

Athenagoras, pp. 42, 95, writes his Apology
177

Pantænus, Missionary to the Indians, pp. 42,
102
189

Demetrius, Bishop of Alexandria, p. 42
189

Clement of Alexandria, Master of the
Catechetical

School, pp. 48-87
189

Theodotus and Artemon, heretics, pp. 22, 35,
114
193

Severus, Emperor, p. 10
193

Victor, Bishop of Rome, pp. 13, 21, 35
197

Quarto-decimans of Asia Minor, p. 15
197

Praxeas, heretic, p. 117
201

Irenæus, Bishop of Lyons, p. 54, martyred
202

Origen, aged 18, Master of the Catechetical
School,

p. 42
203

Tertullian, pp. 138, 188, falls away into
Montanism
204

Philostratus writes the Life of Apollonius
Tyanæus,

p. 109
217

Noetus, heretic, pp. 117, 124
220

Origen converts Gregory Thaumaturgus, p. 66
231

Ammonius the Eclectic, p. 102
232

Gregory Thaumaturgus delivers his panegyric on

Origen, p. 108
239

Plotinus at Rome, pp. 107, 115
244

Babylas, Bishop of Antioch martyred, p. 3
250

Novatian, heresiarch, p. 16
250

Hippolytus, p. 200; martyr
252

Death of Origen, aged 69, p. 107
253

Sabellius, heresiarch, p. 118
255

Dionysius, Bishop of Rome, animadverts on
Dionysius 

of Alexandria, p. 126
260

Paulus of Samosata, heretic, pp. 3, 27, 171,
186, 203
260

Council against Paulus, pp. 27, 128; with
Creed,

pp. 129, 192, 322, 343
264

Death of Dionysius of Alexandria, p. 108
264

Paulus deposed, p. 3
272

Quarto-decimans of the Proconsulate come to an
end, 

p. 14
276

Theonas, Bishop of Alexandria, p. 66
282

Hosius, Bishop of Corduba, pp. 249, 254, 289,
323
295

Meletian Schism in Egypt, pp. 239, 281-2
306

Donatist Schism in Africa, p. 245
306

Constantine's vision of the Labarum, p. 246
312

Lucian, martyred, p. 8
312

Edict of Milan, p. 245
313

Eusebius, Bishop of Nicomedia, p. 260
319

Arius, heresiarch, pp. 205, 237
319

Alexander excommunicates and writes against
Arius,

pp. 217, 238
320

Battle of Hadrianople, pp. 241, 247
323

Constantine writes to Athanasius and Arius, p.
247
324

Ecumenical Council of Nicæa, p. 250
325

Audius, the Quarto-deciman in Mesopotamia, p.
15
325

Athanasius, Bishop of Alexandria, p. 266
326

Arius recalled from exile, p. 266
330

Eustathius, Bishop of Antioch, deposed by the
Arians, 

pp. 280, 360
331

Eusebian Council of Cæsarea, p. 282
333

And of Tyre, \emph{ibid}. Marcellus, Bishop of
Ancyra,

deposed, pp. 280, 313
335

Athanasius banished to Treves, p. 284
335

Death of Arius, p. 269
336

Death of Constantine, who is succeeded in the
East

by Constantius, p. 280
337

Death of Eusebius of Cæsarea, who is succeeded
by

Acacius, p. 275
340

Assemblage of exiled Bishops at Rome, Council
at

Rome, p. 285
340

Eusebian Council of the Dedication at Antioch,
p. 285 

Semi-Arian Creed of Lucian, pp. 286, 322,
343
341

Semi-Arian Creed of Antioch, called the
Macrostich,

p. 287
345

Great Council of Sardica, pp. 289, 313
347

Eusebian Council, p. 289, and Semi-Arian
Creed, 

p. 342, of Philippopolis
347

Council of Milan, p. 292
347

Athanasius returns from exile, pp. 290, 360
348

Formal recantation of Valens and Ursacius, p.
291
349

Death of Constans, p. 303
350

Paulus of Constantinople martyred, p. 311
350

Battle of Mursa, p. 278
351

Eusebian Council, pp. 314, 336, with Semi-Arian

Creed of Sirmium, against Photinus, pp.
314, 322,

343
351

Eusebian Council of Arles, pp. 314, 315
353

Eusebian Council of Milan, p. 316
355

Hilary exiled in Phrygia, p. 300
356

Liberius tempted, p. 318
356

Syrianus and George in Alexandria, p. 329
356

Aetius and Eunomius, Anomœans, p. 337
357

Eusebian or Acacian Conferences and Creeds of

Sirmium; fall of Liberius and Hosius, pp.
323-

326, 341
358

Acacian Council of Antioch, p. 341
358

Semi-Arian Council of Ancyra, pp. 300, 342
358

Acacian Councils of Seleucia (p. 345) and
Ariminum,

p. 348
359

Eudoxius at Constantinople, p. 361
360

Acacian Council at Constantinople, pp. 347,
351, 358
360

Meletius, Bishop of Antioch, p. 361. Death of

Constantius, pp. 344, 352
361

Julian restores the exiled Bishops, p. 353
362

Council of Alexandria, p. 355
362

Schism of Antioch, p. 364
362

Semi-Arian Council of Lampsacus, p. 377
365

Fifty-nine Semi-Arian Bishops accept the \emph{Homoüsion},

p. 378
366

Apollinaris, heresiarch, p. 221
369

Basil, Exarch of Cæsarea, p. 375
370

Death of Eudoxius, p. 379
370

Eighty Catholic Clergy burned at sea, 380
370

Persecution of Catholics, p. 380
371

Athanasius excommunicates one of the dukes of
Lybia

p. 374
371

Death of Athanasius, p. 374
373

Death of Valens, p. 380
378

Theodosius, Emperor, p. 381
379

Gregory Nazianzen at Constantinople, \emph{ibid.}
379

Ecumenical Council of Constantinople, pp. 373,
\&c.
381

Sabbatius, Quarto-deciman, p. 17
395

\xchapter{Appendix}

\textbf{Note 1. The Syrian school of Theology}

(\emph{Vide Supra,}p. 8.)

M\textsc{uch} has been written at home, and more has come to us from
abroad, on the subject of the early Syrian theology, since this Volume
was published. At that time, it was at Oxford considered a paradox to
look to Antioch for the origin of a heresy which takes its name from
an Alexandrian ecclesiastic, and which Mosheim had ruled to be one out
of many instances of the introduction of Neo-Platonic ideas into the
Christian Church. The Divinity Professor of the day, a learned and
kind man, Dr. Burton, in talking with me on the subject, did but
qualify his surprise at the view which I had taken, by saying to me,
``Of course you have a right to your own opinion.'' Since that
time, it has become clear, from the works of Neander and others, that
Arianism was but one out of various errors, traceable to one and the
same mode of theologizing, and that mode, as well as the errors it
originated, the characteristics of the Syrian school.

I have thought it would throw light on the somewhat meagre account
of it at the beginning of this Volume, if I here added a passage on
the same subject, as contained in one of my subsequent works \footnote{``Essay on the Development of Christian
Doctrine,'' pp.
281, 323.}.

———————

The Churches of Syria and Asia Minor were the most intellectual
portion of early Christendom. Alexandria was but one metropolis
in a large region, and contained the philosophy of the whole
Patriarchate; but Syria abounded in wealthy and luxurious cities, the
creation of the Seleucidæ, where the arts and the schools of Greece
had full opportunities of cultivation. For a time too,—for the first
two hundred years, as some think,—Alexandria was the only See as
well as the only School of Egypt; while Syria was divided into small
dioceses, each of which had at first an authority of its own, and
which, even after the growth of the Patriarchal power, received their
respective bishops, not from the See of Antioch, but from their own
metropolitan. In Syria too the schools were private, a circumstance
which would tend both to diversity in religious opinion, and incaution
in the expression of it; but the sole catechetical school of Egypt was
the organ of the Church, and its Bishop could banish Origen for
speculations which developed and ripened with impunity in Syria.

But the immediate source of that fertility in heresy, which is the
unhappy distinction of the Syrian Church, was its celebrated
Exegetical School. The history of that school is summed up in the
broad characteristic fact, on the one hand that it devoted itself to
the literal and critical interpretation of Scripture, and on the other
that it gave rise first to the Arian and then to the Nestorian heresy.
In all ages of the Church, her teachers have shown a disinclination to
confine themselves to the mere literal interpretation of Scripture.
Her most subtle and powerful method of proof, whether in ancient or
modern times, is the mystical sense, which is so frequently used in
doctrinal controversy as on many occasions to supersede any other. In
the early centuries we find this method of interpretation to be the
very ground for receiving as revealed the doctrine of the Holy
Trinity. Whether we betake ourselves to the Ante-Nicene writers or the
Nicene, certain texts will meet us, which do not obviously refer to
that doctrine, yet are put forward as palmary proofs of it. On
the other hand, if evidence be wanted of the connexion of heterodoxy
and biblical criticism in that age, it is found in the fact that, not
long after their contemporaneous appearance in Syria, they are found
combined in the person of Theodore of Heraclea, so called from the
place both of his birth and his bishoprick, an able commentator and an
active enemy of St. Athanasius, though a Thracian unconnected except
by sympathy with the Patriarchate of Antioch. The case had been the
same in a still earlier age;—the Jews clung to the literal sense of
the Old Testament and rejected the Gospel; the Christian Apologists
proved its divinity by means of the allegorical. The formal connexion
of this mode of interpretation with Christian theology is noticed by
Porphyry, who speaks of Origen and others as borrowing it from heathen
philosophy, both in explanation of the Old Testament and in defence of
their own doctrine. It may almost be laid down as an historical fact
that the mystical interpretation and orthodoxy will stand or fall
together.

This is clearly seen, as regards the primitive theology, by a
recent writer, in the course of a Dissertation upon St. Ephrem. After
observing that Theodore of Heraclea, Eusebius, and Diodorus gave a
systematic opposition to the mystical interpretation, which had a sort
of sanction from Antiquity and the orthodox Church, he proceeds;
``Ephrem is not as sober in his interpretations, nor could he be,
since he was a zealous disciple of the orthodox faith. For all those
who are most eminent in such sobriety were as far as possible removed
from the faith of the Councils … On the other hand, all who retained
the faith of the Church never entirely dispensed with the spiritual
sense of the Scriptures. For the Councils watched over the orthodox
faith; nor was it safe in those ages, as we learn especially from the
instance of Theodore of Mopsuestia, to desert the spiritual for an
exclusive cultivation of the literal method. Moreover, the allegorical
interpretation, even when the literal sense was not injured, was also preserved; because in those times, when both heretics and Jews in
controversy were stubborn in their objections to Christian doctrine,
maintaining that the Messiah was yet to come, or denying the
abrogation of the Sabbath and ceremonial law, or ridiculing the
Christian doctrine of the Trinity, and especially that of Christ's
Divine Nature, under such circumstances ecclesiastical writers found
it to their purpose, in answer to such exceptions, violently to refer
every part of Scripture by allegory to Christ and His Church.'' \footnote{Lengerke, de Ephr. S. pp. 78-80.}

The School of Antioch appears to have risen in the middle of the
third century; but there is no evidence to determine whether it was a
local institution, or, as is more probable, a discipline or method
characteristic of the Syrian Church. Dorotheus is one of its earliest
teachers; he is known as a Hebrew scholar, as well as a commentator on
the sacred text, and he was the master of Eusebius of Cæsarea.
Lucian, the friend of the notorious Paul of Samosata, and for three
successive Episcopates after him a seceder from the Church, though
afterwards a martyr in it, was the editor of a new edition of the
Septuagint, and master of the chief original teachers of Arianism.
Eusebius of Cæsarea, Asterius called the Sophist, and Eusebius of
Emesa, Arians of the Nicene period, and Diodorus, a zealous opponent
of Arianism, but the Master of Theodore of Mopsuestia, have all a
place in the Exegetical School. St. Chrysostom and Theodoret, both
Syrians, and the former the pupil of Diodorus, adopted the literal
interpretation, though preserved from its abuse. But the principal
doctor of the School was the master of Nestorius, that Theodore, who
has just been mentioned, and who with his writings, and with the
writings of Theodoret against St. Cyril, and the letter written by
Ibas of Edessa to Maris, was condemned by the fifth Œcumenical
Council. Ibas translated into Syriac, and Maris into Persian, the
books of Theodore and Diodorus \footnote{Asseman. t. 3, p. 30, p. lxviii., \&c.}; and in so doing they became the immediate instruments of the
formation of the great Nestorian school and Church in farther Asia.

As many as ten thousand tracts of Theodore are said in this way to
have been introduced to the knowledge of the Christians of
Mesopotamia, Adiabene, Babylonia, and the neighbouring countries. He
was called by those Churches absolutely ``the Interpreter,''
and it eventually became the very profession of the Nestorian
communion to follow him as such. ``The doctrine of all our Eastern
Churches,'' says the Council under the patriarch Marabas, ``is
founded on the Creed of Nicæa; but in the exposition of the
Scriptures we follow St. Theodore.'' ``We must by all means
remain firm to the commentaries of the great Commentator,'' says
the Council under Sabarjesus; ``whoso shall in any manner oppose
them, or think otherwise, be he anathema.'' \footnote{Assem. t. 3, p. 84, Note 3.} No one since the beginning of Christianity, except Origen and
St. Augustine, has had such great influence on his brethren as
Theodore \footnote{Wegnern, Proleg. in Theod. Opp. p. ix.}.

The original Syrian school had possessed very marked
characteristics, which it did not lose when it passed into a new
country and into strange tongues. Its comments on Scripture seem to
have been clear, natural, methodical, apposite, and logically exact.
``In all Western Aramæa,'' says Lengerke, that is, in Syria,
``there was but one mode of treating whether exegetics or
doctrine, the practical.'' \footnote{De Ephræm Syr. p. 61.}
Thus Eusebius of Cæsarea, whether as a disputant or a commentator, is
confessedly a writer of sense and judgment, and he belongs
historically to the Syrian school, though he does not go so far as to
exclude the mystical interpretation or to deny the verbal inspiration
of Scripture. Again, we see in St. Chrysostom a direct,
straightforward treatment of the sacred text, and a pointed
application of it to things and persons; and Theodoret abounds in
modes of thinking and reasoning which without any great impropriety
may be called English. Again, St. Cyril of Jerusalem, though he does
not abstain from allegory, shows the character of his school by the
great stress he lays upon the study of Scripture, and, I may add, by
the peculiar clearness and neatness of his style, which will be
appreciated by a modern reader.

It would have been well, had the genius of the Syrian theology been
ever in the safe keeping of men such as St. Cyril, St. Chrysostom, and
Theodoret; but in Theodore of Mopsuestia, nay in Diodorus before him,
it developed into those errors, of which Paul of Samosata had been the
omen on its rise. As its attention was chiefly directed to the
examination of the Scriptures, in its interpretation of the Scriptures
was its heretical temper discovered; and though allegory can be made
an instrument of evading Scripture doctrine, criticism may more
readily be turned to the destruction of doctrine and Scripture
together. Bent on ascertaining the literal sense, Theodore was
naturally led to the Hebrew text instead of the Septuagint, and thence
to Jewish commentators. Jewish commentators naturally suggested events
and objects short of evangelical as the fulfilment of the prophetical
announcements, and when it was possible, an ethical sense instead of a
prophetical. The eighth chapter of Proverbs ceased to bear a Christian
meaning, because, as Theodore maintained, the writer of the book had
received the gift, not of prophecy, but of wisdom. The Canticles must
be interpreted literally; and then it was but an easy, or rather a
necessary step, to exclude the book from the Canon. The book of Job
too professed to be historical; yet what was it really but a Gentile
drama? He also gave up the books of Chronicles and Ezra, and, strange
to say, the Epistle of St. James, though it was contained in the
Peschito Version of his Church. He denied that Psalms xxii. and
lxix. applied to our Lord; rather he limited the Messianic passages of
the whole book to four; of which the eighth Psalm was one, and the
forty-fifth another. The rest he explained of Hezekiah and Zerubbabel,
without denying that they might be accommodated to an evangelical
sense \footnote{Lengerke, de Ephræm Syr. pp. 73-75.}. He explained St.
Thomas's words, ``My Lord and my God,'' as a joyful
exclamation; and our Lord's, ``Receive ye the Holy Ghost,'' as
an anticipation of the day of Pentecost. As might be expected, he
denied the verbal inspiration of Scripture. Also, he held that the
deluge did not cover the earth; and, as others before him, he was
heterodox on the doctrine of original sin, and denied the eternity of
punishment.

Maintaining that the real sense of Scripture was, not the scope of
a Divine Intelligence, but the intention of the mere human organ of
inspiration, Theodore was led to hold, not only that that sense was
but one in each text, but that it was continuous and single in a
context; that what was the subject of the composition in one verse,
must be the subject in the next, and that if a Psalm was historical or
prophetical in its commencement, it was the one or the other to its
termination. Even that fulness of meaning, refinement of thought,
subtle versatility of feeling, and delicate reserve or reverent
suggestiveness, which poets exemplify, seem to have been excluded from
his idea of a sacred composition. Accordingly, if a Psalm contained
passages which could not be applied to our Lord, it followed that that
Psalm did not properly apply to Him at all, except by accommodation.
Such at least is the doctrine of Cosmas, a writer of Theodore's
school, who on this ground passes over the twenty-second, sixty-ninth,
and other Psalms, and limits the Messianic to the second, the eighth,
the forty-fifth, and the hundred and tenth. ``David,'' he
says, ``did not make common to the servants what belongs to the
Lord \footnote{[\emph{despotou}], vide La Croze, Thesaur.
Ep. t. 3, § 145.} Christ, but what was
proper to the Lord he spoke of the Lord, and what was proper to
the servants, of servants.'' \footnote{Montf. Coll. Nov. t. 2, p. 227.} Accordingly the twenty-second could not properly belong to
Christ, because in the beginning it spoke of the ``\emph{verba
delictorum meorum}.'' A remarkable consequence would follow
from this doctrine, that as Christ was divided from His Saints, so the
Saints were divided from Christ; and an opening was made for a denial
of the doctrine of their \emph{cultus}, though this denial in the
event has not been developed among the Nestorians. But a more serious
consequence is latently contained in it, and nothing else than the
Nestorian heresy, viz. that our Lord's manhood is not so intimately
included in His Divine Personality that His brethren according to the
flesh may be associated with the Image of the One Christ. Here St.
Chrysostom pointedly contradicts the doctrine of Theodore, though his
fellow-pupil and friend \footnote{Rosenmuller, Hist. Interpr. t. 3, p. 278.};
as does St. Ephræm, though a Syrian also \footnote{Lengerke, de Ephr. Syr. pp. 165-167.}; and St. Basil \footnote{Ernest. de Proph. Mess. p. 462.}.

One other characteristic of the Syrian school, viewed as
independent of Nestorius, should be added:—As it tended to the
separation of the Divine Person of Christ from His manhood, so did it
tend to explain away His Divine Presence in the Sacramental elements.
Ernesti seems to consider that school, in modem language,
Sacramentarian: and certainly some of the most cogent passages brought
by moderns against the Catholic doctrine of the Eucharist are taken
from writers who are connected with that school; as the author, said
to be St. Chrysostom, of the Epistle to Cæsarius, Theodoret in his
Eranistes, and Facundus. Some countenance too is given to the same
view of the Eucharist, at least in some parts of his works, by Origen,
whose language concerning the Incarnation also leans to what was afterwards Nestorianism. To these maybe added Eusebius \footnote{Eccl. Theol. iii. 12.}, who, far removed, as he was, from that heresy, was a disciple
of the Syrian school. The language of the later Nestorian writers
seems to have been of the same character \footnote{Professor Lee's Serm. Oct. 1838, pp.
144-152.}. Such then on the whole is the character of that theology of
Theodore, which passed from Cilicia and Antioch to Edessa first, and
then to Nisibis.

Edessa, the metropolis of Mesopotamia, had remained an Oriental
city till the third century, when it was made a Roman colony by
Caracalla \footnote{Noris. Opp. t. 2, p. 112.}. Its
position on the confines of two empires gave it great ecclesiastical
importance, as the channel by which the theology of Rome and Greece
was conveyed to a family of Christians, dwelling in contempt and
persecution amid a still heathen world. It was the seat of various
schools; apparently of a Greek school, where the classics were studied
as well as theology, where Eusebius of Emesa \footnote{Augusti. Euseb. Em. Opp.} had originally been trained, and where perhaps Protogenes
taught \footnote{Asseman. p. cmxxv.}. There were
Syrian schools attended by heathen and Christian youths in common. The
cultivation of the native language had been an especial object of its
masters since the time of Vespasian, so that the pure and refined
dialect went by the name of the Edessene \footnote{Hoffman, Gram. Syr. Proleg. § 4.}. At Edessa too St. Ephræm formed his own Syrian school, which
lasted long after him; and there too was the celebrated Persian
Christian school, over which Maris presided, who has been already
mentioned as the translator of Theodore into Persian \footnote{The educated Persians were also acquainted
with Syriac.—Assem. t. i. p. 351, Note.}. Even in the time of the predecessor of Ibas in the See
(before A.D. 435) the Nestorianism of this
Persian School was so notorious that Rabbula the Bishop had expelled
its masters and scholars \footnote{Asseman, p. lxx.}; and they, taking refuge in the country with which they were
connected, had introduced the heresy to the Churches subject to the
Persian King.

Something ought to be said of these Churches; though little is
known except what is revealed by the fact, in itself of no slight
value, that they had sustained two persecutions at the hands of the
heathen government in the fourth and fifth centuries. One testimony is
extant as early as the end of the second century, to the effect that
in Parthia, Media, Persia, and Bactria there were Christians who
``were not overcome by evil laws and customs.'' \footnote{Euseb. Præp. vi. 10.} In the early part of the fourth century, a Bishop of Persia
attended the Nicene Council, and about the same time Christianity is
said to have pervaded nearly the whole of Assyria \footnote{Tillemont, Mem. t. 7, p. 77.}. Monachism had been introduced there before the middle of the
fourth century, and shortly after commenced that fearful persecution
in which sixteen thousand Christians are said to have suffered. It
lasted thirty years, and is said to have recommenced at the end of the
century. The second persecution lasted for at least another thirty
years of the next, at the very time when the Nestorian troubles were
in progress in the Empire. Trials such as these show the populousness
as well as the faith of the Churches in those parts; and the number of
the Sees, for the names of twenty-seven Bishops are preserved who
suffered in the former persecution. One of them was apprehended
together with sixteen priests, nine deacons, besides monks and nuns of
his diocese; another with twenty-eight companions, ecclesiastics or
regulars; another with one hundred ecclesiastics of different orders;
another with one hundred and twenty-eight; another with his
chorepiscopus and two hundred and fifty of his clergy. Such was the
Church, consecrated by the blood of so many martyrs, which
immediately after its glorious confession fell a prey to the theology
of Theodore; and which through a succession of ages discovered the
energy, when it had lost the purity of saints.

The members of the Persian school, who had been driven out of
Edessa by Rabbula, found a wide field open for their exertions under
the pagan government with which they had taken refuge. The Persian
monarchs, who had often prohibited by edict \footnote{Gibbon, ch. 47.} the intercommunion of the Church under their sway with the
countries towards the west, readily extended their protection to
exiles, who professed the means of destroying its Catholicity.
Barsumas, the most energetic of them, was placed in the metropolitan
See of Nisibis, where also the fugitive school was settled under the
presidency of another of their party; while Maris was promoted to the
See of Ardaschir. The primacy of the Church had from an early period
belonged to the See of Seleucia in Babylonia. Catholicus was the title
appropriated to its occupant, as well as to the Persian Primate, as
being deputies of the Patriarch of Antioch, and was derived apparently
from the Imperial dignity so called, denoting their function as
Procurators-general, or officers-in-chief for the regions in which
they were placed. Acacius, another of the Edessene party, was put into
this principal See, and suffered, if he did not further, the
innovations of Barsumas. The mode by which the latter effected his
purposes has been left on record by an enemy. ``Barsumas accused
Barbuæus, the Catholicus, before King Pherozes, whispering, 'These
men hold the faith of the Romans, and are their spies. Give me power
against them to arrest them.''' \footnote{Asseman, p. lxxviii.} It is said that in this way he obtained the death of Barbuæus,
whom Acacius succeeded. When a minority resisted \footnote{Gibbon, ibid.} the process of schism, a persecution followed. The death of
seven thousand seven hundred Catholics is said by Monophysite
authorities to have been the price of the severance of the Chaldaic
Churches from Christendom \footnote{Asseman, t. 2, p. 403, t. 3, p. 393.}.
Their loss was compensated in the eyes of the government by the
multitude of Nestorian fugitives, who flocked into Persia from the
Empire, numbers of them industrious artisans, who sought a country
where their own religion was in the ascendant.

The foundation of that religion lay, as we have already seen, in
the literal interpretation of Scripture, of which Theodore was the
principal teacher. The doctrine, in which it formerly consisted, is
known by the name of Nestorius: it lay in the ascription of a human as
well as a Divine Personality to our Lord; and it showed itself in
denying the title of ``Mother of God'' or [\emph{theotokos}],
to St. Mary. As to our Lord's Personality, it is to be observed that
the question of language came in, which always serves to perplex a
subject and make a controversy seem a matter of words. The native
Syrians made a distinction between the word ``Person,'' and
``Prosopon,'' which stands for it in Greek; they allowed that
there was one Prosopon or Parsopa, as they called it, and they held
that there were two Persons. It is asked what they meant by \emph{parsopa}:
the answer seems to be, that they took the word merely in the sense of
\emph{character} or \emph{aspect}, a sense familiar to the Greek \emph{prosopon},
and quite irrelevant as a guarantee of their orthodoxy. It follows
moreover that, since the \emph{aspect} of a thing is its impression
upon the beholder, the personality to which they ascribed unity must
have lain in our Lord's manhood, and not in His Divine Nature. But it
is hardly worth while pursuing the heresy to its limits. Next, as to
the phrase ``Mother of God,'' they rejected it as
unscriptural; they maintained that St. Mary was Mother of the humanity
of Christ, not of the Word, and they fortified themselves by the
Nicene Creed, in which no such title is ascribed to her.

Whatever might be the obscurity or the plausibility of their
original dogma, there is nothing obscure or attractive in the
developments, whether of doctrine or of practice, in which it issued.
The first act of the exiles of Edessa, on their obtaining power in the
Chaldean communion, was to abolish the celibacy of the clergy, or, in
Gibbon's forcible words, to allow ``the public and reiterated
nuptials of the priests, the bishops, and even the patriarch
himself.'' Barsumas, the great instrument of the change of
religion, was the first to set an example of the new usage, and is
even said by a Nestorian writer to have married a nun \footnote{Asseman, t. 3, p. 67.}. He passed a
Canon at Councils, held at Seleucia and elsewhere, that Bishops and
priests might marry, and might renew their wives as often as they lost
them. The Catholic who followed Acacius went so far as to extend the
benefit of the Canon to Monks, that is, to destroy the Monastic order;
and his two successors availed themselves of this liberty, and are
recorded to have been fathers. A restriction, however, was afterwards
placed upon the Catholic, and upon the Episcopal order.

Such were the circumstances, and such the principles, under which
the See of Seleucia became the Rome of the East. In the course of time
the Catholic took on himself the loftier and independent title of
Patriarch of Babylon; and though Seleucia was changed for Ctesiphon
and for Bagdad \footnote{Gibbon, ibid.}, still
the name of Babylon was preserved from first to last as a formal or
ideal Metropolis. In the time of the Caliphs, it was at the head of as
many as twenty-five Archbishops; its Communion extended from China to
Jerusalem; and its numbers, with those of the Monophysites, are said
to have surpassed those of the Greek and Latin Churches together.

\textbf{Note 2. The Doctrine of the Divine Gennesis
according to the early Fathers}\\
(\emph{Vide supra}, p. 240.)

ALREADY in the Notes (Oxf.) on
Athanasius (Ath. Tr. pp. 272-280), and in Dissert. Theolog. iii. I
have explained my difficulty in following Bull and others in the
interpretation they assign to certain statements made in the first age
of the Church concerning the Divine Sonship. Those statements, taken
in their letter, are to the effect that our Lord was the word of God
before He was the Son; that, though, as the Word, He was from
eternity, His \emph{gennesis}is in essential connexion both with the
design and the fact of creation; that He was born indeed of the Father
apart from all time, but still with a definite relation to that
beginning of time when the creation took place, and though born, and
not created, nevertheless born definitely in order to create.

Before the Nicene Council, of the various Schools of the Church,
the Alexandrian alone, is distinctly clear of this doctrine; and even
after the Council it is found in the West, in Upper Italy, Rome, and
Africa; France, as represented by Hilary \footnote{Vide however Hilar. in Matt. xxxi. 3; but he corrects himself,
de Trin. xii.} and Phœbadius, having no part in it. Nay, at Nicæa when that
doctrine lay in the way of the Council to condemn it, it was not
distinctly condemned, though to pass it over was in fact to give it
some countenance. Bull indeed considers it was even recognized
indirectly by the assembled Fathers, in their anathematizing those who
contradicted its distinctive formula, ``He was before He was
born;'' in this (as I have said in the Notes on Athanasius),
I cannot agree with him, but at least it is unaccountable that the
Fathers should not have guarded their anathema from Bull's easy
misinterpretation of it, if the opinion which it seems to countenance
was as much reprobated then, as it rightly is now.

The opinion which I have been describing is, as far as words go,
definitely held by Justin, Tatian, Theophilus, Methodius, in the East;
by Hippolytus, Tertullian, Novatian, Lactantius, Zeno, and Victorinus,
in the West; and that with so plain an identity of view in these
various writers, and with such exact characteristics, that we cannot
explain it away into carelessness of writing, personal idiosyncracy,
or the influence of some particular school; but are forced to consider
it as the common property of them all, so that we may interpret one
writer by the other, and illustrate or supply from the rest what is
obscure or deficient in each.

For instance: Justin says, ``He was begotten, when God at the
beginning through Him created and adorned all things'' (Ap. ii.
6). ``Not a perfect Son, without the flesh, though a perfect
Word,'' says Hippolytus, ``being the Only-begotten, … whom
God called 'Son,' because He was to become such'' (contr. Noet.
15) ... ``There was a time when the Son was not,'' says
Tertullian (adv. Herm. 3); ``He proceeds unto a birth,'' says
Zeno, ``\emph{He} who was, before He was born'' (Tract. ii.
3).

There can be no doubt what the literal sense is of words such as
these, and that in consequence they require some accommodation in
order to reconcile them with the received Catholic teaching \emph{de Deo}
and \emph{de SS. Trinitate}. It is the object of Bull, as of others
after him, to effect this reconciliation. He thinks it a plain duty
both to the authors in question and to the Church, at whatever cost,
to reconcile their statements in all respects with the orthodox
belief; but unless he had felt it a duty, I do not think he would
have ventured upon it. He would have taken them in their literal
sense, had he found them in the writing of some Puritan or Quaker. If
so, his defence of them is but a confirmation of a foregone
conclusion; he starts with the assumption that the words of the early
writers cannot mean what they naturally do mean; and, though this bias
is worthy of all respect, still the fact that it exists is a call on
us to examine closely arguments which without it would not have been
used. And what I have said of Bull applies of course to others, such
as Maran and the Ballerini, who have followed in his track.

Bull then maintains that the terms ``generation,''
``birth,'' and the like, which occur in the passages of the
authors in question, must be accommodated to a literary sense, that
is, taken figuratively, or \emph{impropriè}, to mean merely our Lord's
going forth to create, and the great manifestation of the Sonship made
in and to the universe at its creation; and on these grounds:—1. The
terms used cannot be taken literally, from the fact that in those very
passages, or at least in other passages of the same authors, His
co-eternity with the Father is expressly affirmed. 2. And they must be
taken figuratively, first, because in those passages they actually
stand in connexion with mention of His forthcoming or mission to
create; and next, because unsuspected authors, such as Athanasius,
distinctly connect His creative \emph{office} with His title of
``First-born,'' which belongs to His \emph{nature}.

Now I do not think these arguments will stand; as to the negative
argument, it is true that the Fathers, who speak of the \emph{gennesis}
as having a relation to time and to creation, do in the same passages
or elsewhere speak of the eternity of the Word. Doubtless; but no one
says that these Fathers deny His eternity, as the \emph{Word}, but His
eternity as the \emph{Son}. Bull ought to bring passages in which they
declare the Son and His \emph{gennesis} to be eternal.

As to the positive argument, if they recognized, as he thinks,
any \emph{gennesis} besides that which had a relation to
creation, and which he maintains to be only figuratively a \emph{gennesis},
viz. an eternal \emph{gennesis} from the substance of the Father, why
do they not say so? do they ever compare and contrast the two births
with each other? do they ever recognize them as two, one real and
eternal, the other just before time; the one proper, the other
metaphorical? We know they held a \emph{gennesis} in order to
creation, or with a relation to time; what reason have we for holding
that they held any other? and what reason for saying that the \emph{gennesis}
which they connect with creation was not in their minds a real \emph{gennesis},
that is, such a \emph{gennesis} as we all now hold, all but, as they
expressly state, its not being from eternity?

In other words, what reason have we for saying that the term \emph{gennesis}
is figurative in their use of it? It is true indeed that both the Son's
\emph{gennesis} and also His forthcoming, mission, or manifestation
are sometimes mentioned together by these writers in the same
sentence; but that does not prove they are not in their minds separate
Divine acts; for His creation of the world is mentioned in such
passages too, and as His creation of the world is not His mission,
therefore His mission need not be His \emph{gennesis}; and again, as
His creating is (in their teaching) concurrent with His mission, so
His mission may (in their teaching) be concurrent with His \emph{gennesis}.

Nor are such expositions of the title ``First-born of
creation,'' as Athanasius has so beautifully given us, to the
purpose of Bull. Bull takes it to show that \emph{gennesis} may be
considered to be a mission or forthcoming; whereas Athanasius does not
mean by the ``First-born'' any \emph{gennesis} of our Lord
from the Father at all, but he simply means His coming to the
creature, that is, His exalting the creature into a Divine sonship by
a union with His own Sonship. The \emph{Word} applies His own Sonship
to the creation, and makes Himself, who is the real Son, the first and
the representative of a family of adopted sons \footnote{[(1890, \emph{Add}): adopted, that is,
through the grace of Him, who is in His nature, from eternity, the One
and Only Son of God.]}; the term expresses a relation, not towards God, but towards
the creature. This Athanasius says expressly: ``It is nowhere
written [of the Son] in the Scriptures, 'the First-born \emph{of God},'
nor 'the creature \emph{of God},' but it is 'Only-begotten,' and
'Son,' and 'Word,' and 'Wisdom,' that have relation to the Father. The
same cannot be both Only-begotten and First-born, except in different
relations,—Only-begotten because of His \emph{gennesis}, First-born
because of His condescension.'' Thus Athanasius expressly denies
that, because our Lord is First-born at and to the creation, therefore
He can be said to be begotten at the creation;
``Only-begotten'' is internal to the Divine Essence;
``First-born'' external to It: the one is a word of nature,
the other, of office. If then the authors, whom Bull is defending, had
wished to express a figurative \emph{gennesis}, they would always have
used the word ``First-born,'' never ``Only-begotten:''
and never have associated the generation from the Father with the
coming forth to create. It is true they sometimes associate the Word's
creative office with the term ``First-born;'' but they also
associate it with ``Only-begotten.''

There seems no reason then why the words of Theophilus, Hippolytus,
and the rest should not be taken in their obvious sense; and so far I
agree with Petavius against Bull, Fabricius, Maran, the Ballerini, and
Routh. But, this being granted, still I am not disposed to follow
Petavius in his severe criticism upon those Fathers, and for the
following reasons:—

1. They considered the ``Theos Logos'' to be really
distinct from God, (that is, the Father,) not a mere attribute,
quality, or power, as the Sabellians did, and do.

2. They considered Him to be distinct from God \emph{from everlasting}.

3. Since, as Dionysius says, ``He who speaks is father of his
words,'' they considered the Logos always to be of the \emph{nature}
of a Son. Hence Zeno says He was from everlasting ``Filii
non sine \emph{affectu},'' and Hippolytus, [\emph{teleios logos, \={o}n
monogen\={e}s}].

4. They considered, to use the Scripture term, that He was ``\emph{in
utero Patris}'' before His actual \emph{gennesis}. Victorinus
applies the word ``fœtus'' to Him; ``Non enim fœtus non
est ante partum; sed in occulto est; generatio est manifestatio''
(apud Galland, v. 8, p. 146, col. 2). Zeno says that He ``prodivit
\emph{ex ore Dei} ut rerum naturam fingeret,'' ``\emph{cordis}
ejus nobilis inquilinus,'' and was embraced by the Father ``profundo
suæ sacræ \emph{mentis arcano} sine revelamine.''

5. Hippolytus even considered that the perfection of His Sonship
was not attained till His incarnation, [\emph{teleios logos huios atel\={e}s}];
but even he recognized the identity of the Son with the Logos.

6. Further, this change of the Logos into the Son was internal to
the Divine Mind, Tertull. adv. Prax. 8. contr. Hermog. 18, and
therefore was unlike the \emph{probole} of the Gnostics.

7. Such an opinion was not only not inconsistent with the \emph{Homoüsion},
but implied it. It took for granted that the Son was from the
substance of the Father, and consubstantial with Him; though it
implied a very defective view of the immutability and simplicity of
the Divine Essence.

8. Accordingly, though I cannot allow that it was actually
protected at the Council by the anathema on those who said that our
Lord ``was not before He was born,'' at least it was passed
over on an occasion when the Arian error had to be definitively
reprobated.

This may be said in its favour: but then, on the other hand,—

1. It seriously compromised, as I have said, the simplicity and
immutability of the Divine Essence.

2. It could be resolved, with very little alteration, into
Semi-Arianism on the one hand, or into Sabellianism on the other.

3. On this account it had all along been resisted with definiteness
and earnestness by the Fathers of the Alexandrian School, by whom
finally it was eradicated. Origen urges the doctrine of the [\emph{aeigennes}];
``Perfect Son from Perfect Father,'' says Gregory Thaumaturgus
in his creed; ``The Father being everlasting the Son is
everlasting,'' says Dionysius; ``The Father,'' says
Alexander, ``is ever Father of the ever-present Son,'' and
Athanasius reprobates the [\emph{logos en t\={o}i the\={o}i atel\={e}s,
genn\={e}theis teleios}] (Orat. iv. ii). Hence Gregory Nazianzen
in like manner condemns the [\emph{atel\={e} proteron, eita teleion h\={o}sper
nomos h\={e}meteros genese\={o}s}] (Orat. xx. 9, fin.). And
at length it was classed, and duly, among the heresies. ``Alia (hæresis),''
says Augustine, ``sempiternè natum non intelligens Filium, putat
illam nativitatem sumpsisse à tempore initium; et tamen volens
coæternum Patri Filium confiteri, apud illum fuisse, antequàm de
illo nasceretur, existimat; hoc est, semper eum fuisse, veruntamen
semper eum Filium non fuisse, sed ex quo de illo natus est, Filium
esse cœpisse'' (Hær. 50).

However, this subject should be treated at greater length that I
can allow it here. (Vide Tracts Theol. and
Eccles.)

N.B.—The above addition (page 420) is in consequence of a
misunderstanding, which leads me to repeat, now, 1890, as ever, that
what I have here written is subject to the judgment of Holy Church.

\textbf{Note 3. The Confessions at Sirmium}\\
(\emph{Vide supra}, p. 314.)

\textbf{1. A.D. 351.
Confession against Photinus}\\
(First Sirmian Council)

\footnote{From the Oxford Translation of Athanasius, p. 160.} THIS
Confession was published at a Council of Eastern Bishops (Coustant. in
Hil. p. 1174, Note 1), and was drawn up by the whole body, Hil. de Syn.
37 (according to Sirmond. Diatr. 1. Sirm. p. 366, Petavius de Trin. 1.
9. § 8. Animadv. in Epiph. p. 318 init., and Coustant. in Hil. l.
c.); or by Basil of Ancyra (as Valesius conjectures in Soz. iv. 22,
and Larroquanus, de Liberio, p. 147); or by Mark of Arethusa, Socr.
ii. 30, but Socrates, it is considered, confuses together the dates of
the different Confessions, and this ascription is part of his mistake
(vide Vales. in loc., Coustant. in Hil. de Syn. l. c., Petav. Animad.
in Epiph. l. c.). It was written in Greek.

Till Petavius, Socrates was generally followed in ascribing all
three Sirmian Confessions to this one Council, though at the same time
he was generally considered mistaken as to the year. E.g. Baronius
places them all in 357. Sirmond defended Baronius against Petavius
(though in Facund. x. 6, Note c, he agrees with Petavius); and,
assigning the third Confession to 359, adopted the improbable
conjecture of two Councils, the one Catholic and the other Arian, held
at Sirmium at the same time, putting forth respectively the first and
second Creeds, somewhat after the manner of the contemporary rival
Councils of Sardica. Pagi, Natalis Alexander, Valesius, de Marca,
Tillemont, S. Basnage, Montfaucon, Coustant, Larroquanus agree with
Petavius in placing the Council, at which Photinus was deposed and the
Confession published, in A.D. 351. Mansi dates
it at 358.

Gothofred considers that there were two or three successive
Councils at Sirmium, between A.D. 357 and 359
(in Philostorg. Index, pp. 74, 75; Dissert. pp. 200, 211-214).
Petavius, and Tillemont, speak of three Councils or Conferences held
in A.D. 351, 357, and 359. Mansi, of three in
358, 359; Zaccaria (Dissert. 8) makes in all five, 349 (in which
Photinus was condemned), 351; 357 (in which Hosius lapsed); 357
(following Valesius and Pagi); and 359. Mamachi makes three, 351, 357,
359; Basnage four, 351, 357, 358, 359.

This was the Confession which Pope Liberius signed, according to
Baronius, Natalis Alexander, and Coustant in Hil. Note n. pp.
1335-1337, and as Tillemont thinks probable. Zaccaria says it is the
general opinion, in which he is willing to concur (p. 18).

It would appear (Ath. Tr. p. 114, b.) that Photinus was condemned
at Antioch in the Macrostich, A.D. 345; at
Sardica, 347; at Milan, 348; and at his own See, Sirmium, 351, if not
there, in 349 also;—however, as this is an intricate point on which
there is considerable difference of opinion among critics, it may be
advisable to state here the dates of his condemnation as they are
determined by various writers.

Petavius (de Photino Hæretico, 1) enumerates in all five
condemnations:—1. at Constantinople, A.D. 336,
when Marcellus was deposed. 2. At Sardica, A.D.
347. 3. At Milan, A.D. 347. 4. At Sirmium, A.D.
349. 5. At Sirmium, when he was deposed, A.D.
351. Of these the 4th and 5th were first brought to light by Petavius,
who omits mention of the Macrostich in 345.

Petavius is followed by Natalis Alexander, Montfaucon (vit. Athan.),
and Tillemont; and by De Marca (Diss. de temp. Syn. Sirm.) and S.
Basnage (Annales), and Valesius (in Theod. Hist. ii. 16. p. 23; Socr.
ii. 20), as regards the Council of Milan, except that Valesius places
it with Sirmond in 346; but for the Council of Sirmium in 349, they
substitute a Council of Rome of the same date, while De Marca
considers Photinus condemned again in the Eusebian Council of Milan in
355. De la Roque, on the other hand (Larroquan. Dissert. de Photino
Hær.), considers that Photinus was condemned, 1. in the Macrostich,
344 [345]. 2. At Sardica, 347. 3. At Milan, 348. 4. At Sirmium, 350.
5. At Sirmium, 351. Zaccaria, besides 345 and 347; at Milan, 347; at
Sirmium, 349; at Sirmium again, 351, when he was deposed.

Petavius seems to stand alone in assigning to the Council of
Constantinople, 336, his first condemnation.

\textbf{2. A.D. 357. The
Blasphemy of Potamius and Hosius}\\
(Second Sirmian)

Hilary calls it by the above title, de Syn. 11; vide also Soz. iv.
12, p. 554. He seems also to mean it by the blasphemia Ursacii et
Valentis, contr. Const. 26.

This Confession was the first overt act of disunion between Arians
and Semi-Arians.

Sirmond, De Marca, and Valesius (in Socr. ii. 30), after Phœbadius,
think it put forth by a Council; rather, at a Conference of a few
leading Arians about Constantius, who seems to have been present; e.g.
Ursacius, Valens, and Germinius. Soz. iv. 12. Vide also Hil. Fragm.
vi. 7.

It was written in Latin, Socr. ii. 30. Potamius wrote very
barbarous Latin, judging from the Tract ascribed to him in Dacher.
Spicileg. t. 3. p. 299, unless it be a translation from the
Greek, vide also Galland. Bibl. t. v. p. 96. Petavius thinks the Creed
not written, but merely subscribed by Potamius (de Trin. i. 9. § 8);
and Coustant (in Hil. p. 1155, Note f) that it was written by Ursacius,
Valens, and Potamius. It is remarkable that the Greek in Athanasius is
clearer than the original.

This at first sight is the Creed which Liberius signed, because S.
Hilary speaks of the latter as ``perfidia Ariana,'' Fragm. 6.
Blondel (Prim. dans l'Eglise, p. 484), Larroquanus, \&c., are of
this opinion. And the Roman Breviary, Ed. Ven. 1482, and Ed. Par.
1543, in the Service for S. Eusebius of Rome, August. 14, says that
``Pope Liberius consented to the Arian misbelief,'' Launnoi,
Ep. v. 9. c. 13. Auxilius says the same, Ibid. vi. 14. Animadv. 5. n.
18. Petavius grants that it must be this, \emph{if} any of the three
Sirmian (Animadv. in Epiph. p. 316), but we shall see his own opinion
presently. Zaccaria says that Hosius signed it, but not Liberius (Diss.
8. p. 20, Diss. 7). Zaccaria seems also to consider that there was
another Council or Conference at Sirmium this same year, and it was at
this Conference that Liberius subscribed ``formulæ, quæ contra
Photinum Sirmii edita fuerat, primæ scilicet Sirmiensi, in unum cum
Antiochensi (against Paul of Samosata, also the creed of the
Dedication) libellum conjectæ.'' \emph{Vide infra}. He says he
subscribed it ``iterum,'' the first time being in Berrhœa.

\xsubsection{3. A.D. 357. The foregoing interpolated}

A creed was sent into the East in Hosius's name, Epiph. Hær. 73.
14. Soz. iv. 15, p. 558, of an Anomœan character, which the ``blasphemia''
was not. And St. Hilary may allude to this when he speaks of the
``deliramenta Osii, et \emph{incrementa} Ursacii et Valentis,''
contr. Const. 23. An Anomœan Council of Antioch under Eudoxius
of this date, makes acknowledgments to Ursacius, Valens, and Germinius,
Soz. iv. 12 fin, as being agents in the Arianizing of the West.

Petavius and Tillemont consider this Confession to be the ``blasphemia''
interpolated. Petavius throws out a further conjecture, which seems
gratuitous, that the whole of the latter part of the Creed is a later
addition, and that Liberius only signed the former part. Animadv. in
Epiph. p. 316.

\xsubsection{4. A.D. 358. The Ancyrene Anathemas}

The Semi-Arian party had met in Council at Ancyra in the early
spring of 358 to protest against the ``blasphemia,'' and that
with some kind of correspondence with the Gallic Bishops who had just
condemned it, Phœbadius of Agen writing a Tract against it, which is
still extant. They had drawn up and signed, besides a Synodal Letter,
eighteen anathemas, the last against the ``Consubstantial.''
These, except the last, or the last six, they submitted at the end of
May to the Emperor who was again at Sirmium. Basil, Eustathius,
Eleusius, and another formed the deputation; and their influence
persuaded Constantius to accept the Anathemas, and even to oblige the
party of Valens, at whose ``blasphemia'' they were levelled,
to recant and subscribe them.

\xsubsection{5. A.D. 358. Semi-Arian Digest of Three Confessions}

The Semi-Arian Bishops, pursuing their advantage, composed a Creed
out of three, that of the Dedication, the first Sirmian, and the Creed
of Antioch against Paul, 264-270, in which the
``Consubstantial'' is said to have been omitted or forbidden,
Soz. iv. 15. This Confession was imposed by Imperial authority
on the Arian party, who signed it. So did Liberius, Soz. ibid. Hil.
Fragm. vi. 6, 7; and Petavius considers that this is the subscription
by which he lapsed, de Trin. i. 9. § 5, Animadv. in Epiph. p. 316,
and so Zaccaria, as above, and S. Basnage, in Ann. 358, 13.

It is a point of controversy whether or not the Arians at this time
suppressed the ``blasphemia.'' Socrates and Sozomen say that
they made an attempt to recall the copies they had issued, and even
obtained an edict from the Emperor for this purpose, but without
avail. Socr. ii. 30 fin. Soz. iv. 6, p. 543.

Athanasius, on the other hand, de Syn. 29, relates this in
substance of the third Confession of Sirmium, not of the ``blasphemia''
or second.

Tillemont follows Socrates and Sozomen, considering that Basil's
influence with the Emperor enabled him now to insist on a retraction
of the ``blasphemia.'' And he argues that Germinius in 366,
being suspected of orthodoxy, and obliged to make profession of
heresy, was referred by his party to the formulary of Ariminum, no
notice being taken of the ``blasphemia,'' which looks as if it
were suppressed; whereas Germinius himself appeals to the third
Sirmian, which is a proof that it was not suppressed. Hil. Fragm. 15.
Coustant, in Hil. contr. Const. 26, though he does not adopt the
opinion himself, observes, that the charge brought against Basil, Soz.
iv. 132, Hil. l. c., by the Acacians, of persuading the Africans
against the second Sirmian is an evidence of a great effort on his
part, at a time when he had the Court with him, to suppress it. We
have just seen Basil uniting with the Gallic Bishops against it.

\textbf{6. A.D. 359. The
Confession with a date}\\
(Third Sirmian)

The Semi-Arians, with the hope of striking a further blow at
their opponents by a judgment against the Anomœans, Soz. iv. 16
init., seem to have suggested a general Council, which ultimately
became the Councils of Seleucia and Ariminum. If this was their
measure, they were singularly out-manœuvred by the party of Acacius
and Valens, as may be seen in Athanasius's \emph{de Synodis}. A
preparatory Conference was held at Sirmium at the end of May in this
year, in which the Creed was determined which should be laid before
the great Councils then assembling. Basil and Mark were the chief
Semi-Arians present, and in the event became committed to an almost
Arian Confession. Soz. iv. 16, p. 562. It was finally settled on the
Eve of Pentecost, and the dispute lasted till morning. Epiph. Hær.
73, 22. Mark at length was chosen to draw it up, Soz. iv. 22, p. 573,
yet Valens so managed that Basil could not sign it without an
explanation. It was written in Latin, Socr. ii. 30, Soz. iv. 17, p.
563. Coustant, however, in Hil. p. 1152, note i., seems to consider
this dispute and Mark's confession to belong to the same date (May
22,) in the foregoing year; but p. 1363, note b, he seems to change
his opinion.

Petavius, who, Animadv. in Epiph. p. 318, follows Socrates in
considering that the second Sirmian is the Confession which the Arians
tried to suppress, nevertheless, de Trin. i. 9, § 8, yields to the
testimony of Athanasius in behalf of the third, attributing the
measure to their dissatisfaction with the phrase ``Like in all
things,'' which Constantius had inserted, and with Basil's
explanation on subscribing it, and to the hopes of publishing a bolder
creed which their increasing influence with Constantius inspired. He
does not think it impossible, however, that an attempt was made to
suppress both. Coustant, again, in Hil. p. 1363, note b, asks \emph{when}
it could be that the Eusebians attempted to suppress the second
Confession; and conjectures that the ridicule which followed their
dating of the third and their wish to get rid of the ``Like in all
things,'' were the causes of their anxiety about it. He
observes too with considerable speciousness that Acacius's Second
formulary at Seleucia (Athan. de Syn. 29), and the Confession of Nice
(Ibid. 30), resemble second editions of the third Sirmian. Valesius,
in Socr. ii. 30, and Montfaucon, in Athan. Syn. § 29, take the same
side.

Pagi in Ann. 357. n. 13, supposes that the third Sirmian was the
Creed signed by Liberius. Yet Coustant in Hil. p. 1335, note n,
speaking of Liberius's ``perfidia Ariana,'' as St. Hilary
calls it, says, ``Solus Valesius existimat tertiam [confessionem]
hic memorari:'' whereas Valesius, making four, not to say five,
Sirmian Creeds, understands Liberius to have signed, not the third,
but an intermediate one, between the second and third, as Petavius
does, in Soz. iv. 15 and 16. Moreover, Pagi fixes the date as A.D.
358 ibid.

This Creed, thus drawn up by a Semi-Arian, with an Acacian or Arian
Appendix, then a Semi-Arian insertion, and after all a Semi-Arian
protest on subscription, was proposed at Seleucia by Acacius, Soz. iv.
22, and at Ariminum by Valens, Socr. ii. 37, p. 132.

\textbf{7. A.D. 359. Nicene Editon of the Third
Sirmian}

The third Sirmian was rejected both at Seleucia and Ariminum; but
the Eusebians, dissolving the Council of Seleucia, kept the Fathers at
Ariminum together through the summer and autumn. Meanwhile at Nice in
Thrace they confirmed the third Sirmian, Socr. ii. 37, p. 141, Theod.
Hist. ii. 16, with the additional proscription of the word \emph{hypostasis};
apparently lest the Latins should by means of it evade the
condemnation of the ``consubstantial.'' This Creed, thus
altered, was ultimately accepted at Ariminum and was confirmed in
January 360 at Constantinople; Socr. ii. 41, p. 163. Soz. iv. 24 init.

Liberius retrieved his fault on this occasion; for, whatever was
the confession he had signed, he now refused his assent to the
Ariminian, and, if Socrates is to be trusted, was banished in
consequence, Socr. ii. 37, p. 140.

\textbf{Note 4. The Terms usia and hypostasis,
as used in the early Church}

(Vi\emph{de supra}, p. 186.)

1. \footnote{From the \emph{Atlantis}, July, 1858.} EVEN
before we take into account the effect which would naturally be
produced on the first Christians by the novelty and mysteriousness of
doctrines which depend for their reception simply upon Revelation, we
have reason to anticipate that there would be difficulties and
mistakes in expressing them, when they first came to be set forth by
unauthoritative writers. Even in secular sciences, inaccuracy of
thought and language is but gradually corrected; that is, in
proportion as their subject-matter is thoroughly scrutinized and
mastered by the co-operation of many independent intellects,
successively engaged upon it. Thus, for instance, the word \emph{Person}
requires the rejection of various popular senses, and a careful
definition, before it can serve for philosophical uses. We sometimes
use it for an \emph{individual} as contrasted with a class or
multitude, as when we speak of having ``personal objections''
to another; sometimes for the \emph{body}, in contrast to the soul, as
when we speak of ``beauty of person.'' We sometimes use it in
the abstract, as when we speak of another as ``insignificant in
person;'' sometimes in the concrete, as when we call him ``an
insignificant person.'' How divergent in meaning are the
derivatives, \emph{personable, personalities, personify, personation,
personage, parsonage}! This variety arises partly from our own
carelessness, partly from the necessary developments of language, partly from the exuberance of human thought, partly from the
defects of our vernacular tongue.

Language then requires to be refashioned even for sciences which
are based on the senses and the reason; but much more will this be the
case, when we are concerned with subject-matters, of which, in our
present state, we cannot possibly form any complete or consistent
conception, such as the Catholic doctrines of the Trinity and
Incarnation. Since they are from the nature of the case above our
intellectual reach, and were unknown till the preaching of
Christianity, they required on their first promulgation new words, or
words used in new senses, for their due enunciation; and, since these
were not definitely supplied by Scripture or by tradition, nor, for
centuries, by ecclesiastical authority, variety in the use, and
confusion in the apprehension of them, were unavoidable in the
interval. This conclusion is necessary, admitting the premisses,
antecedently to particular instances in proof.

Moreover, there is a presumption equally strong, that the variety
and confusion that I have anticipated, would in matter of fact issue
here or there in actual heterodoxy, as often as the language of
theologians was misunderstood by hearers or readers, and deductions
were made from it which the teacher did not intend. Thus, for
instance, the word \emph{Person}, used in the doctrine of the Holy
Trinity, would on first hearing suggest Tritheism to one who made the
word synonymous with \emph{individual}; and Unitarianism to another,
who accepted it in the classical sense of a mask or \emph{character}.

Even to this day our theological language is wanting in accuracy:
thus, we sometimes speak of the controversies concerning the \emph{Person}
of Christ, when we mean to include in them those also which belong to
the two \emph{natures} which are predicated of Him.

Indeed, the difficulties of forming a theological phraseology for the whole of Christendom were obviously so great, that we need not
wonder at the reluctance which the first age of Catholic divines
showed in attempting it, even apart from the obstacles caused by the
distraction and isolation of the churches in times of persecution. Not
only had the words to be adjusted and explained which were peculiar to
different schools or traditional in different places, but there was
the formidable necessity of creating a common measure between two, or
rather three languages,—Latin, Greek, and Syriac. The intellect had
to be satisfied, error had to be successfully excluded, parties the
most contrary to each other, and the most obstinate, had to be
convinced. The very confidence which would be felt by Christians in
general that Apostolic truth would never fail,—and that they held it
in each locality themselves and the \emph{orbis terrarum} with them,
in spite of all verbal contrarieties,—would indispose them to define
it, till definition became an imperative duty.

2. I think this plain from the nature of the case; and history
confirms me in the instance of the celebrated word \emph{homoüsion},
which, as one of the first and most necessary steps, so again was
apparently one of the most discouraging, in the attempt to give a
scientific expression to doctrine. This formula, as Athanasius,
Hilary, and Basil affirm, had been disowned, as savouring of
heterodoxy, by the great Council of Antioch in A.D.
264-269; yet, in spite of this disavowal on the part of Bishops of the
highest authority, it was imposed on all the faithful to the end of
time in the Ecumenical Council of Nicæa, A.D.
325, as the one and only safeguard, as it really is, of orthodox
teaching. The misapprehensions and protests which, after such
antecedents, its adoption occasioned for many years, may be easily
imagined. Though above three hundred Bishops had accepted it at Nicæa,
the great body of the Episcopate in the next generation considered it
inexpedient; and Athanasius himself, whose imperishable name is bound
up with it, showed himself most cautious in putting it forward,
though he knew it had the sanction of a General Council. Moreover, the
word does not occur in the \emph{Catecheses} of St. Cyril of Jerusalem, A.D.
347, nor in the recantation made before Pope Julius by Ursacius and
Valens, A.D. 349, nor in the cross-questionings
to which St. Ambrose subjected Palladius and Secundianus, A.D.
381. At Seleucia, A.D. 359, as many as 100
Eastern Bishops, besides the Arian party, were found to abandon it,
while at Ariminum in the same year the celebrated scene took place of
400 Bishops of the West being worried and tricked into a momentary act
of the same character. They had not yet got it deeply fixed into their
minds, as a sort of first principle, that to abandon the formula was
to betray the faith.

3. This disinclination on the part of Catholics to dogmatic
definitions was not confined to the instance of the \emph{homoüsion}.
In the use of the word \emph{hypostasis}, a variation was even allowed
by the authority of a Council [A.D. 362]; and
the circumstances under which it was allowed, and the possibility of
allowing it, without compromising Catholic truth, shall here be
considered.

As to the use of the word. At least in the West, and in St.
Athanasius's day, it was usual to speak of one \emph{hypostasis}, as
of one \emph{usia}, of the Divine Nature. Thus the so-called Sardican
Creed, A.D. 347, speaks of ``one \emph{hypostasis},
which the heretics call \emph{usia}.'' Theod. Hist. ii. 8; the
Roman Council under Damasus, A.D. 371, says that
the Three Persons are of the same \emph{hypostasis} and \emph{usia};
and the Nicene Anathema condemns those who say that the Son ``came
from other \emph{hypostasis} or \emph{usia}.'' Epiphanius too
speaks of ``one \emph{hypostasis},'' \emph{Hær}. 74, 4, \emph{Ancor}.
6 (and though he has the \emph{hypostases}, \emph{Hær}. 62, 3, 72, 1,
yet he is shy of the plural, and prefers ``the \emph{hypostatic}
Father; the \emph{hypostatic} Son,'' \&c., \emph{ibid}. 3 and
4, \emph{Ancor}. 6; and [\emph{tria}], as \emph{Hær}. 74, 4, where he
says ``three \emph{hypostatic} of the same \emph{hypostasis};''
vide also `` in \emph{hypostasis} of perfection,'' \emph{Hær}.
74, 12, \emph{Ancor}. 7 \emph{et alibi}); and Cyril of Jerusalem of
the ``uniform \emph{hypostasis}'' of God, \emph{Cateth}. vi.
7, vide also xvi. 12 and xvii. 9 (though the word may be construed one
out of three in \emph{Cat}. xi. 3); and Gregory Nazianzen, \emph{Orat}.
xxviii, 9, where he is speaking as a Natural, not as a Christian
theologian.

In the preceding century Gregory Thaumaturgus had laid it down that
the Father and the Son were in \emph{hypostasis} one, and the Council
of Antioch, A.D. 264-269, calls the Son in \emph{usia}
and \emph{hypostasis} God, the Son of God. Routh, \emph{Reliq}. t. 2,
p. 466. Accordingly Athanasius expressly tells us, ``\emph{Hypostasis}
is \emph{usia}, and means nothing else but [\emph{auto to on}],'' \emph{ad
Afros}, 4. Jerome says that ``Tota sæcularium litterarum
schola nihil aliud \emph{hypostasin} nisi \emph{usiam} novit,'' \emph{Epist}.
xv. 4; Basil, the Semi-Arian, that ``the Fathers have called \emph{hypostasis
usia},'' Epiph. \emph{Hær}, 73, 12, fin. And Socrates says
that at least it was frequently used for \emph{usia}, when it had
entered into the philosophical schools. \emph{Hist}. iii. 7.

On the other hand the Alexandrians, Origen (\emph{in Joan}. ii. 6 \emph{et
alibi}), Ammonius (\emph{ap. Caten. in Joan}. x. 30, if genuine),
Dionysius (\emph{ap}. Basil \emph{de Sp. S}. n. 72), and Alexander (\emph{ap}.
Theod. \emph{Hist}. i. 4), speak of more \emph{hypostases} than one in
the Divine Nature, that is, of Three; and apparently without the
support of the divines of any other school, unless Eusebius, who is
half an Alexandrian, be an exception. Going down beyond the middle of
the fourth century, we find the Alexandrian Didymus committing himself
to a bold and strong enunciation of the Three \emph{hypostases}, (e.g.
de Trin. 1. 18, \&c.), which is almost without a parallel in
patristical literature.

It was under these circumstances that the Council of Alexandria in A.D.
362, to which I have already referred, a Council in which Athanasius
and Eusebius of Vercellæ were the chief actors, determined to
leave the sense and use of the word open, so that, according to the
custom of their own church or school, Catholics might freely speak of
three \emph{hypostases} or of one.

Thus we are brought to the practice of Athanasius himself. It is
remarkable that he should so far innovate on the custom of his own
Church, as to use the word in each of these two applications of it. In
his \emph{In illud Omnia} he speaks of ``the three perfect \emph{Hypostases}.''
On the other hand, he makes \emph{usia} and \emph{hypostasis}
synonymous in \emph{Orat}. iii. 65, 66, \emph{Orat}. iv. 1 and 33 fin.

There is something more remarkable still in this innovation.
Alexander, his immediate predecessor and master, published, A.D.
320-324, two formal letters against Arius, one addressed to his
namesake of Constantinople, the other encyclical. It is scarcely
possible to doubt that the latter was written by Athanasius; it is so
unlike the former in style and diction, so like the writings of
Athanasius. Now it is observable that in the former the word \emph{hypostasis}
occurs in its Alexandrian sense at least five times; in the latter,
which I attribute to Athanasius, it is dropped, and \emph{usia} is
introduced, which is absent from the former. That is, Athanasius has,
on this supposition, when writing in his Bishop's name a formal
document, pointedly innovated on his Bishop's theological language,
and that the received language of his own Church. I am not supposing
he did this without Alexander's sanction. Indeed the character of the
Arian polemic would naturally lead Alexander, as well as Athanasius,
to be suspicious of their own formula of the ``Three \emph{Hypostases},''
which Arianism was using against them; and the latter would be
confirmed in this feeling by his subsequent familiarity with Latin
theology, and the usage of the Holy See, which, under Pope Damasus, as
we have seen, A.D. 371, spoke of one \emph{hypostasis},
and in the previous century, A.D. 260, protested
by anticipation in the person of Pope Dionysius against the use,
which might be made in the hands of enemies, of the formula of the
Three \emph{Hypostases}. Still it is undeniable that Athanasius does
at least once speak of Three, though his practice is to dispense with
the word and to use others instead of it.

4. Now then we come to the explanation of this difference of usage
in the application of the word. It is difficult to believe that so
accurate a thinker as Athanasius really used an important term in two
distinct, nay contrasted senses; and I cannot but question the fact,
so commonly taken for granted, that the divines of the beginning of
the fourth century had appropriated any word whatever definitely to
express either the idea of \emph{Person} as contrasted with that of \emph{Essence},
or of \emph{Essence} as contrasted with \emph{Person}. I altogether
doubt whether we are correct in saying that they meant by \emph{hypostasis},
in one country \emph{Person}, in another \emph{Essence}. I think such
propositions should be carefully proved, instead of being taken for
granted, as at present is the case. Meanwhile, I have an hypothesis of
my own. I think they used the word both in East and West in one and
the same substantial sense; with some accidental variation or latitude
indeed, but that of so slight a character, as would admit of
Athanasius, or any one else, speaking of one \emph{hypostasis} or
three, without any violence to that sense which remained on the whole
one and the same. What this sense is I proceed to explain:—

The school-men are known to have insisted with great earnestness on
the numerical unity of the Divine Being; each of the three Divine
Persons being one and the same God, unicus, singularis, et totus Deus.
In this, however, they did but follow the recorded doctrine of the
Western theologians of the fifth century, as I suppose will be allowed
by critics generally. So forcible is St. Austin upon the strict unity
of God, that he even thinks it necessary to caution his readers lest
they should suppose that he could allow them to speak of One
Person as well as of Three in the Divine nature \emph{de Trin}., vii.
11. Again, in the (so-called) Athanasian Creed, the same elementary
truth is emphatically insisted on. The neuter \emph{unum} of former
divines is changed into the masculine, in enunciating the mystery.
``Non tres æterni, sed unus æternus.'' I suppose this means,
that each Divine Person is to be received as the one God as entirely
and absolutely as He would be held to be, if we had never heard of the
other Two, and that He is not in any respect less than the one and
only God, because They are each that same one God also; or in other
words, that, as each human individual being has one personality, the
Divine Being has three.

Returning then to Athanasius, I consider that this same mystery is
implied in his twofold application of the word \emph{hypostasis}. The
polytheism and pantheism of the heathen world imagined,—not the God
whom natural reason can discover, conceive, and worship, one
individual, living, and personal,—but a \emph{divinitas}, which was
either a quality, whether energy or life, or an extended substance, or
something else equally inadequate to the real idea which the word
conveys. Such a divinity could not properly be called an \emph{hypostasis}
or said to be \emph{in hypostasi} (except indeed as brute matter may
be called, as in one sense it can be called, an \emph{hypostasis}),
and therefore it was, that that word had some fitness, especially
after the Apostle's adoption of it, \emph{Hebr}. i. 3, to denote the
Christian's God. And this may account for the remark of Socrates, that
it was a new word, strange to the schools of ancient philosophy, which
had seldom professed pure theism or natural theology. ``The
teachers of philosophy among the Greeks,'' he says, ``have
defined \emph{usia} in many ways: but of \emph{hypostasis}, they have
made no mention at all. Irenæus, the grammarian, affirms that the
word is barbarous.''—\emph{Hist}. iii. 7. The better then was it
fitted to express that highest object of thought, of which the
``barbarians'' of Palestine had been the special witnesses.
When the divine \emph{hypostasis} was confessed, the word
expressed or suggested the attributes of individuality,
self-subsistence, self-action, and personality, such as go to form the
idea of the Divine Being to the natural theologian; and, since the
difference between the theist and the Catholic divine in their idea of
His nature is simply this, that, in opposition to the Pantheist, who
cannot understand how the Infinite can be Personal at all, the one
ascribes to him one personality, and the other three, it will be
easily seen how a word, thus characterized and circumstanced, would
admit of being used with but a slight modification of its sense, of
the Trinity as well as of the Unity.

Let us take, by way of illustration, the word \emph{monad}, which
when applied to intellectual beings, includes the idea of personality.
Dionysius of Alexandria, for instance, speaks of the \emph{monad} and
the \emph{triad}: now, would it be very harsh, if, as he has spoken of
``three \emph{hypostases}'' in monad so he had instead spoken
of ``the three monads,'' that is, in the sense of ``thrice
hypostatic \emph{monad},'' as if the intrinsic force of the word \emph{monas}
would preclude the possibility of his use of the plural \emph{monads}
being mistaken to imply that he held more \emph{monads} than one? To
take an analogous case, it would be about the same improper use of
plural for singular, if we said that a martyr by his one act gained
three victories instead of a triple victory, over his three spiritual
foes. And indeed, though Athanasius does not directly speak of three
monads, yet he implies the possibility of such phraseology by teaching
that, though the Father and the Son are two, the \emph{monas} of the
Deity ([\emph{theot\={e}s}]) is indivisible, and that the Deity is
at once Father and Son.

This, then, is what I conceive that he means by sometimes speaking
of one, sometimes of three \emph{hypostases}. The word \emph{hypostasis}
stands neither for \emph{Person} nor for \emph{Essence} exclusively;
but it means the one Personal God of natural theology, the notion of
whom the Catholic corrects and completes as often as he views
him as a Trinity; of which correction Nazianzen's language (\emph{Orat}.
xxviii. 9) contrasted with his usual formula (\emph{vid. Orat}. xx. 6)
of the Three \emph{Hypostases}, is an illustration. The specification
of three \emph{hypostases} does not substantially alter the sense of
the word itself, but is a sort of \emph{catachresis} by which this
Catholic doctrine is forcibly brought out (as it would be by the
phrase ``three monads''), viz. that each of the Divine
Persons is simply the Unus et Singularis Deus. If it be objected, that
by the same mode of reasoning, Athanasius might have said \emph{catachrestically}
not only three \emph{monads} or three \emph{hypostases}, but three
Gods, I deny it, and for this reason, because \emph{hypostasis} is not
equivalent to the simple idea of God, but is rather a definition of
Him, and that in some special elementary points, as essence,
personality, \&c., and because such a mere improper use or varying
application of the term \emph{hypostasis} would not tend to compromise
a truth, which never must even in forms of speech be trifled with, the
absolute numerical unity of the Supreme Being. Though a Catholic could
not say that there are three Gods, he could say, that the \emph{definition}
of God applies to \emph{unus} and \emph{tres}. Perhaps it is for this
reason that Epiphanius speaks of the ``\emph{hypostatic}
Three,'' ``co-\emph{hypostatic},'' ``of the same \emph{hypostasis},''
\emph{Hær}. 74, 4 (\emph{vid}. Jerome, \emph{Ep}. 15, 3), in the
spirit in which St. Thomas, I think, interprets the ``non tres
æterni, sed unus æternus,'' to turn on the contrast of adjective
and substantive.

Petavius makes a remark which is apposite to my present purpose.
``Nomen Dei,'' he says, \emph{de Trin}. iii. 9. § 10, ``cùm
sit ex eorum genere quæ concreta dicuntur, forinam significat, non
abstractam ab individuis proprietatibus, sed in iis subsistentem. Est
enim Deus substantia aliqua divinitatem habens. Sicut homo non humanam
naturam separatam, sed in aliquo individuo subsistentem exponit, ita
tamen ut individuum ac personam, non certam ac determinatam, sed
confuse infiniteque representet, hoc est, \emph{naturam in aliquo}, ut
diximus, \emph{consistentem} ... sic nomen Dei propriè ac
directe divinitatem naturamque divinam indicat, \emph{assignificat autem
eundem, ut in quâpiam personâ subsistentem, nullam de tribus
expresse designans, sed confuse et universe}.'' Here this great
author seems to say, that even the word ``Deus'' may stand,
not barely for the Divine Being, but besides ``in quâpiam
personâ subsistentem,'' without denoting \emph{which} Person; and
in like manner I would understand \emph{hypostasis} to mean the \emph{monas}
with a like indeterminate notion of personality, (without which
attribute the idea of God cannot be,) and thus, according as one \emph{hypostasis}
is spoken of, or three, the word may be roughly translated, in the one
case ``personal substance,'' or ``being with
personality,'' in the other ``substantial person,'' or
``person which is in being.'' In all cases it will be
equivalent to the Deity, to the \emph{monad}, to the divine \emph{usia},
\&c., though with that peculiarity of meaning which I have insisted
on.

5. Since, as has been said above, \emph{hypostasis} is a word more
peculiarly Christian than \emph{usia}, I have judged it best to speak
of it first, that the meaning of it, as it has now been ascertained on
inquiry, may serve as a key for explaining other parallel terms. \emph{Usia}
is one of these the most in use, certainly in the works of Athanasius;
and we have his authority as well as St. Jerome's for stating that it
was once simply synonymous with \emph{hypostasis}. Moreover, in \emph{Orat}.
iii. 65, he uses the two words as equivalent to each other. If this be
so, what has been said above in explanation of the sense he put on the
word \emph{hypostasis}, will apply to \emph{usia} also. This
conclusion is corroborated by the proper meaning of the word \emph{usia}
itself which answers to the English word ``being.'' Now, when
we speak of the Divine Being, we mean to speak of Him, as what he is,
[\emph{ho \={o}n}], including generally His attributes and
characteristics, and among them, at least obscurely, His personality.
By the ``\emph{Divine Being}'' we do not commonly mean a mere \emph{anima
mundi}, or first principle of life or system of laws. \emph{Usia}
then, thus considered, agrees very nearly in sense, from its very
etymology, with \emph{hypostasis}. Further, this was the sense in
which Aristotle used it, viz. for what is ``individuum,'' and
``numero unum;'' and it must not be forgotten that the Neo-platonists,
who exerted so great an influence on the Alexandrian Church, professed
the Aristotelic logic. And so St. Cyril himself, the successor of
Athanasius (Suicer, \emph{Thes. in voce}, [\emph{ousia}].)

This is the word, and not \emph{hypostasis}, which Athanasius
commonly uses in controversy with the Arians, to express the divinity
of the Word. He speaks of the \emph{usia} of the Son as being united
to the Father, and His \emph{usia} being the offspring of the Father's
\emph{usia}. In these and other passages \emph{usia}, I conceive, is
substantially equivalent to \emph{hypostasis}, as I have explained it,
viz. expressing the divine [\emph{monas}] with an obscure intimation
of personality inclusively; and here I think I am able to quote the
words of Father Passaglia, as agreeing (so far) in what I have said.
``Quum \emph{hypostasis},'' he says, \emph{de Trinitate}, p.
1302, ``esse nequeat sine substantiâ, nihil vetabat quominus
trium hypostasum defensores \emph{hypostasim} interdum pro substantiâ
sumerent, præsertim ubi \emph{hypostasis} opponitur rei non
subsistenti ac efficientiæ.'' I should wish to complete the
admission by adding, ``Since an intellectual \emph{usia} naturally
implies an \emph{hypostasis}, there was nothing to hinder \emph{usia}
being used, when \emph{hypostasis} had to be expressed.''

6. After what I have said of \emph{usia} and \emph{hypostasis}, it
will not surprise the reader if I consider that [\emph{physis}]
(nature) also, in the Alexandrian theology, was equally capable of
being applied to the Divine Being viewed as One, or viewed as Three or
each of the Three separately. Thus Athanasius says, One is the Divine
Nature, (\emph{contr. Apoll}. ii. 13, \emph{fin. de Incarn. V}. fin.)
Alexander, on the other hand, calls the Father and Son the ``two \emph{hypostatic}
natures,'' and speaks of the ``only begotten nature,'' (Theod.
Hist. i. 4,) and Clement of ``the Son's nature'' as ``most
intimately near the sole Almighty,'' (Strom. vii. 2,) and
Cyril of a ``generating nature'' and a ``generated'' (Thes.
xi. p. 85) and, in words celebrated in theological history, of
``the Word's One Nature incarnate.''

7. [\emph{Eidos}] is a word of a similar character. As it is found
in \emph{John} v. 37, it may be indifferently interpreted of essence
or of person; the Vulgate translates it ``neque \emph{speciem}
ejus vidistis.'' In Athan. \emph{Orat}. iii. 3, it is synonymous
with \emph{deity} or \emph{usia}; as \emph{ibid}. 6 also; and
apparently in \emph{ibid}. 16, where the Son is said to have the \emph{species}
of the Father. And so in \emph{de Syn}. 52. Athanasius says that there
is only one ``species deitatis.'' Yet, as taken from \emph{Gen}.
xxxii. 31, it is considered to denote the Son; \emph{e.g}. Athan. \emph{Orat}.
i. 20, where it is used as synonymous with Image, [\emph{eikon}]. In
like manner the Son is called ``the very species deitatis.'' \emph{Ep.
Æg}. 17. But again in Athan. \emph{Orat}. iii. 6, it is first said
that the \emph{species} of the Father and Son are one and the same,
then that the Son is the \emph{species} of the Father's (deity), and
then that the Son is the \emph{species} of the Father.

The outcome of this investigation is this:—that we need not by an
officious piety arbitrarily force the language of separate Fathers
into a sense which it cannot bear; nor by an unjust and narrow
criticism accuse them of error; nor impose upon an early age a
distinction of terms belonging to a later. The words \emph{usia} and \emph{hypostasis}
were, naturally and intelligibly, for three or four centuries,
practically synonymous, and were used indiscriminately for two ideas,
which were afterwards respectively denoted by the one and the other.

\textbf{Note 5. The Orthodoxy of the Body of the
Faithful during the Supremacy of Arianism}

(\emph{Vide supra}, p. 358.)
\footnote{From the \emph{Rambler}, July, 1859.}

THE episcopate, whose action was so
prompt and concordant at Nicæa on the rise of Arianism, did not, as a
class or order of men, play a good part in the troubles consequent
upon the Council; and the laity did. The Catholic people, in the
length and breadth of Christendom, were the obstinate champions of
Catholic truth, and the bishops were not. Of course there were great
and illustrious exceptions; first, Athanasius, Hilary, the Latin
Eusebius, and Phœbadius; and after them, Basil, the two Gregories,
and Ambrose; there are others, too, who suffered, if they did nothing
else, as Eustathius, Paulus, Paulinus, and Dionysius; and the Egyptian
bishops, whose weight was small in the Church in proportion to the
great power of their Patriarch. And, on the other hand, as I shall say
presently, there were exceptions to the Christian heroism of the
laity, especially in some of the great towns. And again, in speaking
of the laity, I speak inclusively of their parish-priests (so to call
them), at least in many places; but on the whole, taking a wide view
of the history, we are obliged to say that the governing body of the
Church came short, and the governed were pre-eminent in faith, zeal,
courage, and constancy.

This is a very remarkable fact: but there is a moral in it. Perhaps
it was permitted, in order to impress upon the Church at that very
time passing out of her state of persecution to her long
temporal ascendancy, the great evangelical lesson, that, not the wise
and powerful, but the obscure, the unlearned, and the weak constitute
her real strength. It was mainly by the faithful people that Paganism
was overthrown; it was by the faithful people, under the lead of
Athanasius and the Egyptian bishops, and in some places supported by
their Bishops or priests, that the worst of heresies was withstood and
stamped out of the sacred territory.

The contrast stands as follows:—

1.

1. A.D. 325. The great Council of Nicæa of
318 Bishops, chiefly from the eastern provinces of Christendom, under
the presidency of Hosius of Cordova. It was convoked against Arianism,
which it once for all anathematized; and it inserted the formula of
the ``Consubstantial'' into the Creed, with the view of
establishing the fundamental dogma which Arianism impugned. It is the
first Œcumenical Council, and recognized at the time its own
authority as the voice of the infallible Church. It is so received by
the \emph{orbis terrarum} at this day.

2. A.D. 326. St. Athanasius, the great
champion of the Homoüsion, was elected Bishop of Alexandria.

3. A.D. 334, 335. The Synods of Cæsarea and
Tyre (sixty Bishops) against Athanasius, who was therein accused and
formally condemned of rebellion, sedition, and ecclesiastical tyranny;
of murder, sacrilege, and magic; deposed from his See, forbidden to
set foot in Alexandria for life, and banished to Gaul. Also, they
received Arius into communion.

4. A.D. 341. Council of Rome of fifty
Bishops, attended by the exiles from Thrace, Syria, \&c., by
Athanasius, \&c., in which Athanasius was pronounced innocent.

5. A.D. 341. Great Council of the Dedication
at Antioch, attended by ninety or a hundred Bishops. The council
ratified the proceedings of the Councils of Cæsarea and Tyre, and placed an Arian in the See of Athanasius. Then it proceeded to
pass a dogmatic decree in reversal of the formula of the
``Consubstantial.'' Four or five creeds, instead of the
Nicene, were successively adopted by the assembled Fathers.

Three of these were circulated in the neighbourhood; but as they
wished to send one to Rome, they directed a fourth to be drawn up.
This, too, apparently failed.

6. A.D. 345. Council of the creed called
Macrostich. This Creed suppressed, as did the third, the word
``substance.'' The eastern Bishops sent this to the Bishops of
France, who rejected it.

7. A.D. 347. The great Council of Sardica,
attended by more than 300 Bishops. Before it commenced, a division
between its members broke out on the question whether or not
Athanasius should have a seat in it. In consequence, seventy-six
retired to Philippopolis, on the Thracian side of Mount Hæmus, and
there excommunicated the Pope and the Sardican Fathers. These seceders
published a sixth confession of faith. The Synod of Sardica, including
Bishops from Italy, Gaul, Africa, Egypt, Cyprus, and Palestine,
confirmed the act of the Roman Council, and restored Athanasius and
the other exiles to their Sees. The Synod of Philippopolis, on the
contrary, sent letters to the civil magistrates of those cities,
forbidding them to admit the exiles into them. The Imperial power took
part with the Sardican Fathers, and Athanasius went back to
Alexandria.

8. A.D. 351. The Bishops of the East met at
Sirmium. The semi-Arian Bishops began to detach themselves from the
Arians, and to form a separate party. Under pretence of putting down a
kind of Sabellianism, they drew up a new creed, into which they
introduced the language of some of the ante-Nicene writers on the
subject of our Lord's divinity, and dropped the word
``substance.''

9. A.D. 353. The Council of Arles. The Pope
sent to it several Bishops as legates. The Fathers of the
Council, including the Pope's legate, Vincent, subscribed the
condemnation of Athanasius. Paulinus, Bishop of Treves, was nearly the
only one who stood up for the Nicene faith and for Athanasius. He was
accordingly banished into Phrygia, where he died.

10. A.D. 355. The Council of Milan, of more
than 300 Bishops of the West. Nearly all of them subscribed the
condemnation of Athanasius; whether they generally subscribed the
heretical creed, which was brought forward, does not appear. The Pope's
four legates remained firm, and St. Dionysius of Milan, who died an
exile in Asia Minor. An Arian was put into his See. Saturninus, the
Bishop of Arles, proceeded to hold a Council at Beziers; and its
Fathers banished St. Hilary to Phrygia.

11. A.D. 357-9. The Arians and Semi-Arians
successively drew up fresh creeds at Sirmium.

12. A.D. 357-8. Hosius' fall. ``Constantius
used such violence towards the old man, and confined him so straitly,
that at last, broken by suffering, he was brought, though hardly, to
hold communion with Valens and Ursacius [the Arian leaders], though he
would not subscribe against Athanasius.'' Athan. \emph{Arian. Hist}.
45.

13. A.D. 357-8. And Liberius. ``The
tragedy was not ended in the lapse of Hosius, but in the evil which
befell Liberius, the Roman Pontiff, it became far more dreadful and
mournful, considering that he was Bishop of so great a city, and of
the whole Catholic Church, and that he had so bravely resisted
Constantius two years previously. There is nothing, whether in the
historians and holy fathers, or in his own letters, to prevent our
coming to the conclusion, that Liberius communicated with the Arians,
and confirmed the sentence passed by them against Athanasius; but he
is not at all on that account to be called a heretic.'' Baron.
Ann. 357, 38-45. Athanasius says: ``Liberius, after he had been
in banishment for two years, gave way, and from fear of
threatened death was induced to subscribe. \emph{Arian. Hist}. § 41.
St. Jerome says: ``Liberius, tædio victus exilii, et in
hæreticam pravitatem subscribens, Romam quasi victor intraverat.''
\emph{Chron}. ed. Val. p. 797.

14. A.D. 359. The great Councils of Seleucia
and Ariminum, being one bi-partite Council, representing the East and
West respectively. At Seleucia there were 150 Bishops, of which only
the twelve or thirteen from Egypt were champions of the Nicene
``Consubstantial.'' At Ariminum there were as many as 400
Bishops, who, worn out by the artifice of long delay on the part of
the Arians, abandoned the ``Consubstantial,'' and subscribed
the ambiguous formula which the heretics had substituted for it.

15. About A.D. 360, St. Hilary says: ``I
am not speaking of things foreign to my knowledge; I am not writing
about what I am ignorant of; I have heard and I have seen the
shortcomings of persons who are round about me, not of laymen, but of
Bishops. For, excepting the Bishop Eleusius and a few with him, for
the most part the ten Asian provinces, within whose boundaries I am
situate, are truly ignorant of God.'' \emph{De Syn}. 63. It is
observable, that even Eleusius, who is here spoken of as somewhat
better than the rest, was a Semi-Arian, according to Socrates, and
even a persecutor of Catholics at Constantinople; and, according to
Sozomen, one of those who were active in causing Pope Liberius to give
up the Nicene formula of the ``Consubstantial.'' By the ten
Asian provinces is meant the east and south provinces of Asia Minor,
pretty nearly as cut off by a line passing from Cyzicus to Seleucia
through Synnada.

16. A.D. 360. St. Gregory Nazianzen says,
about this date: ``Surely the pastors have done foolishly; for,
excepting a very few, who either on account of their insignificance
were passed over, or who by reason of their virtue resisted, and who
were to be left as a seed and root for the springing up again
and revival of Israel by the influences of the Spirit, all temporized,
only differing from each other in this, that some succumbed earlier,
and others later; some were foremost champions and leaders in the
impiety, and others joined the second rank of the battle, being
overcome by fear, or by interest, or by flattery, or, what was the
most excusable, by their own ignorance.'' \emph{Orat}. xxi. 24.

17. A.D. 361. About this time, St. Jerome
says: ``Nearly all the churches in the whole world, under the
pretence of peace and of the emperor, are polluted with the communion
of the Arians.'' \emph{Chron}. Of the same date, that is, upon the
Council of Ariminum, are his famous words, ``Ingemuit totus orbis
et se esse Arianum miratus est.'' \emph{In Lucif}. 19. ``The
Catholics of Christendom were strangely surprised to find that the
Council had made Arians of them.''

18. A.D. 362. State of the Church of Antioch
at this time. There were four Bishops or communions of Antioch; first,
the old succession and communion, which had possession before the
Arian troubles; secondly, the Arian succession, which had lately
conformed to orthodoxy in the person of Meletius; thirdly, the new
Latin succession, lately created by Lucifer, whom some have thought
the Pope's legate there; and, fourthly, the new Arian succession,
which was started upon the recantation of Meletius. At length, as
Arianism was brought under, the evil reduced itself to two Episcopal
Successions, that of Meletius and the Latin, which went on for many
years, the West and Egypt holding communion with the latter, and the
East with the former.

19. St. Hilary speaks of the series of ecclesiastical Councils of
that time in the following well-known passage: ``Since the Nicene
Council, we have done nothing but write the Creed. While we fight
about words, inquire about novelties, take advantage of ambiguities,
criticize authors, fight on party questions, have difficulties in
agreeing, and prepare to anathematize each other, there is scarce a
man who belongs to Christ. Take, for instance, last year's
Creed, what alteration is there not in it already? First, we have the
Creed, which bids us not to use the Nicene 'consubstantial;' then
comes another, which decrees and preaches it; next, the third, excuses
the word 'substance,' as adopted by the Fathers in their simplicity;
lastly, the fourth, which instead of excusing, condemns. We determine
creeds by the year or by the month, we change our own determinations,
we prohibit our changes, we anathematize our prohibitions. Thus, we
either condemn others in our own persons, or ourselves in the instance
of others, and while we bite and devour one another, are like to be
consumed one of another.'' \emph{Ad Const}. ii. 4, 5.

20. A.D. 382. St. Gregory writes: ``If I
must speak the truth, I feel disposed to shun every conference of
Bishops: for never saw I Synod brought to a happy issue, and
remedying, and not rather aggravating, existing evils. For rivalry and
ambition are stronger than reason,—do not think me extravagant for
saying so,—and a mediator is more likely to incur some imputation
himself than to clear up the imputations which others lie under.''—\emph{Ep}.
129.

2.

Coming to the opposite side of the contrast, I observe that there
were great efforts made on the part of the Arians to render their
heresy popular. Arius himself, according to the Arian Philostorgius,
``wrote songs for the sea, and for the mill, and for the road, and
then set them to suitable music.'' \footnote{The translations which follow are for the
most part from Bohn's and the Oxford editions, the passages being
abridged.} Hist. ii. 2. Alexander speaks of the ``running about''
of the Arian women, Theod. Hist. i. 4, and of the buffoonery of their
men. Socrates says that ``in the Imperial court, the officers of
the bed-chamber held disputes with the women, and in the city,
in every house, there was a war of dialectics,'' ii. 2. Especially
at Constantinople there were, as Gregory says, ``of Jezebels as
thick a crop as of hemlock in a field,'' Orat. 35, 3; and he
himself suffered from the popular violence there. At Alexandria the
Arian women are described by Athanasius as ``running up and down
like Bacchanals and furies,'' and as ``passing that day in
grief on which they could do no harm.'' \emph{Hist. Arian}. 59.

The controversy was introduced in ridicule into the heathen
theatres, Euseb. v. Const. ii. 6. Socr. i. 6. ``Men of
yesterday,'' says Gregory Nyssen, ``mere mechanics, offhand
dogmatists in theology, servants too and slaves that have been
scourged, run-aways from servile work, and philosophical about things
incomprehensible. Of such the city is full; its entrances, forums,
squares, thoroughfares; the clothes-vendors, the money-lenders, the
victuallers. Ask about pence, and they will discuss the generate and
ingenerate,'' \&c., \&c., tom. ii. p. 898. Socrates, too,
says that the heresy ``ravaged provinces and cities;'' and
Theodoret that, ``quarrels took place in every city and village
concerning the divine dogma, the people looking on, and taking
sides.'' \emph{Hist}. i. 6.

In spite of these attempts, however, on the part of the Arians,
still, viewing Christendom as a whole, we shall find that the Catholic
populations sided with Athanasius; and the fierce disputes above
described evidenced the zeal of the orthodox rather than the strength
of the heretical party. This will appear in the following extracts:—

1. ALEXANDRIA. ``We suppose,'' says
Athanasius, ``you are not ignorant what outrages they [the Arian
Bishops] committed at Alexandria, for they are reported every where.
They attacked the holy virgins and brethren with naked swords; they
beat with scourges their persons, esteemed honourable in God's sight,
so that their feet were lamed by the stripes, whose souls were whole
and sound in purity and all good works.'' Athan. \emph{Ap. c. Arian}.
15.

``Accordingly Constantius writes letters, and commences a
persecution against all. Gathering together a multitude of herdsmen
and shepherds, and dissolute youths belonging to the town, armed with
swords and clubs, they attacked in a body the Church of Quirinus: and
some they slew, some they trampled under foot, others they beat with
stripes and cast into prison or banished. They haled away many women
also, and dragged them openly into the court, and insulted them,
dragging them by the hair. Some they proscribed; from some they took
away their bread, for no other reason but that they might be induced
to join the Arians, and receive Gregory [the Arian Bishop], who had
been sent by the Emperor.'' Athan. \emph{Hist. Arian}. § 10.

``On the week that succeeded the holy Pentecost, when the
people after their fast, had gone out to the cemetery to pray, because
that all refused communion with George [the Arian Bishop], the
commander, Sebastian, straightway with a multitude of soldiers
proceeded to attack the people, though it was the Lord's day; and
finding a few praying (for the greater part had already retired on
account of the lateness of the hour), having lighted a pile, he placed
certain virgins near the fire, and endeavoured to force them to say
that they were of the Arian faith. And having seized on forty men, he
cut some fresh twigs of the palm-tree, with the thorns upon them, and
scourged them on the back so severely that some of them were for a
long time under medical treatment, on account of the thorns which had
entered their flesh, and others, unable to bear up under their
sufferings, died. All those whom they had taken, both the men and the
virgins, they sent away into banishment to the great Oasis. Moreover,
they immediately banished out of Egypt and Libya the following Bishops
[sixteen], and the presbyters, Hierax and Dioscorus; some of them died
on the way, others in the place of their banishment. They caused also
more than thirty Bishops to take to flight.'' Apol. \emph{de Fug}.
7.

2. EGYPT. ``The Emperor Valens having
issued an edict commanding that the orthodox should be expelled both
from Alexandria and the rest of Egypt, depopulation and ruin to an
immense extent immediately followed; some were dragged before the
tribunals, others cast into prison, and many tortured in various ways;
all sorts of punishment being inflicted upon persons who aimed only at
peace and quiet.'' Socr. \emph{Hist}. iv. 24.

3. THE MONKS (1.) \emph{of
Egypt}. ``Antony left the solitude of the desert to go about
every part of the city [Alexandria], warning the inhabitants that the
Arians were opposing the truth, and that the doctrines of the Apostles
were preached only by Athanasius.'' Theod. \emph{Hist}. iv. 27.

``Lucius, the Arian, with a considerable body of troops,
proceeded to the monasteries of Egypt, where he in person assailed the
assemblage of holy men with greater fury than the ruthless soldiery.
When these excellent persons remained unmoved by all the violence, in
despair he advised the military chief to send the fathers of the
monks, the Egyptian Macarius and his namesake of Alexandria, into
exile.'' Socr. iv. 24.

(2.) \emph{Of Constantinople}. ``Isaac, on seeing the emperor
depart at the head of his army, exclaimed, 'You who have declared war
against God cannot gain His aid. Cease from fighting against Him, and
He will terminate the war. Restore the pastors to their flocks, and
then you will obtain a bloodless victory.''' Theod. iv.

(3.) \emph{Of Syria}, \&c. ``That these heretical doctrines
[Apollinarian and Eunomian] did not finally become predominant is
mainly to be attributed to the zeal of the monks of this period; for
all the monks of Syria, Cappadocia, and the neighbouring provinces
were sincerely attached to the Nicene faith. The same fate awaited
them which had been experienced by the Arians; for they incurred the
full weight of the popular odium and aversion, when it was
observed that their sentiments were regarded with suspicion by the
monks.'' Sozom. vi. 27.

(4.) \emph{Of Capadocia}. ``Gregory, the father of Gregory
Theologus, otherwise a most excellent man, and a zealous defender of
the true and Catholic religion, not being on his guard against the
artifices of the Arians, such was his simplicity, received with
kindness certain men who were contaminated with the poison, and
subscribed an impious proposition of theirs. This moved the monks to
such indignation, that they withdrew forthwith from his communion, and
took with them, after their example, a considerable part of his
flock.'' Ed. Bened. Monit. \emph{in Greg. Naz. Orat}. 6.

4. ANTIOCH. ``Whereas he (the Bishop
Leontius) took part in the blasphemy of Arius, he made a point of
concealing this disease, partly for fear of the multitude, partly for
the menaces of Constantius; so those who followed the Apostolical
dogmas gained from him neither patronage nor ordination, but those who
held Arianism were allowed the fullest liberty of speech, and were
placed in the ranks of the sacred ministry. But Flavian and Diodorus,
who had embraced the ascetical life, and maintained the Apostolical
dogmas, openly withstood Leontius's machinations against religious
doctrine. They threatened that they would retire from the communion of
his Church, and would go to the West, and reveal his intrigues. Though
they were not as yet in the sacred ministry, but were in the ranks of
the laity, night and day they used to excite all the people to zeal
for religion. They were the first to divide the singers into two
choirs, and to teach them to sing in alternate parts the strains of
David. They too, assembling the devout at the shrines of the martyrs,
passed the whole night there in hymns to God. These things Leontius
seeing, did not think it safe to hinder them, for he saw that the
multitude was especially well affected towards those excellent
persons. Nothing, however, could persuade Leontius to correct
his wickedness. It follows, that among the clergy were many who were
infected with the heresy: but the mass of the people were champions of
orthodoxy.'' Theodor. \emph{Hist}. ii. 24.

5. EDESSA. ``There is in that city a
magnificent church, dedicated to St. Thomas the Apostle, wherein, on
account of the sanctity of the place, religious assemblies are
continually held. The Emperor Valens wished to inspect this edifice;
when, having learned that all who usually congregated there were
enemies to the heresy which he favoured, he is said to have struck the
prefect with his own hand, because he had neglected to expel them
thence. The prefect, to prevent the slaughter of so great a number of
persons, privately warned them against resorting thither. But his
admonitions and menaces were alike unheeded; for on the following day
they all crowded to the church. When the prefect was going towards it
with a large military force, a poor woman leading her own little child
by the hand, hurried hastily by on her way to the church, breaking
through the ranks of the soldiery. The prefect, irritated at this,
ordered her to be brought to him, and thus addressed her: 'Wretched
woman, whither are you running in so disorderly a manner?' She
replied, 'To the same place that others are hastening.' 'Have you not
heard,' said he, 'that the prefect is about to put to death all that
shall be found there?' 'Yes,' said the woman, 'and therefore I hasten,
that I may be found there.' 'And whither are you dragging that little
child?' said the prefect. The woman answered, 'That he also may be
vouchsafed the honour of martyrdom.' The prefect went back and
informed the Emperor that all were ready to die in behalf of their own
faith; and added that it would be preposterous to destroy so many
persons at one time, and thus succeeded in restraining the Emperor's
wrath.'' Socr. iv. 18. ``Thus was the Christian faith
confessed by the whole city of Edessa.'' Sozom. vi. 18.

6. SAMOSATA. ``The Arians, having
deprived this exemplary flock of their shepherd, elected in his place
an individual with whom none of the inhabitants of the city, whether
poor or rich, servants or mechanics, husbandmen or gardeners, men or
women, young or old, would hold communion. He was left quite alone; no
one even calling to see him, or exchanging a word with him. It is,
however, said that his disposition was extremely gentle; and this is
proved by what I am about to relate. One day, when he went to bathe in
the public baths, the attendants closed the doors; but he ordered the
doors to be thrown open, that the people might be admitted to bathe
with himself. Perceiving that they remained in a standing posture
before him, imagining that great deference towards himself was the
cause of this conduct, he arose and left the bath. These people
believed that the water had been contaminated by his heresy, and
ordered it to be let out, and fresh water to be supplied. When he
heard of this circumstance, he left the city, thinking that he ought
no longer to remain in a place where he was the object of public
aversion and hatred. Upon this retirement of Eunomius, Lucius was
elected as his successor by the Arians. Some young persons were
amusing themselves with playing at ball in the market-place; Lucius
was passing by at the time, and the ball happened to fall beneath the
feet of the ass on which he was mounted. The youths uttered loud
exclamations, believing that the ball was contaminated. They lighted a
fire, and hurled the ball through it, believing that by this process
the ball would be purified. Although this was only a childish deed,
and although it exhibits the remains of ancient superstition, yet it
is sufficient to show the odium which the Arian faction had incurred
in this city. Lucius was far from imitating the mildness of Eunomius,
and he persuaded the heads of the government to exile most of the
clergy.'' Theodor. iv. i5.

7. OSRHOENE. ``Arianism met with similar
opposition at the same period in Osrhoëne and Cappadocia.
Basil, Bishop of Cæsarea, and Gregory, Bishop of Nazianzus, were held
in high admiration and esteem throughout these regions.'' Sozom.
vi. 21.

8. CAPPADOCIA. ``Valens, in passing
through Cappadocia, did all in his power to injure the orthodox, and
to deliver up the churches to the Arians. He thought to accomplish his
designs more easily on account of a dispute which was then pending
between Basil and Eusebius, who governed the Church of Cæsarea. This
dissension had been the cause of Basil's departing to Pontus. The
people, and some of the most powerful and wisest men of the city,
began to regard Eusebius with suspicion, and to meditate a secession
from his communion. The emperor and the Arian Bishops regarded the
absence of Basil and the hatred of the people towards Eusebius, as
circumstances that would tend greatly to the success of their designs.
But their expectations were utterly frustrated. On the first
intelligence of the intention of the emperor to pass through
Cappadocia, Basil returned to Cæsarea, where he effected a
reconciliation with Eusebius. The projects of Valens were thus
defeated, and he returned with his Bishops.'' Sozom. vi. 15.

9. PONTUS. ``It is said that when
Eulalius, Bishop of Amasia in Pontus, returned from exile, he found
that his Church had passed into the hands of an Arian, and that
scarcely fifty inhabitants of the city had submitted to the control of
their new bishop.'' Sozom. vii. 2.

10. ARMENIA. ``That company of Arians,
who came with Eustathius to Nicopolis, had promised that they would
bring over this city to compliance with the commands of the Imperial
vicar. This city had great ecclesiastical importance, both because it
was the metropolis of Armenia, and because it had been ennobled by the
blood of martyrs, and governed hitherto by Bishops of great
reputation, and thus, as Basil calls it, was the nurse of religion and
the metropolis of sound doctrine. Fronto, one of the city
presbyters, who had hitherto shown himself as a champion of the truth,
through ambition gave himself up to the enemies of Christ, and
purchased the bishoprick of the Arians at the price of renouncing the
Catholic faith. This wicked proceeding of Eustathius and the Arians
brought a new glory instead of evil to the Nicopolitans, since it gave
them an opportunity of defending the faith. Fronto, indeed, the Arians
consecrated, but there was a remarkable unanimity of clergy and people
in rejecting him. Scarcely one or two clerks sided with him; on the
contrary, he became the execration of all Armenia.'' \emph{Vita S.
Basil}., Bened. pp. clvii, clviii.

11. NICOMEDIA. ``Eighty pious clergy
proceeded to Nicomedia, and there presented to the emperor a
supplicatory petition complaining of the ill-usage to which they had
been subjected. Valens, dissembling his displeasure in their presence,
gave Modestus, the prefect, a secret order to apprehend these persons
and to put them to death. The prefect, fearing he should excite the
populace to a seditious movement against himself, if he attempted the
public execution of so many, pretended to send them away into
exile,'' \&c. Socr. iv. 16.

12. CAPPADOCIA. St. Basil says, about the
year 372: ``Religious people keep silence, but every blaspheming
tongue is let loose. Sacred things are profaned; those of the laity
who are sound in faith avoid the places of worship as schools of
impiety, and raise their hands in solitudes, with groans and tears to
the Lord in heaven.'' \emph{Ep}. 92. Four years after he writes:
``Matters have come to this pass: the people have left their
houses of prayer, and assemble in deserts,—a pitiable sight; women
and children, old men, and men otherwise infirm, wretchedly faring in
the open air, amid the most profuse rains and snow-storms and winds
and frosts of winter; and again in summer under a scorching sun. To
this they submit, because they will have no part in the wicked
Arian leaven.'' \emph{Ep}. 242. Again: ``Only one offence is
now vigorously punished,—an accurate observance of our fathers'
traditions. For this cause the pious are driven from their countries,
and transported into deserts. The people are in lamentation, in
continual tears at home and abroad. There is a cry in the city, a cry
in the country, in the roads, in the deserts. Joy and spiritual
cheerfulness are no more; our feasts are turned into mourning; our
houses of prayer are shut up, our altars deprived of the spiritual
worship.'' \emph{Ep}. 243.

13. PAPHLAGONIA, \&c. ``I
thought,'' says Julian in one of his Epistles, ``that the
leaders of the Galilæans would feel more grateful to me than to my
predecessor. For in his time they were in great numbers turned out of
their homes, and persecuted, and imprisoned; moreover, multitudes of
so-called heretics'' [the Novatians who were with the Catholics
against the Arians] ``were slaughtered, so that in Samosata,
Paphlagonia, Bithynia, and Galatia, and many other nations, villages
were utterly sacked and destroyed.'' \emph{Ep}. 52.

14. SCYTHIA. ``There are in this country
a great number of cities, of towns, and of fortresses. According to an
ancient custom which still prevails, all the churches of the whole
country are under the sway of one Bishop. Valens [the emperor]
repaired to the Church, and strove to gain over the Bishop to the
heresy of Arius; but this latter manfully opposed his arguments, and
after a courageous defence of the Nicene doctrines, quitted the
emperor, and proceeded to another church, whither he was followed by
the people. Valens was extremely offended at being left alone in a
church with his attendants, and in resentment condemned Vetranio [the
Bishop] to banishment. Not long after, however, he recalled him,
because, I believe, he apprehended insurrection.'' Sozom. vi. 21.

15. CONSTANTINOPLE, ``Those who
acknowledged the doctrine of consubstantiality were not only expelled
from the churches, but also from the cities. But although expulsion at
first satisfied them [the Arians], they soon proceeded to the worse
extremity of inducing compulsory communion with them, caring little
for such a desecration of the churches. They resorted to all kinds of
scourgings, a variety of tortures, and confiscation of property. Many
were punished with exile, some died under the torture, and others were
put to death while being driven from their country. These atrocities
were exercised throughout all the eastern cities, but especially at
Constantinople.'' Socr. ii. 27.

16. ILLYRIA. ``The parents of Theodosius
were Christians and were attached to the Nicene doctrine, hence he
took pleasure in the ministration of Ascholius [Bishop of
Thessalonica]. He also rejoiced at finding that the Arian heresy had
not been received in Illyria.'' Sozom, vii. 4.

17. NEIGHBOURHOOD OF MACEDONIA.
``Theodosius inquired concerning the religious sentiments which
were prevalent in the other provinces, and ascertained that, as far as
Macedonia, one form of belief was universally predominant,''
\&c. Ibid.

18. ROME. ``With respect to the doctrine
no dissension arose either at Rome or in any other of the Western
Churches; the people unanimously adhered to the form of belief
established at Nicæa.'' Sozom. vi. 23.

``Liberius, returning to Rome, found the mind of the mass of
men alienated from him, because he had so shamefully yielded to
Constantius. And thus it came to pass, that those persons who had
hitherto kept aloof from Felix [the rival Pope], and had avoided his
communion in favour of Liberius, on hearing what had happened, left
him for Felix, who raised the Catholic standard.'' Baron. \emph{Ann}.
357. 56. He tells us besides (57), that the people would not even go
to the public baths, lest they should bathe with the party of Liberius.

19. MILAN. At the Council of Milan, Eusebius
of Vercellæ, when it was proposed to draw up a declaration against
Athanasius, ``said that the Council ought first to be sure of the
faith of the Bishops attending it, for he had found out that some of
them were polluted with heresy. Accordingly he brought before the
Fathers the Nicene Creed, and said he was willing to comply with all
their demands, after they had subscribed that confession. Dionysius,
Bishop of Milan, at once took up the paper and began to write his
assent; but Valens [the Arian] violently pulled pen and paper out of
his hands, crying out that such a course of proceeding was impossible.
Whereupon, after much tumult, the question came before the people, and
great was the distress of all of them; the faith of the Church was
attacked by the Bishops. They then, dreading the judgment of the
people, transfer their meeting from the church to the Imperial
palace.'' Hilar. \emph{ad Const}. i. 8.

Again: ``As the feast of Easter approached, the empress sent to
St. Ambrose to ask a church of him, where the Arians who attended her
might meet together. He replied, that a Bishop could not give up the
temple of God. The pretorian prefect came into the church, where St.
Ambrose was attended by the people, and endeavoured to persuade him to
yield up at least the Portian Basilica. The people were clamorous
against the proposal; and the prefect retired to report how matters
stood to the emperor. The Sunday following St. Ambrose was explaining
the creed, when he was informed that the officers were hanging up the
Imperial hangings in the Portian Basilica, and that upon this news the
people were repairing thither. While he was offering up the holy
sacrifice, a second message came that the people had seized an Arian
priest as he was passing through the street. He despatched a number of
his clergy to the spot to rescue the Arian from his danger. The court
looked on this resistance of the people as seditious, and immediately
laid considerable fines upon the whole body of the tradesmen of
the city. Several were thrown into prison. In three days' time these
tradesmen were fined two hundred pounds weight of gold, and they said
that they were ready to give as much again on condition that they
might retain their faith. The prisons were filled with tradesmen; all
the officers of the household, secretaries, agents of the emperor, and
dependent officers who served under various counts, were kept within
doors, and were forbidden to appear in public, under pretence that
they should bear no part in sedition. Men of higher rank were menaced
with severe consequences, unless the Basilica were surrendered ...

``Next morning the Basilica was surrounded by soldiers; but it
was reported, that these soldiers had sent to the Emperor to tell him,
that if he wished to come abroad he might, and that they would attend
him, if he was going to the assembly of the Catholics: otherwise, that
they would go to that which would be held by St. Ambrose. Indeed, the
soldiers were all Catholics, as well as the citizens of Milan: there
were so few heretics there, except a few officers of the emperor and
some Goths.

``St. Ambrose was continuing his discourse, when he was told
that the Emperor had withdrawn the soldiers from the Basilica, and
that he had restored to the tradesmen the fines which he had exacted
from them. This news gave joy to the people, who expressed their
delight with applauses and thanksgivings; the soldiers themselves were
eager to bring the news, throwing themselves on the altars, and
kissing them in token of peace. ``Fleury's \emph{Hist}. xviii. 41,
42, Oxf. trans.

20. CHRISTENDOM GENERALLY.
St. Hilary to Constantius: ``Not only in words, but in tears, we
beseech you to save the Catholic Churches from any longer continuance
of these most grievous injuries, and of their present intolerable
persecutions and insults, which moreover they are enduring, monstrous as it is, from our brethren. Surely your clemency should
listen to the voice of those who cry out so loudly, 'I am a Catholic,
I have no wish to be a heretic.' It should seem equitable to your
sanctity, most glorious Augustus, that they who fear the Lord God and
His judgment should not be polluted and contaminated with execrable
blasphemies, but should have liberty to follow those Bishops and
prelates who both observe inviolate the laws of charity, and who
desire a perpetual and sincere peace. It is impossible, it is
unreasonable, to mix true and false, to confuse light and darkness,
and bring into union, of whatever kind, night and day. Give permission
to the populations to hear the teaching of the pastors whom they have
wished, whom they fixed on, whom they have chosen, to attend their
celebration of the divine mysteries, to offer prayers through them for
your safety and prosperity.'' \emph{ad Const}. i. 1, 2.

In drawing out this comparison between the conduct of the Catholic
Bishops and that of their flocks during the Arian troubles, I must not
be understood as intending any conclusion inconsistent with the
infallibility of the Ecclesia docens, (that is, the Church when
teaching) and with the claim of the Pope and the Bishops to constitute
the Church in that aspect. I am led to give this caution, because, for
the want of it, I was seriously misunderstood in some quarters on my
first writing on the above subject in the \emph{Rambler} Magazine of
May, 1859. But on that occasion I was writing simply historically, not
doctrinally, and, while it is historically true, it is in no sense
doctrinally false, that a Pope, as a private doctor, and much more
Bishops, when not teaching formally, may err, as we find they did err
in the fourth century. Pope Liberius might sign a Eusebian formula at
Sirmium, and the mass of Bishops at Ariminum or elsewhere, and yet
they might, in spite of this error, be infallible in their \emph{ex
cathedrâ} decisions.

The reason of my being misunderstood arose from two or three
clauses or expressions which occurred in the course of my remarks,
which I should not have used had I anticipated how they would be
taken, and which I avail myself of this opportunity to explain and
withdraw. First, I will quote the passage which bore a meaning which I
certainly did not intend, and then I will note the phrases which seem
to have given this meaning to it. It will be seen how little, when
those phrases are withdrawn, the sense of the passage, as I intended
it, is affected by the withdrawal. I said then:—``It is not a
little remarkable, that, though, historically speaking, the fourth
century is the age of doctors, illustrated, as it is, by the Saints
Athanasius, Hilary, the two Gregories, Basil, Chrysostom, Ambrose,
Jerome, and Augustine, (and all those saints bishops also), except
one, nevertheless in that very day the Divine tradition committed to
the infallible Church was proclaimed and maintained far more by the
faithful than by the Episcopate.

``Here of course I must explain:—in saying this then,
undoubtedly I am not denying that the great body of the Bishops were
in their internal belief orthodox; nor that there were numbers of
clergy who stood by the laity and acted as their centres and guides;
nor that the laity actually received their faith, in the first
instance, from the Bishops and clergy; nor that some portions of the
laity were ignorant, and other portions were at length corrupted by
the Arian teachers, who got possession of the sees, and ordained an
heretical clergy:—but I mean still, that in that time of immense
confusion the divine dogma of our Lord's divinity was proclaimed,
enforced, maintained, and (humanly speaking) preserved, far more by
the ``Ecclesia docta'' than by the ``Ecclesia docens;''
that the body of the Episcopate was unfaithful to its commission,
while the body of the laity was faithful to its baptism; that at one
time the pope, at other times a patriarchal, metropolitan, or other
great see, at other times general councils, said what they
should not have said, or did what obscured and compromised revealed
truth; while, on the other hand, it was the Christian people, who,
under Providence, were the ecclesiastical strength of Athanasius,
Hilary, Eusebius of Vercellæ, and other great solitary confessors,
who would have failed without them ...

``On the one hand, then, I say, that there was a temporary
suspense of the functions of the 'Ecclesia docens.' The body of
Bishops failed in their confession of the faith. They spoke variously,
one against another; there was nothing, after Nicæa, of firm,
unvarying, consistent testimony, for nearly sixty years ...

``We come secondly to the proofs of the fidelity of the laity,
and the effectiveness of that fidelity, during that domination of
Imperial heresy, to which the foregoing passages have related.''

The three clauses which furnished matter of objection were these:—I
said, (1), that ``there was a temporary suspense of the functions
of the 'Ecclesia docens;''' (2), that ``the body of Bishops
failed in their confession of the faith.'' (3), that ``general
councils, \&c., said what they should not have said, or did what
obscured and compromised revealed truth.''

(1). That ``there was a temporary \emph{suspense} of the
functions of the Ecclesia docens'' is not true, if by saying so is
meant that the Council of Nicæa held in 325 did not sufficiently
define and promulgate for all times and all places the dogma of our
Lord's divinity, and that the notoriety of that Council and the voices
of its great supporters and maintainers, as Athanasius, Hilary,
\&c., did not bring home the dogma to the intelligence of the
faithful in all parts of Christendom. But what I meant by
``suspense'' (I did not say ``suspension,''
purposely,) was only this, that there was no authoritative utterance
of the Church's infallible voice in matter of fact between the Nicene
Council, A.D. 325, and the Council of
Constantinople, A.D. 381, or, in the words which
I actually used, ``there was nothing after Nicæa of firm,
unvarying, consistent testimony for nearly sixty years.'' As
writing before the Vatican Definition of 1870, I did not lay stress
upon the Roman Councils under Popes Julius and Damasus \footnote{A distinguished theologian infers from my
words that I deny that ``the Church is in every time the activum
instrumentum docendi.'' But I do not admit the fairness of this
inference. Distinguo: activum instrumentum docendi virtuale, C.
Actuale, N. The Ecumenical Council of 325 was an effective authority
in 341, 351, and 359, though at those dates the Arians were in the
seats of teaching. Fr. Perrone agrees with me. 1. He reckons the
``fidelium sensus'' among the ``instrumenta traditionis.''
(\emph{Immac. Concept}. p. 139.) 2. He contemplates, nay he instances,
the case in which the ``sensus fidelium'' supplies, as the
``instrumentum,'' the absence of the other instruments, the \emph{magisterium}
of the Church, as exercised at Nicæa, being always supposed. One of
his instances is that of the dogma de visione Dei beatificâ. He says:
``Certe quidem in Ecclesiâ non deerat quoad hunc fidei articulum
divina traditio; alioquin, nunquam is definiri potuisset: verum non
omnibus illa erat comperta: divina eloquia haud satis in re sunt
conspicua; Patres, ut vidimus, in varias abierunt sententias;
liturgiæ ipsæ non modicam præ se ferunt difficultatem. His omnibus
succurrit juge Ecclesiæ magisterium; communis præterea fidelium
sensus.'' p. 148.}.

(2). That ``the \emph{body} of Bishops failed in their
confession of the faith,'' p. 17. Here, if the word
``body'' is used in the sense of the Latin ``corpus,''
as ``corpus'' is used in theological treatises, and as it
doubtless would be translated for the benefit of readers ignorant of
the English language, certainly this would be a heretical statement.
But I meant nothing of the kind. I used it in the vague, familiar,
genuine sense of which Johnson gives instances in his dictionary, as
meaning ``the great preponderance,'' or, ``the mass''
of Bishops, viewing them in the main or the gross, as a \emph{cumulus} of
individuals. Thus Hooker says, ``Life and death have divided
between them the whole body of mankind;'' Clarendon, after
speaking of the van of the king's army, says, ``in the body was
the king and the prince:'' and Addison speaks of
``navigable rivers, which ran up into the body of Italy.'' In
this sense it is true historically that the body of Bishops failed in
their confession. Tillemont, quoting from St. Gregory Nazianzen, says,
``La souscription (Arienne) etait une des dispositions necessaires
pour entrer et pour se conserver dans l'episcopat. L'encre était
toujours toute prête, et l'accusateur aussi. Ceux qui avaient paru
invincibles jusques alors, céderent à cette tempête. Si leur esprit
ne tomba pas dans l'heresie, leur main néanmoins y consentit ... Peu
d'Evêques s'exemterent de ce malheur, n' y ayant eu que ceux que leur
propre bassesse faisait negliger, ou que leur vertu fit resister
genereusement, et que Dieu conserva afin qu'il restât encore quelque
semence et quelque racine pour faire refleurir Israel.'' T. vi. p.
499. In St. Gregory's own words, [\emph{pl\={e}n olig\={o}n agan,
pantes tou kairou gegonasi; tosouton all\={e}l\={o}n
dienenkontes, hoson tous men proteron, tous de husteron touto pathein}].
Orat. xxi. 24. p. 401. \emph{Ed. Bened}.

(3). That ``\emph{general} councils said what they should not
have said, and did what obscured and compromised revealed truth.''
Here again the question to be determined is what is meant by the word
``general.'' If I meant by ``general'' ecumenical, I
should have spoken as no Catholic can speak; but ecumenical Councils
there were none between 325 and 381, and so I could not be referring
to any; and in matter of fact I used the word ``general'' in \emph{contrast}
to ``ecumenical,'' as I had used it in Tract No. 90, and as
Bellarmine uses the word. He makes a fourfold division of
``general Councils,'' viz., those which are approbata;
reprobata; partim confirmata, partim reprobata; and nec manifeste
probata nec manifeste reprobata. Among the ``reprobata'' he
placed the Arian Councils. They were quite large enough to be called
``generalia;'' the twin Councils of Seleucia and Ariminum
numbering as many as 540 Bishops. When I spoke then of ``general
councils compromising revealed truth,'' I spoke of the Arian or
Eusebian Councils, not of the Catholic.

I hope this is enough to observe on this subject.

\textbf{Note 6. Chronology of the Councils}

(\emph{Vide supra}, p. 271.)

A\textsc{s} the direct object of the foregoing Volume was to exhibit
the doctrine, temper, and conduct of the Arians in the fourth century
rather than to write their history, there is much incidental confusion
in the order in which the events which it includes are brought before
the reader. However, in truth, the chronology of the period is by no
means clear, and the author may congratulate himself that, by the
scope of his work, he is exempt from the necessity of deciding
questions relative to it, on which ancient testimonies and modern
critics are in hopeless variance both with themselves and with each
other.

Accordingly, he has chosen one authority, the accurate Tillemont,
and followed him almost throughout. Here, however, he thinks it well
to subjoin some tables on the subject, taken from the Oxford Library
of the Fathers, which delineate the main outline of the history, while
they vividly illustrate the difficulty of determining in detail the
succession of dates.

PRINCIPAL EVENTS
BETWEEN A.D. 325 AND
A.D. 381,

IN CHRONOLOGICAL ORDER

\textbf{1. From 325 to 337}

(Mainly from Tillemont.)

\textbf{A.D.}

325.
(From June 19 to August 25.) COUNCIL
OF NICÆA.

Arius and his partisans anathematized and banished, 

Arius to Illyricum. The Eusebians subscribe to the

\emph{Homoüsion}.

326.
Athanasius raised to the See of Alexandria at
the age

of about 30.

328-9.
Eusebius of Nicomedia in favour with
Constantine.

330.
An Arian priest gains the ear of Constantine,
who

recalls Arius from exile to Alexandria.

331.
Athanasius refuses to restore him to communion.

Eustathius deposed by the Eusebians on a charge

of Sabellianism; other Bishops deposed.

334.
Council of Cæsarea against Athanasius, who
refuses

to attend it.

335.
Council of Tyre and Jerusalem, in which Arius
and

the Arians are formerly readmitted. Athanasius,

forced by the emperor to attend, abruptly leaves it

in order to appeal to Constantine. THE EUSEBIANS

DEPOSE ATHANASIUS,
AND CONSTANTINE BANISHES

HIM TO TREVES.

336.
Eusebians hold a Council at Constantinople to
condemn

Marcellus on the ground of his Sabellianism; and to

recognize Arius. DEATH OF
ARIUS.

337.
DEATH OF
CONSTANTINE. The Eusebian Constantius

succeeds him in the East, the orthodox Constans and

Constantine in the West.

\textbf{2. From 337 to 342}

338
Exiles
recalled by the three new Emperors.

(End of June.) Athanasius leaves Treves for Alexandria.

(\emph{From Valesius,
Schel- strate, Pagi, Montfau- con and S. Basnage}.)
(From \emph{Baronius and Petavius}.)
(\emph{From Tillemont and Papebroke}.)

339
Eusebius sends to Pope

Julius for a Council.
Eusebius, \&c.

COUNCIL OF
ALEX-

ANDRIA DEFENDS

ATHANASIUS TO

THE POPE.
Eusebius, \&c.

COUNCIL OF ALEX-

ANDRIA, \&c.

(Sept.) \emph{Athanasius goes to Rome}.

\footnote{The events in italics are grounded on an hypothesis of
the authors who introduce them, that Athanasius made two journeys to
Rome, which they adopt in order to lighten the difficulties of the
chronology.}

340
COUNCIL OF
ALEX-

ANDRIA DEFENDS

ATHANASIUS TO

THE POPE.
Papal Legates sent to

Antioch from Rome.

(Early in year) \emph{Athana- sius goes to Rome}.
Papal Legates, \&c.

(End of year) \emph{Athana- sius returns to Alex- andria.} 341
(Christmas or before

Sept.)

COUNCIL OF THE

DEDICATION AT

ANTIOCH (Eusebian),

not in order to

anticipate the Council

at Rome.

(Lent) THE ARIAN

GREGORY IN

ALEXANDRIA.

(March-May.) ATHA-

NASIUS ESCAPES TO

ROME, after the

Council of the

Dedication,

immediately before

or after the Papal

Legates set out

from Rome.
COUNCIL OF

DEDICATION, \&c.,

In order to anticipate

the Council at Rome.

The Papal Legates leave

Antioch.

\emph{A Roman Council}.
(Christmas or before

Sept.) COUNCIL.,

\&c.

(Lent.) THE ARIAN

GREGORY, \&c.

ATHANASIUS

ESCAPES, \&c.

The Papal Legates, \&c.

The Papal Legates

arrive at Rome during

the Council there.

(June till Aug. or Sept.)

COUNCIL OF ROME.

THE POPE'S
LETTER

TO THE EUSEBIANS

immediately after the

Council.

342
(April or June.)

The Papal Legates arrive

at Antioch.

(Jan.) The Papal Legates

leave Antioch.

(March or April.) The

Papal Legates arrive at

Rome.

COUNCIL OF ROME.

THE POPE’S
LETTER

TO THE EUSEBIANS.
(End of year) \emph{Athanasius returns to Alexandria.}\\
(Or beginning Lent.)

THE ARIAN
GREGORY

IN ALEXANDRIA.

The Papal Legates arrive

at Rome.

ATHANASIUS ESCAPES

TO ROME
shortly after

the Roman Council

there.

COUNCIL OF ROME.

THE POPE’S
LETTER

TO THE EUSEBIANS,

\&c.

\textbf{3. From 342 to 351}

(Mainly from Tillemont.)

345.
COUNCIL OF
ANTIOCH (Eusebian), at which the

Macrostich is drawn up.

347.
G\textsc{reat} C\textsc{ouncil}
OF S\textsc{ardica}, at the
instance

of the orthodox Constans. Council of Milan against

Photinus. Ursacius and Valens sue for reconciliation

to the Church.

349.
Council of Jerusalem, at which Athanasius is
present.

Athanasius returns to Alexandria. Ursacius and

Valens recant, and are reconciled at Rome. Council

at Sirmium or at Rome against Photinus.

350.
DEATH OF
CONSTANS. The Eusebian Constantius

sole Emperor.

351.
GREAT COUNCIL
OF SIRMIUM, at which
Photinus

is deposed. First Sirmian creed, \&c.

\textbf{4. From 351 to 361}

B

a

r

o

n

i

u

s
P

e

t

a

v

i

u

s
V

a

l

e

s

i

u

s
P

a

g

i
B

a

s

n

a

g

e
T

i

l

l

e

m

o

n

t
N.

A

l

e

x

a

n

d

e

r
C

o

u

s

t

a

n

t
M

o

n

t

f

o

u

c

o

n
M

a

n

s

i
M

a

m

a

c

h

i
Z

a

c

c

a

r

i

a

1. Great Council of Sirmium
357
351
351
351
351
351
351
351
351
357/8
351
351

2. Photinus deposed
357
351

351
351
351
351
351
351
358
351
351

3. First Sirmian Creed

(Semi-Arian)
357
351

351
351
351
351

351
358
351

4. Signed by Pope Liberius

with a condemnation of

Athanasius
357
o
o
o
o
357/8
358
357
o
358
o
357

5. Council of Arles (Eusebian)

Athanasius condemned
353
353
353
353
353
353/4
353
353

354
353

6. Great Council of Milan

(Eusebian) Athanasius

condemned
355
(communiter)
355

7. Rise of the Eunomians
356
356

356

8. Syrianus in Alexandria,
and

George in Cappadocia
356
356
356
356
356
356
355
356
356
356/7
356

9. Council of Beziers.
Hilary

deposed and banished
356
355

356
356
356
356
356?

355
356

10. Fresh Council or 

Conference at Sirmium
o
357
357
357
357
357
357
357
357
359
357
357

11. Second Sirmian Creed, the

blasphemy of Potamius

and Hosius (Homœan, if

not Anomœan)
357
357

357
357
357
357
357
357
359
357

12. Signed by Hosius, but

without condemning

Athanasius
357
357
357
357
357
357
357
357
357
355
357
357

13. Signed by Liberius, with a

condemnation of

Athanasius
o
357

357
o
o
o

14. Another or an altered

Creed signed by Liberius

with condemnation of

Athanasius
o
357
357
o
o
o
o
o
o
o
o
o

15. Council of Antioch in

favour of Eunomius

358
358

358

16. Its Creed (Anomœan)

358
358

358

17. Council of Ancyra of 12

Bishops

18. Its Creed (Semi-Arian)

against both the

Homoüsian and the

Anomœan, signed by

Liberius
357

357
358

358

358

358
358

358
358

358
358
358

358

359

359

358

358

19. Fresh Council or

Conference at Sirmium
o

359
358
358/9
358
359
359
359
359
359

20. Third Sirmian. Creed

(Homœan) drawn up by

Semi-Arians
357
358
359
358
358/9
359
359
359
359
359
359
359

21. Signed by Liberius
o
o
358?
358
358
o

o
o

o
o

22. BI-PARTITE
COUNCIL OF

ARIMINUM (Homœan) and

of Seleucia (Semi-Arian)
359
(communiter)

23. Council of Constantinople

(Homœan)

360
359/

60

360
359

360
359
359

24. Council of Antioch

(Anomœan)

361
360

361
361
361

25. DEATH OF
CONSTANTIUS
361

(communiter)

\textbf{5. From 361 to 381}

(From Tillemont.)

362.
COUNCIL OF
ALEXANDRIA.

365.
Council of Lampsacus (Semi-Arian or Macedonian).

366.
Macedonian Bishops reconciled to the Church at
Rome.

367.
Council of Tyre for the same purpose.

373.
DEATH OF
ATHANASIUS.

381.
SECOND ŒCUMENICAL
COUNCIL AT CONSTANTINOPLE.

\textbf{Note 7. Omissions in the Text of the
Third Edition}

(\emph{Vide Advertisement}.)

Here follow the two sentences, which, as was stated in the
Advertisement to this Edition, have forfeited their place in the text:—

1. Supra. p. 11 (p. 12, 1st Ed.), after ``external
observers,'' the text proceeded. ``Presenting then the
characters of a religion, sufficiently correct in the main articles of
faith to satisfy the reason, and yet indulgent to the carnal nature of
man, Judaism occupied that place in the Christian world, which has
since been filled by a corruption of Christianity itself. While its
adherents manifested a rancorous malevolence,'' \&c.

2. Supra, p. 393 (p. 421, 1st Ed.), after ``his place could
nowhere be found,'' the text proceeded. ``Even the Papal
Apostasy, which seems at first sight an exception to this rule, has
lasted but the same proportion of the whole duration of Christianity,
which Arianism occupied in its day; that is, if we date it, as in
fairness we ought, from the fatal Council of Trent. And, as to the
present perils,'' \&c.
